% Generated mechanically from frozen input inputs/p05/ja/flat.tex.
% Do not edit this generated file; edit the bound locale source instead.

% OK, start here.
%

\setcounter{chapter}{37}
\chapter{平坦性についてさらに}


\phantomsection
\label{flat-section-phantom}




\section{序論}
\label{flat-section-introduction}

\noindent
本章では、平坦加群とスキームの平坦射に関するいくつかの高度な結果、
およびその応用を論じる。平坦性に関する結果の大部分は Raynaud と Gruson の
論文 \cite{GruRay} に見いだすことができる。

\medskip\noindent
本章を読む前に、読者には次の結果に目を通すことを勧める
（この一覧は以前の結果への案内も兼ねている）。
\begin{enumerate}
\item 平坦加群の一般論については
Algebra, Section \ref{algebra-section-flat} を参照せよ。
\item $\text{Tor}$ 群と平坦性の関係については
Algebra, Section \ref{algebra-section-tor} を参照せよ。
\item 平坦性の判定法については
Algebra, Section \ref{algebra-section-criteria-flatness}
（Noetherian の場合）、
Algebra, Section \ref{algebra-section-flatness-artinian}
（Artinian の場合）、
Algebra, Section \ref{algebra-section-more-flatness-criteria}
（非 Noetherian の場合）、そして最後に
More on Morphisms, Section \ref{more-morphisms-section-criterion-flat-fibres}
を参照せよ。
\item 総称平坦性については
Algebra, Section \ref{algebra-section-generic-flatness}
および Morphisms, Section \ref{morphisms-section-generic-flatness} を参照せよ。
\item 平坦軌跡の開性については
Algebra, Section \ref{algebra-section-open-flat}
および More on Morphisms, Section \ref{more-morphisms-section-open-flat}
を参照せよ。
\item 平坦化については More on Algebra, Sections
\ref{more-algebra-section-flattening},
\ref{more-algebra-section-flattening-artinian},
\ref{more-algebra-section-flattening-local-base},
\ref{more-algebra-section-flattening-local-source-base}, および
\ref{more-algebra-section-flattening-Noetherian-complete-local} を参照せよ。
\item さらに More on Algebra, Sections
\ref{more-algebra-section-descent-flatness-integral},
\ref{more-algebra-section-torsion-flat},
\ref{more-algebra-section-flat-finite}, および
\ref{more-algebra-section-blowup-flat} に追加の結果がある。
\end{enumerate}
平坦性に関する内容の応用として、次の話題を論じる：Grothendieck の存在定理の
非 Noetherian 版、吹き上げと平坦性、コンパクト化に関する Nagata の定理、
h 位相、吹き上げ正方形と降下、弱正規化、正標数におけるベクトル束の降下、
および h 位相における完全複体の局所構造である。






\section{エタール局所化に関する補題}
\label{flat-section-etale-localization}

\noindent
本節では、本章の後半で有用となるエタール局所化に関するいくつかの補題を
列挙する。初読時には本節を飛ばしてよい。

\begin{lemma}
\label{flat-lemma-lift-etale}
$i : Z \to X$ をアフィンスキームの閉埋め込みとする。
$Z' \to Z$ をエタール射とし、$Z'$ はアフィンであるとする。
このとき、$X'$ がアフィンであり、$Z$ 上のスキームとして
$Z' \cong Z \times_X X'$ となるエタール射 $X' \to X$ が存在する。
\end{lemma}

\begin{proof}
Algebra, Lemma \ref{algebra-lemma-lift-etale} を参照せよ。
\end{proof}

\begin{lemma}
\label{flat-lemma-etale-at-point}
次の可換図式を考える：
$$
\xymatrix{
X \ar[d] & X' \ar[l] \ar[d] \\
S & S' \ar[l]
}
$$
$X' \to X$ および $S' \to S$ はエタール射であるとする。
$s' \in S'$ を点とする。このとき
$$
X' \times_{S'} \Spec(\mathcal{O}_{S', s'})
\longrightarrow
X \times_S \Spec(\mathcal{O}_{S', s'})
$$
はエタールである。
\end{lemma}

\begin{proof}
$X$ 上エタールなスキーム間の射として $X' \to X_{S'}$ はエタールである。
Morphisms, Lemma \ref{morphisms-lemma-etale-permanence} を参照せよ。
また、エタール射の基底変換はエタールである。
Morphisms, Lemma \ref{morphisms-lemma-base-change-etale} を参照せよ。
\end{proof}

\begin{lemma}
\label{flat-lemma-etale-flat-up-down}
$X \to T \to S$ をスキームの射とし、$T \to S$ はエタールであるとする。
$\mathcal{F}$ を準連接 $\mathcal{O}_X$-加群とし、$x \in X$ を点とする。
このとき
$$
\mathcal{F}\text{ は }S\text{ 上 }x\text{ で平坦}
\Leftrightarrow
\mathcal{F}\text{ は }T\text{ 上 }x\text{ で平坦}
$$
特に、$\mathcal{F}$ が $S$ 上平坦であることと、$\mathcal{F}$ が
$T$ 上平坦であることは同値である。
\end{lemma}

\begin{proof}
エタール射は平坦射なので
（Morphisms, Lemma \ref{morphisms-lemma-etale-flat} を参照せよ）、含意
``$\Leftarrow$'' は Algebra, Lemma \ref{algebra-lemma-composition-flat}
から従う。逆に、$\mathcal{F}$ が $S$ 上 $x$ で平坦であると仮定する。
$X$ では $x$ の上にあり、$T$ では $x$ の $T$ における像の上にある点を
$\tilde x \in X \times_S T$ と記す。このとき
$(X \times_S T \to X)^*\mathcal{F}$ は
$\text{pr}_2 : X \times_S T \to T$ によって $T$ 上 $\tilde x$ で平坦である。
Morphisms, Lemma \ref{morphisms-lemma-base-change-module-flat} を参照せよ。
対角射 $\Delta_{T/S} : T \to T \times_S T$ は開埋め込みである。これは
Morphisms, Lemmas \ref{morphisms-lemma-diagonal-unramified-morphism} および
\ref{morphisms-lemma-etale-smooth-unramified} を組み合わせればよい。
したがって $X$ は $X \times_S T$ の開部分スキームと同一視される。
この開部分への $\text{pr}_2$ の制限は与えられた射 $X \to T$ であり、
$\tilde x$ はこの開部分の点 $x$ に対応し、
$(X \times_S T \to X)^*\mathcal{F}$ のこの開部分への制限は
$\mathcal{F}$ である。ゆえに $\mathcal{F}$ は $T$ 上 $x$ で平坦である。
\end{proof}

\begin{lemma}
\label{flat-lemma-etale-flat-up-down-local-ring}
$T \to S$ をエタール射とし、$t \in T$ の像を $s \in S$ とする。
$M$ を $\mathcal{O}_{T, t}$-加群とする。このとき
$$
M\text{ は }\mathcal{O}_{S, s}\text{ 上平坦}
\Leftrightarrow
M\text{ は }\mathcal{O}_{T, t}\text{ 上平坦}.
$$
\end{lemma}

\begin{proof}
$S$ を $s$ のアフィン近傍で置き換え、次に $T$ を $t$ のアフィン近傍で
置き換えてよい。
$\mathcal{F} = (\Spec(\mathcal{O}_{T, t}) \to T)_*\widetilde M$ とおく。
これは $T$ 上の準連接層であり
（Schemes, Lemma \ref{schemes-lemma-push-forward-quasi-coherent} を参照するか、
直接論じればよい）、$t$ における茎は $M$ である（詳細は省略する）。
Lemma \ref{flat-lemma-etale-flat-up-down} を適用せよ。
\end{proof}

\begin{lemma}
\label{flat-lemma-flat-up-down-henselization}
$S$ をスキーム、$s \in S$ を点とする。$\mathcal{O}_{S, s}^h$
（それぞれ $\mathcal{O}_{S, s}^{sh}$）を Hensel 化
（それぞれ狭義 Hensel 化）と記す。Algebra, Definition
\ref{algebra-definition-henselization} を参照せよ。
$M^{sh}$ を $\mathcal{O}_{S, s}^{sh}$-加群とする。次は同値である：
\begin{enumerate}
\item $M^{sh}$ は $\mathcal{O}_{S, s}$ 上平坦である。
\item $M^{sh}$ は $\mathcal{O}_{S, s}^h$ 上平坦である。
\item $M^{sh}$ は $\mathcal{O}_{S, s}^{sh}$ 上平坦である。
\end{enumerate}
$M^{sh} = M^h \otimes_{\mathcal{O}_{S, s}^h} \mathcal{O}_{S, s}^{sh}$
ならば、これらは次とも同値である：
\begin{enumerate}
\item[(4)] $M^h$ は $\mathcal{O}_{S, s}$ 上平坦である。
\item[(5)] $M^h$ は $\mathcal{O}_{S, s}^h$ 上平坦である。
\end{enumerate}
$M^h = M \otimes_{\mathcal{O}_{S, s}} \mathcal{O}_{S, s}^h$ ならば、
これらはさらに次とも同値である：
\begin{enumerate}
\item[(6)] $M$ は $\mathcal{O}_{S, s}$ 上平坦である。
\end{enumerate}
\end{lemma}

\begin{proof}
More on Algebra, Lemma
\ref{more-algebra-lemma-dumb-properties-henselization}
により、局所環の射
$\mathcal{O}_{S, s} \to \mathcal{O}_{S, s}^h \to \mathcal{O}_{S, s}^{sh}$
は忠実平坦である。したがって (3) $\Rightarrow$ (2) $\Rightarrow$ (1) および
(5) $\Rightarrow$ (4) は
Algebra, Lemma \ref{algebra-lemma-composition-flat} から従う。
忠実平坦性により、同値 (6) $\Leftrightarrow$ (5) および
(5) $\Leftrightarrow$ (3) は
Algebra, Lemma \ref{algebra-lemma-flatness-descends} から従う。
ゆえに次を示せば十分である：
(1) $\Rightarrow$ (2) $\Rightarrow$ (3) および
(4) $\Rightarrow$ (5).
これらを証明するにあたり、$S$ はアフィンスキームであると仮定してよい。

\medskip\noindent
(1) を仮定する。Lemma \ref{flat-lemma-etale-flat-up-down-local-ring} により、
任意のエタール近傍 $(T, t) \to (S, s)$ に対して $M^{sh}$ は
$\mathcal{O}_{T, t}$ 上平坦である。$\mathcal{O}_{S, s}^h$ および
$\mathcal{O}_{S, s}^{sh}$ は $\mathcal{O}_{T, t}$ の形の局所環の有向余極限である
（次を参照せよ：
Algebra, Lemmas \ref{algebra-lemma-henselization-different}
および \ref{algebra-lemma-strict-henselization-different}）。したがって
Algebra, Lemma \ref{algebra-lemma-colimit-rings-flat} により、$M^{sh}$ は
$\mathcal{O}_{S, s}^h$ および $\mathcal{O}_{S, s}^{sh}$ 上平坦である。
ゆえに (1) は (2) と (3) を含意する。もちろん、
$\mathcal{O}_{S, s}$ を $\mathcal{O}_{S, s}^h$ で置き換えれば、これは
(2) $\Rightarrow$ (3) も含意する。同じ議論で
(4) $\Rightarrow$ (5) が証明される。
\end{proof}

\begin{lemma}
\label{flat-lemma-tor-amplitude-up-down-henselization}
$S$ をスキーム、$s \in S$ を点とする。$\mathcal{O}_{S, s}^h$
（それぞれ $\mathcal{O}_{S, s}^{sh}$）を Hensel 化
（それぞれ狭義 Hensel 化）と記す。Algebra, Definition
\ref{algebra-definition-henselization} を参照せよ。
$M^{sh}$ を $D(\mathcal{O}_{S, s}^{sh})$ の対象とし、
$a, b \in \mathbf{Z}$ とする。次は同値である：
\begin{enumerate}
\item $M^{sh}$ は $\mathcal{O}_{S, s}$ 上 $[a, b]$ 内の Tor 振幅をもつ。
\item $M^{sh}$ は $\mathcal{O}_{S, s}^h$ 上 $[a, b]$ 内の Tor 振幅をもつ。
\item $M^{sh}$ は $\mathcal{O}_{S, s}^{sh}$ 上 $[a, b]$ 内の Tor 振幅をもつ。
\end{enumerate}
$M^{sh} =
M^h \otimes_{\mathcal{O}_{S, s}^h}^\mathbf{L} \mathcal{O}_{S, s}^{sh}$
で $M^h \in D(\mathcal{O}_{S, s}^h)$ ならば、これらは次とも同値である：
\begin{enumerate}
\item[(4)] $M^h$ は $\mathcal{O}_{S, s}$ 上 $[a, b]$ 内の Tor 振幅をもつ。
\item[(5)] $M^h$ は $\mathcal{O}_{S, s}^h$ 上 $[a, b]$ 内の Tor 振幅をもつ。
\end{enumerate}
$M^h = M \otimes_{\mathcal{O}_{S, s}}^\mathbf{L} \mathcal{O}_{S, s}^h$
で $M \in D(\mathcal{O}_{S, s})$ ならば、これらはさらに次とも同値である：
\begin{enumerate}
\item[(6)] $M$ は $\mathcal{O}_{S, s}$ 上 $[a, b]$ 内の Tor 振幅をもつ。
\end{enumerate}
\end{lemma}

\begin{proof}
More on Algebra, Lemma
\ref{more-algebra-lemma-dumb-properties-henselization}
により、局所環の射
$\mathcal{O}_{S, s} \to \mathcal{O}_{S, s}^h \to \mathcal{O}_{S, s}^{sh}$
は忠実平坦である。したがって (3) $\Rightarrow$ (2) $\Rightarrow$ (1) および
(5) $\Rightarrow$ (4) は More on Algebra, Lemma
\ref{more-algebra-lemma-flat-push-tor-amplitude} から従う。
忠実平坦性により、同値 (6) $\Leftrightarrow$ (5) および
(5) $\Leftrightarrow$ (3) は More on Algebra, Lemma
\ref{more-algebra-lemma-flat-descent-tor-amplitude} から従う。
ゆえに (1) $\Rightarrow$ (3)、(2) $\Rightarrow$ (3)、および
(4) $\Rightarrow$ (5) を示せば十分である。

\medskip\noindent
(1) を仮定する。特に $M^{sh}$ のコホモロジーは次数 $< a$ および
$> b$ で消える。したがって、$i > b$ に対して $P^i = 0$ となる自由
% P05-EN-CH38-0001
$\mathcal{O}_{X, x}^{sh}$-加群の複体 $P^\bullet$ によって $M^{sh}$ を表せる
（例えば、非常に一般的な Derived Categories, Lemma
\ref{derived-lemma-subcategory-left-resolution} を参照せよ）。
すべての $n$ に対して $P^n$ は $\mathcal{O}_{S, s}$ 上平坦である。
$\Coker(d_P^{a - 1})$ を考える。More on Algebra, Lemma
\ref{more-algebra-lemma-last-one-flat}
により、これは平坦な $\mathcal{O}_{S, s}$-加群である。したがって
Lemma \ref{flat-lemma-flat-up-down-henselization} により、これは平坦な
$\mathcal{O}_{S, s}^{sh}$-加群である。ゆえに
$\tau_{\geq a}P^\bullet$ は $M^{sh}$ を
% P05-EN-CH38-0002
$D(\mathcal{O}_{S, s}^{sh}$ で表す平坦 $\mathcal{O}_{S, s}^{sh}$-加群の
複体であり、$M^{sh}$ は $[a, b]$ 内の Tor 振幅をもつことが分かる。
More on Algebra, Lemma \ref{more-algebra-lemma-tor-amplitude} を参照せよ。
ゆえに (1) は (3) を含意する。もちろん、$\mathcal{O}_{S, s}$ を
$\mathcal{O}_{S, s}^h$ で置き換えれば、これは (2) $\Rightarrow$ (3) も
含意する。同じ議論で (4) $\Rightarrow$ (5) が証明される。
\end{proof}

\begin{lemma}
\label{flat-lemma-finite-flat-weak-assassin-up-down}
$g : T \to S$ をスキームの有限平坦射とする。
$\mathcal{G}$ を準連接 $\mathcal{O}_S$-加群とし、$t \in T$ の像を
$s \in S$ とする。このとき
$$
t \in \text{WeakAss}(g^*\mathcal{G})
\Leftrightarrow
s \in \text{WeakAss}(\mathcal{G})
$$
\end{lemma}

\begin{proof}
含意 ``$\Leftarrow$'' は Divisors, Lemma
\ref{divisors-lemma-weakly-ass-pullback} から直ちに従う。
$t \in \text{WeakAss}(g^*\mathcal{G})$ と仮定する。
$\Spec(A) \subset S$ を $s$ のアフィン開近傍とし、$\mathcal{G}$ を
$A$-加群 $M$ に付随する準連接層とする。$s$ に対応する素イデアルを
$\mathfrak p \subset A$ とする。$g$ は有限平坦なので、ある有限平坦
$A$-代数 $B$ に対して $g^{-1}(\Spec(A)) = \Spec(B)$ である。
$g^*\mathcal{G}$ は $B$-加群 $M \otimes_A B$ に付随する準連接
$\mathcal{O}_{\Spec(B)}$-加群であり、$g_*g^*\mathcal{G}$ は
$A$-加群 $M \otimes_A B$ に付随する準連接
$\mathcal{O}_{\Spec(A)}$-加群である。Algebra, Lemma
\ref{algebra-lemma-finite-flat-local} により、ある整数 $n \geq 0$ に対して
$B_{\mathfrak p} \cong A_{\mathfrak p}^{\oplus n}$ である。
$s$ の上に $T$ の点が少なくとも一つ存在すると仮定したので $n \geq 1$ である。
したがって茎を調べれば
$$
s \in \text{WeakAss}(\mathcal{G})
\Leftrightarrow
s \in \text{WeakAss}(g_*g^*\mathcal{G})
$$
である。いま、仮定 $t \in \text{WeakAss}(g^*\mathcal{G})$ は Divisors,
Lemma \ref{divisors-lemma-weakly-associated-finite} により
$s \in \text{WeakAss}(g_*g^*\mathcal{G})$ を含意する。したがって上の同値から
$s \in  \text{WeakAss}(\mathcal{G})$ である。
\end{proof}

\begin{lemma}
\label{flat-lemma-etale-weak-assassin-up-down}
$h : U \to S$ をスキームのエタール射とする。
$\mathcal{G}$ を準連接 $\mathcal{O}_S$-加群とし、$u \in U$ の像を
$s \in S$ とする。このとき
$$
u \in \text{WeakAss}(h^*\mathcal{G})
\Leftrightarrow
s \in \text{WeakAss}(\mathcal{G})
$$
\end{lemma}

\begin{proof}
$S$ と $U$ をそれぞれ $s$ と $u$ のアフィン近傍で置き換えることにより、
% P05-EN-CH38-0003
$g$ はアフィン間の標準エタール射であると仮定してよい。Morphisms, Lemma
\ref{morphisms-lemma-etale-locally-standard-etale} を参照せよ。
したがって $S = \Spec(A)$ かつ
% P05-EN-CH38-0004
$X = \Spec(A[x, 1/g]/(f))$ と仮定してよい。ここで $f$ はモニックで、
$f'$ は $A[x, 1/g]$ で可逆である。
$A[x, 1/g]/(f) = (A[x]/(f))_g$ は有限自由 $A$-代数 $A[x]/(f)$ の
局所化でもある。したがって、$U$ を $S$ 上有限局所自由なスキーム
$T = \Spec(A[x]/(f))$ の開部分スキームとみなせる。これにより、上の
Lemma \ref{flat-lemma-finite-flat-weak-assassin-up-down} に帰着する。
\end{proof}

\begin{lemma}
\label{flat-lemma-weakly-associated-henselization}
$S$ をスキーム、$s \in S$ を点とする。$\mathcal{O}_{S, s}^h$
（それぞれ $\mathcal{O}_{S, s}^{sh}$）を Hensel 化
（それぞれ狭義 Hensel 化）と記す。Algebra, Definition
\ref{algebra-definition-henselization} を参照せよ。
$\mathcal{F}$ を準連接 $\mathcal{O}_S$-加群とする。次は同値である：
\begin{enumerate}
\item $s$ は $\mathcal{F}$ の弱随伴点である。
\item $\mathfrak m_s$ は $\mathcal{F}_s$ の弱随伴素イデアルである。
\item $\mathfrak m_s^h$ は
$\mathcal{F}_s \otimes_{\mathcal{O}_{S, s}} \mathcal{O}_{S, s}^h$ の
弱随伴素イデアルである。
\item $\mathfrak m_s^{sh}$ は
$\mathcal{F}_s \otimes_{\mathcal{O}_{S, s}} \mathcal{O}_{S, s}^{sh}$ の
弱随伴素イデアルである。
\end{enumerate}
\end{lemma}

\begin{proof}
(1) と (2) の同値は定義である。Divisors, Definition
\ref{divisors-definition-weakly-associated} を参照せよ。
% P05-EN-CH38-0005
含意 (2) $\Rightarrow$ (3) $\Rightarrow$ (4) は、平坦射
（More on Algebra, Lemma
\ref{more-algebra-lemma-dumb-properties-henselization}）である
$$
\Spec(\mathcal{O}_{S, s}^{sh}) \to
\Spec(\mathcal{O}_{S, s}^h) \to
\Spec(\mathcal{O}_{S, s})
$$
と閉点に Divisors, Lemma \ref{divisors-lemma-weakly-ass-pullback} を
適用すれば従う。(4) $\Rightarrow$ (2) を証明するため、$S$ をアフィン近傍で
置き換えてよい。元
$x \in \mathcal{F}_s \otimes_{\mathcal{O}_{S, s}} \mathcal{O}_{S, s}^{sh}$
をとり、その零化イデアルの根基は $\mathfrak m_s^{sh}$ に等しいとする。
Algebra, Lemma \ref{algebra-lemma-weakly-ass-local} を参照せよ。
% P05-EN-CH38-0006
Algebra, Lemma
\ref{algebra-lemma-strict-henselization-different}
により、$\mathcal{O}_{S, s}^{sh}$ はエタール近傍
$f : (U, u) \to (S, s)$ 上の $\mathcal{O}_{U, u}$ の極限に等しいので、
$x$ はある
$x' \in \mathcal{F}_s \otimes_{\mathcal{O}_{S, s}} \mathcal{O}_{U, u}$
の像であると仮定してよい。局所環の射
$\mathcal{O}_{U, u} \to \mathcal{O}_{S, s}^{sh}$ は
（狭義 Hensel 化であるから）忠実平坦であり、したがって普遍単射である
（Algebra, Lemma
\ref{algebra-lemma-faithfully-flat-universally-injective}）。
ゆえに $x'$ の零化イデアルは $x$ の零化イデアルの逆像である。
したがって、この零化イデアルの根基は $\mathfrak m_u$ に等しい。
よって $u$ は $f^*\mathcal{F}$ の弱随伴点である。Lemma
\ref{flat-lemma-etale-weak-assassin-up-down} により、$s$ は
$\mathcal{F}$ の弱随伴点である。
\end{proof}





\section{有限型加群の局所構造}
\label{flat-section-local-structure-module}

\noindent
本章の多くの議論を機能させる鍵となる技術的補題は、幾何学的な
Lemma \ref{flat-lemma-elementary-devissage} である。

\begin{lemma}
\label{flat-lemma-sheaf-lives-on-subscheme}
$f : X \to S$ をアフィンスキームの有限型射とする。
$\mathcal{F}$ を有限型準連接 $\mathcal{O}_X$-加群とする。
$x \in X$ の $S$ における像を $s = f(x)$ とし、
$\mathcal{F}_s = \mathcal{F}|_{X_s}$ とおく。
このとき、有限表示の閉埋め込み $i : Z \to X$ と、有限型準連接
$\mathcal{O}_Z$-加群 $\mathcal{G}$ で、$i_*\mathcal{G} = \mathcal{F}$ および
$Z_s = \text{Supp}(\mathcal{F}_s)$ を満たすものが存在する。
\end{lemma}

\begin{proof}
射 $f : X \to S$ は環準同型 $A \to B$ で与えられ、$\mathcal{F}$ は
$B$-加群 $M$ に付随する準連接層であるとする。Morphisms, Lemma
\ref{morphisms-lemma-locally-finite-type-characterize} により
$A \to B$ は有限型環準同型であり、Properties, Lemma
\ref{properties-lemma-finite-type-module} により $M$ は有限 $B$-加群である。
特に、$\mathcal{F}$ の台は $M$ の零化イデアル
$I = \{x \in B \mid xm = 0\ \forall m \in M\}$ で切り出される
$\Spec(B)$ の閉部分スキームである。Algebra, Lemma
\ref{algebra-lemma-support-closed} を参照せよ。
$x$ に対応する素イデアルを $\mathfrak q \subset B$、$s$ に対応する素イデアルを
$\mathfrak p \subset A$ とする。$X_s = \Spec(B \otimes_A \kappa(\mathfrak p))$
であり、$\mathcal{F}_s$ は $B \otimes_A \kappa(\mathfrak p)$-加群
$M \otimes_A \kappa(\mathfrak p)$ に付随する準連接層である。
Morphisms, Lemma \ref{morphisms-lemma-support-finite-type} により、
$\mathcal{F}_s$ の台は $V(I(B \otimes_A \kappa(\mathfrak p)))$ に等しい。
$B \otimes_A \kappa(\mathfrak p)$ は $\kappa(\mathfrak p)$ 上有限型なので、
次を満たす有限個の元 $f_1, \ldots, f_m \in I$ が存在する：
$$
% P05-EN-CH38-0007
I(B \otimes_A \kappa(\mathfrak p)) =
(f_1, \ldots, f_n)(B \otimes_A \kappa(\mathfrak p)).
$$
$(f_1, \ldots, f_m)$ で切り出される閉部分スキームを $i : Z \to X$ と記す。
式では $Z = \Spec(B/(f_1, \ldots, f_m))$ である。$M$ は $I$ で零化されるため、
$M$ を $B/(f_1, \ldots, f_m)$-加群とみなせる。言い換えれば、
$\mathcal{F}$ は $Z$ 上の有限型加群の押し出しである。構成により
$Z_s = \text{Supp}(\mathcal{F}_s)$ なので、補題が証明された。
\end{proof}

\begin{lemma}
\label{flat-lemma-elementary-devissage}
% P05-EN-CH38-0008
$f : X \to S$ を局所有限型なスキームの射とする。
$\mathcal{F}$ を有限型準連接 $\mathcal{O}_X$-加群とする。
$x \in X$ の $S$ における像を $s = f(x)$ とし、
$\mathcal{F}_s = \mathcal{F}|_{X_s}$ および
$n = \dim_x(\text{Supp}(\mathcal{F}_s))$ とおく。このとき次を構成できる：
\begin{enumerate}
\item 基本エタール近傍 $g : (X', x') \to (X, x)$ および
$e : (S', s') \to (S, s)$。
\item 可換図式
$$
\xymatrix{
X \ar[dd]_f & X' \ar[dd] \ar[l]^g & Z' \ar[l]^i \ar[d]^\pi \\
& & Y' \ar[d]^h \\
S & S' \ar[l]_e & S' \ar@{=}[l]
}
$$
\item $i(z') = x'$、$y' = \pi(z')$、$h(y') = s'$ を満たす点
$z' \in Z'$。
\item 有限型準連接 $\mathcal{O}_{Z'}$-加群 $\mathcal{G}$。
\end{enumerate}
これらは次の性質を満たす：
\begin{enumerate}
\item $X'$、$Z'$、$Y'$、$S'$ はアフィンスキームである。
\item $i$ は有限表示の閉埋め込みである。
\item $i_*(\mathcal{G}) \cong g^*\mathcal{F}$ である。
\item $\pi$ は有限であり、$\pi^{-1}(\{y'\}) = \{z'\}$ である。
\item 拡大 $\kappa(y')/\kappa(s')$ は純超越的である。
\item $h$ は相対次元 $n$ の平滑射であり、幾何学的に整なファイバーをもつ。
\end{enumerate}
\end{lemma}

\begin{proof}
$V \subset S$ を $s$ のアフィン近傍、$U \subset f^{-1}(V)$ を $x$ の
アフィン近傍とする。このとき $f|_U : U \to V$ と $\mathcal{F}|_U$ に対して
補題を証明すれば十分である。したがって証明の残りでは $X$ と $S$ は
アフィンであると仮定する。

\medskip\noindent
まず $X_s = \text{Supp}(\mathcal{F}_s)$ と仮定する。特に
$n = \dim_x(X_s)$ である。More on Morphisms, Lemmas
\ref{more-morphisms-lemma-local-local-structure-finite-type} および
\ref{more-morphisms-lemma-local-local-structure-finite-type-affine} を適用する。
これにより可換図式
$$
\xymatrix{
X \ar[dd] & X' \ar[l]^g \ar[d]^\pi \\
& Y' \ar[d]^h  \\
S & S' \ar[l]_e
}
$$
と点 $x' \in X'$ を得る。$Z' = X'$、$i = \text{id}$、および
$\mathcal{G} = g^*\mathcal{F}$ とおけば、この場合の解を得る。

\medskip\noindent
一般の場合には、Lemma \ref{flat-lemma-sheaf-lives-on-subscheme} のような
閉埋め込み $Z \to X$ と $Z$ 上の層 $\mathcal{G}$ を選ぶ。
前段落の結果を $Z \to S$ と $\mathcal{G}$ に適用すると、図式
$$
\xymatrix{
X \ar[dd]_f & Z \ar[l] \ar[dd]_{f|_Z} & Z' \ar[l]^g \ar[d]^\pi \\
& & Y' \ar[d]^h \\
S & S \ar@{=}[l] & S' \ar[l]_e
}
$$
と、要求されるすべての性質を満たす点 $z' \in Z'$ を得る。
Lemma \ref{flat-lemma-lift-etale} を用いて $Z'$ を $X$ 上エタールなスキームへ
埋め込む。$X'$ を $S'$ 上のスキームにしたいので、この補題を直接適用する
ことはできない。代わりに射
$$
\xymatrix{
Z' \ar[r] & Z \times_S S' \ar[r] & X \times_S S'
}
$$
を考える。最初の射は Morphisms, Lemma
\ref{morphisms-lemma-etale-permanence} によりエタールである。
第二の射は閉埋め込みの基底変換なので閉埋め込みである。最後に、
$X$、$S$、$S'$、$Z$、$Z'$ はすべてアフィンなので、Lemma
\ref{flat-lemma-lift-etale} を適用して、次を満たすアフィンスキームのエタール射
$X' \to X \times_S S'$ を得る：
$$
Z' = (Z \times_S S') \times_{(X \times_S S')} X' = Z \times_X X'.
$$
$Z \to X$ は有限表示の閉埋め込みなので、$Z' \to X'$ もそうである。
$x' \in X'$ を $z' \in Z'$ に対応する点とする。このとき完成した図式
$$
\xymatrix{
X \ar[dd] & X' \ar[dd] \ar[l] & Z' \ar[l]^i \ar[d]^\pi \\
& & Y' \ar[d]^h \\
S & S' \ar[l]_e & S' \ar@{=}[l]
}
$$
は元の問題の解である。
\end{proof}

\begin{lemma}
\label{flat-lemma-devissage-finite-presentation}
仮定と記法は Lemma \ref{flat-lemma-elementary-devissage} と同じとする。
$f$ が局所有限表示ならば $\pi$ は有限表示である。
この場合、次は同値である：
\begin{enumerate}
\item $\mathcal{F}$ は $x$ の近傍で有限表示の $\mathcal{O}_X$-加群である。
\item $\mathcal{G}$ は $z'$ の近傍で有限表示の $\mathcal{O}_{Z'}$-加群である。
\item $\pi_*\mathcal{G}$ は $y'$ の近傍で有限表示の
$\mathcal{O}_{Y'}$-加群である。
\end{enumerate}
引き続き $f$ が局所有限表示であると仮定すると、次も互いに同値である：
\begin{enumerate}
\item[(a)] $\mathcal{F}_x$ は有限表示の $\mathcal{O}_{X, x}$-加群である。
\item[(b)] $\mathcal{G}_{z'}$ は有限表示の $\mathcal{O}_{Z', z'}$-加群である。
\item[(c)] $(\pi_*\mathcal{G})_{y'}$ は有限表示の
$\mathcal{O}_{Y', y'}$-加群である。
\end{enumerate}
\end{lemma}

\begin{proof}
$f$ は局所有限表示であると仮定する。このとき $Z' \to S$ はそのような射の
合成なので局所有限表示である。Morphisms, Lemma
\ref{morphisms-lemma-composition-finite-presentation} を参照せよ。
$Y' \to S$ も平滑射とエタール射の合成なので局所有限表示である。したがって
Morphisms, Lemma \ref{morphisms-lemma-finite-presentation-permanence} により
$\pi$ は局所有限表示である。$\pi$ は有限なので分離的かつ準コンパクトでもあり、
ゆえに $\pi$ は実際に有限表示である。

\medskip\noindent
(1)、(2)、(3) の同値を証明するため、次も考える：
(4) $g^*\mathcal{F}$ は $x'$ の近傍で有限表示の
$\mathcal{O}_{X'}$-加群である。有限表示加群の引き戻しは有限表示である。
Modules, Lemma \ref{modules-lemma-pullback-finite-presentation} を参照せよ。
したがって (1) $\Rightarrow$ (4) である。エタール射 $g$ は開である。
Morphisms, Lemma \ref{morphisms-lemma-etale-open} を参照せよ。
したがって、$x'$ の任意の開近傍 $U' \subset X'$ に対し、像 $g(U')$ は
$x$ の開近傍であり、写像 $\{U' \to g(U')\}$ はエタール被覆である。
ゆえに Descent, Lemma \ref{descent-lemma-finite-presentation-descends} により
(4) $\Rightarrow$ (1) である。Descent, Lemma
\ref{descent-lemma-finite-finitely-presented-module} と、いくつかの容易な
位相的議論（More on Morphisms, Lemma
\ref{more-morphisms-lemma-finite-morphism-single-point-in-fibre} を参照せよ）を
用いると、(4) $\Leftrightarrow$ (2) $\Leftrightarrow$ (3) が分かる。

\medskip\noindent
(a)、(b)、(c) の同値を証明するため、環準同型
$$
\mathcal{O}_{X, x} \to
\mathcal{O}_{X', x'} \to
\mathcal{O}_{Z', z'} \leftarrow
\mathcal{O}_{Y', y'}
$$
を考える。最初の環準同型は忠実平坦である。したがって
$\mathcal{F}_x$ が $\mathcal{O}_{X, x}$ 上有限表示であることと、
% P05-EN-CH38-0010
$g^*\mathcal{F}_{x'}$ が $\mathcal{O}_{X', x'}$ 上有限表示であることは
同値である。Algebra, Lemma
\ref{algebra-lemma-descend-properties-modules} を参照せよ。
第二の環準同型は全射（したがって有限）であり、仮定により有限表示である。
したがって $g^*\mathcal{F}_{x'}$ が $\mathcal{O}_{X', x'}$ 上有限表示で
あることと、$\mathcal{G}_{z'}$ が $\mathcal{O}_{Z', z'}$ 上有限表示で
あることは同値である。Algebra, Lemma
\ref{algebra-lemma-finite-finitely-presented-extension} を参照せよ。
$\pi$ は有限かつ有限表示であり、
% P05-EN-CH38-0009
$\pi^{-1}(\{y'\}) = \{x'\}$ なので、環準同型
$\mathcal{O}_{Y', y'} \leftarrow \mathcal{O}_{Z', z'}$ は有限かつ有限表示である。
More on Morphisms, Lemma
\ref{more-morphisms-lemma-finite-morphism-single-point-in-fibre} を参照せよ。
したがって $\mathcal{G}_{z'}$ が $\mathcal{O}_{Z', z'}$ 上有限表示で
あることと、
% P05-EN-CH38-0011
$\pi_*\mathcal{G}_{y'}$ が $\mathcal{O}_{Y', y'}$ 上有限表示であることは
同値である。Algebra, Lemma
\ref{algebra-lemma-finite-finitely-presented-extension} を参照せよ。
\end{proof}

\begin{lemma}
\label{flat-lemma-devissage-flat}
仮定と記法は Lemma \ref{flat-lemma-elementary-devissage} と同じとする。
次は同値である：
\begin{enumerate}
\item $\mathcal{F}$ は $x$ の近傍で $S$ 上平坦である。
\item $\mathcal{G}$ は $z'$ の近傍で $S'$ 上平坦である。
\item $\pi_*\mathcal{G}$ は $y'$ の近傍で $S'$ 上平坦である。
\end{enumerate}
次も同値である：
\begin{enumerate}
\item[(a)] $\mathcal{F}_x$ は $\mathcal{O}_{S, s}$ 上平坦である。
\item[(b)] $\mathcal{G}_{z'}$ は $\mathcal{O}_{S', s'}$ 上平坦である。
\item[(c)] $(\pi_*\mathcal{G})_{y'}$ は $\mathcal{O}_{S', s'}$ 上平坦である。
\end{enumerate}
\end{lemma}

\begin{proof}
(1)、(2)、(3) の同値を証明するため、次も考える：
(4) $g^*\mathcal{F}$ は $x'$ の近傍で $S$ 上平坦である。
以下では、Lemma \ref{flat-lemma-etale-flat-up-down} を用いて $S$ 上の平坦性と
$S'$ 上の平坦性を同一視し、その都度断らない。エタール射 $g$ は平坦かつ
開である。Morphisms, Lemma \ref{morphisms-lemma-etale-open} を参照せよ。
したがって、$x'$ の任意の開近傍 $U' \subset X'$ に対して、像 $g(U')$ は
$x$ の開近傍であり、射 $U' \to g(U')$ は全射かつ平坦である。
ゆえに Morphisms, Lemma \ref{morphisms-lemma-flat-permanence} により
(4) $\Leftrightarrow$ (1) である。また
$$
\Gamma(X', g^*\mathcal{F}) =
\Gamma(Z', \mathcal{G}) =
\Gamma(Y', \pi_*\mathcal{G})
$$
である。したがって $g^*\mathcal{F}$、$\mathcal{G}$、$\pi_*\mathcal{G}$ の
$S'$ 上の平坦性はすべて同値である（ここでは $X'$、$Z'$、$Y'$、$S'$ が
すべてアフィンであることを用いる）。アフィン近傍に関する省略した位相的議論
（More on Morphisms, Lemma
\ref{more-morphisms-lemma-finite-morphism-single-point-in-fibre} と比較せよ）から、
(4) $\Leftrightarrow$ (2) $\Leftrightarrow$ (3) が分かる。

\medskip\noindent
(a)、(b)、(c) の同値を証明するため、局所環準同型の可換図式
$$
\xymatrix{
\mathcal{O}_{X', x'} \ar[r]_\iota &
\mathcal{O}_{Z', z'} &
\mathcal{O}_{Y', y'} \ar[l]^\alpha &
\mathcal{O}_{S', s'} \ar[l]^\beta \\
\mathcal{O}_{X, x} \ar[u]^\gamma & & &
\mathcal{O}_{S, s} \ar[lll]_\varphi \ar[u]_\epsilon
}
$$
を考える。以下では、Lemma \ref{flat-lemma-etale-flat-up-down-local-ring} を用いて
$\mathcal{O}_{S, s}$ 上の平坦性と $\mathcal{O}_{S', s'}$ 上の平坦性を
同一視し、その都度断らない。写像 $\gamma$ は忠実平坦である。したがって
$\mathcal{F}_x$ が $\mathcal{O}_{S, s}$ 上平坦であることと、
% P05-EN-CH38-0010
$g^*\mathcal{F}_{x'}$ が $\mathcal{O}_{S', s'}$ 上平坦であることは同値である。
Algebra, Lemma \ref{algebra-lemma-flatness-descends-more-general} を参照せよ。
$\mathcal{O}_{S', s'}$-加群として、加群 $g^*\mathcal{F}_{x'}$、
$\mathcal{G}_{z'}$、および
% P05-EN-CH38-0011
$\pi_*\mathcal{G}_{y'}$ はすべて同型である。More on Morphisms, Lemma
\ref{more-morphisms-lemma-finite-morphism-single-point-in-fibre} を参照せよ。
これで証明が完了する。
\end{proof}









\section{一段階デヴィサージュ}
\label{flat-section-one-step-devissage}

\noindent
本節では、加群の一段階デヴィサージュとは何かを説明する。
% P05-EN-CH38-0012
一段階デヴィサージュは基底と終域上エタール局所的に存在する。
一段階デヴィサージュの基底変換、Zariski 縮小、およびエタール局所化を論じる。

\begin{definition}
\label{flat-definition-one-step-devissage}
$S$ をスキームとする。$X$ は $S$ 上局所有限型であるとする。
$\mathcal{F}$ を有限型準連接 $\mathcal{O}_X$-加群とし、
$s \in S$ を点とする。$s$ 上の {\it $\mathcal{F}/X/S$ の一段階
デヴィサージュ}とは、$S$ 上のスキームの射
$$
\xymatrix{
X & Z \ar[l]_i \ar[r]^\pi & Y
}
$$
と、次を満たす有限型準連接 $\mathcal{O}_Z$-加群 $\mathcal{G}$ によって
与えられる：
\begin{enumerate}
\item $X$、$S$、$Z$、$Y$ はアフィンである。
\item $i$ は有限表示の閉埋め込みである。
\item $\mathcal{F} \cong i_*\mathcal{G}$ である。
\item $\pi$ は有限である。
\item 構造射 $Y \to S$ は平滑であり、次元
$\dim(\text{Supp}(\mathcal{F}_s))$ の幾何学的既約ファイバーをもつ。
\end{enumerate}
この場合、$(Z, Y, i, \pi, \mathcal{G})$ は $s$ 上の
$\mathcal{F}/X/S$ の一段階デヴィサージュであるという。
\end{definition}

\noindent
このような一段階デヴィサージュは $X$ と $S$ がアフィンの場合にしか
存在しないことに注意せよ。上の定義では $X$ が $S$ 上（局所）有限型である
ことだけを要求し、以下でもこの設定で作業する。\cite{GruRay} の著者は、
$X \to S$ が局所有限表示であり、
% P05-EN-CH38-0013
$\mathcal{F}$ が有限型準連接 $\mathcal{O}_X$-加群である設定を一貫して用いる。
この選択の利点は、$\mathcal{F}$ が $\mathcal{O}_X$-加群として有限表示であるかを
問うことが「意味をもつ」一方、我々の設定では「意味をもたない」ことである。
議論については More on Morphisms, Section
\ref{more-morphisms-section-finite-type-finite-presentation} を参照せよ。
そこでの観察により、我々の設定では $\mathcal{F}$ が「$S$ に相対的に
局所有限表示」であるという条件を考え、この概念を一貫して用いることもできる。
しかし代わりに、Lemma \ref{flat-lemma-devissage-finite-presentation} の結果と
Remark \ref{flat-remark-finite-presentation} の観察に頼り、この問題が現れるたびに
個別的に処理する。

\begin{definition}
\label{flat-definition-one-step-devissage-at-x}
$S$ をスキームとする。$X$ は $S$ 上局所有限型であるとする。
$\mathcal{F}$ を有限型準連接 $\mathcal{O}_X$-加群とし、
$x \in X$ の $S$ における像を $s$ とする。
$x$ における {\it $\mathcal{F}/X/S$ の一段階デヴィサージュ}とは、系
$(Z, Y, i, \pi, \mathcal{G}, z, y)$ であって、
$(Z, Y, i, \pi, \mathcal{G})$ が $s$ 上の $\mathcal{F}/X/S$ の
一段階デヴィサージュであり、かつ次を満たすものをいう：
\begin{enumerate}
\item $\dim_x(\text{Supp}(\mathcal{F}_s)) = \dim(\text{Supp}(\mathcal{F}_s))$。
\item $z \in Z$ は $i(z) = x$ および $\pi(z) = y$ を満たす点である。
\item $\pi^{-1}(\{y\}) = \{z\}$ である。
\item 拡大 $\kappa(y)/\kappa(s)$ は純超越的である。
\end{enumerate}
\end{definition}

\noindent
$x$ における $\mathcal{F}/X/S$ の一段階デヴィサージュは、$X$ と $S$ が
アフィンの場合にしか存在しない。条件 (1) と Definition
\ref{flat-definition-one-step-devissage} の条件 (5) により、$Y \to S$ の相対次元は
$\dim_x(\text{Supp}(\mathcal{F}_s))$ に等しい。

\begin{lemma}
\label{flat-lemma-elementary-devissage-variant}
% P05-EN-CH38-0014
$f : X \to S$ を局所有限型なスキームの射とする。
$\mathcal{F}$ を有限型準連接 $\mathcal{O}_X$-加群とし、$x \in X$ の
$S$ における像を $s = f(x)$ とする。このとき、点付きスキームの可換図式
$$
\xymatrix{
(X, x) \ar[d]_f & (X', x') \ar[l]^g \ar[d] \\
(S, s) & (S', s') \ar[l] \\
}
$$
が存在し、$(S', s') \to (S, s)$ および $(X', x') \to (X, x)$ は
基本エタール近傍であり、$g^*\mathcal{F}/X'/S'$ は $x'$ における
一段階デヴィサージュをもつ。
\end{lemma}

\begin{proof}
Definition \ref{flat-definition-one-step-devissage-at-x} と
Lemma \ref{flat-lemma-elementary-devissage} から直ちに従う。
\end{proof}

\begin{lemma}
\label{flat-lemma-base-change-one-step}
$S$、$X$、$\mathcal{F}$、$s$ は Definition
\ref{flat-definition-one-step-devissage} と同じとする。
$(Z, Y, i, \pi, \mathcal{G})$ を $s$ 上の $\mathcal{F}/X/S$ の
一段階デヴィサージュとする。$(S', s') \to (S, s)$ を点付きスキームの
任意の射とする。このデータに対し、$X', Z', Y', i', \pi'$ を
$S' \to S$ による $X, Z, Y, i, \pi$ の基底変換とする。
$\mathcal{F}'$ を $\mathcal{F}$ の $X'$ への引き戻し、$\mathcal{G}'$ を
$\mathcal{G}$ の $Z'$ への引き戻しとする。$S'$ がアフィンならば、
$(Z', Y', i', \pi', \mathcal{G}')$ は $s'$ 上の
$\mathcal{F}'/X'/S'$ の一段階デヴィサージュである。
\end{lemma}

\begin{proof}
アフィンのファイバー積はアフィンである。Schemes, Lemma
\ref{schemes-lemma-fibre-product-affines} を参照せよ。基底変換は、閉埋め込み、
有限表示の射、有限射、平滑射、幾何学的既約ファイバーをもつ射、および
相対次元 $n$ の射を保つ。次を参照せよ：
Morphisms, Lemmas \ref{morphisms-lemma-base-change-closed-immersion},
\ref{morphisms-lemma-base-change-finite-presentation},
\ref{morphisms-lemma-base-change-finite},
\ref{morphisms-lemma-base-change-smooth},
\ref{morphisms-lemma-base-change-relative-dimension-d}, and
More on Morphisms, Lemma
\ref{more-morphisms-lemma-base-change-fibres-geometrically-irreducible}.
有限射 $i$ に沿う押し出しは基底変換と可換なので
$i'_*\mathcal{G}' \cong \mathcal{F}'$ である。Cohomology of Schemes,
Lemma \ref{coherent-lemma-affine-base-change} を参照せよ。また
$\dim(\text{Supp}(\mathcal{F}_s)) = \dim(\text{Supp}(\mathcal{F}'_{s'}))$
である。実際、
$$
\text{Supp}(\mathcal{F}_s) \times_s s' = \text{Supp}(\mathcal{F}'_{s'}).
$$
であり、Morphisms, Lemma
\ref{morphisms-lemma-dimension-fibre-after-base-change} を適用できる。
これで補題が証明された。
\end{proof}

\begin{lemma}
\label{flat-lemma-base-change-one-step-at-x}
$S$、$X$、$\mathcal{F}$、$x$、$s$ は Definition
\ref{flat-definition-one-step-devissage-at-x} と同じとする。
$(Z, Y, i, \pi, \mathcal{G}, z, y)$ を $x$ における
$\mathcal{F}/X/S$ の一段階デヴィサージュとする。
$(S', s') \to (S, s)$ を、同型 $\kappa(s) = \kappa(s')$ を誘導する
点付きスキームの射とする。$(Z', Y', i', \pi', \mathcal{G}')$ は
Lemma \ref{flat-lemma-base-change-one-step} で構成したものとし、
$x' \in X'$（それぞれ $z' \in Z'$、$y' \in Y'$）を、
$x \in X$（それぞれ $z \in Z$、$y \in Y$）と $s' \in S'$ の双方へ写る
一意な点とする。$S'$ がアフィンならば、
$(Z', Y', i', \pi', \mathcal{G}', z', y')$ は $x'$ における
$\mathcal{F}'/X'/S'$ の一段階デヴィサージュである。
\end{lemma}

\begin{proof}
Lemma \ref{flat-lemma-base-change-one-step} により、
$(Z', Y', i', \pi', \mathcal{G}')$ は $s'$ 上の
$\mathcal{F}'/X'/S'$ の一段階デヴィサージュである。仮定
$\kappa(s) = \kappa(s')$ により、ファイバー $X'_{s'}$、$Z'_{s'}$、
$Y'_{s'}$ はそれぞれ $X_s$、$Z_s$、$Y_s$ と同型である。したがって
Definition \ref{flat-definition-one-step-devissage-at-x} の性質 (1) -- (4) は
$(Z', Y', i', \pi', \mathcal{G}', z', y')$ に対して成り立つ。
\end{proof}

\begin{definition}
\label{flat-definition-shrink}
$S$、$X$、$\mathcal{F}$、$x$、$s$ は Definition
\ref{flat-definition-one-step-devissage-at-x} と同じとする。
$(Z, Y, i, \pi, \mathcal{G}, z, y)$ を $x$ における
$\mathcal{F}/X/S$ の一段階デヴィサージュとする。この状況の
{\it 標準縮小}を、$s \in S'$、$x \in X'$、$z \in Z'$、$y \in Y'$ を
満たす標準開 $S' \subset S$、$X' \subset X$、$Z' \subset Z$、
$Y' \subset Y$ であって、
$$
(Z', Y', i|_{Z'}, \pi|_{Z'}, \mathcal{G}|_{Z'}, z, y)
$$
が $x$ における $\mathcal{F}|_{X'}/X'/S'$ の一段階デヴィサージュと
なるものによって定める。
\end{definition}

\begin{lemma}
\label{flat-lemma-shrink}
仮定と記法を Definition \ref{flat-definition-shrink} と同じとする。このとき：
\begin{enumerate}
\item
\label{flat-item-shrink-base}
$S' \subset S$ が $s$ の標準開近傍ならば、$X' = X_{S'}$、
$Z' = Z_{S'}$、$Y' = Y_{S'}$ とおくことで標準縮小を得る。
\item
\label{flat-item-shrink-on-Y}
$W \subset Y$ を $y$ の標準開近傍とする。このとき
$Y' = W \times_S S'$ となる標準縮小が存在する。
\item
\label{flat-item-shrink-on-X}
$U \subset X$ を $x$ の開近傍とする。このとき $X' \subset U$ となる
標準縮小が存在する。
\end{enumerate}
\end{lemma}

\begin{proof}
(1) は Lemma \ref{flat-lemma-base-change-one-step-at-x} と、アフィンスキームの射に
よる標準開の逆像が標準開であるという事実から直ちに従う。Algebra, Lemma
\ref{algebra-lemma-spec-functorial} を参照せよ。

\medskip\noindent
$W \subset Y$ を (2) のようにとる。$Y \to S$ は平滑なので開である。
Morphisms, Lemma \ref{morphisms-lemma-smooth-open} を参照せよ。
したがって、$W$ の像に含まれる $s$ の標準開近傍 $S'$ を見つけられる。
すると $W_{S'} \to S'$ のファイバーは、$S'$ 上の $Y \to S$ のファイバーの
空でない開部分スキームであり、したがって幾何学的既約でもある。
$Y' = W_{S'}$ および $Z' = \pi^{-1}(Y')$ とおくと、$Z' \subset Z$ は
$z$ の標準開近傍である。$Z' = D(\overline{h})$ を満たす関数
$\overline{h} \in \Gamma(Z, \mathcal{O}_Z)$ をとる。$i : Z \to X$ は
閉埋め込みなので、$i^\sharp(h) = \overline{h}$ を満たす関数
$h \in \Gamma(X, \mathcal{O}_X)$ が存在する。$X' = D(h) \subset X$ と
とる。このようにして (2) の標準縮小を得る。

\medskip\noindent
$U \subset X$ を (3) のようにとる。$U$ を縮小して、$U$ は標準開であると
仮定してよい。More on Morphisms, Lemma
\ref{more-morphisms-lemma-finite-morphism-single-point-in-fibre} により、
$\pi^{-1}(W) \subset i^{-1}(U)$ を満たす $y$ の標準開近傍
$W \subset Y$ が存在する。(2) を適用し、$Y' = W_{S'}$ を満たす標準縮小
$X', S', Z', Y'$ を得る。$Z' \subset \pi^{-1}(W) \subset i^{-1}(U)$ なので、
標準縮小を定めるどの条件も損なわずに $X'$ を $X' \cap U$ で置き換えられる
（$U$ も標準開なので、これは依然として標準開である）。これでよい。
\end{proof}

\begin{lemma}
\label{flat-lemma-elementary-etale-neighbourhood}
$S$、$X$、$\mathcal{F}$、$x$、$s$ は Definition
\ref{flat-definition-one-step-devissage-at-x} と同じとする。
$(Z, Y, i, \pi, \mathcal{G}, z, y)$ を $x$ における
$\mathcal{F}/X/S$ の一段階デヴィサージュとする。次の図式を考える：
$$
\xymatrix{
(Y, y) \ar[d] & (Y', y') \ar[l] \ar[d] \\
(S, s) & (S', s') \ar[l]
}
$$
これは点付きスキームの可換図式で、水平矢印は基本エタール近傍であるとする。
このとき可換図式
$$
\xymatrix{
& & (X'', x'') \ar[lld] \ar[d] & (Z'', z'') \ar[l] \ar[lld] \ar[d] \\
(X, x) \ar[d] & (Z, z) \ar[l] \ar[d] &
(S'', s'') \ar[lld] & (Y'', y'') \ar[lld] \ar[l] \\
(S, s) & (Y, y) \ar[l]
}
$$
が存在し、次の性質をもつ：
\begin{enumerate}
\item $(S'', s'') \to (S', s')$ は基本エタール近傍であり、射
$S'' \to S$ は合成 $S'' \to S' \to S$ である。
\item $Y''$ は $Y' \times_{S'} S''$ の開部分スキームである。
\item $Z'' = Z \times_Y Y''$ である。
\item $(X'', x'') \to (X, x)$ は基本エタール近傍である。
\item $(Z'', Y'', i'', \pi'', \mathcal{G}'', z'', y'')$ は層
$\mathcal{F}''$ の $x''$ における一段階デヴィサージュである。
\end{enumerate}
ここで $\mathcal{F}''$（それぞれ $\mathcal{G}''$）は射 $X'' \to X$
（それぞれ $Z'' \to Z$）による $\mathcal{F}$（それぞれ $\mathcal{G}$）の
引き戻しであり、$i'' : Z'' \to X''$ と $\pi'' : Z'' \to Y''$ は図式の
とおりである。
\end{lemma}

\begin{proof}
$(S'', s'') \to (S', s')$ を、$S''$ がアフィンである任意の基本エタール近傍と
する。$Y'' \subset Y' \times_{S'} S''$ を、点 $y'' = (y', s'')$ を含む任意の
アフィン開近傍とする。このとき (3) によりアフィンな $(Z'', z'')$ を得る。
さらに $Z_{S''} \to X_{S''}$ は閉埋め込みであり、$Z'' \to Z_{S''}$ は
エタール射である。したがって Lemma \ref{flat-lemma-lift-etale} を適用でき、
% P05-EN-CH38-0015
$Z'' \cong X'' \times_{X_{S'}} Z_{S'}$ を満たすアフィンのエタール射
$X'' \to X_{S'}$ を見つけられる。対応する閉埋め込みを
$i'' : Z'' \to X''$ と記す。$x'' = i''(z'')$ とおけば、補題のような可換図式を
得る。性質 (1)、(2)、(3)、(4) は構成により成り立つ。したがって、
$(S'', s'') \to (S', s')$ と $Y''$ を適切に選んだとき (5) が成り立つことを
示せば十分である。

\medskip\noindent
まず、第一段落のような $(S'', s'') \to (S', s')$ と $Y''$ の任意の選択に
対して成り立つ性質を列挙する。構成により $Z'' = X'' \times_X Z$ なので、
$i''_*\mathcal{G}'' = \mathcal{F}''$ である（記法は補題の主張のとおり）。
Cohomology of Schemes, Lemma \ref{coherent-lemma-affine-base-change} を参照せよ。
$n = \dim(\text{Supp}(\mathcal{F}_s)) = \dim_x(\text{Supp}(\mathcal{F}_s))$
とおく。射 $Y'' \to S''$ は相対次元 $n$ の平滑射である。実際、
$Y' \to S'$ は、エタール射と相対次元 $n$ の平滑射の合成
$Y' \to Y_{S'} \to S'$ として相対次元 $n$ の平滑射であり、基底変換は
相対次元 $n$ の平滑射を保つ。また $\kappa(y'') = \kappa(y)$ および
$\kappa(s) = \kappa(s'')$ なので、$\kappa(y'')$ は $\kappa(s'')$ の
純超越拡大である。ファイバーの射 $X''_{s''} \to X_s$ は
$\kappa(s) = \kappa(s'')$ 上のアフィンスキームのエタール射であり、点
$x''$ を点 $x$ へ写し、$\mathcal{F}_s$ を $\mathcal{F}''_{s''}$ へ引き戻す。
したがって
$$
\dim(\text{Supp}(\mathcal{F}''_{s''})) =
\dim(\text{Supp}(\mathcal{F}_s)) = n =
\dim_x(\text{Supp}(\mathcal{F}_s)) =
\dim_{x''}(\text{Supp}(\mathcal{F}''_{s''}))
$$
である。次元がエタール局所化で不変であることを用いた。Descent, Lemma
\ref{descent-lemma-dimension-at-point-local} を参照せよ。
$\pi'' : Z'' \to Y''$ は $\pi$ の基底変換なので $\pi''$ は有限であり、
$\kappa(y) = \kappa(y'')$ なので
% P05-EN-CH38-0016
$\pi^{-1}(\{y''\}) = \{z''\}$ である。

\medskip\noindent
この時点で、$Y'' \to S''$ が幾何学的既約ファイバーをもつこと以外は
Definition \ref{flat-definition-one-step-devissage} のすべての条件を確認した。
もちろん一般にはこれは成り立たず、ここで $\kappa(s) \subset \kappa(y)$ が
純超越的であることを用いる。すなわち、$T \subset Y'_{s'}$ を
$y' = (y, s')$ を含む $Y'_{s'}$ の既約成分とする。このスキームは
$\kappa(s')$ 上平滑なので、$T$ は $Y'_{s'}$ の開部分スキームである。
$\kappa(s') = \kappa(s)$ は $\kappa(y') = \kappa(y)$ の中で代数的に閉じて
いるので、Varieties, Lemma
\ref{varieties-lemma-geometrically-connected-if-connected-and-point} により
$T$ は幾何学的連結である。$T$ は平滑なので、$T$ は幾何学的既約である。
したがって More on Morphisms, Lemma
\ref{more-morphisms-lemma-normal-morphism-irreducible} を適用でき、
$(S'', s'') \to (S', s')$ という基本エタール射とアフィン開
$Y'' \subset Y'_{S''}$ で、$Y'' \to S''$ のすべてのファイバーが
幾何学的既約であり、$T = Y''_{s''}$ となるものを見つけられる。
縮小することにより（まず $Y''$、次に $S''$）、$Y''$ と $S''$ はともに
アフィンであると仮定してよい。これで補題の証明が完了する。
\end{proof}

\begin{lemma}
\label{flat-lemma-existence-alpha}
$S$、$X$、$\mathcal{F}$、$s$ は Definition
\ref{flat-definition-one-step-devissage} と同じとする。
$(Z, Y, i, \pi, \mathcal{G})$ を $s$ 上の $\mathcal{F}/X/S$ の
一段階デヴィサージュとし、$\xi \in Y_s$ を（一意な）総称点とする。
このとき、整数 $r > 0$ と $\mathcal{O}_Y$-加群の射
$$
\alpha : \mathcal{O}_Y^{\oplus r} \longrightarrow \pi_*\mathcal{G}
$$
であって、
$$
\alpha :
\kappa(\xi)^{\oplus r}
\longrightarrow
(\pi_*\mathcal{G})_\xi \otimes_{\mathcal{O}_{Y, \xi}} \kappa(\xi)
$$
が同型となるものが存在する。さらに、この場合
$$
\dim(\text{Supp}(\Coker(\alpha)_s)) < \dim(\text{Supp}(\mathcal{F}_s)).
$$
\end{lemma}

\begin{proof}
仮定によりスキーム $S$ と $Y$ はアフィンである。
$S = \Spec(A)$ および $Y = \Spec(B)$ と書く。$\pi$ は有限なので、
$\mathcal{O}_Y$-加群 $\pi_*\mathcal{G}$ は有限型準連接
$\mathcal{O}_Y$-加群である。したがって、ある有限 $B$-加群 $N$ に対して
$\pi_*\mathcal{G} = \widetilde{N}$ である。$\xi$ に対応する素イデアルを
$\mathfrak p \subset B$ とする。$\alpha$ を得るため、
$r = \dim_{\kappa(\mathfrak p)} N \otimes_B \kappa(\mathfrak p)$ とおき、
$N \otimes_B \kappa(\mathfrak p)$ の基底をなす
$x_1, \ldots, x_r \in N$ を選ぶ。$\alpha : B^{\oplus r} \to N$ を、式
$\alpha(b_1, \ldots, b_r) = \sum b_ix_i$ で与えられる写像とする。
望みどおり
$\alpha : \kappa(\mathfrak p)^{\oplus r} \to N \otimes_B \kappa(\mathfrak p)$
は同型である。最後に、$\alpha$ をこの性質をもつ任意の写像とする。
すると $N' = \Coker(\alpha)$ は $N' \otimes \kappa(\mathfrak p) = 0$ を
満たす有限 $B$-加群である。Nakayama の補題
（Algebra, Lemma \ref{algebra-lemma-NAK}）により
$N'_{\mathfrak p} = 0$ である。ファイバー $Y_s$ は総称点 $\xi$ をもつ
% P05-EN-CH38-0017
次元 $n$ の幾何学的既約スキームであり、$\xi$ は $\Coker(\alpha)$ の台に
含まれないことを今示したので、補題の最後の主張が成り立つ。
\end{proof}


\section{完全デヴィサージュ}
\label{flat-section-complete-devissage}

\noindent
% P05-EN-CH38-0018: uniquely determined grammar repair; see sidecar.
本節では、加群の完全デヴィサージュとは何かを説明し、そのような
デヴィサージュが存在することを証明する。本節の内容は主として
データの整理に関するものである。

\begin{definition}
\label{flat-definition-complete-devissage}
$S$ をスキームとする。
$X$ は $S$ 上局所有限型であるとする。
$\mathcal{F}$ を有限型準連接 $\mathcal{O}_X$ 加群とする。
$s \in S$ を点とする。
{\it $s$ 上の $\mathcal{F}/X/S$ の完全デヴィサージュ}とは、次の図式
$$
\xymatrix{
X & Z_1 \ar[l]^{i_1} \ar[d]^{\pi_1} \\
& Y_1 & Z_2 \ar[l]^{i_2} \ar[d]^{\pi_2} \\
& & Y_2 & Z_3 \ar[l] \ar[d] \\
& & & ... & ... \ar[l] \ar[d] \\
& & & & Y_n
}
$$
をなす $S$ 上のスキーム、有限型準連接 $\mathcal{O}_{Z_k}$ 加群
$\mathcal{G}_k$、および $\mathcal{O}_{Y_k}$ 加群の写像
$$
\alpha_k :
\mathcal{O}_{Y_k}^{\oplus r_k}
\longrightarrow
\pi_{k, *}\mathcal{G}_k,
\quad
k = 1, \ldots, n
$$
であって、次の性質を満たすものとして与えられる：
\begin{enumerate}
\item $(Z_1, Y_1, i_1, \pi_1, \mathcal{G}_1)$ は $s$ 上の
$\mathcal{F}/X/S$ の一段階デヴィサージュである、
\item 写像 $\alpha_k$ は同型
$$
\kappa(\xi_k)^{\oplus r_k} \longrightarrow
(\pi_{k, *}\mathcal{G}_k)_{\xi_k}
\otimes_{\mathcal{O}_{Y_k, \xi_k}} \kappa(\xi_k)
$$
を誘導する。ここで $\xi_k \in (Y_k)_s$ は唯一の総称点である、
\item $k = 2, \ldots, n$ に対し、系
$(Z_k, Y_k, i_k, \pi_k, \mathcal{G}_k)$
は $s$ 上の $\Coker(\alpha_{k - 1})/Y_{k - 1}/S$ の
一段階デヴィサージュである、
\item $\Coker(\alpha_n) = 0$.
\end{enumerate}
このとき、
$(Z_k, Y_k, i_k, \pi_k, \mathcal{G}_k, \alpha_k)_{k = 1, \ldots, n}$
を $s$ 上の $\mathcal{F}/X/S$ の完全デヴィサージュという。
\end{definition}

\begin{definition}
\label{flat-definition-complete-devissage-at-x}
$S$ をスキームとする。
$X$ は $S$ 上局所有限型であるとする。
$\mathcal{F}$ を有限型準連接 $\mathcal{O}_X$ 加群とする。
$x \in X$ を像が $s \in S$ である点とする。
{\it $x$ における $\mathcal{F}/X/S$ の完全デヴィサージュ}とは、系
$$
(Z_k, Y_k, i_k, \pi_k, \mathcal{G}_k, \alpha_k, z_k, y_k)_{k = 1, \ldots, n}
$$
であって、$(Z_k, Y_k, i_k, \pi_k, \mathcal{G}_k, \alpha_k)$ が
$s$ 上の $\mathcal{F}/X/S$ の完全デヴィサージュであり、かつ次を満たすものとして与えられる：
\begin{enumerate}
\item $(Z_1, Y_1, i_1, \pi_1, \mathcal{G}_1, z_1, y_1)$ は
$x$ における $\mathcal{F}/X/S$ の一段階デヴィサージュである、
\item $k = 2, \ldots, n$ に対し、系
$(Z_k, Y_k, i_k, \pi_k, \mathcal{G}_k, z_k, y_k)$
は $y_{k - 1}$ における $\Coker(\alpha_{k - 1})/Y_{k - 1}/S$ の
一段階デヴィサージュである。
\end{enumerate}
\end{definition}

\noindent
ここでも、完全デヴィサージュが存在し得るのは $X$ と $S$ が
アフィンである場合に限ることを注意しておく。

\begin{lemma}
\label{flat-lemma-base-change-complete}
$S$, $X$, $\mathcal{F}$, $s$ は
Definition \ref{flat-definition-complete-devissage} のとおりとする。
$(S', s') \to (S, s)$ を点付きスキームの任意の射とする。
$(Z_k, Y_k, i_k, \pi_k, \mathcal{G}_k, \alpha_k)_{k = 1, \ldots, n}$
を $s$ 上の $\mathcal{F}/X/S$ の完全デヴィサージュとする。
このデータに対し、$X', Z'_k, Y'_k, i'_k, \pi'_k$ をそれぞれ
$S' \to S$ による $X, Z_k, Y_k, i_k, \pi_k$ の基底変換とする。
$\mathcal{F}'$ を $\mathcal{F}$ の $X'$ への引き戻しとし、
$\mathcal{G}'_k$ を $\mathcal{G}_k$ の $Z'_k$ への引き戻しとする。
$\alpha'_k$ を $\alpha_k$ の $Y'_k$ への引き戻しとする。
$S'$ がアフィンならば、
$(Z'_k, Y'_k, i'_k, \pi'_k, \mathcal{G}'_k, \alpha'_k)_{k = 1, \ldots, n}$
は $s'$ 上の $\mathcal{F}'/X'/S'$ の完全デヴィサージュである。
\end{lemma}

\begin{proof}
Lemma \ref{flat-lemma-base-change-one-step} により、
一段階デヴィサージュの基底変換は一段階デヴィサージュである。
したがって、$\Coker(\alpha_k)$ の形成が基底変換と可換であること、
および Definition \ref{flat-definition-complete-devissage} の条件 (2) が
基底変換によって保たれることを示せば十分である。前者は、
$\pi'_{k, *}\mathcal{G}'_k$ が $\pi_{k, *}\mathcal{G}_k$ の引き戻しであり
（Cohomology of Schemes, Lemma \ref{coherent-lemma-affine-base-change}）、
かつ $\otimes$ が右完全であることから従う。後者についても、
同じ理由により
$$
(\pi_{k, *}\mathcal{G}_k)_{\xi_k}
\otimes_{\mathcal{O}_{Y_k, \xi_k}} \kappa(\xi_k)
\otimes_{\kappa(\xi_k)} \kappa(\xi'_k)
\cong
(\pi'_{k, *}\mathcal{G}'_k)_{\xi'_k}
\otimes_{\mathcal{O}_{Y'_k, \xi'_k}} \kappa(\xi'_k)
$$
が成り立つ。記号の意味は明らかである。
\end{proof}

\begin{lemma}
\label{flat-lemma-base-change-complete-at-x}
$S$, $X$, $\mathcal{F}$, $x$, $s$ は
Definition \ref{flat-definition-complete-devissage-at-x} のとおりとする。
$(S', s') \to (S, s)$ を、同型 $\kappa(s) = \kappa(s')$ を誘導する
点付きスキームの射とする。また、
$(Z_k, Y_k, i_k, \pi_k, \mathcal{G}_k, \alpha_k, z_k, y_k)_{k = 1, \ldots, n}$
を $x$ における $\mathcal{F}/X/S$ の完全デヴィサージュとする。
さらに、
$(Z'_k, Y'_k, i'_k, \pi'_k, \mathcal{G}'_k, \alpha'_k)_{k = 1, \ldots, n}$
を Lemma \ref{flat-lemma-base-change-complete} で構成したものとする。
% P05-EN-CH38-0019: source-faithful unindexed target schemes; see sidecar.
$x' \in X'$（それぞれ $z'_k \in Z'$, $y'_k \in Y'$）を、
$x \in X$（それぞれ $z_k \in Z_k$, $y_k \in Y_k$）と
$s' \in S'$ の双方へ写る唯一の点とする。
$S'$ がアフィンならば、
$(Z'_k, Y'_k, i'_k, \pi'_k, \mathcal{G}'_k, \alpha'_k,
z'_k, y'_k)_{k = 1, \ldots, n}$
は $x'$ における $\mathcal{F}'/X'/S'$ の完全デヴィサージュである。
\end{lemma}

\begin{proof}
Lemma \ref{flat-lemma-base-change-complete} と
Lemma \ref{flat-lemma-base-change-one-step-at-x} を組み合わせればよい。
\end{proof}

\begin{definition}
\label{flat-definition-shrink-complete}
$S$, $X$, $\mathcal{F}$, $x$, $s$ は
Definition \ref{flat-definition-complete-devissage-at-x} のとおりとする。
$x$ における $\mathcal{F}/X/S$ の完全デヴィサージュ
$(Z_k, Y_k, i_k, \pi_k, \mathcal{G}_k, \alpha_k, z_k, y_k)_{k = 1, \ldots, n}$
を考える。この状況の{\it 標準縮小}を次のように定義する。
% P05-EN-CH38-0020: source-faithful undefined point indices and unindexed opens; see sidecar.
標準開集合 $S' \subset S$, $X' \subset X$,
$Z'_k \subset Z_k$, $Y'_k \subset Y_k$ であって、$s_k \in S'$,
$x_k \in X'$, $z_k \in Z'$, $y_k \in Y'$ を満たし、かつ
$$
(Z'_k, Y'_k, i'_k, \pi'_k,
\mathcal{G}'_k, \alpha'_k, z_k, y_k)_{k = 1, \ldots, n}
$$
% P05-EN-CH38-0021: source-faithful one-step/complete mismatch; see sidecar.
が $x$ における $\mathcal{F}'/X'/S'$ の一段階デヴィサージュとなるものとする。ここで
$\mathcal{G}'_k = \mathcal{G}_k|_{Z'_k}$ かつ
$\mathcal{F}' = \mathcal{F}|_{X'}$ である。
\end{definition}

\begin{lemma}
\label{flat-lemma-shrink-complete}
Definition \ref{flat-definition-shrink-complete} の仮定と記号のもとで、
次が成り立つ：
\begin{enumerate}
\item
\label{flat-item-shrink-base-complete}
$S' \subset S$ が $s$ の標準開近傍ならば、
$X' = X_{S'}$, $Z'_k = Z_{S'}$, $Y'_k = Y_{S'}$ とおくことで
標準縮小を得る。
\item
\label{flat-item-shrink-on-Y-complete}
% P05-EN-CH38-0022: source-faithful undefined point y; see sidecar.
$W \subset Y_n$ を $y$ の標準開近傍とする。
このとき、$Y'_n = W \times_S S'$ となる標準縮小が存在する。
\item
\label{flat-item-shrink-on-X-complete}
$U \subset X$ を $x$ の開近傍とする。
このとき、$X' \subset U$ となる標準縮小が存在する。
\end{enumerate}
\end{lemma}

\begin{proof}
(1) は Lemmas \ref{flat-lemma-base-change-complete-at-x} および
\ref{flat-lemma-shrink} から直ちに従う。

\medskip\noindent
(2) の証明。便宜上 $X = Y_0$ と書く。
Lemma \ref{flat-lemma-shrink} (\ref{flat-item-shrink-on-Y}) を適用して、
$y_{n - 1}$ における $\Coker(\alpha_{n - 1})/Y_{n - 1}/S$ の
一段階デヴィサージュの標準縮小
$S', Y'_{n - 1}, Z'_n, Y'_n$
で $Y'_n = W \times_S S'$ を満たすものを得る。この手続きを繰り返し、
$y_{n - 2}$ における $\Coker(\alpha_{n - 2})/Y_{n - 2}/S$ の
一段階デヴィサージュの標準縮小
$S'', Y''_{n - 2}, Z''_{n - 1}, Y''_{n - 1}$
で $Y''_{n - 1} = Y'_{n - 1} \times_S S''$ を満たすものを得られる。
このように続けて、最終的に
$S^{(n)}, Y^{(n)}_0, Z^{(n)}_1, Y^{(n)}_1$
を得る。この時点で、$S^{(n)}$, $X^{(n)}$,
$Z_k^{(n - k)} \times_S S^{(n)}$, $Y_k^{(n - k)} \times_S S^{(n)}$
を取れば、所望の性質をもつ標準縮小が得られることは明らかである。

\medskip\noindent
(3) の証明。完全デヴィサージュの長さに関する帰納法を用いる。
まず Lemma \ref{flat-lemma-shrink} (\ref{flat-item-shrink-on-X}) を適用し、
$x$ における $\mathcal{F}/X/S$ の一段階デヴィサージュの標準縮小
$S', X', Z'_1, Y'_1$
で $X' \subset U$ を満たすものを得る。$n = 1$ ならば証明は終わる。
$n > 1$ ならば、帰納法により、
% P05-EN-CH38-0023: source-faithful wrong distinguished point; see sidecar.
$x$ における $\Coker(\alpha_1)/Y_1/S$ の完全デヴィサージュ
$(Z_k, Y_k, i_k, \pi_k, \mathcal{G}_k, \alpha_k, z_k, y_k)_{k = 2, \ldots, n}$
の標準縮小
$S''$, $Y''_1$, $Z''_k$, $Y''_k$
で $Y''_1 \subset Y'_1$ を満たすものを得られる。さらに
Lemma \ref{flat-lemma-shrink} (\ref{flat-item-shrink-on-Y}) を用いると、
$S''' \subset S'$, $X''' \subset X'$, $Z'''_1$ および
$Y'''_1 = Y''_1 \times_S S'''$ からなる標準縮小を得られる。
したがって、求めるものは
$$
S''', X''', Z'''_1, Y'''_1, Z''_2 \times_S S''',
Y''_2 \times_S S''', \ldots, Z''_n \times_S S''', Y''_n \times_S S'''
$$
を取ることで得られる。これで補題の証明が終わる。
\end{proof}

\begin{proposition}
\label{flat-proposition-existence-complete-at-x}
$S$ をスキームとする。
$X$ は $S$ 上局所有限型であるとする。
$x \in X$ を像が $s \in S$ である点とする。
% P05-EN-CH38-0024: source-faithful missing declaration of F; see sidecar.
可換図式
$$
\xymatrix{
(X, x) \ar[d] & (X', x') \ar[l]^g \ar[d] \\
(S, s) & (S', s') \ar[l]
}
$$
であって、横向きの射が基本エタール近傍であり、かつ
% P05-EN-CH38-0025: source-faithful x/x' mismatch; see sidecar.
$g^*\mathcal{F}/X'/S'$ が $x$ において完全デヴィサージュをもつものが存在する。
\end{proposition}

\begin{proof}
整数 $d = \dim_x(\text{Supp}(\mathcal{F}_s))$ に関する帰納法で証明する。
Lemma \ref{flat-lemma-elementary-devissage-variant} により、図式
$$
\xymatrix{
(X, x) \ar[d] & (X', x') \ar[l]^g \ar[d] \\
(S, s) & (S', s') \ar[l]
}
$$
であって、横向きの射が基本エタール近傍であり、かつ
$g^*\mathcal{F}/X'/S'$ が $x'$ において一段階デヴィサージュをもつものが存在する。
問題は局所的なので、$(X, x) \to (S, s)$ を
$(X', x') \to (S', s')$ で置き換えてよい。したがって、置き換えた後には
$x$ における $\mathcal{F}/X/S$ の一段階デヴィサージュ
$(Z_1, Y_1, i_1, \pi_1, \mathcal{G}_1)$ が存在すると仮定してよい。

\medskip\noindent
Lemma \ref{flat-lemma-existence-alpha} を適用し、写像
$$
\alpha_1 :
\mathcal{O}_{Y_1}^{\oplus r_1}
\longrightarrow
\pi_{1, *}\mathcal{G}_1
$$
を得る。これは $\kappa(\xi_1)$ 上のベクトル空間の同型を誘導する。
ここで $\xi_1 \in Y_1$ は $s$ 上の $Y_1$ のファイバーの唯一の総称点である。さらに
$\dim_{y_1}(\text{Supp}(\Coker(\alpha_1)_s)) < d$.
$y_1$ における $\Coker(\alpha_1)_s$ の茎が零となる場合がある。
この場合、Lemma \ref{flat-lemma-shrink} (\ref{flat-item-shrink-on-Y}) により
$Y_1$ を縮小して $\Coker(\alpha_1) = 0$ と仮定できるので、
% P05-EN-CH38-0026: source-faithful zero/one length mismatch; see sidecar.
長さ零の完全デヴィサージュを得る。

\medskip\noindent
次に、$y_1$ における $\Coker(\alpha_1)_s$ の茎が零でないと仮定する。
この場合、帰納法により可換図式
\begin{equation}
\label{flat-equation-overcome-this}
\vcenter{
\xymatrix{
(Y_1, y_1) \ar[d] & (Y'_1, y'_1) \ar[l]^h \ar[d] \\
(S, s) & (S', s') \ar[l]
}
}
\end{equation}
であって、横向きの射が基本エタール近傍であり、かつ
$h^*\Coker(\alpha_1)/Y'_1/S'$ が $y'_1$ において完全デヴィサージュ
$$
(Z_k, Y_k, i_k, \pi_k, \mathcal{G}_k, \alpha_k, z_k, y_k)_{k = 2, \ldots, n}
$$
をもつものが存在する。（特に、
% P05-EN-CH38-0027: source-faithful wrong closed-immersion target; see sidecar.
$i_2 : Z_2 \to Y'_1$ は $Y'_2$ への閉埋め込みである。）ここで
Lemma \ref{flat-lemma-elementary-etale-neighbourhood} を
$S, X, \mathcal{F}, x, s$、系
$(Z_1, Y_1, i_1, \pi_1, \mathcal{G}_1)$
および図式 (\ref{flat-equation-overcome-this}) に適用する。すると図式
$$
\xymatrix{
& & (X'', x'') \ar[lld] \ar[d] & (Z''_1, z''_1) \ar[l] \ar[lld] \ar[d] \\
(X, x) \ar[d] & (Z_1, z_1) \ar[l] \ar[d] &
(S'', s'') \ar[lld] & (Y''_1, y''_1) \ar[lld] \ar[l] \\
(S, s) & (Y_1, y_1) \ar[l]
}
$$
で、引用した補題に列挙されたすべての性質をもつものを得る。
特に $Y''_1 \subset Y'_1 \times_{S'} S''$ である。
$X_1 = Y'_1 \times_{S'} S''$ とおき、$\mathcal{F}_1$ を
$\Coker(\alpha_1)$ の引き戻しとする。
Lemma \ref{flat-lemma-base-change-complete-at-x} により、系
\begin{equation}
\label{flat-equation-shrink-this}
(Z_k \times_{S'} S'',
Y_k \times_{S'} S'', i''_k, \pi''_k, \mathcal{G}''_k,
\alpha''_k, z''_k, y''_k)_{k = 2, \ldots, n}
\end{equation}
% P05-EN-CH38-0028: source-faithful malformed/incomplete relative datum; see sidecar.
は $\mathcal{F}_1$ から $X_1$ への完全デヴィサージュである。
再び、問題の性質から $(X, x) \to (S, s)$ を
$(X'', x'') \to (S'', s'')$ で置き換えてよい。
以上により、次を仮定してよいことが分かる：
\begin{enumerate}
\item[(a)] $x$ における $\mathcal{F}/X/S$ の一段階デヴィサージュ
$(Z_1, Y_1, i_1, \pi_1, \mathcal{G}_1)$ が存在する、
\item[(b)] $\alpha_1 : \mathcal{O}_{Y_1}^{\oplus r_1}
\to \pi_{1, *}\mathcal{G}_1$ であって、
% P05-EN-CH38-0029: source-faithful alpha/alpha_1 mismatch; see sidecar.
$\alpha \otimes \kappa(\xi_1)$ が同型となるものが存在する、
\item[(c)] $Y_1 \subset X_1$ は開で、$y_1 = x_1$ かつ
$\mathcal{F}_1|_{Y_1} \cong \Coker(\alpha_1)$ である、
\item[(d)] $x_1$ における $\mathcal{F}_1/X_1/S$ の完全デヴィサージュ
$(Z_k, Y_k, i_k, \pi_k, \mathcal{G}_k, \alpha_k, z_k, y_k)_{k = 2, \ldots, n}$
が存在する。
\end{enumerate}
証明を終えるには、一段階デヴィサージュと完全デヴィサージュを縮小して、
両者が一つの完全デヴィサージュとしてつながるようにすればよい。
（以下の証明を読む代わりに、Lemmas \ref{flat-lemma-shrink} と
\ref{flat-lemma-shrink-complete} を用いて読者自身でこれを行うことを勧める。）
$Y_1 \subset X_1$ は $x_1$ の開近傍なので、
Lemma \ref{flat-lemma-shrink-complete} (\ref{flat-item-shrink-on-X-complete})
を適用し、データ (d) の標準縮小
$S', X'_1, Z'_2, Y'_2, \ldots, Y'_n$ で $X'_1 \subset Y_1$ を満たすものを得る。
$X'_1$ はアフィンスキーム $Y_1$ の標準開集合でもあることに注意する。
次に、データ (a) を次のように縮小する。まず
Lemma \ref{flat-lemma-shrink} (\ref{flat-item-shrink-base}) により基底 $S$ を $S'$ に縮小し、
次いで Lemma \ref{flat-lemma-shrink} (\ref{flat-item-shrink-on-Y}) を用いて
その結果を $S''$, $X''$, $Z''_1$, $Y''_1$ に縮小し、最終的に
$Y''_1 = X'_1 \times_S S''$ かつ $S'' \subset S'$ となるようにする。
すると、
$$
Z''_1, Y''_1, Z'_2 \times_{S'} S'', Y'_2 \times_{S'} S'', \ldots,
Y'_n \times_{S'} S''
$$
が求めていた完全デヴィサージュを与えることが分かる。
\end{proof}

\noindent
さらにデータを整理すると、次の帰結を得る。

\begin{lemma}
\label{flat-lemma-existence-complete}
$X \to S$ をスキームの有限型射とする。
$\mathcal{F}$ を有限型準連接 $\mathcal{O}_X$ 加群とする。
$s \in S$ を点とする。
基本エタール近傍 $(S', s') \to (S, s)$ およびエタール射
$h_i : Y_i \to X_{S'}$, $i = 1, \ldots, n$ であって、各 $i$ に対し
$s'$ 上の $\mathcal{F}_i/Y_i/S'$ の完全デヴィサージュが存在し、かつ
$X_s = (X_{S'})_{s'} \subset \bigcup h_i(Y_i)$ を満たすものが存在する。
ここで $\mathcal{F}_i$ は $\mathcal{F}$ の $Y_i$ への引き戻しである。
\end{lemma}

\begin{proof}
各点 $x \in X_s$ に対し、図式
$$
\xymatrix{
(X, x) \ar[d] & (X', x') \ar[l]^g \ar[d] \\
(S, s) & (S', s') \ar[l]
}
$$
であって、横向きの射が基本エタール近傍であり、かつ
$g^*\mathcal{F}/X'/S'$ が $x'$ において完全デヴィサージュをもつものを得られる。
$X \to S$ は有限型なのでファイバー $X_s$ は準コンパクトであり、
上の各 $g : X' \to X$ は開写像だから、$X_s$ は有限個の
$g(X'_{s'})$ で被覆できる。したがって、そのような図式の有限族
$$
\vcenter{
\xymatrix{
% P05-EN-CH38-0030: source-faithful unindexed distinguished point; see sidecar.
(X, x) \ar[d] & (X'_i, x'_i) \ar[l]^{g_i} \ar[d] \\
(S, s) & (S'_i, s'_i) \ar[l]
}
}
\quad i = 1, \ldots, n
$$
% P05-EN-CH38-0031: source-faithful equality/covering-containment mismatch; see sidecar.
で $X_s = \bigcup g_i(X'_i)$ を満たすものを得られる。
$S' = S'_1 \times_S \ldots \times_S S'_n$ とおく。
% P05-EN-CH38-0032: source-faithful undefined X_i; see sidecar.
$Y_i = X_i \times_{S'_i} S'$ を $X'_i$ の $S'$ への基底変換とする。
Lemma \ref{flat-lemma-base-change-complete} により、$\mathcal{F}$ の $Y_i$ への引き戻しは
% P05-EN-CH38-0033: source-faithful s/s' mismatch; see sidecar.
$s$ 上の完全デヴィサージュをもつことが分かり、証明が完了する。
\end{proof}



\section{代数への翻訳}
\label{flat-section-translation}

\noindent
完全デヴィサージュをもつことの意味を代数的に明記しておくと有用かもしれない。
そこで、次の概念を導入する（それほど有用な概念ではないので、
途方もなく長い名前を付けることにする）。

\begin{definition}
\label{flat-definition-elementary-etale-neighbourhood}
$R \to S$ を環準同型とする。$\mathfrak q$ を、$R$ の素イデアル
$\mathfrak p$ の上にある $S$ の素イデアルとする。
{\it $\mathfrak q$ における環準同型 $R \to S$ の基本エタール局所化}とは、
環の可換図式およびそれに付随する素イデアル
$$
\xymatrix{
S \ar[r] & S' \\
R \ar[u] \ar[r] & R' \ar[u]
}
\quad\quad
\xymatrix{
\mathfrak q \ar@{-}[r] & \mathfrak q' \\
\mathfrak p \ar@{-}[u] \ar@{-}[r] & \mathfrak p' \ar@{-}[u]
}
$$
であって、$R \to R'$ と $S \to S'$ がエタール環準同型であり、かつ
$\kappa(\mathfrak p) = \kappa(\mathfrak p')$ および
$\kappa(\mathfrak q) = \kappa(\mathfrak q')$ を満たすものとして与えられる。
\end{definition}

\begin{definition}
\label{flat-definition-complete-devissage-algebra}
$R \to S$ を有限型環準同型とする。
$\mathfrak r$ を $R$ の素イデアルとする。
$N$ を有限 $S$ 加群とする。
{\it $\mathfrak r$ 上の $N/S/R$ の完全デヴィサージュ}とは、
$R$ 代数準同型
$$
\xymatrix{
& A_1 & & A_2 & & ... & & A_n \\
S \ar[ru] & & B_1 \ar[lu] \ar[ru] & & ... \ar[lu] \ar[ru] & &
... \ar[lu] \ar[ru] & & B_n \ar[lu]
}
$$
、有限 $A_i$ 加群 $M_i$、および $B_i$ 加群準同型
$\alpha_i : B_i^{\oplus r_i} \to M_i$ であって次を満たすものとして与えられる：
\begin{enumerate}
\item $S \to A_1$ は全射かつ有限表示である、
% P05-EN-CH38-0034: source-faithful missing i<n range in items (2) and (6); see sidecar.
\item $B_i \to A_{i + 1}$ は全射かつ有限表示である、
\item $B_i \to A_i$ は有限である、
\item $R \to B_i$ は滑らかで、幾何学的既約なファイバーをもつ、
\item $S$ 加群として $N \cong M_1$ である、
\item $B_i$ 加群として $\Coker(\alpha_i) \cong M_{i + 1}$ である、
\item $\mathfrak p_i = \mathfrak rB_i$ とおくと、$\alpha_i : \kappa(\mathfrak p_i)^{\oplus r_i}
\to M_i \otimes_{B_i} \kappa(\mathfrak p_i)$ は同型である、
\item $\Coker(\alpha_n) = 0$.
\end{enumerate}
この状況で、
$(A_i, B_i, M_i, \alpha_i)_{i = 1, \ldots, n}$
を $\mathfrak r$ 上の $N/S/R$ の完全デヴィサージュという。
\end{definition}

\begin{remark}
\label{flat-remark-finite-presentation}
すべての $i$ に対する $R$ 代数 $B_i$ と、$i \geq 2$ に対する $A_i$ は、
$R$ 上有限表示であることに注意する。$S$ が $R$ 上有限表示ならば、
$A_1$ も $R$ 上有限表示である。この場合、完全デヴィサージュに現れる
すべての環準同型は有限表示である。Algebra, Lemma
\ref{algebra-lemma-compose-finite-type} を参照せよ。
引き続き $S$ が $R$ 上有限表示であると仮定すると、次は同値である：
\begin{enumerate}
% P05-EN-CH38-0035: source-faithful undefined M/N mismatch; see sidecar.
\item $M$ は $S$ 上有限表示である、
\item $M_1$ は $A_1$ 上有限表示である、
\item $M_1$ は $B_1$ 上有限表示である、
\item 各 $M_i$ は $A_i$ 加群としても $B_i$ 加群としても有限表示である。
\end{enumerate}
(1) $\Leftrightarrow$ (2) および (2) $\Leftrightarrow$ (3) は
Algebra, Lemma \ref{algebra-lemma-finite-finitely-presented-extension}
から従う。$M_1$ が有限表示ならば $\Coker(\alpha_1)$ も有限表示であり
（Algebra, Lemma \ref{algebra-lemma-extension} を参照）、したがって
$M_2$ なども有限表示である。
\end{remark}

\begin{definition}
\label{flat-definition-complete-devissage-at-x-algebra}
$R \to S$ を有限型環準同型とする。
$\mathfrak q$ を、$R$ の素イデアル $\mathfrak r$ の上にある $S$ の素イデアルとする。
$N$ を有限 $S$ 加群とする。
{\it $\mathfrak q$ における $N/S/R$ の完全デヴィサージュ}とは、
$\mathfrak r$ 上の $N/S/R$ の完全デヴィサージュ
$(A_i, B_i, M_i, \alpha_i)_{i = 1, \ldots, n}$ と、
$\mathfrak r$ の上にある素イデアル $\mathfrak q_i \subset B_i$ であって、
次を満たすものとして与えられる：
\begin{enumerate}
\item $\kappa(\mathfrak r) \subset \kappa(\mathfrak q_i)$ は純超越拡大である、
\item $\mathfrak q_i \subset B_i$ の上にある素イデアル
$\mathfrak q'_i \subset A_i$ は一意である、
% P05-EN-CH38-0036: source-faithful ill-typed contraction to A_i; see sidecar.
\item $\mathfrak q = \mathfrak q'_1 \cap S$ かつ
$\mathfrak q_i = \mathfrak q'_{i + 1} \cap A_i$ である、
\item $R \to B_i$ の相対次元は
$\dim_{\mathfrak q_i}(\text{Supp}(M_i \otimes_R \kappa(\mathfrak r)))$
である。
\end{enumerate}
\end{definition}

\begin{remark}
\label{flat-remark-same-notion}
$A \to B$ を有限型環準同型とし、$N$ を有限 $B$ 加群とする。
$\mathfrak q$ を、$A$ の素イデアル $\mathfrak r$ の上にある $B$ の素イデアルとする。
$X = \Spec(B)$, $S = \Spec(A)$ とおき、$X$ 上で
$\mathcal{F} = \widetilde{N}$ とする。$x$ を $\mathfrak q$ に対応する点とし、
% P05-EN-CH38-0037: source-faithful undefined p/r mismatch throughout the remark; see sidecar.
$s \in S$ を $\mathfrak p$ に対応する点とする。このとき、
\begin{enumerate}
\item $s$ 上の $\mathcal{F}/X/S$ の完全デヴィサージュが存在するならば、
$\mathfrak p$ 上の $N/B/A$ の完全デヴィサージュが存在する、
\item $x$ における $\mathcal{F}/X/S$ の完全デヴィサージュが存在することと、
$\mathfrak q$ における $N/B/A$ の完全デヴィサージュが存在することは同値である。
\end{enumerate}
ただし、「$\mathfrak p$ 上の $N/B/A$ の完全デヴィサージュ」の定式化では
相対次元に関する条件を省略したという小さな違いがあり、そのため (1) の含意は
一方向にしか成り立たない。$\mathfrak q$ における完全デヴィサージュという概念には、
この条件が組み込まれている。いずれにせよ、ここで用いるのは
$\mathcal{F}/X/S$ に対する存在が $N/B/A$ に対する存在を導くことだけである。
\end{remark}

\begin{lemma}
\label{flat-lemma-existence-algebra}
$R \to S$ を有限型環準同型とする。
$M$ を有限 $S$ 加群とする。
$\mathfrak q$ を $S$ の素イデアルとする。
$\mathfrak q$ における環準同型 $R \to S$ の基本エタール局所化
$R' \to S', \mathfrak q', \mathfrak p'$ であって、
$\mathfrak q'$ における $(M \otimes_S S')/S'/R'$ の
完全デヴィサージュが存在するものがある。
\end{lemma}

\begin{proof}
これは Proposition \ref{flat-proposition-existence-complete-at-x} を
Remark \ref{flat-remark-same-notion} によって言い換えたものである。
\end{proof}




\section{局所化と普遍単射写像}
\label{flat-section-localize-universally-injective}


\begin{lemma}
\label{flat-lemma-homothety-spectrum}
$R \to S$ を環準同型とする。
$N$ を $S$ 加群とする。次を仮定する：
\begin{enumerate}
\item $R$ は極大イデアル $\mathfrak m$ をもつ局所環である、
\item $\overline{S} = S/\mathfrak m S$ は Noetherian である、
% P05-EN-CH38-0038: source-faithful undefined m_R/m mismatch; see sidecar.
\item $\overline{N} = N/\mathfrak m_R N$ は有限 $\overline{S}$ 加群である。
\end{enumerate}
$\Sigma \subset S$ を、$\overline{N}$ 上の零因子でない元からなる
乗法的部分集合とする。このとき $\Sigma^{-1}S$ は半局所環であり、
そのスペクトルは $\text{Ass}_S(\overline{N})$ のある元に含まれる
素イデアル $\mathfrak q \subset S$ からなる。さらに、
$\Sigma^{-1}S$ の任意の極大イデアルは、$\overline{S}$ 上の
$\overline{N}$ の随伴素イデアルに対応する。
\end{lemma}

\begin{proof}
Algebra, Lemma \ref{algebra-lemma-ass-quotient-ring} により、
$\text{Ass}_S(\overline{N}) = \text{Ass}_{\overline{S}}(\overline{N})$
であることに注意する。Algebra, Lemma \ref{algebra-lemma-finite-ass}
により、これは有限集合である。
$\{\mathfrak q_1, \ldots, \mathfrak q_r\} = \text{Ass}_S(\overline{N})$
と書く。Algebra, Lemma \ref{algebra-lemma-ass-zero-divisors} により
$\Sigma = S \setminus (\bigcup \mathfrak q_i)$ である。
Algebra, Lemma \ref{algebra-lemma-spec-localization} における
$\Spec(\Sigma^{-1}S)$ の記述と Algebra, Lemma \ref{algebra-lemma-silly}
により、$\Sigma^{-1}S$ の素イデアルは、いずれかの $\mathfrak q_i$ に
含まれる $S$ の素イデアルに対応する。したがって、$\Sigma^{-1}S$ の
極大イデアルは、集合 $\{\mathfrak q_1, \ldots, \mathfrak q_r\}$ の
包含関係に関する極大元と一対一に対応する。これで補題が証明された。
\end{proof}

\begin{lemma}
\label{flat-lemma-homothety-universally-injective}
仮定と記号は Lemma \ref{flat-lemma-homothety-spectrum} のとおりとする。
さらに次を仮定する：
\begin{enumerate}
\item $S$ は局所環で、$R \to S$ は局所準同型である、
\item $S$ は $R$ 上本質的有限表示である、
\item $N$ は $S$ 上有限表示である、
\item $N$ は $R$ 上平坦である。
\end{enumerate}
このとき、各 $s \in \Sigma$ は普遍単射な $R$ 加群準同型
$s : N \to N$ を定め、写像 $N \to \Sigma^{-1}N$ は
$R$-普遍単射である。
\end{lemma}

\begin{proof}
Algebra, Lemma \ref{algebra-lemma-mod-injective-general} により、列
$0 \to N \to N \to N/sN \to 0$ は完全で、$N/sN$ は $R$ 上平坦である。
したがって $s : N \to N$ は普遍単射である。Algebra, Lemma
\ref{algebra-lemma-flat-tor-zero} を参照せよ。写像
$N \to \Sigma^{-1}N$ は、写像 $s : N \to N$ の有向余極限なので普遍単射である。
\end{proof}

\begin{lemma}
\label{flat-lemma-base-change-universally-flat-local}
$R \to S$ を環準同型とする。
$N$ を $S$ 加群とする。
$S \to S'$ を環準同型とする。次を仮定する：
\begin{enumerate}
% P05-EN-CH38-0039: uniquely determined missing list punctuation; see sidecar.
\item $R \to S$ は局所環の局所準同型である、
\item $S$ は $R$ 上本質的有限表示である、
\item $N$ は $S$ 上有限表示である、
\item $N$ は $R$ 上平坦である、
\item $S \to S'$ は平坦である、
\item $\Spec(S') \to \Spec(S)$ の像は、$\mathfrak m_R$ の上にある
$S$ の素イデアル $\mathfrak q$ で、$\mathfrak q$ が
$N/\mathfrak m_R N$ の随伴素イデアルとなるものをすべて含む。
\end{enumerate}
このとき $N \to N \otimes_S S'$ は $R$-普遍単射である。
\end{lemma}

\begin{proof}
% P05-EN-CH38-0040: source-faithful R/S tensor mismatch; see sidecar.
$N' = N \otimes_R S'$ とおく。可換図式
$$
\xymatrix{
N \ar[d] \ar[r] & N' \ar[d] \\
\Sigma^{-1}N \ar[r] & \Sigma^{-1}N'
}
$$
を考える。ここで $\Sigma \subset S$ は $N/\mathfrak m_R N$ 上の
零因子でない元の集合である。写像 $N \to \Sigma^{-1}N'$ が
普遍単射であることを示せば、$N \to N'$ も普遍単射となる
（Algebra, Lemma \ref{algebra-lemma-universally-injective-permanence} を参照）。

\medskip\noindent
Lemma \ref{flat-lemma-homothety-spectrum} により、環 $\Sigma^{-1}S$ は
半局所環で、その極大イデアルは $N/\mathfrak m_R N$ の随伴素イデアルに対応する。
したがって仮定により、$\Spec(\Sigma^{-1}S') \to \Spec(\Sigma^{-1}S)$
の像はこれらの極大イデアルをすべて含む。
Algebra, Lemma \ref{algebra-lemma-ff-rings} により、環準同型
$\Sigma^{-1}S \to \Sigma^{-1}S'$ は忠実平坦である。したがって、写像
$\Sigma^{-1}N \to \Sigma^{-1}N'$、すなわち
$$
N \otimes_S \Sigma^{-1}S \longrightarrow N \otimes_S \Sigma^{-1}S'
$$
は普遍単射である。Algebra, Lemmas
\ref{algebra-lemma-faithfully-flat-universally-injective} および
\ref{algebra-lemma-universally-injective-tensor} を参照せよ。
最後に Lemma \ref{flat-lemma-homothety-universally-injective} を適用すると、
$N \to \Sigma^{-1}N$ は普遍単射であることが分かる。普遍単射な加群準同型の
合成は普遍単射なので（Algebra, Lemma
\ref{algebra-lemma-composition-universally-injective} を参照）、
$N \to \Sigma^{-1}N'$ は普遍単射であり、証明が完了する。
\end{proof}

\begin{lemma}
\label{flat-lemma-base-change-universally-flat}
$R \to S$ を環準同型とする。
$N$ を $S$ 加群とする。
$S \to S'$ を環準同型とする。次を仮定する：
\begin{enumerate}
\item $R \to S$ は有限表示で、$N$ は $S$ 上有限表示である、
\item $N$ は $R$ 上平坦である、
\item $S \to S'$ は平坦である、
\item $\Spec(S') \to \Spec(S)$ の像は、素イデアル $\mathfrak q$ であって、
$\mathfrak q$ が $N \otimes_R \kappa(\mathfrak p)$ の随伴素イデアルとなるものを
すべて含む。ここで $\mathfrak p$ は $\mathfrak q$ の $R$ における逆像である。
\end{enumerate}
このとき $N \to N \otimes_S S'$ は $R$-普遍単射である。
\end{lemma}

\begin{proof}
Algebra, Lemma \ref{algebra-lemma-universally-injective-check-stalks}
により、$R$ の $\mathfrak p$ の上にある $S$ の任意の素イデアル
$\mathfrak q$ に対し、
% P05-EN-CH38-0041: source-faithful R/S tensor mismatch; see sidecar.
$N_{\mathfrak q} \to (N \otimes_R S')_{\mathfrak q}$ が
% P05-EN-CH38-0042: uniquely determined article repair; see sidecar.
$R_{\mathfrak p}$-普遍単射であることを示せば十分である。したがって、環準同型
$R_{\mathfrak p} \to S_{\mathfrak q} \to S'_{\mathfrak q}$
および加群 $N_{\mathfrak q}$ に Lemma
\ref{flat-lemma-base-change-universally-flat-local} を適用できる。
\end{proof}

\noindent
読者は次の補題を Algebra, Lemmas \ref{algebra-lemma-mod-injective},
\ref{algebra-lemma-mod-injective-general} および Section
\ref{flat-section-variants-mod-injective} の結果と比較するとよいだろう。
いずれの場合にも、結論は写像 $u : M \to N$ が平坦な余核をもつ
普遍単射であるというものである。

\begin{lemma}
\label{flat-lemma-universally-injective-local}
$(R, \mathfrak m)$ を局所環とする。$u : M \to N$ を $R$ 加群準同型とする。
$M$ が射影的 $R$ 加群、$N$ が平坦 $R$ 加群であり、
$\overline{u} : M/\mathfrak mM \to N/\mathfrak mN$ が単射ならば、
$u$ は普遍単射である。
\end{lemma}

\begin{proof}
Algebra, Theorem \ref{algebra-theorem-projective-free-over-local-ring}
により、加群 $M$ は自由である。$M$ の有限生成直和因子のそれぞれについて
結果を示せば補題が従う。したがって $M$ は有限自由であると仮定してよい。
$N = \colim_i N_i$ を有限自由加群の有向余極限として表す。Algebra, Theorem
\ref{algebra-theorem-lazard} を参照せよ。$M$ は有限自由なので、
$u : M \to N$ はある $i$ に対して $N_i$ を経由する。対応する $R$ 加群準同型を
$u_i : M \to N_i$ と書く。$\overline{u}$ は単射なので、
$\overline{u_i} : M/\mathfrak mM \to N_i/\mathfrak mN_i$ は単射であり、
任意の $i' \geq i$ に対して写像
$N_i/\mathfrak mN_i \to N_{i'}/\mathfrak mN_{i'}$ と合成した後も単射である。
$M$ と $N_{i'}$ は局所環 $R$ 上有限自由なので、任意の $i' \geq i$ に対し
$M \to N_{i'}$ は分裂単射である。したがって、任意の $R$ 加群 $Q$ に対し
$M \otimes_R Q \to N_{i'} \otimes_R Q$ はすべての $i' \geq i$ について単射である。
$- \otimes_R Q$ は余極限と可換なので、
% P05-EN-CH38-0043: source-faithful N_{i'}/N colimit-target mismatch; see sidecar.
$M \otimes_R Q \to N_{i'} \otimes_R Q$ は所望のとおり単射である。
\end{proof}

\begin{lemma}
\label{flat-lemma-invert-universally-injective}
仮定と記号は Lemma \ref{flat-lemma-homothety-spectrum} のとおりとする。
さらに $N$ は $R$ 加群として射影的であると仮定する。
このとき、各 $s \in \Sigma$ は普遍単射な $R$ 加群準同型
$s : N \to N$ を定め、写像 $N \to \Sigma^{-1}N$ は
$R$-普遍単射である。
\end{lemma}

\begin{proof}
$s \in \Sigma$ を取る。Lemma \ref{flat-lemma-universally-injective-local}
により、写像 $s : N \to N$ は普遍単射である。写像
$N \to \Sigma^{-1}N$ は、写像 $s : N \to N$ の有向余極限なので普遍単射である。
\end{proof}






\section{完備化と Mittag--Leffler 加群}
\label{flat-section-completion-ML}


\begin{lemma}
\label{flat-lemma-completed-direct-sum-ML}
$R$ を環、$I \subset R$ をイデアル、$A$ を集合とする。
$R$ は Noetherian で $I$ に関して完備であると仮定する。このとき完備化
$(\bigoplus\nolimits_{\alpha \in A} R)^\wedge$ は平坦かつ Mittag--Leffler である。
\end{lemma}

\begin{proof}
More on Algebra, Lemma
\ref{more-algebra-lemma-ui-completion-direct-sum-into-product} により、写像
$(\bigoplus\nolimits_{\alpha \in A} R)^\wedge
\to \prod_{\alpha \in A} R$ は普遍単射である。したがって Algebra, Lemmas
\ref{algebra-lemma-ui-flat-domain} および
\ref{algebra-lemma-pure-submodule-ML} により、
$\prod_{\alpha \in A} R$ が平坦かつ Mittag--Leffler であることを示せば十分である。
Algebra, Proposition \ref{algebra-proposition-characterize-coherent}
（および Algebra, Lemma \ref{algebra-lemma-Noetherian-coherent}）により、
$\prod_{\alpha \in A} R$ は平坦である。また、$R$ のコピーの積は
Mittag--Leffler なので結論を得る。Algebra, Lemma
\ref{algebra-lemma-product-over-Noetherian-ring} を参照せよ。
\end{proof}

\begin{lemma}
\label{flat-lemma-lift-ML}
$R$ を環、$I \subset R$ をイデアルとする。
$M$ を $R$ 加群とする。次を仮定する：
\begin{enumerate}
\item $R$ は Noetherian で $I$ 進完備である、
\item $M$ は $R$ 上平坦である、
\item $M/IM$ は射影的 $R/I$ 加群である。
\end{enumerate}
このとき $I$ 進完備化 $M^\wedge$ は平坦な Mittag--Leffler $R$ 加群である。
\end{lemma}

\begin{proof}
$F$ が自由 $R$ 加群となる全射 $F \to M$ を選ぶ。Algebra, Lemma
\ref{algebra-lemma-split-completed-sequence} により、加群 $M^\wedge$ は
加群 $F^\wedge$ の直和因子である。したがって $F$ について補題を証明すれば十分であり、
この場合は Lemma \ref{flat-lemma-completed-direct-sum-ML} から従う。
\end{proof}

\noindent
Lemmas \ref{flat-lemma-universally-injective-to-completion} および
\ref{flat-lemma-universally-injective-to-completion-flat} において、
$R \to S$ が有限型ならば $S$ が Noetherian であるという仮定は満たされる。
Algebra, Lemma \ref{algebra-lemma-Noetherian-permanence} を参照せよ。

\begin{lemma}
\label{flat-lemma-universally-injective-to-completion}
$R$ を環、$I \subset R$ をイデアルとする。
$R \to S$ を環準同型、$N$ を $S$ 加群とする。次を仮定する：
\begin{enumerate}
\item $R$ は Noetherian 環である、
\item $S$ は Noetherian 環である、
\item $N$ は有限 $S$ 加群である、
\item 任意の有限 $R$ 加群 $Q$ と任意の
$\mathfrak q \in \text{Ass}_S(Q \otimes_R N)$ に対し
$IS + \mathfrak q \not = S$ が成り立つ。
\end{enumerate}
このとき、$N$ から $N$ の $I$ 進完備化への写像 $N \to N^\wedge$ は、
$R$ 加群の写像として普遍単射である。
\end{lemma}

\begin{proof}
任意の有限 $R$ 加群 $Q$ に対し、写像
$Q \otimes_R N \to Q \otimes_R N^\wedge$ が単射であることを示す必要がある。
Algebra, Theorem \ref{algebra-theorem-universally-exact-criteria} を参照せよ。
標準写像 $Q \otimes_R N^\wedge \to (Q \otimes_R N)^\wedge$ が存在するので、
標準写像 $Q \otimes_R N \to (Q \otimes_R N)^\wedge$ が単射であることを
証明すれば十分である。したがって $N$ を $Q \otimes_R N$ で置き換えてよく、
写像 $N \to N^\wedge$ の単射性を証明すれば十分である。

\medskip\noindent
$K = \Ker(N \to N^\wedge)$ とおく。$N$ は
$\prod_{\mathfrak q \in \text{Ass}(N)} N_{\mathfrak q}$ の部分加群なので、
$\mathfrak q \in \text{Ass}(N)$ に対して $K_{\mathfrak q} = 0$ を示せば十分である。
Algebra, Lemma \ref{algebra-lemma-zero-at-ass-zero} を参照せよ。
$\mathfrak q \in \text{Ass}(N)$ を取る。最後の仮定により、
$\mathfrak q' \supset IS + \mathfrak q$ となる素イデアルが存在する。
$K_{\mathfrak q}$ は $K_{\mathfrak q'}$ の局所化なので、$K_{\mathfrak q'}$ の
消滅を証明すれば十分である。$K = \bigcap I^nN$ であるから
$K_{\mathfrak q'} \subset \bigcap I^nN_{\mathfrak q'}$ である。よって
Algebra, Lemma \ref{algebra-lemma-intersect-powers-ideal-module-zero} により
$K_{\mathfrak q'} = 0$ である。
\end{proof}

\begin{lemma}
\label{flat-lemma-universally-injective-to-completion-flat}
$R$ を環、$I \subset R$ をイデアルとする。
$R \to S$ を環準同型、$N$ を $S$ 加群とする。次を仮定する：
\begin{enumerate}
\item $R$ は Noetherian 環である、
\item $S$ は Noetherian 環である、
\item $N$ は有限 $S$ 加群である、
\item $N$ は $R$ 上平坦である、
\item $N \otimes_R \kappa(\mathfrak p)$ の随伴素イデアルとなる任意の
素イデアル $\mathfrak q \subset S$ に対し、ここで
% P05-EN-CH38-0044: source-faithful contraction/intersection ambiguity; see sidecar.
$\mathfrak p = R \cap \mathfrak q$ とおくと、$IS + \mathfrak q \not = S$ が成り立つ。
\end{enumerate}
このとき、$N$ から $N$ の $I$ 進完備化への写像 $N \to N^\wedge$ は、
$R$ 加群の写像として普遍単射である。
\end{lemma}

\begin{proof}
Algebra, Lemma \ref{algebra-lemma-bourbaki-fibres} および
Remark \ref{algebra-remark-bourbaki} により、テンソル積
$N \otimes_R Q$ の随伴素イデアルの集合は、加群
$N \otimes_R \kappa(\mathfrak p)$ の随伴素イデアルの集合に含まれる。
% P05-EN-CH38-0045: uniquely determined subject-agreement repair; see sidecar.
したがって Lemma \ref{flat-lemma-universally-injective-to-completion} から従う。
\end{proof}




\section{射影加群}
\label{flat-section-projective}

\noindent
次の補題を用いると、基底に関する Noetherian 帰納法によって射影性を証明できる。
Lemma \ref{flat-lemma-fibres-irreducible-flat-projective} を参照せよ。

\begin{lemma}
\label{flat-lemma-flat-pure-over-complete-projective}
$R$ を環、$I \subset R$ をイデアルとする。
$R \to S$ を環準同型、$N$ を $S$ 加群とする。次を仮定する：
\begin{enumerate}
\item $R$ は Noetherian で $I$ 進完備である、
\item $R \to S$ は有限型である、
\item $N$ は有限 $S$ 加群である、
\item $N$ は $R$ 上平坦である、
\item $N/IN$ は射影的 $R/I$ 加群である、
\item $N \otimes_R \kappa(\mathfrak p)$ の随伴素イデアルとなる任意の
素イデアル $\mathfrak q \subset S$ に対し、ここで
% P05-EN-CH38-0046: source-faithful contraction/intersection ambiguity; see sidecar.
$\mathfrak p = R \cap \mathfrak q$ とおくと、$IS + \mathfrak q \not = S$ が成り立つ。
\end{enumerate}
このとき $N$ は射影的 $R$ 加群である。
\end{lemma}

\begin{proof}
Lemma \ref{flat-lemma-universally-injective-to-completion-flat} により、写像
$N \to N^\wedge$ は普遍単射である。Lemma \ref{flat-lemma-lift-ML} により、
加群 $N^\wedge$ は Mittag--Leffler である。Algebra, Lemma
\ref{algebra-lemma-pure-submodule-ML} により、$N$ も Mittag--Leffler である。
したがって $N$ は $R$ 加群として可算生成、平坦、かつ Mittag--Leffler であり、
Algebra, Lemma \ref{algebra-lemma-countgen-projective} により射影的である。
\end{proof}

\begin{lemma}
\label{flat-lemma-fibres-irreducible-flat-projective}
$R$ を環、$R \to S$ を環準同型とする。次を仮定する：
\begin{enumerate}
\item $R$ は Noetherian である、
\item $R \to S$ は有限型かつ平坦である、
\item 各ファイバー環 $S \otimes_R \kappa(\mathfrak p)$ は
$\kappa(\mathfrak p)$ 上幾何学的整である。
\end{enumerate}
このとき $S$ は射影的 $R$ 加群である。
\end{lemma}

\begin{proof}
集合
$$
\{I \subset R \mid S/IS\text{ は }R/I\text{-加群として射影的でない}\}
$$
を考える。この集合が空であることを示さなければならない。矛盾を導くため、
空でないと仮定する。このとき極大元 $I$ をもつ。その根基を
$J = \sqrt{I}$ とする。$I \not = J$ ならば、$S/JS$ は射影的 $R/J$ 加群であり、
$S/IS$ は $R/I$ 上平坦で、$J/I$ は $R/I$ の冪零イデアルである。
Algebra, Lemma \ref{algebra-lemma-lift-projective} を適用すると、
$S/IS$ は射影的 $R/I$ 加群となり、矛盾する。したがって $I$ は根基イデアルと
仮定してよい。言い換えれば、$R$ が被約環で、$R$ の任意の非零イデアル $I$ に対し
$S/IS$ が射影的 $R/I$ 加群である場合に補題を証明すればよい。

\medskip\noindent
$R$ は被約環で、$R$ の任意の非零イデアル $I$ に対して $S/IS$ は
射影的 $R/I$ 加群であると仮定する。一般平坦性、すなわち
Algebra, Lemma \ref{algebra-lemma-generic-flatness-Noetherian}
（整域となる局所化 $R_g$ に適用する）または、より一般的な
Algebra, Lemma \ref{algebra-lemma-generic-flatness-reduced} により、
$S_f$ が自由 $R_f$ 加群となる非零元 $f \in R$ が存在する。
$R$ の $(f)$ 進完備化を $R^\wedge = \lim R/(f^n)$ と書く。環準同型
$$
R \longrightarrow R_f \times R^\wedge
$$
は忠実平坦である。Algebra, Lemma \ref{algebra-lemma-completion-flat}
を参照せよ。したがって射影性の忠実平坦降下、すなわち Algebra, Theorem
\ref{algebra-theorem-ffdescent-projectivity} により、
$S \otimes_R R^\wedge$ が射影的 $R^\wedge$ 加群であることを示せば十分である。
そのために Lemma \ref{flat-lemma-flat-pure-over-complete-projective} の判定法を用いる。
まず、$S/fS = (S \otimes_R R^\wedge)/f(S \otimes_R R^\wedge)$ は
射影的 $R/(f)$ 加群であり、$S \otimes_R R^\wedge$ はその基底変換として
$R^\wedge$ 上平坦かつ有限型であることに注意する。次に、
$\mathfrak p^\wedge$ を $R^\wedge$ の素イデアルとする。
$\mathfrak p \subset R$ を対応する $R$ の素イデアルとする。
$R \to S$ は幾何学的整なファイバー環をもつので、任意の基底変換の
ファイバー環についても同じことが成り立つ。したがって
$\mathfrak q^\wedge = \mathfrak p^\wedge(S \otimes_R R^\wedge)$,
% P05-EN-CH38-0047: uniquely determined singular/plural and punctuation repair; see sidecar.
は $\mathfrak p^\wedge$ の上にある素イデアルであり、
$S \otimes_R \kappa(\mathfrak p^\wedge)$ の唯一の随伴素イデアルである。
よって $f(S \otimes_R R^\wedge) + \mathfrak q^\wedge
\not = S \otimes_R R^\wedge$ ならば証明は終わる。これは成り立つ。実際、
$f$ は $f$ 進完備環 $R^\wedge$ の Jacobson 根基に属するので
$\mathfrak p^\wedge + fR^\wedge \not = R^\wedge$ であり、また
$R^\wedge \to S \otimes_R R^\wedge$ のファイバーは空でない
（既約空間は空でない）ため、この射はスペクトル上全射である。
\end{proof}

\begin{lemma}
\label{flat-lemma-fibres-irreducible-flat-projective-nonnoetherian}
$R$ を環、$R \to S$ を環準同型とする。次を仮定する：
\begin{enumerate}
\item $R \to S$ は有限表示かつ平坦である、
\item 各ファイバー環 $S \otimes_R \kappa(\mathfrak p)$ は
$\kappa(\mathfrak p)$ 上幾何学的整である。
\end{enumerate}
このとき $S$ は射影的 $R$ 加群である。
\end{lemma}

\begin{proof}
環の余 Cartesian 図式
$$
\xymatrix{
S_0 \ar[r] & S \\
R_0 \ar[u] \ar[r] & R \ar[u]
}
$$
であって、$R_0$ が $\mathbf{Z}$ 上有限型であり、$R_0 \to S_0$ が
幾何学的整なファイバーをもつ有限型平坦射となるものを得られる。
More on Morphisms, Lemmas
\ref{more-morphisms-lemma-Noetherian-approximation-flat},
\ref{more-morphisms-lemma-Noetherian-approximation-geometrically-reduced},
\ref{more-morphisms-lemma-Noetherian-approximation-geometrically-irreducible},
\ref{more-morphisms-lemma-Noetherian-approximation-combine} を参照せよ。
Lemma \ref{flat-lemma-fibres-irreducible-flat-projective} により、$S_0$ は
射影的 $R_0$ 加群である。したがって $S = S_0 \otimes_{R_0} R$ は
射影的 $R$ 加群である。Algebra, Lemma
\ref{algebra-lemma-ascend-properties-modules} を参照せよ。
\end{proof}

\begin{remark}
\label{flat-remark-how-in-RG}
Lemma \ref{flat-lemma-fibres-irreducible-flat-projective-nonnoetherian} は、
本章の諸結果を展開するうえで鍵となる段階である。\cite{GruRay} における
この補題の類似物は \cite[I Proposition 3.3.1]{GruRay} である。すなわち、
$R \to S$ が幾何学的整なファイバーをもつ滑らかな射ならば、$S$ は
射影的 $R$ 加群である。これは Lemma
\ref{flat-lemma-fibres-irreducible-flat-projective-nonnoetherian} の特殊な場合であるが、
後でいずれこの補題を改良するので、現時点でより強い結果を得ても大きな利点はない。
\cite{GruRay} で与えられているこの証明を簡潔に概説する。
\begin{enumerate}
\item まず、上と同様に $R$ が Noetherian である場合へ帰着する。
\item 射影性は忠実平坦な環準同型に沿って降下するので（Algebra, Theorem
\ref{algebra-theorem-ffdescent-projectivity} を参照）、$R$ 上の fppf 位相で
局所的に作業してよい。したがって $R \to S$ は切断
$\sigma : S \to R$ をもつと仮定してよい（$S$ へ基底変換する通常の方法による）。
$I = \Ker(S \to R)$ とおく。
\item $R$ 上でもう少し局所化すると、$I/I^2$ が自由 $R$ 加群であり、
$I$ に関する $S$ の完備化 $S^\wedge$ が $R[[t_1, \ldots, t_n]]$ と
同型であると仮定してよい。Morphisms, Lemma
\ref{morphisms-lemma-section-smooth-morphism} を参照せよ。ここでは
$R \to S$ が滑らかであることを用いている。
\item $S$ が射影的 $R$ 加群であることを証明するには、$S$ が $R$ 加群として
平坦、可算生成、かつ Mittag--Leffler であることを証明すれば十分である。
Algebra, Lemma \ref{algebra-lemma-countgen-projective} を参照せよ。
最初の二つの性質は明らかである。したがって $S$ が $R$ 加群として
Mittag--Leffler であることを証明すれば十分である。Algebra, Lemma
\ref{algebra-lemma-power-series-ML} により、加群 $R[[t_1, \ldots, t_n]]$ は
$R$ 上 Mittag--Leffler である。よって Algebra, Lemma
\ref{algebra-lemma-pure-submodule-ML} により、
% P05-EN-CH38-0048: uniquely determined spurious-article repair; see sidecar.
$S \to S^\wedge$ が $R$ 加群の写像として普遍単射であることを示せば十分である。
\item Lemma \ref{flat-lemma-base-change-universally-flat} を適用すると、
$S \to S^\wedge$ は $R$-普遍単射であることが分かる。実際、
$R \to S$ は幾何学的整なファイバーをもつので、任意のファイバー環の
任意の随伴点はそのファイバー環の総称点にほかならず、これは
$\Spec(S^\wedge) \to \Spec(S)$ の像に含まれる。
\end{enumerate}
今概説した証明と、Lemma
\ref{flat-lemma-fibres-irreducible-flat-projective-nonnoetherian} の証明へ至る
議論の展開との間には類似がある。
% P05-EN-CH38-0049: uniquely determined punctuation repair; see sidecar.
いずれにおいても完備化が本質的な役割を果たし、またいずれにおいても
幾何学的整なファイバーをもつという仮定により、$S$ から $S$ の完備化への写像が
$R$-普遍単射であることが保証される。
\end{remark}













\section{有限型平坦加群 I}
\label{flat-section-finite-type-flat-I}

\noindent
有限表示な環準同型 $R \to S$ と有限 $S$-加群 $N$ が与えられたとき、
場合によっては $N$ が $R$ 上平坦であることから $N$ が有限表示であることが従う。
本節では、この主張が「各点で」成り立つことを証明する。
なお、命題 \ref{flat-proposition-finite-type-flat-at-point} の第1の証明は
節 \ref{flat-section-local-structure-module} の幾何学的結果を用いるが、
完全デヴィサージュの存在は用いない。

\begin{lemma}
\label{flat-lemma-induction-step}
$(R, \mathfrak m)$ を局所環とする。$R \to S$ を、幾何学的整ファイバーをもつ
有限表示平坦環準同型とする。$\mathfrak p = \mathfrak mS$ と書く。
$\mathfrak q \subset S$ を $\mathfrak m$ の上にある素イデアルとし、
$N$ を有限 $S$-加群とする。このとき、ある $r \geq 0$ と $S$-加群準同型
$$
\alpha : S^{\oplus r} \longrightarrow N
$$
が存在し、
$\alpha : \kappa(\mathfrak p)^{\oplus r} \to N \otimes_S \kappa(\mathfrak p)$
は同型である。このような任意の $\alpha$ について、次は同値である。
\begin{enumerate}
\item $N_{\mathfrak q}$ は $R$ 上平坦である。
\item $\alpha$ は $R$-普遍単射であり、
$\Coker(\alpha)_{\mathfrak q}$ は $R$ 上平坦である。
\item $\alpha$ は単射であり、
$\Coker(\alpha)_{\mathfrak q}$ は $R$ 上平坦である。
\item $\alpha_{\mathfrak p}$ は同型であり、
$\Coker(\alpha)_{\mathfrak q}$ は $R$ 上平坦である。
\item $\alpha_{\mathfrak q}$ は単射であり、
$\Coker(\alpha)_{\mathfrak q}$ は $R$ 上平坦である。
\end{enumerate}
\end{lemma}

\begin{proof}
$\alpha$ を得るため、
$r = \dim_{\kappa(\mathfrak p)} N \otimes_S \kappa(\mathfrak p)$ とおき、
$N \otimes_S \kappa(\mathfrak p)$ の基底をなす
$x_1, \ldots, x_r \in N$ を選ぶ。
$\alpha(s_1, \ldots, s_r) = \sum s_i x_i$ と定めればよい。これで存在が示された。

\medskip\noindent
$\alpha$ を一つ固定する。最も重要な含意は (1) $\Rightarrow$ (2) なので、
これを最初に証明する。(1) を仮定する。
$S/\mathfrak mS$ は分数体 $\kappa(\mathfrak p)$ をもつ整域であるから、
$(S/\mathfrak mS)^{\oplus r} \to
N_{\mathfrak p}/\mathfrak mN_{\mathfrak p} = N \otimes_S \kappa(\mathfrak p)$
は単射である。したがって、
% P05-EN-CH38-0129: uniquely determined sentence-boundary repair; see sidecar.
Lemmas \ref{flat-lemma-universally-injective-local} and
\ref{flat-lemma-fibres-irreducible-flat-projective-nonnoetherian} により、
準同型 $S^{\oplus r} \to N_{\mathfrak p}$ は $R$-普遍単射である。
よって $S^{\oplus r} \to N$ も $R$-普遍単射である（Algebra,
Lemma \ref{algebra-lemma-universally-injective-permanence} 参照）。
さらに、その局所化 $\alpha_{\mathfrak q}$ も $R$-普遍単射である
（Algebra, Lemma \ref{algebra-lemma-universally-injective-localize} 参照）。
Algebra, Lemma \ref{algebra-lemma-ui-flat-domain} により、
$\Coker(\alpha)_{\mathfrak q}$ は $R$ 上平坦であると結論できる。

\medskip\noindent
(2) $\Rightarrow$ (3) は明らかである。(3) が成り立つなら、
$\alpha_{\mathfrak p}$ は単射加群準同型の局所化なので単射である。
中山の補題（Algebra, Lemma \ref{algebra-lemma-NAK}）により、
$\alpha_{\mathfrak p}$ は全射でもある。よって (3) $\Rightarrow$ (4) である。
(4) が成り立つなら $\alpha_{\mathfrak p}$ は同型であり、
$S_{\mathfrak q} \to S_{\mathfrak p}$ が単射なので $\alpha$ は単射である。
実際、Lemmas \ref{flat-lemma-invert-universally-injective} and
\ref{flat-lemma-fibres-irreducible-flat-projective-nonnoetherian} を合わせると、
$S \setminus \mathfrak p$ の元は $S$ 上の非零因子である。
したがって (4) $\Rightarrow$ (5) である。最後に (5) が成り立つなら、
$N_{\mathfrak q}$ は平坦加群の拡大として $R$ 上平坦である
（Algebra, Lemma \ref{algebra-lemma-flat-ses} 参照）。
ゆえに (5) $\Rightarrow$ (1) であり、証明が完了した。
\end{proof}

\begin{lemma}
\label{flat-lemma-complete-devissage-flat-finite-type-module}
$(R, \mathfrak m)$ を局所環とする。
$R \to S$ を有限表示な環準同型とし、$N$ を有限 $S$-加群とする。
$\mathfrak q$ を $\mathfrak m$ の上にある $S$ の素イデアルとする。
$N_{\mathfrak q}$ は $R$ 上平坦であり、かつ $\mathfrak q$ における
$N/S/R$ の完全デヴィサージュが存在すると仮定する。
このとき $N$ は有限表示 $S$-加群であり、$R$-加群として自由であり、
同型
$$
N \cong B_1^{\oplus r_1} \oplus \ldots \oplus B_n^{\oplus r_n}
$$
が $R$-加群の同型として存在する。ここで各 $B_i$ は幾何学的既約ファイバーをもつ
滑らかな $R$-代数である。
\end{lemma}

\begin{proof}
与えられた完全デヴィサージュを
$(A_i, B_i, M_i, \alpha_i, \mathfrak q_i)_{i = 1, \ldots, n}$ とする。
$n$ に関する帰納法で補題を証明する。Remark \ref{flat-remark-finite-presentation}
によれば、$N$ が有限表示 $S$-加群であることと、$M_1$ が有限表示
% P05-EN-CH38-0130: uniquely determined article repair; see sidecar.
$B_1$-加群であることは同値である。また、$R$-加群として
$N_{\mathfrak q} \cong (M_1)_{\mathfrak q_1}$ である。実際、
(a) $N_{\mathfrak q} \cong (M_1)_{\mathfrak q'_1}$ であり、ここで
$\mathfrak q'_1$ は $\mathfrak q_1$ の上にある $A_1$ の一意な素イデアルである。
(b) Algebra, Lemma \ref{algebra-lemma-unique-prime-over-localize-below} により
$(A_1)_{\mathfrak q'_1} = (A_1)_{\mathfrak q_1}$ である。
したがって (c) $(M_1)_{\mathfrak q'_1} \cong (M_1)_{\mathfrak q_1}$ である。
ゆえに $(M_1)_{\mathfrak q_1}$ は平坦 $R$-加群である。したがって補題を
証明するために $(S, N)$ を $(B_1, M_1)$ で置き換えてよい。
Lemma \ref{flat-lemma-induction-step} により、準同型
$\alpha_1 : B_1^{\oplus r_1} \to M_1$ は $R$-普遍単射であり、
% P05-EN-CH38-0131: source-faithful q/q_1 mismatch; see sidecar.
$\Coker(\alpha_1)_{\mathfrak q}$ は $R$ 上平坦である。
さらに $(A_i, B_i, M_i, \alpha_i, \mathfrak q_i)_{i = 2, \ldots, n}$ は、
$\mathfrak q_1$ における $\Coker(\alpha_1)/B_1/R$ の完全デヴィサージュである。
したがって帰納法の仮定により、$\Coker(\alpha_1)$ は有限表示
$B_1$-加群であり、$R$-加群として自由であり、補題に述べた形の分解をもつ。
Algebra, Lemma \ref{algebra-lemma-extension} によって、これは $M_1$ が
有限表示 $B_1$-加群であることを含意する。さらに、$R$-加群として
$M_1 \cong B_1^{\oplus r_1} \oplus \Coker(\alpha_1)$
であり、したがって補題に述べた形の分解を得る。最後に、
Lemma \ref{flat-lemma-fibres-irreducible-flat-projective-nonnoetherian} により
$B_1$ は射影 $R$-加群であり、Algebra,
Theorem \ref{algebra-theorem-projective-free-over-local-ring} により
$R$-加群として自由である。これで証明が完了した。
\end{proof}

\begin{proposition}
\label{flat-proposition-finite-type-flat-at-point}
$f : X \to S$ をスキームの射とし、$\mathcal{F}$ を $X$ 上の準連接層とする。
$x \in X$ の像を $s \in S$ とする。次を仮定する。
\begin{enumerate}
\item $f$ は局所有限表示である。
\item $\mathcal{F}$ は有限型である。
\item $\mathcal{F}$ は $S$ 上 $x$ で平坦である。
\end{enumerate}
このとき、基本エタール近傍 $(S', s') \to (S, s)$ と開部分スキーム
$$
V \subset X \times_S \Spec(\mathcal{O}_{S', s'})
$$
が存在する。これは $x$ に写る
$X \times_S \Spec(\mathcal{O}_{S', s'})$ の一意な点を含み、
$\mathcal{F}$ の $V$ への引き戻しは有限表示 $\mathcal{O}_V$-加群であり、
$\mathcal{O}_{S', s'}$ 上平坦である。
\end{proposition}

\begin{proof}[第1の証明]
この証明は長いが、完全デヴィサージュの存在を用いない。
問題は $x$ と $s$ の近傍で局所的なので、$X$ と $S$ はアフィンであると
仮定してよい。証明の途中で有限回、
% P05-EN-CH38-0132: uniquely determined word-order repair; see sidecar.
$S$ を $(S, s)$ の基本エタール近傍で置き換える。このような置換の後に、
$x$ を含む開集合 $V \subset X \times_S \Spec(\mathcal{O}_{S, s})$ で、
$\mathcal{F}|_V$ が $S$ 上平坦かつ有限表示となるものを見つけることが目標である。
もちろん証明の任意の段階で $S$ を $\Spec(\mathcal{O}_{S, s})$ に
置き換えてもよい。すなわち $S$ は局所スキームであると仮定してよい。
整数 $n = \dim_x(\text{Supp}(\mathcal{F}_s))$ に関する帰納法で命題を証明する。

\medskip\noindent
Lemma \ref{flat-lemma-elementary-devissage} に従って、次を選べる。
\begin{enumerate}
\item 基本エタール近傍 $g : (X', x') \to (X, x)$ および
$e : (S', s') \to (S, s)$。
\item 可換図式
$$
\xymatrix{
X \ar[dd]_f & X' \ar[dd] \ar[l]^g & Z' \ar[l]^i \ar[d]^\pi \\
& & Y' \ar[d]^h \\
S & S' \ar[l]_e & S' \ar@{=}[l]
}
$$
\item $i(z') = x'$、$y' = \pi(z')$、$h(y') = s'$ を満たす点 $z' \in Z'$。
\item 有限型準連接 $\mathcal{O}_{Z'}$-加群 $\mathcal{G}$。
\end{enumerate}
証明の第1段落の注意に従い、$S$ を $\Spec(\mathcal{O}_{S', s'})$ で
置き換えることにする。次の図式を考える。
$$
\xymatrix{
X_{\mathcal{O}_{S', s'}} \ar[ddr]_f &
X'_{\mathcal{O}_{S', s'}} \ar[dd] \ar[l]^g &
Z'_{\mathcal{O}_{S', s'}} \ar[l]^i \ar[d]^\pi \\
& & Y'_{\mathcal{O}_{S', s'}} \ar[dl]^h \\
& \Spec(\mathcal{O}_{S', s'})
}
$$
ここでは $S'$ 上のスキーム $X', Z', Y'$ を
$\Spec(\mathcal{O}_{S', s'}) \to S'$ に沿って、また $S$ 上のスキーム $X$ を
$\Spec(\mathcal{O}_{S', s'}) \to S$ に沿って基底変換した。
Lemma \ref{flat-lemma-etale-at-point} により、$g$ はなおエタールである。
$X$ を $X_{\mathcal{O}_{S', s'}}$ で、$X'$ を
$X'_{\mathcal{O}_{S', s'}}$ で、$Z'$ を $Z'_{\mathcal{O}_{S', s'}}$ で、
$Y'$ を $Y'_{\mathcal{O}_{S', s'}}$ で置き換えると、
Lemma \ref{flat-lemma-elementary-devissage} の図式をもち、さらに
$S = S'$ は閉点 $s$ をもつ局所スキームであると仮定できる。
Lemmas \ref{flat-lemma-devissage-finite-presentation} and
\ref{flat-lemma-devissage-flat} により、$Y' \to S$、層 $\pi_*\mathcal{G}$、点 $y'$
に対する結果から、$X \to S$、$\mathcal{F}$、$x$ に対する結果が従う。
したがって $S$ は局所的であり、$X \to S$ は次元 $n$ の幾何学的既約
ファイバーをもつアフィンスキーム間の滑らかな射であると仮定してよい。

\medskip\noindent
帰納法の基底である $n = 0$ の場合を考える。
$X \to S$ は次元 $0$ の幾何学的既約ファイバーをもつ滑らかな射なので、
Descent, Lemma \ref{descent-lemma-universally-injective-etale-open-immersion}
により $X \to S$ は開埋め込みである。$S$ は局所的であり、その閉点が
$X \to S$ の像に入るので、$X = S$ と結論できる。したがって
$\mathcal{F}$ は有限平坦 $\mathcal{O}_{S, s}$-加群に対応する。
この場合、Algebra, Lemma \ref{algebra-lemma-finite-flat-local} により
$\mathcal{F}$ は実際に有限自由であり、結果が従う。

\medskip\noindent
次に帰納段階を扱う。閉ファイバー内の支持の次元が $< n$ なら結果が
成り立つと仮定する。$S = \Spec(A)$、$X = \Spec(B)$ と書き、ある
$B$-加群 $N$ に対して $\mathcal{F} = \widetilde{N}$ と書く。
$A$ は局所環であり、その極大イデアルを $\mathfrak m$ とする。
$X \to S$ は幾何学的既約ファイバーをもつので、
$\mathfrak p = \mathfrak mB$ は $\mathfrak m$ の上にある一意な極小素イデアルである。
最後に、$\mathfrak q \subset B$ を $x$ に対応する素イデアルとする。
Lemma \ref{flat-lemma-induction-step} により、準同型
$$
\alpha : B^{\oplus r} \to N
$$
を選べ、$\kappa(\mathfrak p)^{\oplus r} \to N \otimes_B \kappa(\mathfrak p)$
は同型である。さらに $N_{\mathfrak q}$ は $A$ 上平坦なので、同じ補題から
$\alpha$ は単射であり、$\Coker(\alpha)_{\mathfrak q}$ は $A$ 上平坦である。
$Q = \Coker(\alpha)$ とおく。$Q/\mathfrak mQ$ の支持は
$\mathfrak p$ を含まない。したがって確かに
$\dim_{\mathfrak q}(\text{Supp}(Q/\mathfrak mQ)) < n$ である。
$Q$ について判明したことを合わせると、帰納法の仮定を $Q$ に適用できる。
ゆえに、結論が $B \otimes_A A'$ 上の $Q \otimes_A A'$ について成り立つような
% P05-EN-CH38-0133: source-faithful s/s' pointed-neighbourhood mismatch; see sidecar.
基本エタール射 $(S', s) \to (S, s)$ が存在する。ここで
$A' = \mathcal{O}_{S', s'}$ である。$A$ を $A'$ で置き換えると、完全列
$$
0 \to B^{\oplus r} \to N \to Q \to 0
$$
を得る（ここでは、上で述べた $\alpha$ の単射性を用いた）。これは有限
$B$-加群の完全列であり、さらに $g \in B$, $g \not \in \mathfrak q$ で、
$Q_g$ が $B_g$ 上有限表示かつ $A$ 上平坦となる元も得られる。
局所化は完全なので、
$$
0 \to B_g^{\oplus r} \to N_g \to Q_g \to 0
$$
も完全である。$B_g$ と $Q_g$ は $A$ 上平坦なので、Algebra,
Lemma \ref{algebra-lemma-flat-ses} により $N_g$ は $A$ 上平坦である。
また $B_g$ と $Q_g$ は $B_g$ 上有限表示なので、Algebra,
Lemma \ref{algebra-lemma-extension} により $N_g$ も同様である。
\end{proof}

\begin{proof}[第2の証明]
Proposition \ref{flat-proposition-existence-complete-at-x} を適用して、
点付きスキームの可換図式
$$
\xymatrix{
(X, x) \ar[d] & (X', x') \ar[l]^g \ar[d] \\
(S, s) & (S', s') \ar[l]
}
$$
を得る。ここで横の矢印は基本エタール近傍であり、
% P05-EN-CH38-0134: source-faithful x/x' complete-devissage point mismatch; see sidecar.
$g^*\mathcal{F}/X'/S'$ は $x$ において完全デヴィサージュをもつ。
（特に $S'$ と $X'$ はアフィンである。）Morphisms,
Lemma \ref{morphisms-lemma-flat-permanence} により、$g^*\mathcal{F}$ は
$S$ 上 $x'$ で平坦であり、Lemma \ref{flat-lemma-etale-flat-up-down} により
$S'$ 上 $x'$ で平坦である。Remark \ref{flat-remark-same-notion} を通じて、
$$
\Gamma(X', g^*\mathcal{F})/
\Gamma(X', \mathcal{O}_{X'})/
\Gamma(S', \mathcal{O}_{S'})
$$
は $x'$ に対応する $\Gamma(X', \mathcal{O}_{X'})$ の素イデアルにおいて
完全デヴィサージュをもつ。この完全デヴィサージュを、$s'$ に対応する
$\Gamma(S', \mathcal{O}_{S'})$ の素イデアルにおける局所環
$\mathcal{O}_{S', s'}$ へ基底変換できる。したがって
Lemma \ref{flat-lemma-complete-devissage-flat-finite-type-module} により、
% P05-EN-CH38-0135: source-faithful undefined F' notation; see sidecar.
$$
\Gamma(X', \mathcal{F}')
\otimes_{\Gamma(S', \mathcal{O}_{S'})}
\mathcal{O}_{S', s'}
$$
は $\mathcal{O}_{S', s'}$ 上平坦であり、
$\Gamma(X', \mathcal{O}_{X'})
\otimes_{\Gamma(S', \mathcal{O}_{S'})} \mathcal{O}_{S', s'}$ 上有限表示である。
言い換えれば、$\mathcal{F}$ の
$X' \times_{S'} \Spec(\mathcal{O}_{S', s'})$ への制限は有限表示であり、
$\mathcal{O}_{S', s'}$ 上平坦である。射
$X' \times_{S'} \Spec(\mathcal{O}_{S', s'})
\to X \times_S \Spec(\mathcal{O}_{S', s'})$
はエタールなので（Lemma \ref{flat-lemma-etale-at-point}）、その像
$V \subset X \times_S \Spec(\mathcal{O}_{S', s'})$ は開部分スキームである。
さらにエタール降下により、$\mathcal{F}$ の $V$ への制限は有限表示であり、
$\mathcal{O}_{S', s'}$ 上平坦である。ここでは Morphisms,
Lemma \ref{morphisms-lemma-etale-open}、Descent,
Lemma \ref{descent-lemma-finite-presentation-descends}、および Morphisms,
Lemma \ref{morphisms-lemma-flat-permanence} を用いた。
\end{proof}

\begin{lemma}
\label{flat-lemma-open-in-fibre-where-flat}
$f : X \to S$ を局所有限型なスキームの射とし、$\mathcal{F}$ を有限型
準連接 $\mathcal{O}_X$-加群とする。$s \in S$ とする。このとき集合
$$
\{x \in X_s \mid \mathcal{F}\text{ は }S\text{ 上 }x\text{ で平坦}\}
$$
はファイバー $X_s$ において開である。
\end{lemma}

\begin{proof}
% P05-EN-CH38-0136: source-faithful undefined U notation; see sidecar.
$x \in U$ と仮定する。Proposition \ref{flat-proposition-finite-type-flat-at-point}
における基本エタール近傍 $(S', s') \to (S, s)$ と開集合
$V \subset X \times_S \Spec(\mathcal{O}_{S', s'})$ を選ぶ。
$\kappa(s) = \kappa(s')$ なので $X_{s'} = X_s$ であることに注意する。
$x' \in V \cap X_{s'}$ なら、$\mathcal{F}$ の $X \times_S S'$ への
引き戻しは $S'$ 上 $x'$ で平坦である。したがって Morphisms,
Lemma \ref{morphisms-lemma-flat-permanence} により、$\mathcal{F}$ は
$S$ 上 $x'$ で平坦である。言い換えれば、$X_s \cap V \subset U$ は
$U$ における $x$ の開近傍である。
\end{proof}

\begin{lemma}
\label{flat-lemma-finite-type-flat-at-point}
$f : X \to S$ をスキームの射とし、$\mathcal{F}$ を $X$ 上の準連接層とする。
$x \in X$ の像を $s \in S$ とする。次を仮定する。
\begin{enumerate}
\item $f$ は局所有限型である。
\item $\mathcal{F}$ は有限型である。
\item $\mathcal{F}$ は $S$ 上 $x$ で平坦である。
\end{enumerate}
このとき、基本エタール近傍 $(S', s') \to (S, s)$ と開部分スキーム
$$
V \subset X \times_S \Spec(\mathcal{O}_{S', s'})
$$
が存在する。これは $x$ に写る
$X \times_S \Spec(\mathcal{O}_{S', s'})$ の一意な点を含み、
$\mathcal{F}$ の $V$ への引き戻しは $\mathcal{O}_{S', s'}$ 上平坦である。
\end{lemma}

\begin{proof}
% P05-EN-CH38-0146: source-faithful missing locally qualifier; see sidecar.
（これと Proposition \ref{flat-proposition-finite-type-flat-at-point} の唯一の違いは、
$f$ が有限表示であると仮定しないことである。）問題は $X$ と $S$ 上局所的なので、
$X$ と $S$ はアフィンであると仮定してよい。$X = \Spec(B)$、
$S = \Spec(A)$ と書き、$B = A[x_1, \ldots, x_n]/I$ と書く。
言い換えれば閉埋め込み $i : X \to \mathbf{A}^n_S$ を得る。
$t = i(x) \in \mathbf{A}^n_S$ とおく。
$\mathbf{A}^n_S \to S$、層 $i_*\mathcal{F}$、点 $t$ に
Proposition \ref{flat-proposition-finite-type-flat-at-point} を適用すると、
基本エタール近傍 $(S', s') \to (S, s)$ と開部分スキーム
$$
W \subset \mathbf{A}^n_{\mathcal{O}_{S', s'}}
$$
を得て、$i_*\mathcal{F}$ の $W$ への引き戻しは
$\mathcal{O}_{S', s'}$ 上平坦である。これは、
$V := W \cap \big(X \times_S \Spec(\mathcal{O}_{S', s'})\big)$
が求める開部分スキームであることを意味する。
\end{proof}

\begin{lemma}
\label{flat-lemma-finite-type-flat-along-fibre}
$f : X \to S$ をスキームの射とし、$\mathcal{F}$ を $X$ 上の準連接層とする。
$s \in S$ とする。次を仮定する。
\begin{enumerate}
\item $f$ は有限表示である。
\item $\mathcal{F}$ は有限型である。
\item $\mathcal{F}$ はファイバー $X_s$ の各点で $S$ 上平坦である。
\end{enumerate}
このとき、基本エタール近傍 $(S', s') \to (S, s)$ と開部分スキーム
$$
V \subset X \times_S \Spec(\mathcal{O}_{S', s'})
$$
が存在する。これはファイバー $X_s = X \times_S s'$ を含み、
$\mathcal{F}$ の $V$ への引き戻しは有限表示 $\mathcal{O}_V$-加群であり、
$\mathcal{O}_{S', s'}$ 上平坦である。
\end{lemma}

\begin{proof}
各点 $x \in X_s$ に対し、
Proposition \ref{flat-proposition-finite-type-flat-at-point} を用いると、
基本エタール近傍 $(S_x, s_x) \to (S, s)$ と開集合
$V_x \subset X \times_S \Spec(\mathcal{O}_{S_x, s_x})$ を得る。
ここで $x \in X_s = X \times_S s_x$ は $V_x$ に含まれ、$\mathcal{F}$ の
$V_x$ への引き戻しは有限表示 $\mathcal{O}_{V_x}$-加群であり、
$\mathcal{O}_{S_x, s_x}$ 上平坦である。特にファイバー $(V_x)_{s_x}$ を
$X_s$ における $x$ の開近傍とみなせる。$X_s$ は準コンパクトなので、
$X_s$ が $(V_{x_i})_{s_{x_i}}$ の合併となるような有限個の点
$x_1, \ldots, x_n \in X_s$ を選べる。More on Morphisms,
Lemma \ref{more-morphisms-lemma-elementary-etale-neighbourhoods} により、
各近傍 $(S_{x_i}, s_{x_i})$ を支配する基本エタール近傍
$(S' , s') \to (S, s)$ を選ぶ。
% P05-EN-CH38-0137: uniquely determined singular/plural agreement repair; see sidecar.
$V = \bigcup V_i$ とおく。ここで $V_i$ は、射
$$
X \times_S \Spec(\mathcal{O}_{S', s'})
\longrightarrow
X \times_S \Spec(\mathcal{O}_{S_{x_i}, s_{x_i}})
$$
による開集合 $V_{x_i}$ の逆像である。構成より $V$ は $X_s$ を含み、
$\mathcal{F}$ の $V$ への引き戻しは有限表示 $\mathcal{O}_V$-加群であり、
$\mathcal{O}_{S', s'}$ 上平坦である。
\end{proof}

\begin{lemma}
\label{flat-lemma-finite-type-flat-along-fibre-variant}
$f : X \to S$ をスキームの射とし、$\mathcal{F}$ を $X$ 上の準連接層とする。
$s \in S$ とする。次を仮定する。
\begin{enumerate}
\item $f$ は有限型である。
\item $\mathcal{F}$ は有限型である。
\item $\mathcal{F}$ はファイバー $X_s$ の各点で $S$ 上平坦である。
\end{enumerate}
このとき、基本エタール近傍 $(S', s') \to (S, s)$ と開部分スキーム
$$
V \subset X \times_S \Spec(\mathcal{O}_{S', s'})
$$
が存在する。これはファイバー $X_s = X \times_S s'$ を含み、
$\mathcal{F}$ の $V$ への引き戻しは $\mathcal{O}_{S', s'}$ 上平坦である。
\end{lemma}

\begin{proof}
（これと Lemma \ref{flat-lemma-finite-type-flat-along-fibre} の唯一の違いは、
$f$ が有限表示であると仮定しないことである。）各点 $x \in X_s$ に対し、
Lemma \ref{flat-lemma-finite-type-flat-at-point} を用いると、基本エタール近傍
$(S_x, s_x) \to (S, s)$ と開集合
$V_x \subset X \times_S \Spec(\mathcal{O}_{S_x, s_x})$ を得る。
ここで $x \in X_s = X \times_S s_x$ は $V_x$ に含まれ、$\mathcal{F}$ の
$V_x$ への引き戻しは $\mathcal{O}_{S_x, s_x}$ 上平坦である。
特にファイバー $(V_x)_{s_x}$ を $X_s$ における $x$ の開近傍とみなせる。
$X_s$ は準コンパクトなので、$X_s$ が $(V_{x_i})_{s_{x_i}}$ の合併となる
ような有限個の点 $x_1, \ldots, x_n \in X_s$ を選べる。
More on Morphisms,
Lemma \ref{more-morphisms-lemma-elementary-etale-neighbourhoods} により、
各近傍 $(S_{x_i}, s_{x_i})$ を支配する基本エタール近傍
$(S' , s') \to (S, s)$ を選ぶ。
% P05-EN-CH38-0138: uniquely determined singular/plural agreement repair; see sidecar.
$V = \bigcup V_i$ とおく。ここで $V_i$ は、射
$$
X \times_S \Spec(\mathcal{O}_{S', s'})
\longrightarrow
X \times_S \Spec(\mathcal{O}_{S_{x_i}, s_{x_i}})
$$
による開集合 $V_{x_i}$ の逆像である。構成より $V$ は $X_s$ を含み、
$\mathcal{F}$ の $V$ への引き戻しは $\mathcal{O}_{S', s'}$ 上平坦である。
\end{proof}

\begin{lemma}
\label{flat-lemma-finite-type-flat-at-point-X}
$S$ をスキームとし、$X$ を $S$ 上局所有限型とする。
$x \in X$ の像を $s \in S$ とする。$X$ が $S$ 上 $x$ で平坦なら、
基本エタール近傍 $(S', s') \to (S, s)$ と開部分スキーム
$$
V \subset X \times_S \Spec(\mathcal{O}_{S', s'})
$$
が存在する。これは $x$ に写る
$X \times_S \Spec(\mathcal{O}_{S', s'})$ の一意な点を含み、
$V \to \Spec(\mathcal{O}_{S', s'})$ は平坦かつ有限表示である。
\end{lemma}

\begin{proof}
問題は $X$ と $S$ 上局所的なので、$X$ と $S$ はアフィンであると仮定してよい。
$X = \Spec(B)$、$S = \Spec(A)$ と書き、$B = A[x_1, \ldots, x_n]/I$ と書く。
言い換えれば閉埋め込み $i : X \to \mathbf{A}^n_S$ を得る。
$t = i(x) \in \mathbf{A}^n_S$ とおく。
$\mathbf{A}^n_S \to S$、層 $\mathcal{F} = i_*\mathcal{O}_X$、点 $t$ に
Proposition \ref{flat-proposition-finite-type-flat-at-point} を適用すると、
基本エタール近傍 $(S', s') \to (S, s)$ と開部分スキーム
$$
W \subset \mathbf{A}^n_{\mathcal{O}_{S', s'}}
$$
を得て、$i_*\mathcal{O}_X$ の引き戻しは平坦かつ有限表示である。これは、
$V := W \cap \big(X \times_S \Spec(\mathcal{O}_{S', s'})\big)$
が求める開部分スキームであることを意味する。
\end{proof}

\begin{lemma}
\label{flat-lemma-finite-type-flat-at-point-local}
$f : X \to S$ を局所有限表示な射とし、$\mathcal{F}$ を有限型準連接
$\mathcal{O}_X$-加群とする。$x \in X$ であり、$\mathcal{F}$ が $S$ 上
$x$ で平坦なら、$\mathcal{F}_x$ は有限表示 $\mathcal{O}_{X, x}$-加群である。
\end{lemma}

\begin{proof}
$s = f(x)$ とおく。Proposition \ref{flat-proposition-finite-type-flat-at-point} により、
基本エタール近傍 $(S', s') \to (S, s)$ が存在し、$\mathcal{F}$ の
$X \times_S \Spec(\mathcal{O}_{S', s'})$ への引き戻しは、$x$ に対応する点
$x' \in X_{s'} = X_s$ の近傍で有限表示である。環準同型
$$
\mathcal{O}_{X, x} \longrightarrow
\mathcal{O}_{X \times_S \Spec(\mathcal{O}_{S', s'}), x'}
=
\mathcal{O}_{X \times_S S', x'}
$$
はエタール環準同型の局所化として平坦かつ局所的である。したがって降下により、
$\mathcal{F}_x$ は $\mathcal{O}_{X, x}$ 上有限表示である。Algebra,
Lemma \ref{algebra-lemma-descend-properties-modules} を参照せよ。また、平坦な
局所環準同型が忠実平坦であることについては Algebra,
Lemma \ref{algebra-lemma-local-flat-ff} を参照せよ。
\end{proof}

\begin{lemma}
\label{flat-lemma-finite-type-flat-at-point-local-X}
$f : X \to S$ を局所有限型な射とする。$x \in X$ の像を $s \in S$ とする。
$f$ が $S$ 上 $x$ で平坦なら、$\mathcal{O}_{X, x}$ は
$\mathcal{O}_{S, s}$ 上本質的有限表示である。
\end{lemma}

\begin{proof}
$X$ と $S$ はアフィンであると仮定してよい。$X = \Spec(B)$、
$S = \Spec(A)$ と書き、$B = A[x_1, \ldots, x_n]/I$ と書く。
言い換えれば閉埋め込み $i : X \to \mathbf{A}^n_S$ を得る。
$t = i(x) \in \mathbf{A}^n_S$ とおく。
$\mathbf{A}^n_S \to S$、層 $\mathcal{F} = i_*\mathcal{O}_X$、点 $t$ に
Lemma \ref{flat-lemma-finite-type-flat-at-point-local} を適用すると、
$\mathcal{O}_{X, x}$ は $\mathcal{O}_{\mathbf{A}^n_S, t}$ 上有限表示である。
これは求める主張を含意する。
\end{proof}







\section{開集合からの性質の延長}
\label{flat-section-extending-properties}

\noindent
本節では次の形の結果をいくつか集める。$f : X \to S$ がスキームの平坦射で、
$f$ が $S$ の稠密開集合上である性質を満たすなら、$f$ は $S$ 全体上で同じ性質を満たす。

\begin{lemma}
\label{flat-lemma-flat-finite-type-finitely-presented-over-dense-open}
\begin{slogan}
（良い）開集合上の有限表示加群の、$S$-平坦かつ有限型な延長は、
$X$ 上でも有限表示である。
\end{slogan}
$f : X \to S$ をスキームの射とし、$\mathcal{F}$ を準連接
$\mathcal{O}_X$-加群とする。$U \subset S$ を開集合とする。次を仮定する。
\begin{enumerate}
\item $f$ は局所有限表示である。
\item $\mathcal{F}$ は有限型であり、$S$ 上平坦である。
\item $U \subset S$ は逆コンパクトかつスキーム論的に稠密である。
\item $\mathcal{F}|_{f^{-1}U}$ は有限表示である。
\end{enumerate}
このとき $\mathcal{F}$ は有限表示である。
\end{lemma}

\begin{proof}
問題は $X$ と $S$ 上局所的なので、$X$ と $S$ はアフィンであると仮定してよい。
$S = \Spec(A)$、$X = \Spec(B)$ と書く。$\mathcal{F}$ が $N$ に付随する
準連接層となるような有限 $B$-加群 $N$ をとる。Algebra,
Lemma \ref{algebra-lemma-qc-open} により、ある $f_i \in A$ に対して
$U = D(f_1) \cup \ldots \cup D(f_n)$ である。$U$ はスキーム論的に稠密なので、
準同型 $A \to A_{f_1} \times \ldots \times A_{f_n}$ は単射である。
$s \in S$ に写る $x \in X$ に対応する $\mathfrak p \subset A$ の上にある
素イデアル $\mathfrak q \subset B$ を選ぶ。
Lemma \ref{flat-lemma-finite-type-flat-at-point-local} により、$N_\mathfrak q$ は
$B_\mathfrak q$ 上有限表示である。$B$-加群の全射
$\varphi : B^{\oplus m} \to N$ を選ぶ。$k_1, \ldots, k_t \in \Ker(\varphi)$
を選び、$N' = B^{\oplus m}/\sum Bk_j$ とおく。標準的全射 $N' \to N$ があり、
$N$ はこのように構成される $B$-加群 $N'$ のフィルター余極限である。
したがって、(a) $N'_{f_i} \cong N_{f_i}$, $i = 1, \ldots, n$ および
(b) $N'_\mathfrak q \cong N_\mathfrak q$ を満たすように
$k_1, \ldots, k_t$ を選べる。特に $N'_\mathfrak q$ は $A$ 上平坦である。
平坦性の開性（Algebra, Theorem \ref{algebra-theorem-openness-flatness}）により、
$g \in B$, $g \not \in \mathfrak q$ で $N'_g$ が $A$ 上平坦となるものが存在する。
次の可換図式を考える。
$$
\xymatrix{
N'_g \ar[r] \ar[d] & N_g \ar[d] \\
\prod N'_{gf_i} \ar[r] & \prod N_{gf_i}
}
$$
下の矢印は $k_1, \ldots, k_t$ の選び方により同型である。
$A \to \prod A_{f_i}$ は単射であり $N'_g$ は $A$ 上平坦なので、
左の縦矢印は単射である。したがって上の横矢印は単射であり、ゆえに同型である。
これにより $N_g$ は $B_g$ 上有限表示である。最後に Algebra,
Lemma \ref{algebra-lemma-cover} を適用すればよい。
\end{proof}

\begin{lemma}
\label{flat-lemma-flat-finite-type-finitely-presented-over-dense-open-X}
$f : X \to S$ をスキームの射とし、$U \subset S$ を開集合とする。次を仮定する。
\begin{enumerate}
\item $f$ は局所有限型かつ平坦である。
\item $U \subset S$ は逆コンパクトかつスキーム論的に稠密である。
\item $f|_{f^{-1}U} : f^{-1}U \to U$ は局所有限表示である。
\end{enumerate}
% P05-EN-CH38-0139: uniquely determined spurious-preposition repair; see sidecar.
このとき $f$ は局所有限表示である。
\end{lemma}

\begin{proof}
問題は $X$ と $S$ 上局所的なので、$X$ と $S$ はアフィンであると仮定してよい。
閉埋め込み $i : X \to \mathbf{A}^n_S$ を選び、$i_*\mathcal{O}_X$ に
Lemma \ref{flat-lemma-flat-finite-type-finitely-presented-over-dense-open} を適用する。
% P05-EN-CH38-0140: source-faithful explicit proof-detail omission; see sidecar.
一部の詳細は省略する。
\end{proof}

\begin{lemma}
\label{flat-lemma-flat-finite-presentation-dimension-over-dense-open}
$f : X \to S$ を平坦かつ局所有限型なスキームの射とする。
$U \subset S$ を、$X_U \to U$ が相対次元 $\leq e$ をもつ稠密開集合とする
（Morphisms, Definition \ref{morphisms-definition-relative-dimension-d} 参照）。
さらに次のいずれかを仮定する。
\begin{enumerate}
\item $f$ は局所有限表示である。
\item $U \subset S$ は逆コンパクトである。
\end{enumerate}
このとき $f$ は相対次元 $\leq e$ をもつ。
\end{lemma}

\begin{proof}
まず (1) の場合を証明する。More on Morphisms, Lemmas
\ref{more-morphisms-lemma-flat-finite-presentation-CM-open} and
\ref{more-morphisms-lemma-flat-finite-presentation-CM-pieces} で構成・考察された
開部分スキームを $W \subset X$ とする。各ファイバーの各総称点は $W$ に
含まれるので、$W$ について結果を証明すれば十分である。
$W = \bigcup_{d \geq 0} U_d$ なので、$d > e$ に対して
$U_d = \emptyset$ を示せば十分である。$f$ は平坦かつ局所有限表示なので開であり、
$f(U_d)$ は開である（Morphisms, Lemma \ref{morphisms-lemma-fppf-open}）。
したがって $U_d$ が空でなければ、求める通り
$f(U_d) \cap U \not = \emptyset$ である。

\medskip\noindent
次に (2) の場合を証明する。$S$ をその被約化で置き換えてよい。
すると $U$ はスキーム論的に稠密である。したがって
Lemma \ref{flat-lemma-flat-finite-type-finitely-presented-over-dense-open-X} により、
$f$ は局所有限表示である。こうして (1) の場合に帰着する。
\end{proof}

\begin{lemma}
\label{flat-lemma-proper-flat-finite-over-dense-open}
$f : X \to S$ を平坦かつ固有なスキームの射とする。
$U \subset S$ を、$X_U \to U$ が有限となる稠密開集合とする。
さらに $f$ が局所有限表示であるか、または $U \subset S$ が逆コンパクトなら、
$f$ は有限である。
\end{lemma}

\begin{proof}
Lemma \ref{flat-lemma-flat-finite-presentation-dimension-over-dense-open} により、
$f$ のファイバーは次元 0 である。したがって $f$ は準有限であり
（Morphisms, Lemma \ref{morphisms-lemma-locally-quasi-finite-rel-dimension-0}）、
ゆえに有限ファイバーをもつ
（Morphisms, Lemma \ref{morphisms-lemma-quasi-finite}）。したがって
More on Morphisms, Lemma \ref{more-morphisms-lemma-characterize-finite} により
$f$ は有限である。
\end{proof}

\begin{lemma}
\label{flat-lemma-zariski}
$f : X \to S$ をスキームの射とし、$U \subset S$ を開集合とする。次を仮定する。
\begin{enumerate}
\item $f$ は分離的、局所有限型、かつ平坦である。
\item $f^{-1}(U) \to U$ は同型である。
\item $U \subset S$ は逆コンパクトかつスキーム論的に稠密である。
\end{enumerate}
このとき $f$ は開埋め込みである。
\end{lemma}

\begin{proof}
Lemma \ref{flat-lemma-flat-finite-type-finitely-presented-over-dense-open-X} により、
射 $f$ は局所有限表示である。像 $f(X) \subset S$ は開なので
（Morphisms, Lemma \ref{morphisms-lemma-fppf-open}）、$S$ を $f(X)$ で
置き換えてよい。したがって $f$ が同型であることを示せばよい。
$S$ はアフィンであると仮定してよい。また、像が $S$ である任意の
準コンパクト開集合 $X' \subset X$ が $S$ に同型に写ることを示せば十分なので、
$X$ が準コンパクトな場合に帰着できる。よって $f$ は準コンパクトであると
仮定してよい。Lemma \ref{flat-lemma-flat-finite-presentation-dimension-over-dense-open}
により、$f$ のすべてのファイバーは次元 $0$ である。したがって Morphisms,
Lemma \ref{morphisms-lemma-locally-quasi-finite-rel-dimension-0} により
$f$ は準有限である。$s \in S$ とする。$V \to T$ が有限かつ
$W_t = \emptyset$ である分解 $X \times_S T = V \amalg W$ を与える
基本エタール近傍 $g : (T, t) \to (S, s)$ を選ぶ。More on Morphisms,
Lemma \ref{more-morphisms-lemma-etale-splits-off-quasi-finite-part} を参照せよ。
与えられた射を $\pi : V \amalg W \to T$ と書く。$\pi$ は平坦かつ
局所有限表示なので、$\pi(V)$ は $T$ で開である
（Morphisms, Lemma \ref{morphisms-lemma-fppf-open}）。$T$ を縮小して、
$T = \pi(V)$ と仮定してよい。$f$ は $U$ 上同型なので、$\pi$ は
$g^{-1}U$ 上同型である。$\pi(V) = T$ なので、これは
$\pi^{-1}g^{-1}U$ が $V$ に含まれることを含意する。Morphisms,
Lemma \ref{morphisms-lemma-flat-morphism-scheme-theoretically-dense-open} により、
$\pi^{-1}g^{-1}U \subset V \amalg W$ はスキーム論的に稠密である。
したがって $W = \emptyset$ である。ゆえに $X \times_S T = V$ は $T$ 上有限である。
これは、例えば Descent, Lemma \ref{descent-lemma-descending-property-finite} により、
$S$ を $s$ の開近傍で置き換えた後に $f$ が有限であることを含意する。
すると $f$ は有限局所自由である
（Morphisms, Lemma \ref{morphisms-lemma-finite-flat}）。さらに $S$ を $s$ の
より小さい開近傍に縮めると、$f$ はある次数 $d$ の有限局所自由射となる
（Morphisms, Lemma \ref{morphisms-lemma-finite-locally-free}）。
しかし稠密開集合 $U$ 上で次数が $1$ であることから明らかなように、
$d = 1$ である。したがって $f$ は同型である。
\end{proof}





\section{有限表示平坦加群}
\label{flat-section-finitely-presented-flat}

\noindent
有限表示な環準同型 $R \to S$ と有限表示 $S$-加群 $N$ が与えられたとき、
場合によっては、少なくとも $S$ をエタール拡大で置き換えた後、
$N$ が $R$ 上平坦であることから $N$ が射影 $R$-加群であることが従う。
% P05-EN-CH38-0141: uniquely determined article repair; see sidecar.
本節では、この種の結果をいくつか集める。

\begin{lemma}
\label{flat-lemma-induction-step-fp}
$R$ を環とする。$R \to S$ を、幾何学的整ファイバーをもつ有限表示平坦環準同型とする。
$\mathfrak q \subset S$ を素イデアル $\mathfrak r \subset R$ の上にある
素イデアルとする。$\mathfrak p = \mathfrak r S$ とおく。
$N$ を有限表示 $S$-加群とする。このとき、ある $r \geq 0$ と $S$-加群準同型
$$
\alpha : S^{\oplus r} \longrightarrow N
$$
が存在し、
$\alpha : \kappa(\mathfrak p)^{\oplus r} \to N \otimes_S \kappa(\mathfrak p)$
は同型である。このような任意の $\alpha$ について、次は同値である。
\begin{enumerate}
\item $N_{\mathfrak q}$ は $R$ 上平坦である。
\item $f \in R$, $f \not \in \mathfrak r$ で、
$\alpha_f : S_f^{\oplus r} \to N_f$ が $R_f$-普遍単射となるものと、
$g \in S$, $g \not \in \mathfrak q$ で、$\Coker(\alpha)_g$ が
$R$ 上平坦となるものが存在する。
\item $\alpha_{\mathfrak r}$ は $R_{\mathfrak r}$-普遍単射であり、
% P05-EN-CH38-0142: uniquely determined missing list punctuation; see sidecar.
$\Coker(\alpha)_{\mathfrak q}$ は $R$ 上平坦である。
\item $\alpha_{\mathfrak r}$ は単射であり、
$\Coker(\alpha)_{\mathfrak q}$ は $R$ 上平坦である。
\item $\alpha_{\mathfrak p}$ は同型であり、
$\Coker(\alpha)_{\mathfrak q}$ は $R$ 上平坦である。
\item $\alpha_{\mathfrak q}$ は単射であり、
$\Coker(\alpha)_{\mathfrak q}$ は $R$ 上平坦である。
\end{enumerate}
\end{lemma}

\begin{proof}
$\alpha$ を得るため、
$r = \dim_{\kappa(\mathfrak p)} N \otimes_S \kappa(\mathfrak p)$ とおき、
$N \otimes_S \kappa(\mathfrak p)$ の基底をなす
$x_1, \ldots, x_r \in N$ を選ぶ。
$\alpha(s_1, \ldots, s_r) = \sum s_i x_i$ と定めればよい。これで存在が示された。

\medskip\noindent
$\alpha$ の選択を一つ固定する。準同型
$\alpha_{\mathfrak r} : S_{\mathfrak r}^{\oplus r} \to N_{\mathfrak r}$ に
Lemma \ref{flat-lemma-induction-step} を適用できる。したがって
(1), (3), (4), (5), (6) はすべて同値である。また (2) が (3) を
含意することも明らかなので、(1) が (2) を含意することだけを示せばよい。

\medskip\noindent
(1) を仮定する。平坦性の開性（Algebra,
Theorem \ref{algebra-theorem-openness-flatness}）により、集合
$$
U_1 = \{\mathfrak q' \subset S \mid N_{\mathfrak q'}\text{ は }R\text{ 上平坦}\}
$$
は $\Spec(S)$ で開である。仮定により $\mathfrak q$ を含み、したがって
$\mathfrak p$ も含む。$S^{\oplus r}$ と $N$ は有限表示 $S$-加群なので、集合
$$
U_2 = \{\mathfrak q' \subset S \mid
\alpha_{\mathfrak q'}\text{ は同型}\}
$$
は $\Spec(S)$ で開である（Algebra,
Lemma \ref{algebra-lemma-map-between-finitely-presented} 参照）。
(5) により、これは $\mathfrak p$ を含む。$R \to S$ は有限表示かつ平坦なので、
写像 $\Phi : \Spec(S) \to \Spec(R)$ は開である（Algebra,
Proposition \ref{algebra-proposition-fppf-open} 参照）。任意の素イデアル
$\mathfrak r' \in \Phi(U_1 \cap U_2)$ に対し、$\mathfrak r'$ の上にある素イデアル
$\mathfrak q'$ で、$N_{\mathfrak q'}$ が平坦かつ $\alpha_{\mathfrak q'}$ が同型となる
ものが存在する。これは、$\mathfrak p' = \mathfrak r' S$ とおけば
$\alpha \otimes \kappa(\mathfrak p')$ が同型であることを含意する。
したがって含意 (1) $\Rightarrow$ (3) により、$\alpha_{\mathfrak r'}$ は
$R_{\mathfrak r'}$-普遍単射である。よって
$f \in R$, $f \not \in \mathfrak r$ で $D(f) \subset \Phi(U_1 \cap U_2)$
となるものを選べば、$\alpha_f$ は $R_f$-普遍単射である
（Algebra, Lemma \ref{algebra-lemma-universally-injective-check-stalks} 参照）。
同じ議論により、任意の
$\mathfrak q' \in U_1 \cap \Phi^{-1}(\Phi(U_1 \cap U_2))$ に対し、
$\Coker(\alpha)_{\mathfrak q'}$ は $R$ 上平坦である。また
$\mathfrak q \in U_1 \cap \Phi^{-1}(\Phi(U_1 \cap U_2))$ である。
したがって $g \in S$, $g \not \in \mathfrak q$ で
$D(g) \subset U_1 \cap \Phi^{-1}(\Phi(U_1 \cap U_2))$ となるものを選べ、
証明が完了する。
\end{proof}

\begin{lemma}
\label{flat-lemma-complete-devissage-flat-finitely-presented-module}
$R \to S$ を有限表示な環準同型とする。$N$ を $R$ 上平坦な有限表示
$S$-加群とする。$\mathfrak r \subset R$ を素イデアルとする。
$\mathfrak r$ 上の $N/S/R$ の完全デヴィサージュが存在すると仮定する。
このとき $f \in R$, $f \not \in \mathfrak r$ で、
$$
N_f \cong B_1^{\oplus r_1} \oplus \ldots \oplus B_n^{\oplus r_n}
$$
が $R$-加群の同型となるものが存在する。ここで各 $B_i$ は幾何学的既約
ファイバーをもつ滑らかな $R_f$-代数である。さらに $N_f$ は射影
$R_f$-加群である。
\end{lemma}

\begin{proof}
与えられた完全デヴィサージュを
$(A_i, B_i, M_i, \alpha_i)_{i = 1, \ldots, n}$ とする。
$n$ に関する帰納法で補題を証明する。補題の主張はすべて $R$-加群としての
$N$ の構造に関する。したがって $N$ を $M_1$ で置き換え、$M_1$ を
$B_1$-加群とみなしてよい。$M_1$ が有限表示 $B_1$-加群である理由については
Remark \ref{flat-remark-finite-presentation} を参照せよ。
Lemma \ref{flat-lemma-induction-step-fp} により、ある
$f \in R$, $f \not \in \mathfrak r$ に対して $R$ を $R_f$ で置き換えた後、
準同型 $\alpha_1 : B_1^{\oplus r_1} \to M_1$ が $R$-普遍単射であると
仮定してよい。$M_1$ と $B_1^{\oplus r_1}$ は $R$ 上平坦で、かつ有限表示
$B_1$-加群なので、$\Coker(\alpha_1)$ は $R$ 上平坦
（Algebra, Lemma \ref{algebra-lemma-ui-flat-domain}）かつ有限表示
$B_1$-加群である。また
$(A_i, B_i, M_i, \alpha_i)_{i = 2, \ldots, n}$ は
$\Coker(\alpha_1)$ の完全デヴィサージュである。したがって帰納法の仮定により、
ある $f \in R$, $f \not \in \mathfrak r$ に対して $R$ を $R_f$ で
置き換えた後、$\Coker(\alpha_1)$ は補題に述べた形の分解をもち、
かつ射影的であると仮定してよい。特に
$M_1 = B_1^{\oplus r_1} \oplus \Coker(\alpha_1)$ である。
これで分解に関する主張が示された。射影性に関する主張は、
Lemma \ref{flat-lemma-fibres-irreducible-flat-projective-nonnoetherian} により
$B_1$ が射影 $R$-加群であることから従う。
\end{proof}

\begin{remark}
\label{flat-remark-complete-devissage-flat-finitely-presented-module}
Lemma \ref{flat-lemma-complete-devissage-flat-finitely-presented-module} には
次の変形がある。平坦性の条件を弱めて、$\mathfrak r$ の上にある与えられた
素イデアル $\mathfrak q$ において $N$ が平坦であることだけを仮定する一方、
デヴィサージュの条件を強めて、{\it $\mathfrak q$ における}完全デヴィサージュの
存在を仮定する。Lemma \ref{flat-lemma-complete-devissage-flat-finite-type-module} と比較せよ。
\end{remark}

\noindent
次が本節の主結果である。

\begin{proposition}
\label{flat-proposition-finite-presentation-flat-at-point}
$f : X \to S$ をスキームの射とし、$\mathcal{F}$ を $X$ 上の準連接層とする。
$x \in X$ の像を $s \in S$ とする。次を仮定する。
\begin{enumerate}
\item $f$ は局所有限表示である。
\item $\mathcal{F}$ は有限表示である。
\item $\mathcal{F}$ は $S$ 上 $x$ で平坦である。
\end{enumerate}
このとき点付きスキームの可換図式
$$
\xymatrix{
(X, x) \ar[d] & (X', x') \ar[l]^g \ar[d] \\
(S, s) & (S', s') \ar[l]
}
$$
が存在する。横の矢印は基本エタール近傍であり、$X'$ と $S'$ はアフィンであり、
$\Gamma(X', g^*\mathcal{F})$ は射影 $\Gamma(S', \mathcal{O}_{S'})$-加群である。
\end{proposition}

\begin{proof}
平坦性の開性（More on Morphisms,
Theorem \ref{more-morphisms-theorem-openness-flatness}）により、$X$ を $x$ の
開近傍で置き換え、$\mathcal{F}$ は $S$ 上平坦であると仮定してよい。
次に Proposition \ref{flat-proposition-existence-complete-at-x} を適用して、
命題に述べた図式で、$g^*\mathcal{F}/X'/S'$ が $s'$ 上完全デヴィサージュを
もつものを得る。（特に $S'$ と $X'$ はアフィンである。）Morphisms,
Lemma \ref{morphisms-lemma-flat-permanence} により、$g^*\mathcal{F}$ は
$S$ 上平坦であり、Lemma \ref{flat-lemma-etale-flat-up-down} により
$S'$ 上平坦である。Remark \ref{flat-remark-same-notion} を通じて、
$$
\Gamma(X', g^*\mathcal{F})/
\Gamma(X', \mathcal{O}_{X'})/
\Gamma(S', \mathcal{O}_{S'})
$$
は $s'$ に対応する $\Gamma(S', \mathcal{O}_{S'})$ の素イデアル上で
完全デヴィサージュをもつ。したがって
Lemma \ref{flat-lemma-complete-devissage-flat-finitely-presented-module} により、
$S'$ を $s'$ の標準開近傍で置き換えた後に命題の結論が成り立つ。
\end{proof}

\noindent
本節の残りでは、この結果の変形をいくつか証明する。最初は「大域的」な版である。

\begin{lemma}
\label{flat-lemma-finite-presentation-flat-along-fibre}
$f : X \to S$ をスキームの射とし、$\mathcal{F}$ を $X$ 上の準連接層とする。
$s \in S$ とする。次を仮定する。
\begin{enumerate}
\item $f$ は有限表示である。
\item $\mathcal{F}$ は有限表示である。
\item $\mathcal{F}$ はファイバー $X_s$ の各点で $S$ 上平坦である。
\end{enumerate}
このとき、基本エタール近傍 $(S', s') \to (S, s)$ とスキームの可換図式
$$
\xymatrix{
X \ar[d] & X' \ar[l]^g \ar[d] \\
S & S' \ar[l]
}
$$
が存在する。ここで $g$ はエタール、$X_s \subset g(X')$ であり、スキーム
$X'$ と $S'$ はアフィンであり、$\Gamma(X', g^*\mathcal{F})$ は
射影 $\Gamma(S', \mathcal{O}_{S'})$-加群である。
\end{lemma}

\begin{proof}
各点 $x \in X_s$ に対し、
Proposition \ref{flat-proposition-finite-presentation-flat-at-point} を用いて可換図式
$$
\xymatrix{
(X, x) \ar[d] & (Y_x, y_x) \ar[l]^{g_x} \ar[d] \\
(S, s) & (S_x, s_x) \ar[l]
}
$$
を得る。横の矢印は基本エタール近傍であり、$Y_x$ と $S_x$ はアフィンで、
$\Gamma(Y_x, g_x^*\mathcal{F})$ は射影
$\Gamma(S_x, \mathcal{O}_{S_x})$-加群である。特に
$g_x(Y_x) \cap X_s$ は $X_s$ における $x$ の開近傍である。
$X_s$ は準コンパクトなので、$X_s$ が $g_{x_i}(Y_{x_i}) \cap X_s$ の合併と
なるような有限個の点 $x_1, \ldots, x_n \in X_s$ を選べる。
More on Morphisms,
Lemma \ref{more-morphisms-lemma-elementary-etale-neighbourhoods} により、
各近傍 $(S_{x_i}, s_{x_i})$ を支配する基本エタール近傍
$(S' , s') \to (S, s)$ を選ぶ。$S'$ はアフィンであると仮定してもよい。
$X' = \coprod Y_{x_i} \times_{S_{x_i}} S'$ とおき、明らかな射
$g : X' \to X$ を与える。構成より $g(X')$ は $X_s$ を含み、
$$
\Gamma(X', g^*\mathcal{F})
=
\bigoplus \Gamma(Y_{x_i}, g_{x_i}^*\mathcal{F})
\otimes_{\Gamma(S_{x_i}, \mathcal{O}_{S_{x_i}})}
\Gamma(S', \mathcal{O}_{S'}).
$$
これは射影 $\Gamma(S', \mathcal{O}_{S'})$-加群である。Algebra,
Lemma \ref{algebra-lemma-ascend-properties-modules} を参照せよ。
\end{proof}

\noindent
次の二つの補題は、$\mathcal{F} = \mathcal{O}_X$ の場合に上の結果を
言い換えたものである。

\begin{lemma}
\label{flat-lemma-finite-presentation-flat-at-point-X}
$f : X \to S$ を局所有限表示とする。$x \in X$ の像を $s \in S$ とする。
$f$ が $S$ 上 $x$ で平坦なら、点付きスキームの可換図式
$$
\xymatrix{
(X, x) \ar[d] & (X', x') \ar[l]^g \ar[d] \\
(S, s) & (S', s') \ar[l]
}
$$
が存在する。横の矢印は基本エタール近傍であり、$X'$ と $S'$ はアフィンであり、
$\Gamma(X', \mathcal{O}_{X'})$ は射影 $\Gamma(S', \mathcal{O}_{S'})$-加群である。
\end{lemma}

\begin{proof}
これは Proposition \ref{flat-proposition-finite-presentation-flat-at-point} の特別な場合である。
\end{proof}

\begin{lemma}
\label{flat-lemma-finite-presentation-flat-along-fibre-X}
$f : X \to S$ を有限表示とし、$s \in S$ とする。
$X$ が $X_s$ のすべての点で $S$ 上平坦なら、基本エタール近傍
$(S', s') \to (S, s)$ とスキームの可換図式
$$
\xymatrix{
X \ar[d] & X' \ar[l]^g \ar[d] \\
S & S' \ar[l]
}
$$
が存在する。ここで $g$ はエタール、$X_s \subset g(X')$ であり、
$X'$ と $S'$ はアフィンであり、$\Gamma(X', \mathcal{O}_{X'})$ は
射影 $\Gamma(S', \mathcal{O}_{S'})$-加群である。
\end{lemma}

\begin{proof}
これは Lemma \ref{flat-lemma-finite-presentation-flat-along-fibre} の特別な場合である。
\end{proof}

\noindent
次の補題は、射と層が有限型であることだけを仮定し（有限表示とは限らない）、
Proposition \ref{flat-proposition-finite-presentation-flat-at-point} から得られる
帰結を説明する。

\begin{lemma}
\label{flat-lemma-finite-type-flat-at-point-free}
$f : X \to S$ をスキームの射とし、$\mathcal{F}$ を $X$ 上の準連接層とする。
$x \in X$ の像を $s \in S$ とする。次を仮定する。
\begin{enumerate}
\item $f$ は局所有限表示である。
\item $\mathcal{F}$ は有限型である。
\item $\mathcal{F}$ は $S$ 上 $x$ で平坦である。
\end{enumerate}
このとき、基本エタール近傍 $(S', s') \to (S, s)$ と点付きスキームの可換図式
$$
\xymatrix{
(X, x) \ar[d] & (X', x') \ar[l]^g \ar[d] \\
(S, s) & (\Spec(\mathcal{O}_{S', s'}), s') \ar[l]
}
$$
が存在する。ここで $X' \to X \times_S \Spec(\mathcal{O}_{S', s'})$ は
エタールで、$\kappa(x) = \kappa(x')$ であり、スキーム $X'$ は
$\mathcal{O}_{S', s'}$ 上有限表示なアフィンスキームであり、層
$g^*\mathcal{F}$ は有限表示 $\mathcal{O}_{X'}$-加群であり、
$\Gamma(X', g^*\mathcal{F})$ は自由 $\mathcal{O}_{S', s'}$-加群である。
\end{lemma}

\begin{proof}
補題を証明するために $(S, s)$ を任意の基本エタール近傍で置き換えてよく、
さらに $S$ を $\Spec(\mathcal{O}_{S, s})$ で置き換えてもよい。
したがって Proposition \ref{flat-proposition-finite-type-flat-at-point} により、
$\mathcal{F}$ は $x$ の近傍で有限表示かつ $S$ 上平坦であると仮定してよい。
この場合、Algebra, Theorem \ref{algebra-theorem-projective-free-over-local-ring}
により局所環上では射影 $=$ 自由であるから、結果は
Proposition \ref{flat-proposition-finite-presentation-flat-at-point} から従う。
\end{proof}

\begin{lemma}
\label{flat-lemma-finite-type-flat-at-point-free-variant}
$f : X \to S$ をスキームの射とし、$\mathcal{F}$ を $X$ 上の準連接層とする。
$x \in X$ の像を $s \in S$ とする。次を仮定する。
\begin{enumerate}
\item $f$ は局所有限型である。
\item $\mathcal{F}$ は有限型である。
\item $\mathcal{F}$ は $S$ 上 $x$ で平坦である。
\end{enumerate}
このとき、基本エタール近傍 $(S', s') \to (S, s)$ と点付きスキームの可換図式
$$
\xymatrix{
(X, x) \ar[d] & (X', x') \ar[l]^g \ar[d] \\
(S, s) & (\Spec(\mathcal{O}_{S', s'}), s') \ar[l]
}
$$
が存在する。ここで $X' \to X \times_S \Spec(\mathcal{O}_{S', s'})$ は
エタールで、$\kappa(x) = \kappa(x')$ であり、スキーム $X'$ はアフィンであり、
$\Gamma(X', g^*\mathcal{F})$ は自由 $\mathcal{O}_{S', s'}$-加群である。
\end{lemma}

\begin{proof}
% P05-EN-CH38-0147: source-faithful missing locally qualifier; see sidecar.
（これと Lemma \ref{flat-lemma-finite-type-flat-at-point-free} の唯一の違いは、
$f$ が有限表示であると仮定しないことである。）問題は $X$ と $S$ 上局所的なので、
$X$ と $S$ はアフィンであると仮定してよい。$X = \Spec(B)$、
$S = \Spec(A)$ と書く。$B$ は有限型 $A$-代数なので、全射
$A[x_1, \ldots, x_n] \to B$ をとれる。言い換えれば、閉埋め込み
$i : X \to \mathbf{A}^n_S$ を選べる。$t = i(x)$、
$\mathcal{G} = i_*\mathcal{F}$ とおく。
% P05-EN-CH38-0143: uniquely determined subject-agreement repair; see sidecar.
$\mathcal{G}_t \cong \mathcal{F}_x$ は $\mathcal{O}_{S, s}$-加群の同型である。
したがって $\mathcal{G}$ は $S$ 上 $t$ で平坦である。射
$\mathbf{A}^n_S \to S$、点 $t$、層 $\mathcal{G}$ に
Lemma \ref{flat-lemma-finite-type-flat-at-point-free} を適用する。
すると基本エタール近傍 $(S', s') \to (S, s)$ と点付きスキームの可換図式
$$
\xymatrix{
(\mathbf{A}^n_S, t) \ar[d] & (Y, y) \ar[l]^h \ar[d] \\
(S, s) & (\Spec(\mathcal{O}_{S', s'}), s') \ar[l]
}
$$
を得る。ここで $Y \to \mathbf{A}^n_{\mathcal{O}_{S', s'}}$ はエタール、
$\kappa(t) = \kappa(y)$ であり、スキーム $Y$ はアフィンであり、
% P05-EN-CH38-0144: source-faithful projective/free conclusion mismatch; see sidecar.
$\Gamma(Y, h^*\mathcal{G})$ は射影 $\mathcal{O}_{S', s'}$-加群である。
元の問題の解は、$Y$ の閉部分スキーム
$X' = Y \times_{\mathbf{A}^n_S} X$ によって与えられる。
\end{proof}

\begin{lemma}
\label{flat-lemma-finite-type-flat-along-fibre-free}
$f : X \to S$ をスキームの射とし、$\mathcal{F}$ を $X$ 上の準連接層とする。
$s \in S$ とする。次を仮定する。
\begin{enumerate}
\item $f$ は有限表示である。
\item $\mathcal{F}$ は有限型である。
\item $\mathcal{F}$ は $X_s$ のすべての点で $S$ 上平坦である。
\end{enumerate}
このとき、基本エタール近傍 $(S', s') \to (S, s)$ とスキームの可換図式
$$
\xymatrix{
X \ar[d] & X' \ar[l]^g \ar[d] \\
S & \Spec(\mathcal{O}_{S', s'}) \ar[l]
}
$$
が存在する。ここで $X' \to X \times_S \Spec(\mathcal{O}_{S', s'})$ は
エタール、$X_s = g((X')_{s'})$ であり、スキーム $X'$ は
$\mathcal{O}_{S', s'}$ 上有限表示なアフィンスキームであり、層
$g^*\mathcal{F}$ は有限表示 $\mathcal{O}_{X'}$-加群であり、
$\Gamma(X', g^*\mathcal{F})$ は自由 $\mathcal{O}_{S', s'}$-加群である。
\end{lemma}

\begin{proof}
各点 $x \in X_s$ に対し、Lemma \ref{flat-lemma-finite-type-flat-at-point-free}
を用いて、基本エタール近傍 $(S_x , s_x) \to (S, s)$ と可換図式
$$
\xymatrix{
(X, x) \ar[d] & (Y_x, y_x) \ar[l]^{g_x} \ar[d] \\
(S, s) & (\Spec(\mathcal{O}_{S_x, s_x}), s_x) \ar[l]
}
$$
を得る。ここで $Y_x \to X \times_S \Spec(\mathcal{O}_{S_x, s_x})$ はエタール、
$\kappa(x) = \kappa(y_x)$ であり、スキーム $Y_x$ は
$\mathcal{O}_{S_x, s_x}$ 上有限表示なアフィンスキームであり、層
$g_x^*\mathcal{F}$ は有限表示 $\mathcal{O}_{Y_x}$-加群であり、
$\Gamma(Y_x, g_x^*\mathcal{F})$ は自由 $\mathcal{O}_{S_x, s_x}$-加群である。
特に $g_x((Y_x)_{s_x})$ は $X_s$ における $x$ の開近傍である。
$X_s$ は準コンパクトなので、$X_s$ が $g_{x_i}((Y_{x_i})_{s_{x_i}})$ の
合併となるような有限個の点 $x_1, \ldots, x_n \in X_s$ を選べる。
More on Morphisms,
Lemma \ref{more-morphisms-lemma-elementary-etale-neighbourhoods} により、
各近傍 $(S_{x_i}, s_{x_i})$ を支配する基本エタール近傍
$(S' , s') \to (S, s)$ を選ぶ。次のようにおく。
$$
X' = \coprod Y_{x_i} \times_{\Spec(\mathcal{O}_{S_{x_i}, s_{x_i}})}
\Spec(\mathcal{O}_{S', s'})
$$
さらに明らかな射 $g : X' \to X$ を与える。構成より $X_s = g(X'_{s'})$ であり、
$$
\Gamma(X', g^*\mathcal{F})
=
\bigoplus \Gamma(Y_{x_i}, g_{x_i}^*\mathcal{F})
\otimes_{\mathcal{O}_{S_{x_i}, s_{x_i}}}
\mathcal{O}_{S', s'}.
$$
これは自由加群の基底変換の直和として、自由 $\mathcal{O}_{S', s'}$-加群である。
% P05-EN-CH38-0145: source-faithful explicit proof-detail omission; see sidecar.
細部をいくつか省略する。
\end{proof}

\begin{lemma}
\label{flat-lemma-finite-type-flat-along-fibre-free-variant}
$f : X \to S$ をスキームの射とし、$\mathcal{F}$ を $X$ 上の準連接層とする。
$s \in S$ とする。次を仮定する。
\begin{enumerate}
\item $f$ は有限型である。
\item $\mathcal{F}$ は有限型である。
\item $\mathcal{F}$ は $X_s$ のすべての点で $S$ 上平坦である。
\end{enumerate}
このとき、基本エタール近傍 $(S', s') \to (S, s)$ とスキームの可換図式
$$
\xymatrix{
X \ar[d] & X' \ar[l]^g \ar[d] \\
S & \Spec(\mathcal{O}_{S', s'}) \ar[l]
}
$$
が存在する。ここで $X' \to X \times_S \Spec(\mathcal{O}_{S', s'})$ はエタール、
$X_s = g((X')_{s'})$ であり、スキーム $X'$ はアフィンであり、
$\Gamma(X', g^*\mathcal{F})$ は自由 $\mathcal{O}_{S', s'}$-加群である。
\end{lemma}

\begin{proof}
（これと Lemma \ref{flat-lemma-finite-type-flat-along-fibre-free} の唯一の違いは、
$f$ が有限表示であると仮定しないことである。）各点 $x \in X_s$ に対し、
Lemma \ref{flat-lemma-finite-type-flat-at-point-free-variant} を用いて、
基本エタール近傍 $(S_x , s_x) \to (S, s)$ と可換図式
$$
\xymatrix{
(X, x) \ar[d] & (Y_x, y_x) \ar[l]^{g_x} \ar[d] \\
(S, s) & (\Spec(\mathcal{O}_{S_x, s_x}), s_x) \ar[l]
}
$$
を得る。ここで $Y_x \to X \times_S \Spec(\mathcal{O}_{S_x, s_x})$ はエタール、
$\kappa(x) = \kappa(y_x)$ であり、スキーム $Y_x$ はアフィンであり、
$\Gamma(Y_x, g_x^*\mathcal{F})$ は自由 $\mathcal{O}_{S_x, s_x}$-加群である。
特に $g_x((Y_x)_{s_x})$ は $X_s$ における $x$ の開近傍である。
$X_s$ は準コンパクトなので、$X_s$ が $g_{x_i}((Y_{x_i})_{s_{x_i}})$ の
合併となるような有限個の点 $x_1, \ldots, x_n \in X_s$ を選べる。
More on Morphisms,
Lemma \ref{more-morphisms-lemma-elementary-etale-neighbourhoods} により、
各近傍 $(S_{x_i}, s_{x_i})$ を支配する基本エタール近傍
$(S' , s') \to (S, s)$ を選ぶ。次のようにおく。
$$
X' = \coprod Y_{x_i} \times_{\Spec(\mathcal{O}_{S_{x_i}, s_{x_i}})}
\Spec(\mathcal{O}_{S', s'})
$$
さらに明らかな射 $g : X' \to X$ を与える。構成より $X_s = g(X'_{s'})$ であり、
$$
\Gamma(X', g^*\mathcal{F})
=
\bigoplus \Gamma(Y_{x_i}, g_{x_i}^*\mathcal{F})
\otimes_{\mathcal{O}_{S_{x_i}, s_{x_i}}}
\mathcal{O}_{S', s'}.
$$
これは自由加群の基底変換の直和として、自由 $\mathcal{O}_{S', s'}$-加群である。
\end{proof}




\section{有限型平坦加群 II}
\label{flat-section-finite-type-flat-II}

\noindent
以下の補題が必要となる。

\begin{lemma}
\label{flat-lemma-weak-bourbaki-pre-pre}
$R \to S$ を有限表示な環準同型とする。$N$ を有限表示 $S$-加群とする。
$\mathfrak q \subset S$ を $\mathfrak p \subset R$ の上にある素イデアルとする。次のようにおく。
$\overline{S} = S \otimes_R \kappa(\mathfrak p)$、
$\overline{\mathfrak q} = \mathfrak q \overline{S}$、および
$\overline{N} = N \otimes_R \kappa(\mathfrak p)$。このとき、
$g \in S$ であって $g \not \in \mathfrak q$ であり、かつ
$\overline{g} \in \mathfrak r$ が、
$\mathfrak r \in \text{Ass}_{\overline{S}}(\overline{N})$ かつ
$\mathfrak r \not \subset \overline{\mathfrak q}$ となるすべての場合に
成り立つようなものを見いだせる。
\end{lemma}

\begin{proof}
実際、$\text{Ass}_{\overline{S}}(\overline{N}) =
\{\mathfrak r_1, \ldots, \mathfrak r_n\}$
とする（有限性は Algebra, Lemma \ref{algebra-lemma-finite-ass} による）。
番号を付け替えることにより、次を仮定してよい。
$$
\mathfrak r_1 \subset \overline{\mathfrak q},
\ldots,
\mathfrak r_r \subset \overline{\mathfrak q}, \quad
\mathfrak r_{r + 1} \not \subset \overline{\mathfrak q},
\ldots,
\mathfrak r_n \not \subset \overline{\mathfrak q}
$$
ここで $\overline{\mathfrak q}$ は素イデアルなので、積
$\mathfrak r_{r + 1} \ldots \mathfrak r_n$ は
$\overline{\mathfrak q}$ に含まれない。したがって、
$\mathfrak r_{r + 1}, \ldots, \mathfrak r_n$ のすべてに含まれるが
$\overline{\mathfrak q}$ には含まれない $\overline{S}$ の元 $a$ を選べる。
$a$ に写る $g \in S$ が存在すれば、その $g$ が所望のものである。一般には、
$\lambda \in \kappa(\mathfrak p)$
が非零で、$\lambda a$ がある $g \in S$ の像となるように選べる。
\end{proof}

\noindent
次の補題には、後出の Lemma \ref{flat-lemma-weak-bourbaki} という
少し強い変種がある。

\begin{lemma}
\label{flat-lemma-weak-bourbaki-pre}
$R \to S$ を有限表示な環準同型とする。
$N$ を有限表示 $S$-加群であって $R$-加群として平坦なものとする。
$M$ を $R$-加群とする。$\mathfrak q$ を $\mathfrak p \subset R$ の上にある
$S$ の素イデアルとする。このとき
$$
\mathfrak q \in \text{WeakAss}_S(M \otimes_R N)
\Leftrightarrow
\Big(
\mathfrak p \in \text{WeakAss}_R(M)
\text{ かつ }
\overline{\mathfrak q} \in \text{Ass}_{\overline{S}}(\overline{N})
\Big)
$$
ここで $\overline{S} = S \otimes_R \kappa(\mathfrak p)$、
$\overline{\mathfrak q} = \mathfrak q \overline{S}$、および
$\overline{N} = N \otimes_R \kappa(\mathfrak p)$ とする。
\end{lemma}

\begin{proof}
Lemma \ref{flat-lemma-weak-bourbaki-pre-pre} のように $g \in S$ を選ぶ。
スキームの射 $\Spec(S_g) \to \Spec(R)$、$N_g$ に対応する準連接加群、
および素イデアル $\mathfrak qS_g$ と $\mathfrak p$ に対応する点に
Proposition \ref{flat-proposition-finite-presentation-flat-at-point} を適用する。
代数へ翻訳すると、次の環の可換図式を得る。
$$
\xymatrix{
S \ar[r] & S_g \ar[r] & S' \\
& R \ar[lu] \ar[u] \ar[r] & R' \ar[u]
}
\quad\quad
\xymatrix{
\mathfrak q \ar@{-}[r] \ar@{-}[rd] &
\mathfrak qS_g \ar@{-}[d] \ar@{-}[r] & \mathfrak q' \ar@{-}[d] \\
& \mathfrak p \ar@{-}[r] & \mathfrak p'
}
$$
ここでは図示した素イデアルが付随し、水平矢印はエタールであり、
$N \otimes_S S'$ は $R'$-加群として射影的である。次のようにおく。
$N' = N \otimes_S S'$, $M' = M \otimes_R R'$,
% P05-EN-CH38-0148: source-faithful residue-field/base-ring mismatch; see sidecar.
$\overline{S}' = S' \otimes_{R'} \kappa(\mathfrak q')$,
$\overline{\mathfrak q}' = \mathfrak q' \overline{S}'$,
および
$$
\overline{N}' = N' \otimes_{R'} \kappa(\mathfrak p') =
\overline{N} \otimes_{\overline{S}} \overline{S}'
$$
Lemma \ref{flat-lemma-etale-weak-assassin-up-down} により
\begin{align*}
\text{WeakAss}_{S'}(M' \otimes_{R'} N') & =
(\Spec(S') \to \Spec(S))^{-1}\text{WeakAss}_S(M \otimes_R N) \\
\text{WeakAss}_{R'}(M') & =
(\Spec(R') \to \Spec(R))^{-1}\text{WeakAss}_R(M) \\
\text{Ass}_{\overline{S}'}(\overline{N}') & =
(\Spec(\overline{S}') \to \Spec(\overline{S}))^{-1}
\text{Ass}_{\overline{S}}(\overline{N})
\end{align*}
$\overline{N}$ と $\overline{N}'$ に Algebra,
Lemma \ref{algebra-lemma-ass-weakly-ass} を用いる。特に
\begin{align*}
\mathfrak q \in \text{WeakAss}_S(M \otimes_R N)
& \Leftrightarrow
\mathfrak q' \in \text{WeakAss}_{S'}(M' \otimes_{R'} N') \\
\mathfrak p \in \text{WeakAss}_R(M)
& \Leftrightarrow
\mathfrak p' \in \text{WeakAss}_{R'}(M') \\
\overline{\mathfrak q} \in \text{Ass}_{\overline{S}}(\overline{N})
& \Leftrightarrow
\overline{\mathfrak q}' \in \text{WeakAss}_{\overline{S}'}(\overline{N}')
\end{align*}
$g$ を慎重に選んだことと、上の
$\text{Ass}_{\overline{S}'}(\overline{N}')$ の公式から次が従う。
\begin{equation}
\label{flat-equation-key-observation}
\text{もし }\mathfrak r' \in \text{Ass}_{\overline{S}'}(\overline{N}')
\text{ が }\mathfrak r \subset \overline{S}\text{ の上にあれば }
\mathfrak r \subset \overline{\mathfrak q}
\end{equation}
これは後の証明で鍵となる観察である。以下では、Algebra,
Lemma \ref{algebra-lemma-weakly-ass-local} による弱随伴素イデアルの特徴づけを
断りなく用いる。

\medskip\noindent
$\overline{\mathfrak q} \not \in \text{Ass}_{\overline{S}}(\overline{N})$
と仮定する。このとき
$\overline{\mathfrak q}' \not \in \text{Ass}_{\overline{S}'}(\overline{N}')$.
したがって
Algebra, Lemmas \ref{algebra-lemma-ass-zero-divisors},
\ref{algebra-lemma-finite-ass}、および
\ref{algebra-lemma-silly}
により、$\overline{N}'$ 上の零因子でない元
$\overline{a}' \in \overline{\mathfrak q}'$ が存在する。
ある非零元 $\lambda \in \kappa(\mathfrak p)$ によって
$\overline{a}'$ を $\lambda \overline{a}'$ で置き換えれば、
$\overline{a}'$ に写る $a' \in \mathfrak q'$ を見いだせる。
Lemma \ref{flat-lemma-invert-universally-injective} により、写像
$a' : N'_{\mathfrak p'} \to N'_{\mathfrak p'}$ は
$R'_{\mathfrak p'}$ 上普遍単射である。特に、
$a' : M' \otimes_{R'} N' \to M' \otimes_{R'} N'$ は
$\mathfrak p'$ で局所化した後に単射であり、したがって
$\mathfrak q'$ で局所化した後にも単射である。これは明らかに
$\mathfrak q' \not \in \text{WeakAss}_{S'}(M' \otimes_{R'} N')$.
を意味する。ゆえに $\mathfrak q \in \text{WeakAss}_S(M \otimes_R N)$ ならば
$\overline{\mathfrak q} \in \text{Ass}_{\overline{S}}(\overline{N})$.
である。

\medskip\noindent
$\mathfrak q \in \text{WeakAss}_S(M \otimes_R N)$ と仮定する。ここで
% P05-EN-CH38-0149: source-faithful module/base-ring namespace mismatch; see sidecar.
$\mathfrak p \in \text{WeakAss}_S(M)$ を示したい。
$\mathfrak q$ が $J = \text{Ann}_S(z)$ 上極小となるような元
$z \in M \otimes_R N$ をとる。イデアル $\mathfrak p$ の生成元の集合を
$f_i \in \mathfrak p$, $i \in I$ とする。$\mathfrak q$ は
$\mathfrak p$ の上にあるので、各 $i$ に対して、
$g_i \not \in \mathfrak q$ かつ $g_i f_i^{n_i} \in J$、すなわち
$g_i f_i^{n_i} z = 0$ となる $n_i \geq 1$ と $g_i \in S$ を選べる。
$z$ の像を $z' \in (M' \otimes_{R'} N')_{\mathfrak p'}$ とする。
$z$ の $(M \otimes_R N)_\mathfrak q$ における像が非零であり、かつ
$S_\mathfrak q \to S'_{\mathfrak q'}$ が忠実平坦なので、$z'$ は非零である。
$f_i^{n_i} z' = 0$ であると主張する。

\medskip\noindent
主張の証明：$g_i$ の像を $g'_i \in S'$ とする。
鍵となる観察 (\ref{flat-equation-key-observation}) により、像
$\overline{g}'_i \in \overline{S}'$ は、任意の $\mathfrak r'$ であって
% P05-EN-CH38-0150: source-faithful missing-prime module mismatch; see sidecar.
$\mathfrak r' \in \text{Ass}_{\overline{S}'}(\overline{N})$ となるものに
含まれない。したがって Lemma \ref{flat-lemma-invert-universally-injective} により、
$g'_i : N'_{\mathfrak p'} \to N'_{\mathfrak p'}$ は
$R'_{\mathfrak p'}$ 上普遍単射である。特に、
% P05-EN-CH38-0151: uniquely determined spelling repair; see sidecar.
$g'_i : M' \otimes_{R'} N' \to M' \otimes_{R'} N'$ は
$\mathfrak p'$ で局所化した後に単射である。
$g_i f_i^{n_i} z' = 0$ なので主張が従う。

\medskip\noindent
この主張から、$R'_{\mathfrak p'}$ における $z'$ の零化イデアルは
元 $f_i^{n_i}$ を含む。$R \to R'$ はエタールなので、Algebra,
Lemma \ref{algebra-lemma-etale-at-prime} により
$\mathfrak p'R'_{\mathfrak p'} = \mathfrak pR'_{\mathfrak p'}$
である。したがって、$R'_{\mathfrak p'}$ における $z'$ の零化イデアルの根基は
% P05-EN-CH38-0152: source-faithful missing-prime localization mismatch; see sidecar.
$\mathfrak p' R_{\mathfrak p'}$
に等しい（ここで $z'$ が非零であることを用いた）。一方、
$$
z' \in (M' \otimes_{R'} N')_{\mathfrak p'} =
M'_{\mathfrak p'} \otimes_{R'_{\mathfrak p'}} N'_{\mathfrak p'}
$$
加群 $N'_{\mathfrak p'}$ は局所環 $R'_{\mathfrak p'}$ 上射影的なので自由である
（Algebra, Theorem \ref{algebra-theorem-projective-free-over-local-ring}）。
したがって、$z' \in M'_{\mathfrak p'} \otimes_{R'_{\mathfrak p'}} F'$
となる有限自由直和因子 $F' \subset N'_{\mathfrak p'}$ を見いだせる。
$F'$ の階数を $n$ とすれば、
$\mathfrak p' R'_{\mathfrak p'} \in
\text{WeakAss}_{R'_{\mathfrak p'}}({M'_{\mathfrak p'}}^{\oplus n})$.
が従う。これは、たとえば Algebra, Lemma \ref{algebra-lemma-weakly-ass} により
$\mathfrak p'R'_{\mathfrak p'} \in \text{WeakAss}(M'_{\mathfrak p'})$
を意味する。ゆえに $\mathfrak p' \in \text{WeakAss}_{R'}(M')$ であり、
これはさらに $\mathfrak p \in \text{WeakAss}_R(M)$ を与える。
以上で補題の同値における含意 ``$\Rightarrow$'' の証明が完了した。

\medskip\noindent
$\mathfrak p \in \text{WeakAss}_R(M)$ かつ
$\overline{\mathfrak q} \in \text{Ass}_{\overline{S}}(\overline{N})$
と仮定する。$\mathfrak q$ が $M \otimes_R N$ の弱随伴素イデアルであることを
示したい。$\overline{\mathfrak q}'$ は
$\text{Ass}_{\overline{S}'}(\overline{N}')$ の極大元であることに注意する。
これは (\ref{flat-equation-key-observation}) と、$\overline{\mathfrak q}$ の上にある
$\overline{S}'$ の素イデアルの間には包含関係がないという事実から従う
（エタール射のファイバーは離散的である。Morphisms,
Lemma \ref{morphisms-lemma-etale-over-field}）。したがって、
$R, S, \mathfrak p, \mathfrak q, M, N$ をそれぞれ
$R', S', \mathfrak p', \mathfrak q', M', N'$ で置き換えることにより、
補題の仮定に加えて次を仮定してよい。
\begin{enumerate}
\item $\mathfrak p \in \text{WeakAss}_R(M)$ である。
\item $\overline{\mathfrak q} \in \text{Ass}_{\overline{S}}(\overline{N})$ である。
\item $N$ は $R$-加群として射影的である。
\item $\overline{\mathfrak q}$ は
$\text{Ass}_{\overline{S}}(\overline{N})$ において極大である。
\end{enumerate}
もう一つ簡約がある。すなわち $R, S, M, N$ をそれぞれ
$\mathfrak p$ における局所化で置き換えてよい。これにより、さらに次の条件を得る。
\begin{enumerate}
\item[(5)] $R$ は極大イデアル $\mathfrak p$ をもつ局所環である。
\end{enumerate}
(1) -- (5) から $\mathfrak q \in \text{WeakAss}(M \otimes_R N)$ が
従うことを示して証明を終える。

\medskip\noindent
$R$ は局所環であり $\mathfrak p \in \text{WeakAss}_R(M)$ なので、
その零化イデアル $I$ の根基が $\mathfrak p$ に等しくなるような
$y \in M$ を選べる。ある $\overline{g}_i \in \overline{S}$ を用いて
$\overline{\mathfrak q} = (\overline{g}_1, \ldots, \overline{g}_n)$
と書く。$\overline{g}_i$ に写る $g_i \in S$ を選ぶ。このとき
$\mathfrak q = \mathfrak pS + g_1S + \ldots + g_nS$ である。写像
$$
\Psi : N/IN \longrightarrow (N/IN)^{\oplus n}, \quad
z \longmapsto (g_1z, \ldots, g_nz).
$$
を考える。これは射影 $R/I$-加群の準同型である。
$\mathfrak p/I$ は局所冪零なので、局所環 $R/I$ は自己随伴的である
（More on Algebra, Definition \ref{more-algebra-definition-auto-ass}）。
$\overline{\mathfrak q} \in \text{Ass}_{\overline{S}}(\overline{N})$ なので、
写像 $\Psi \otimes \kappa(\mathfrak p)$ は単射でない。
したがって More on Algebra, Lemma \ref{more-algebra-lemma-P-fPD-zero} により
$\Psi$ は単射でない。$\Psi$ の核の非零元 $z \in N/IN$ を選ぶ。
零化イデアル $J = \text{Ann}_S(z)$ は、構成により $IS$ と $g_i$ を含む。
したがって $\sqrt{J} \subset S$ は $\mathfrak q$ を含む。
$J$ 上極小な素イデアル $\mathfrak s \subset S$ をとる。このとき
$\mathfrak q \subset \mathfrak s$ であり、$\mathfrak s$ は
$\mathfrak p$ の上にあり、かつ
$\mathfrak s \in \text{WeakAss}_S(N/IN)$ である。
% P05-EN-CH38-0153: uniquely determined missing-verb repair; see sidecar.
最後の事実は弱随伴素イデアルの定義から従う。
環準同型 $R \to S$ と加群 $R/I$, $N$ に、すでに証明した補題の
``$\Rightarrow$'' の部分を適用すると、
$\overline{\mathfrak s} \in \text{Ass}_{\overline{S}}(\overline{N})$.
を得る。$\overline{\mathfrak q} \subset \overline{\mathfrak s}$ なので、
上の条件 (4) にある $\overline{\mathfrak q}$ の極大性から
$\overline{\mathfrak q} = \overline{\mathfrak s}$ が従う。
% P05-EN-CH38-0154: uniquely determined spelling repair; see sidecar.
よって $\mathfrak q = \mathfrak s$ となり、所望の結論を得る。
\end{proof}

\begin{lemma}
\label{flat-lemma-bourbaki-finite-type-general-base-at-point}
$S$ をスキームとする。$f : X \to S$ を局所有限型とする。
$x \in X$ の像を $s \in S$ とする。
$\mathcal{F}$ を $X$ 上の有限型準連接層とし、
$\mathcal{G}$ を $S$ 上の準連接層とする。
$\mathcal{F}$ が $x$ において $S$ 上平坦ならば、
$$
x \in \text{WeakAss}_X(\mathcal{F} \otimes_{\mathcal{O}_X} f^*\mathcal{G})
\Leftrightarrow
s \in \text{WeakAss}_S(\mathcal{G})
\text{ かつ }
x \in \text{Ass}_{X_s}(\mathcal{F}_s).
$$
\end{lemma}

\begin{proof}
この段落では、$f$ が有限表示である場合に帰着する。
問題は $X$ と $S$ 上局所的なので、$X$ と $S$ はアフィンであると仮定してよい。
$X = \Spec(B)$、$S = \Spec(A)$ と書き、
$B = A[x_1, \ldots, x_n]/I$ と書く。言い換えると、$S$ 上の閉埋め込み
$i : X \to \mathbf{A}^n_S$ を得る。$t = i(x) \in \mathbf{A}^n_S$ とおく。
$i_*\mathcal{F}$ は $\mathbf{A}^n_S$ 上の有限型準連接層であり、
$t$ において $S$ 上平坦であることに注意する。また
$$
i_*(\mathcal{F} \otimes_{\mathcal{O}_X} f^*\mathcal{G}) =
i_*\mathcal{F} \otimes_{\mathcal{O}_{\mathbf{A}^n_S}} p^*\mathcal{G}
$$
である。ここで $p : \mathbf{A}^n_S \to S$ は射影である。
Divisors, Lemma \ref{divisors-lemma-weakly-associated-finite} により、
$t$ が $i_*(\mathcal{F} \otimes_{\mathcal{O}_X} f^*\mathcal{G})$ の
弱随伴点であることと、$x$ が
$\mathcal{F} \otimes_{\mathcal{O}_X} f^*\mathcal{G}$ の
弱随伴点であることは同値である。同様に、
$x \in \text{Ass}_{X_s}(\mathcal{F}_s)$ であることと
$t \in \text{Ass}_{\mathbf{A}^n_s}((i_*\mathcal{F})_s)$ であることも同値である
（Algebra, Lemma \ref{algebra-lemma-ass-quotient-ring}）。
したがって $X = \mathbf{A}^n_S$ の場合に補題を証明すれば十分である。
ゆえに $X \to S$ は有限表示であると仮定してよい。

\medskip\noindent
この段落では、$\mathcal{F}$ が有限表示で $S$ 上平坦である場合に帰着する。
Proposition \ref{flat-proposition-finite-type-flat-at-point} のように、
基本エタール近傍 $e : (S', s') \to (S, s)$ と開部分
$V \subset X \times_S \Spec(\mathcal{O}_{S', s'})$ を選ぶ。
$x$ と $s'$ に写る一意な点を $x' \in X' = X \times_S S'$ とする。
Lemma \ref{flat-lemma-etale-weak-assassin-up-down} により、
$X' \to S'$、$x'$、$s'$、$(X' \to X)^*\mathcal{F}$、および
$e^*\mathcal{G}$ に対して主張を証明すれば十分である。
% P05-EN-CH38-0155: uniquely determined missing-verb repair; see sidecar.
$x'$ に写る一意な点を $v \in V$ とし、
$s' \in \Spec(\mathcal{O}_{S', s'})$ を閉点とする。このとき
$\mathcal{O}_{V, v} = \mathcal{O}_{X', x'}$ および
$\mathcal{O}_{\Spec(\mathcal{O}_{S', s'}), s'} =
\mathcal{O}_{S', s'}$ であり、$\mathcal{F}$ と $\mathcal{G}$ の
引き戻しの茎についても同様である。また
$V_{s'} \subset X'_{s'}$ は開部分スキームである。
% P05-EN-CH38-0156: uniquely determined subject-agreement repair; see sidecar.
弱随伴点であるという条件は層の茎のみに依存するので、
$X' \to S'$、$x'$、$s'$、$(X' \to X)^*\mathcal{F}$、$e^*\mathcal{G}$ を、
$V \to \Spec(\mathcal{O}_{S', s'})$, $v$, $s'$, $(V \to X)^*\mathcal{F}$,
および $(\Spec(\mathcal{O}_{S', s'}) \to S)^*\mathcal{G}$ で
置き換えてよい。したがって $f$ は有限表示であり、
$\mathcal{F}$ は有限表示で $S$ 上平坦であると仮定してよい。

\medskip\noindent
$f$ は有限表示であり、$\mathcal{F}$ は有限表示で $S$ 上平坦であると仮定する。
$X$ と $S$ を $x$ と $s$ のアフィン近傍へ縮小すれば、この場合は
Lemma \ref{flat-lemma-weak-bourbaki-pre} から従う。
\end{proof}

\begin{lemma}
\label{flat-lemma-weak-bourbaki}
$R \to S$ を本質的有限型な環準同型とする。
$N$ を $R$ 上平坦な有限 $S$-加群の局所化とする。
$M$ を $R$-加群とする。このとき
$$
\text{WeakAss}_S(M \otimes_R N)
=
\bigcup\nolimits_{\mathfrak p \in \text{WeakAss}_R(M)}
\text{Ass}_{S \otimes_R \kappa(\mathfrak p)}(N \otimes_R \kappa(\mathfrak p))
$$
\end{lemma}

\begin{proof}
この補題は Lemma \ref{flat-lemma-bourbaki-finite-type-general-base-at-point} を
代数へ翻訳したものである。翻訳の詳細は省略する。
\end{proof}

\begin{lemma}
\label{flat-lemma-bourbaki-finite-type-general-base}
$f : X \to S$ を局所有限型な射とする。
$\mathcal{F}$ を $X$ 上の有限型準連接層であって $S$ 上平坦なものとする。
$\mathcal{G}$ を $S$ 上の準連接層とする。このとき
$$
\text{WeakAss}_X(\mathcal{F} \otimes_{\mathcal{O}_X} f^*\mathcal{G}) =
\bigcup\nolimits_{s \in \text{WeakAss}_S(\mathcal{G})}
\text{Ass}_{X_s}(\mathcal{F}_s)
$$
\end{lemma}

\begin{proof}
Lemma \ref{flat-lemma-bourbaki-finite-type-general-base-at-point} の
直ちに得られる帰結である。
\end{proof}

\begin{theorem}
\label{flat-theorem-finite-type-flat}
\begin{slogan}
平坦軌跡は開である（非ネーター版）。
\end{slogan}
$f : X \to S$ をスキームの射とする。
$\mathcal{F}$ を準連接 $\mathcal{O}_X$-加群とする。次を仮定する。
\begin{enumerate}
\item $X \to S$ は局所有限表示である。
\item $\mathcal{F}$ は有限型 $\mathcal{O}_X$-加群である。
\item $S$ の弱随伴点の集合は $S$ において局所有限である。
\end{enumerate}
このとき $U = \{x \in X \mid \mathcal{F}\text{ は }S\text{ 上 }x\text{ で平坦}\}$
は $X$ において開であり、$\mathcal{F}|_U$ は有限表示
$\mathcal{O}_U$-加群であって $S$ 上平坦である。
\end{theorem}

\begin{proof}
$\mathcal{F}$ が $x$ において $S$ 上平坦となる $x \in X$ をとる。
$x$ のある開近傍上で、$\mathcal{F}$ の制限が有限表示かつ $S$ 上平坦な
加群となることを示さなければならない。問題は $X$ と $S$ 上局所的なので、
$X$ と $S$ はアフィンであると仮定してよい。
Lemma \ref{flat-lemma-finite-type-flat-at-point-local} により
$\mathcal{F}_x$ は有限表示 $\mathcal{O}_{X, x}$-加群である。
したがって Algebra, Lemma \ref{algebra-lemma-construct-fp-module-from-stalk} により、
有限表示 $\mathcal{O}_X$-加群 $\mathcal{F}'$ と、同型
$\varphi_x : \mathcal{F}'_x \to \mathcal{F}_x$ を誘導する写像
$\varphi : \mathcal{F}' \to \mathcal{F}$ が存在する。特に
$\mathcal{F}'$ は $x$ において $S$ 上平坦である。よって平坦性の開性
More on Morphisms, Theorem \ref{more-morphisms-theorem-openness-flatness} により、
$X$ を縮小した後、$\mathcal{F}'$ は $S$ 上平坦であると仮定してよい。
$\mathcal{F}$ は有限型なので、さらに $X$ を縮小すれば、$\varphi$ は全射であると
仮定してよい。Modules, Lemma \ref{modules-lemma-finite-type-surjective-on-stalk}
を参照せよ。あるいは中山の補題
（Algebra, Lemma \ref{algebra-lemma-NAK}）を用いてもよい。
Lemma \ref{flat-lemma-bourbaki-finite-type-general-base} により
$$
\text{WeakAss}_X(\mathcal{F}') \subset
\bigcup\nolimits_{s \in \text{WeakAss}(S)} \text{Ass}_{X_s}(\mathcal{F}'_s)
$$
である。仮定により $\text{WeakAss}(S)$ は有限であり、
Divisors, Lemma \ref{divisors-lemma-finite-ass} により
$\text{Ass}_{X_s}(\mathcal{F}'_s)$ も有限なので、
$\text{WeakAss}_X(\mathcal{F}')$ は有限である。Algebra,
Lemma \ref{algebra-lemma-silly} を用い、$X$ をもう一度縮小すれば、
$\text{WeakAss}_X(\mathcal{F}')$ は $x$ の一般化全体に含まれると
仮定してよい。ここで $\mathcal{K} = \Ker(\varphi)$ を考える。すると
$\text{WeakAss}_X(\mathcal{K}) \subset \text{WeakAss}_X(\mathcal{F}')$
である（
Divisors, Lemma \ref{divisors-lemma-ses-weakly-ass})
による）。一方、$\varphi_x$ は同型なので、
% P05-EN-CH38-0157: uniquely determined sentence-connection repair; see sidecar.
すべての $x' \leadsto x$ に対して $\varphi_{x'}$ も同型である。したがって
$\text{WeakAss}_X(\mathcal{K}) = \emptyset$
であり、Divisors, Lemma \ref{divisors-lemma-weakly-ass-zero} により
$\mathcal{K} = 0$ である。
\end{proof}

\begin{lemma}
\label{flat-lemma-finite-type-flat-algebra}
$R \to S$ を有限表示な環準同型とする。$M$ を有限 $S$-加群とする。
$\text{WeakAss}_R(R)$ は有限であると仮定する。このとき
$$
U = \{\mathfrak q \subset S \mid M_{\mathfrak q}\text{ は }R\text{ 上平坦}\}
$$
は $\Spec(S)$ において開である。また、$D(g) \subset U$ となる
すべての $g \in S$ に対して、局所化 $M_g$ は有限表示
$S_g$-加群であって $R$ 上平坦である。
\end{lemma}

\begin{proof}
Theorem \ref{flat-theorem-finite-type-flat} から直ちに従う。
\end{proof}

\begin{lemma}
\label{flat-lemma-finite-type-flat-X}
$f : X \to S$ を局所有限型なスキームの射とする。$S$ の弱随伴点の集合は
$S$ において局所有限であると仮定する。このとき、$f$ が平坦である点
$x \in X$ の集合は開部分スキーム $U \subset X$ であり、
$U \to S$ は平坦かつ局所有限表示である。
\end{lemma}

\begin{proof}
問題は $X$ と $S$ 上局所的なので、$X$ と $S$ はアフィンであると仮定してよい。
このとき $X \to S$ は有限型環準同型 $A \to B$ に対応する。
全射 $A[x_1, \ldots, x_n] \to B$ を選び、$B$ を
$A[x_1, \ldots, x_n]$-加群とみなす。
Lemma \ref{flat-lemma-finite-type-flat-algebra} を適用すれば証明が完了する。
\end{proof}

\begin{lemma}
\label{flat-lemma-finite-type-flat-over-integral}
$f : X \to S$ を局所有限型かつ平坦なスキームの射とする。
$S$ が整ならば、$f$ は局所有限表示である。
\end{lemma}

\begin{proof}
Lemma \ref{flat-lemma-finite-type-flat-X} の特別な場合である。
\end{proof}

\begin{proposition}
\label{flat-proposition-flat-finite-type-finite-presentation-domain}
$R$ を整域とする。$R \to S$ を有限型環準同型とする。
$M$ を有限 $S$-加群とする。
\begin{enumerate}
\item $S$ が $R$ 上平坦ならば、$S$ は有限表示 $R$-代数である。
\item $M$ が $R$-加群として平坦ならば、$M$ は有限表示 $S$-加群である。
\end{enumerate}
\end{proposition}

\begin{proof}
(1) は Lemma \ref{flat-lemma-finite-type-flat-over-integral} の特別な場合である。
(2) について、全射 $R[x_1, \ldots, x_n] \to S$ を選ぶ。
Lemma \ref{flat-lemma-finite-type-flat-algebra} により、$M$ は有限表示
$R[x_1, \ldots, x_n]$-加群である。Algebra, Lemma
\ref{algebra-lemma-finitely-presented-over-subring} により結論を得る。
\end{proof}

\begin{lemma}[Theorem \ref{flat-theorem-finite-type-flat} の有限型版]
\label{flat-lemma-finite-type-flat}
$f : X \to S$ をスキームの射とする。
$\mathcal{F}$ を準連接 $\mathcal{O}_X$-加群とする。次を仮定する。
\begin{enumerate}
\item $X \to S$ は局所有限型である。
\item $\mathcal{F}$ は有限型 $\mathcal{O}_X$-加群である。
\item $S$ の弱随伴点の集合は $S$ において局所有限である。
\end{enumerate}
このとき $U = \{x \in X \mid \mathcal{F}\text{ は }S\text{ 上 }x\text{ で平坦}\}$
は $X$ において開であり、$\mathcal{F}|_U$ は $S$ 上平坦かつ
$S$ に関して局所有限表示である（More on Morphisms, Definition
\ref{more-morphisms-definition-relatively-finitely-presented-sheaf} を参照）。
\end{lemma}

\begin{proof}
問題は $X$ と $S$ 上局所的である。したがって $X$ と $S$ は
アフィンであると仮定してよい。このとき閉埋め込み
$i : X \to \mathbf{A}^n_S$ を選べる。
$X' = \mathbf{A}^n_S \to S$ と有限型準連接加群
$\mathcal{F}' = i_*\mathcal{F}$ に
Theorem \ref{flat-theorem-finite-type-flat} を適用すると、
$$
U' = \{x' \in X' \mid \mathcal{F}'\text{ は }S\text{ 上 }x'\text{ で平坦}\}
$$
は $X'$ において開であり、$\mathcal{F}'|_{U'}$ は有限表示であることが分かる。
$\mathcal{F}'$ は $X' \setminus i(X)$ 上で零に制限され、かつすべての
$x \in X$ に対して $\mathcal{F}'_{i(x)} \cong \mathcal{F}_x$ なので、
$$
U' = i(U) \amalg (X' \setminus i(X))
$$
である。したがって $U = i^{-1}(U')$ は開である。さらに明らかに
$\mathcal{F}'|_{U'} = (i|_U)_*(\mathcal{F}|_U)$.
である。よって More on Morphisms, Lemmas
\ref{more-morphisms-lemma-relative-finite-presentation} および
\ref{more-morphisms-lemma-finite-morphism-relative-finite-presentation} により、
$\mathcal{F}|_U$ は $S$ に関して有限表示である。
\end{proof}

\begin{lemma}
\label{flat-lemma-finite-type-flat-algebra-bis}
$R \to S$ を有限型環準同型とする。$M$ を有限 $S$-加群とする。
$\text{WeakAss}_R(R)$ は有限であると仮定する。このとき
$$
U = \{\mathfrak q \subset S \mid M_{\mathfrak q}\text{ は }R\text{ 上平坦}\}
$$
は $\Spec(S)$ において開である。また、$D(g) \subset U$ となる
すべての $g \in S$ に対して、局所化 $M_g$ は $R$ 上平坦で、
$R$ に関して有限表示な $S_g$-加群である（More on Algebra, Definition
\ref{more-algebra-definition-relatively-finitely-presented} を参照）。
\end{lemma}

\begin{proof}
これは Lemma \ref{flat-lemma-finite-type-flat} を代数へ翻訳したものである。
\end{proof}








\section{相対的純粋加群の例}
\label{flat-section-examples-pure-modules}

\noindent
% P05-EN-CH38-0158: uniquely determined missing-demonstrative repair; see sidecar.
この短い節では、後に導入する {\it 相対的純粋加群} の概念と
{\it 非純性} の概念の動機となる結果の例をいくつか論じる。
各例は補題として述べる。随伴素イデアルに関する条件と、次の補題に現れる条件との
類似性に注意せよ。
Lemmas \ref{flat-lemma-base-change-universally-flat},
\ref{flat-lemma-universally-injective-to-completion},
\ref{flat-lemma-universally-injective-to-completion-flat}、および
\ref{flat-lemma-flat-pure-over-complete-projective}.
この点の議論については Algebra,
Lemma \ref{algebra-lemma-compare-relative-assassins} も参照せよ。

\begin{lemma}
\label{flat-lemma-explain-why-pure}
$R$ を極大イデアル $\mathfrak m$ をもつ局所環とする。
$R \to S$ を環準同型とし、$N$ を $S$-加群とする。次を仮定する。
\begin{enumerate}
\item $N$ は $R$-加群として射影的である。
\item $S/\mathfrak mS$ はネーター環であり、$N/\mathfrak mN$ は有限
$S/\mathfrak mS$-加群である。
\end{enumerate}
このとき、$\mathfrak p = R \cap \mathfrak q$ としたときに
$N \otimes_R \kappa(\mathfrak p)$ の随伴素イデアルとなる任意の素イデアル
$\mathfrak q \subset S$ に対して、$\mathfrak q + \mathfrak m S \not = S$ である。
\end{lemma}

\begin{proof}
Lemmas \ref{flat-lemma-homothety-spectrum} および
\ref{flat-lemma-invert-universally-injective} の仮定が満たされることに注意する。
以下ではこれらの補題の結論を断りなく用いる。
$N/\mathfrak mN$ 上の零因子でない元からなる乗法的集合を
$\Sigma \subset S$ とする。写像 $N \to \Sigma^{-1}N$ は
$R$ 上普遍単射である。したがって、
$N \otimes_R \kappa(\mathfrak p)$ の随伴素イデアルである任意の
$\mathfrak q \subset S$ は、
$\Sigma^{-1}N \otimes_R \kappa(\mathfrak p)$ の随伴素イデアルでもある。
これは明らかに、$\mathfrak q$ が $\Sigma^{-1}S$ の素イデアルに対応することを
意味する。よって $\mathfrak q \subset \mathfrak q'$ となる
$\mathfrak q'$ が存在し、それは $N/\mathfrak mN$ の随伴素イデアルに対応する。
これで結論を得る。
\end{proof}

\noindent
次の補題は、この現象のもう一つの（少しばかり自明な）例を与える。

\begin{lemma}
\label{flat-lemma-explain-why-pure-complete}
$R$ を環とし、$I \subset R$ をイデアルとする。
$R \to S$ を環準同型とし、$N$ を $S$-加群とする。
$N$ が $I$ 進完備ならば、任意の $R$-加群 $M$ と、
$N \otimes_R M$ の随伴素イデアルである任意の素イデアル
$\mathfrak q \subset S$ に対して $\mathfrak q + I S \not = S$ である。
\end{lemma}

\begin{proof}
$S$ の $I$ 進完備化を $S^\wedge$ と書く。
$N$ は $S^\wedge$-加群であり、したがって
$N \otimes_R M$ も $S^\wedge$-加群であることに注意する。
$\mathfrak q = \text{Ann}_S(z)$ となる元 $z \in N \otimes_R M$ をとる。
$z \not = 0$ なので $\text{Ann}_{S^\wedge}(z) \not = S^\wedge$ である。
したがって $\mathfrak q S^\wedge \not = S^\wedge$ である。ゆえに
$\mathfrak q S^\wedge \subset \mathfrak m$ となる極大イデアル
$\mathfrak m \subset S^\wedge$ が存在する。Algebra,
Lemma \ref{algebra-lemma-radical-completion} により
$IS^\wedge \subset \mathfrak m$ なので、結論を得る。
\end{proof}

\noindent
局所環上の射影加群は自由であるから、次の補題は
Lemma \ref{flat-lemma-explain-why-pure} の別証明を与えることに注意する。
Algebra, Theorem \ref{algebra-theorem-projective-free-over-local-ring} を参照せよ。

\begin{lemma}
\label{flat-lemma-explain-why-pure-direct-sum-finite-modules}
$R$ を極大イデアル $\mathfrak m$ をもつ局所環とする。
$R \to S$ を環準同型とし、$N$ を $S$-加群とする。
$N$ は $R$-加群として有限 $R$-加群の直和と同型であると仮定する。
このとき、任意の $R$-加群 $M$ と、$N \otimes_R M$ の随伴素イデアルである
任意の素イデアル $\mathfrak q \subset S$ に対して
$\mathfrak q + \mathfrak m S \not = S$ である。
\end{lemma}

\begin{proof}
各 $M_i$ が有限 $R$-加群となるように
$N = \bigoplus_{i \in I} M_i$ と書く。$M$ を $R$-加群とし、
$\mathfrak q \subset S$ を $N \otimes_R M$ の随伴素イデアルであって
$\mathfrak q + \mathfrak m S = S$ となるものとする。
$\mathfrak q = \text{Ann}_S(z)$ となる元 $z \in N \otimes_R M$ をとる。
直和分解を少し変更すれば、ある元 $1 \in I$ に対して
$z \in M_1 \otimes_R M$ であると仮定してよい。
ある $f \in \mathfrak q$、$x_j \in \mathfrak m$、および $g_j \in S$ を用いて
$1 = f + \sum x_j g_j$ と書く。任意の $g \in S$ に対して、
% P05-EN-CH38-0159: uniquely determined missing-preposition repair; see sidecar.
次の $R$-線形写像を $g'$ と書く。
$$
M_1 \to N \xrightarrow{g} N \to M_1
$$
ここで第一の矢印は包含写像、第二の矢印は $g$ による乗法、
第三の矢印は射影写像である。各 $x_j \in R$ なので、等式
$$
f' + \sum x_j g'_j = \text{id}_{M_1} \in \text{End}_R(M_1)
$$
を得る。中山の補題（Algebra, Lemma \ref{algebra-lemma-NAK}）により
$f'$ は全射である。したがって Algebra, Lemma \ref{algebra-lemma-fun} により
$f'$ は同型である。特に写像
$$
M_1 \otimes_R M \to N \otimes_R M \xrightarrow{f} N \otimes_R M
\to M_1 \otimes_R M
$$
は同型である。これは $fz = 0$ という仮定に矛盾する。
\end{proof}

\begin{lemma}
\label{flat-lemma-explain-why-pure-ML}
$R$ を極大イデアル $\mathfrak m$ をもつ Hensel 局所環とする。
$R \to S$ を環準同型とし、$N$ を $S$-加群とする。
$N$ は $R$-加群として可算生成かつ Mittag--Leffler であると仮定する。
このとき、任意の $R$-加群 $M$ と、$N \otimes_R M$ の随伴素イデアルである
任意の素イデアル $\mathfrak q \subset S$ に対して
$\mathfrak q + \mathfrak m S \not = S$ である。
\end{lemma}

\begin{proof}
Algebra, Lemma \ref{algebra-lemma-split-ML-henselian} により、この補題は
Lemma \ref{flat-lemma-explain-why-pure-direct-sum-finite-modules} に帰着する。
\end{proof}

\noindent
$f : X \to S$ をスキームの射とし、$\mathcal{F}$ を $X$ 上の
準連接加群とする。$\xi \in \text{Ass}_{X/S}(\mathcal{F})$ とし、
$Z = \overline{\{\xi\}}$ とする。次の図を考える。
$$
\xymatrix{
\xi \ar@{|->}[d] & Z \ar[r] \ar[rd] & X \ar[d]^f \\
f(\xi) & & S
}
$$
$f(Z) \subset \overline{\{f(\xi)\}}$ であり、$f(Z)$ が閉であることと
等号が成り立つこと、すなわち $f(Z) = \overline{\{f(\xi)\}}$ であることは
同値である。Lemma \ref{flat-lemma-explain-why-pure} により、$S$, $X$ がアフィンで、
ファイバー $X_s$ がネーター的であり、$\mathcal{F}$ が有限型であり、かつ
$\Gamma(X, \mathcal{F})$ が射影 $\Gamma(S, \mathcal{O}_S)$-加群ならば、
$f(Z) = \overline{\{f(\xi)\}}$ は閉部分集合である。
% P05-EN-CH38-0160: uniquely determined subject-agreement repair; see sidecar.
次の補題で記述される場合にも、少し異なる類似の主張が成り立つ。
Lemmas \ref{flat-lemma-explain-why-pure-complete},
\ref{flat-lemma-explain-why-pure-direct-sum-finite-modules}、および
\ref{flat-lemma-explain-why-pure-ML}.




\section{非純性}
\label{flat-section-impure}

\noindent
Section \ref{flat-section-examples-pure-modules} で例示した現象を、
$\text{Ass}_{X/S}(\mathcal{F})$ の点の特殊化を用いて形式化したい。
また、点 $s \in S$ の周りで局所的に議論したい。そのため次の定義を置く。

\begin{situation}
\label{flat-situation-pre-pure}
ここで $S$, $X$ はスキームであり、$f : X \to S$ は有限型な射である。
また、$\mathcal{F}$ は有限型準連接 $\mathcal{O}_X$-加群である。
最後に $s$ は $S$ の点である。
\end{situation}

\noindent
この状況で、射 $g : T \to S$、$g(t) = s$ となる点 $t \in T$、
特殊化 $t' \leadsto t$、および $t'$ の上にある $X$ の基底変換内の点
$\xi \in X_T$ を考える。図示すると
\begin{equation}
\label{flat-equation-impurity}
\vcenter{
\xymatrix{
\xi \ar@{|->}[d] & \\
t' \ar@{~>}[r] & t \ar@{|->}[r] & s
}
}
\quad\quad
\vcenter{
\xymatrix{
X_T \ar[d] \ar[r] & X \ar[d] \\
T \ar[r]^g & S
}
}
\end{equation}
% P05-EN-CH38-0161: uniquely determined missing-preposition repair; see sidecar.
さらに、$\mathcal{F}$ の $X_T$ への引き戻しを $\mathcal{F}_T$ と書く。

\begin{definition}
\label{flat-definition-impurity}
Situation \ref{flat-situation-pre-pure} において、
$\xi \in \text{Ass}_{X_T/T}(\mathcal{F}_T)$ かつ
$\overline{\{\xi\}} \cap X_t = \emptyset$ であるとき、図式
(\ref{flat-equation-impurity}) は
{\it $s$ 上の $\mathcal{F}$ の非純性}を定めるという。
これは「$(g : T \to S, t' \leadsto t, \xi)$ を
$s$ 上の $\mathcal{F}$ の非純性とする」と表す。
\end{definition}

\begin{lemma}
\label{flat-lemma-impure-finite-presentation}
Situation \ref{flat-situation-pre-pure} において、$s$ 上の $\mathcal{F}$ の
非純性が存在するならば、$g$ が局所有限表示であり、かつ
% P05-EN-CH38-0162: uniquely determined missing-verb repair; see sidecar.
$t$ が $s$ 上の $g$ のファイバーの閉点となるような、$s$ 上の
$\mathcal{F}$ の非純性 $(g : T \to S, t' \leadsto t, \xi)$ が存在する。
\end{lemma}

\begin{proof}
$(g : T \to S, t' \leadsto t, \xi)$ を $s$ 上の $\mathcal{F}$ の
任意の非純性とする。$t \in T$ と $Z = \overline{\{\xi\}}$ に
Limits, Lemma \ref{limits-lemma-separate} を適用すると、
$t$ の開近傍 $V \subset T$、可換図式
$$
\xymatrix{
V \ar[d] \ar[r]_a & T' \ar[d]^b \\
T \ar[r]^g & S,
}
$$
および、次を満たす閉部分スキーム $Z' \subset X_{T'}$ を得る。
\begin{enumerate}
\item 射 $b : T' \to S$ は局所有限表示である。
\item $Z' \cap X_{a(t)} = \emptyset$ である。
\item $Z \cap X_V$ は射 $X_V \to X_{T'}$ によって $Z'$ 内へ写る。
\end{enumerate}
$t'$ は $t$ に特殊化するので、$T$ を $t$ の開近傍 $V$ で置き換えてよい。
したがって次の可換図式を得る。
$$
\xymatrix{
X_T \ar[d] \ar[r] &
X_{T'} \ar[d] \ar[r] &
X \ar[d] \\
T \ar[r]^a & T' \ar[r]^b & S
}
$$
ここで $b \circ a = g$ である。$\xi$ の像を $\xi' \in X_{T'}$ と書く。
Divisors, Lemma \ref{divisors-lemma-base-change-relative-assassin} により
$\xi' \in \text{Ass}_{X_{T'}/T'}(\mathcal{F}_{T'})$ である。
% P05-EN-CH38-0163: uniquely determined redundant-closure repair; see sidecar.
さらに構成により、$\overline{\{\xi'\}}$ は、ファイバー $X_{a(t)}$ と
交わらない閉部分集合 $Z'$ に含まれる。したがって
$(T' \to S, a(t') \leadsto a(t), \xi')$ は
$s$ 上の $\mathcal{F}$ の非純性である。

\medskip\noindent
以上により $g : T \to S$ は局所有限表示であると仮定してよい。
$Z = \overline{\{\xi\}}$ とおく。仮定により $Z_t = \emptyset$ である。
More on Morphisms, Lemma \ref{more-morphisms-lemma-empty-generic-fibre} により、
これは $\overline{\{t\}}$ のある開部分集合内の $t''$ に対して
$Z_{t''} = \emptyset$ となることを意味する。$s$ 上の $T \to S$ の
ファイバーは Jacobson スキームなので（Morphisms,
Lemma \ref{morphisms-lemma-ubiquity-Jacobson-schemes}）、
% P05-EN-CH38-0164: uniquely determined subject-agreement repair; see sidecar.
$Z_{t''} = \emptyset$ となる閉点 $t'' \in \overline{\{t\}}$ が存在する。
このとき $(g : T \to S, t' \leadsto t'', \xi)$ が所望の非純性である。
\end{proof}

\begin{lemma}
\label{flat-lemma-impure-limit}
Situation \ref{flat-situation-pre-pure} において、
$(g : T \to S, t' \leadsto t, \xi)$ を $s$ 上の $\mathcal{F}$ の
非純性とする。$T = \lim_{i \in I} T_i$ は $S$ 上のアフィンスキームの
有向極限であると仮定する。このとき、ある $i$ に対して三つ組
$(T_i \to S, t'_i \leadsto t_i, \xi_i)$ は
$s$ 上の $\mathcal{F}$ の非純性である。
\end{lemma}

\begin{proof}
主張の記法は次の意味である。射影射を $p_i : T \to T_i$ とし、
$t_i = p_i(t)$ および $t'_i = p_i(t')$ とする。最後に
$\xi_i \in X_{T_i}$ は $\xi$ の像である。Divisors,
Lemma \ref{divisors-lemma-base-change-relative-assassin} により、
$\xi_i$ は $T_i$ 上の $\mathcal{F}_{T_i}$ の相対随伴点である。
したがって、ある $i$ に対して
$\overline{\{\xi_i\}} \cap X_{t_i} = \emptyset$ であることだけを示せばよい。

\medskip\noindent
第一の証明。$Z_i = \overline{\{\xi_i\}} \subset X_{T_i}$ および
$Z = \overline{\{\xi\}} \subset X_T$ に被約誘導スキーム構造を入れる。
Limits, Lemma \ref{limits-lemma-inverse-limit-irreducibles} により
$Z = \lim Z_i$ である。体 $k$ と、像が $t$ となる射
$\Spec(k) \to T$ を選ぶ。このとき
$$
\emptyset =
Z \times_T \Spec(k) = (\lim Z_i) \times_{(\lim T_i)} \Spec(k)
= \lim Z_i \times_{T_i} \Spec(k)
$$
である。これは極限がファイバー積と可換だからである（極限は極限と可換する）。
$X_{T_i} \to T_i$ は有限型であり、したがって $Z_i \to T_i$ も有限型なので、
各 $Z_i \times_{T_i} \Spec(k)$ は準コンパクトである。よって Limits,
Lemma \ref{limits-lemma-limit-nonempty} により、ある $i$ に対して
$Z_i \times_{T_i} \Spec(k)$ は空である。合成
$\Spec(k) \to T \to T_i$ の像は $t_i$ なので、所望の結論を得る。

\medskip\noindent
第二の証明。$Z = \overline{\{\xi\}}$ とおく。この状況に
Limits, Lemma \ref{limits-lemma-separate} を適用すると、
$t$ の開近傍 $V \subset T$、可換図式
$$
\xymatrix{
V \ar[d] \ar[r]_a & T' \ar[d]^b \\
T \ar[r]^g & S,
}
$$
および、次を満たす閉部分スキーム $Z' \subset X_{T'}$ を得る。
\begin{enumerate}
\item 射 $b : T' \to S$ は局所有限表示である。
\item $Z' \cap X_{a(t)} = \emptyset$ である。
\item $Z \cap X_V$ は射 $X_V \to X_{T'}$ によって $Z'$ 内へ写る。
\end{enumerate}
$V$ は $T$ のアフィン開部分であると仮定してよい。したがって Limits,
Lemmas \ref{limits-lemma-descend-opens} および
\ref{limits-lemma-limit-affine} により、ある $i$ とアフィン開部分
$V_i \subset T_i$ を、
% P05-EN-CH38-0165: source-faithful undefined-projection notation; see sidecar.
$V = f_i^{-1}(V_i)$ となるように見いだせる。Limits,
Proposition \ref{limits-proposition-characterize-locally-finite-presentation} により、
必要なら $i$ を少し大きくした後、$a = a_i \circ f_i|_V$ となる射
$a_i : V_i \to T'$ を見いだせる。誘導される射
$X_{V_i} \to X_{T'}$ は $\xi_i$ を $Z'$ 内へ写す。
$Z' \cap X_{a(t)} = \emptyset$ なので、
$(T_i \to S, t'_i \leadsto t_i, \xi_i)$ は
$s$ 上の $\mathcal{F}$ の非純性である。
\end{proof}

\begin{lemma}
\label{flat-lemma-quasi-finite-impurity-elementary}
Situation \ref{flat-situation-pre-pure} において、$g$ が $t$ において
準有限となる $s$ 上の $\mathcal{F}$ の非純性
$(g : T \to S, t' \leadsto t, \xi)$ が存在するならば、
$(T, t) \to (S, s)$ が基本エタール近傍となるような非純性
$(g : T \to S, t' \leadsto t, \xi)$ が存在する。
\end{lemma}

\begin{proof}
$(g : T \to S, t' \leadsto t, \xi)$ を、$g$ が $t$ において準有限となる
$s$ 上の $\mathcal{F}$ の非純性とする。$T$ を縮小した後、
$g$ は局所有限型であると仮定してよい。$T \to S$ と $t \mapsto s$ に
More on Morphisms,
Lemma \ref{more-morphisms-lemma-etale-makes-quasi-finite-finite-at-point}
を適用すると、次の図式を得る。
$$
\xymatrix{
T \ar[d] & T \times_S U \ar[l] \ar[d] & V \ar[l] \ar[ld] \\
S & U \ar[l]
}
$$
ここで $(U, u) \to (S, s)$ は基本エタール近傍であり、
$V \subset T \times_S U$ は $v = (t, u)$ の開近傍であって、
$V \to U$ は有限であり、$v$ は $u$ の上にある $V$ の一意な点である。
射 $V \to T$ はエタール、したがって平坦なので、$v' \mapsto t'$ となる
特殊化 $v' \leadsto v$ が存在する。$\kappa(t') \subset \kappa(v')$ は
有限分離拡大であることに注意する。$\xi \in X_{t'}$ に写る任意の点
$\zeta \in X_{v'}$ を選ぶ。Divisors,
Lemma \ref{divisors-lemma-base-change-relative-assassin} により
$\zeta \in \text{Ass}_{X_V/V}(\mathcal{F}_V)$ である。さらに仮定により
閉包 $\overline{\{\xi\}}$ は $X_t$ と交わらないので、閉包
$\overline{\{\zeta\}}$ はファイバー $X_v$ と交わらない。言い換えると、
% P05-EN-CH38-0166: source-faithful base-point/scheme mismatch; see sidecar.
$(V \to S, v' \leadsto v, \zeta)$ は $S$ 上の $\mathcal{F}$ の非純性である。

\medskip\noindent
次に、$v'$ の像を
% P05-EN-CH38-0167: source-faithful undefined U' notation; see sidecar.
$u' \in U'$ とし、$\zeta$ の像を $\theta \in X_U$ とする。
このとき $\theta \mapsto u'$ かつ $u' \leadsto u$ である。
Divisors, Lemma \ref{divisors-lemma-base-change-relative-assassin} により
% P05-EN-CH38-0168: source-faithful missing base-change subscript; see sidecar.
$\theta \in \text{Ass}_{X_U/U}(\mathcal{F})$ である。さらに
$\pi : X_V \to X_U$ は有限なので、
$\pi\big(\overline{\{\zeta\}}\big) = \overline{\{\pi(\zeta)\}}$ である。
$v$ は $u$ の上にある $V$ の一意な点なので、
$X_v \cap \overline{\{\zeta\}} = \emptyset$ から
$X_u \cap \overline{\{\pi(\zeta)\}} = \emptyset$ が従う。したがって
$(U \to S, u' \leadsto u, \theta)$ は $s$ 上の $\mathcal{F}$ の非純性であり、
結論を得る。
\end{proof}

\begin{lemma}
\label{flat-lemma-Noetherian-impurity-quasi-finite}
Situation \ref{flat-situation-pre-pure} において、$S$ は局所ネーター的であると
仮定する。$s$ 上の $\mathcal{F}$ の非純性が存在するならば、
$g$ が $t$ において準有限となるような $s$ 上の $\mathcal{F}$ の非純性
$(g : T \to S, t' \leadsto t, \xi)$ が存在する。
\end{lemma}

\begin{proof}
$S$ を $s$ のアフィン近傍で置き換えてよい。
Lemma \ref{flat-lemma-impure-finite-presentation} により、
% P05-EN-CH38-0169: uniquely determined corrupted-clause repair; see sidecar.
$g$ は局所有限型であり、$t$ は $s$ 上の $g$ のファイバーの閉点となるような
非純性 $(g : T \to S, t' \leadsto t, \xi)$ をもつと仮定してよい。
$T$ を $\overline{\{t'\}}$ 上の被約誘導スキーム構造で置き換えてよい。
$Z = \overline{\{\xi\}} \subset X_T$ とおく。仮定により
$Z_t = \emptyset$ であり、$Z \to T$ の像は $t'$ を含む。
More on Morphisms,
Lemma \ref{more-morphisms-lemma-relative-assassin-in-neighbourhood} により、
任意の $w \in f(V)$ に対して $V_w$ の任意の生成点 $\xi'$ が
$\text{Ass}_{X_T/T}(\mathcal{F}_T)$ に属するような非空開部分
$V \subset Z$ が存在する。More on Morphisms,
Lemma \ref{more-morphisms-lemma-nonempty-generic-fibre} により、
$W \subset f(V)$ となる非空開部分 $W \subset T$ が存在する。
More on Morphisms, Lemma
\ref{more-morphisms-lemma-quasi-finite-quasi-section-meeting-nearby-open-X} により、
$t \in T'$ であり、$T' \to S$ は $t$ において準有限であり、かつ
$s$ へ写らない点 $z \in T' \cap W$, $z \leadsto t$ が存在するような
閉部分スキーム $T' \subset T$ が存在する。非空スキーム $V_z$ の
任意の生成点 $\xi'$ を選ぶ。このとき
$(T' \to S, z \leadsto t, \xi')$ が所望の非純性である。
\end{proof}

\noindent
以下では、$s$ における $S$ の Hensel 化
$S^h = \Spec(\mathcal{O}_{S, s}^h)$
を用いる。\'Etale Cohomology,
Definition \ref{etale-cohomology-definition-etale-local-rings} を参照せよ。
% P05-EN-CH38-0170: uniquely determined article/word-order repair; see sidecar.
$S^h \to S$ は $S^h$ の閉点を $s$ へ写し、剰余体の同型を誘導するので、
% P05-EN-CH38-0171: uniquely determined denote-construction repair; see sidecar.
この閉点も $s \in S^h$ と表す。したがって
$(S^h, s) \to (S, s)$ は点付きスキームの射である。

\begin{lemma}
\label{flat-lemma-impurity-on-henselization}
Situation \ref{flat-situation-pre-pure} において、$s$ 上の $\mathcal{F}$ の非純性
$(S^h \to S, s' \leadsto s, \xi)$ が存在するならば、
$(T, t) \to (S, s)$ が基本エタール近傍となるような、$s$ 上の
$\mathcal{F}$ の非純性 $(T \to S, t' \leadsto t, \xi)$ が存在する。
\end{lemma}

\begin{proof}
$S$ を $s$ のアフィン近傍で置き換えてよい。
$S = \Spec(A)$ とし、$s$ は素イデアル $\mathfrak p \subset A$ に
対応するとする。このとき
$\mathcal{O}_{S, s}^h = \colim_{(T, t)} \Gamma(T, \mathcal{O}_T)$
である。ここで極限は、$(S, s)$ のアフィン基本エタール近傍
$(T, t)$ からなる余フィルター圏の反対圏上でとる。More on Morphisms,
Lemma \ref{more-morphisms-lemma-describe-henselization} とその証明を参照せよ。
したがって $S^h = \lim_i T_i$ であり、Lemma \ref{flat-lemma-impure-limit} により
結論を得る。
\end{proof}

\begin{lemma}
\label{flat-lemma-pure-along-X-s}
Situation \ref{flat-situation-pre-pure} において、次は同値である。
\begin{enumerate}
\item $S^h$ を $s$ における $S$ の Hensel 化とするとき、
$s$ 上の $\mathcal{F}$ の非純性 $(S^h \to S, s' \leadsto s, \xi)$ が存在する。
\item $(T, t) \to (S, s)$ が基本エタール近傍となるような、$s$ 上の
$\mathcal{F}$ の非純性 $(T \to S, t' \leadsto t, \xi)$ が存在する。
\item $T \to S$ が $t$ において準有限となるような、$s$ 上の
$\mathcal{F}$ の非純性 $(T \to S, t' \leadsto t, \xi)$ が存在する。
\end{enumerate}
\end{lemma}

\begin{proof}
エタール射は局所準有限なので、(2) が (3) を含意することは明らかである。
Lemma \ref{flat-lemma-quasi-finite-impurity-elementary} で (3) が (2) を含意することを
示した。Lemma \ref{flat-lemma-impurity-on-henselization} で (1) が (2) を
含意することを示した。最後に、$(T \to S, t' \leadsto t, \xi)$ を
$s$ 上の $\mathcal{F}$ の非純性であって、
$(T, t) \to (S, s)$ が基本エタール近傍となるものとする。このとき構造射
$S^h \to S$ の分解 $S^h \to T \to S$ を選べる。
$t'$ に写る任意の点 $s' \in S^h$ と、$\xi \in X_{t'}$ に写る任意の
$\xi' \in X_{s'}$ を選ぶ。このとき
$(S^h \to S, s' \leadsto s, \xi')$ は $s$ 上の $\mathcal{F}$ の非純性である。
% P05-EN-CH38-0172: source-faithful explicit proof-detail omission; see sidecar.
詳細は省略する。
\end{proof}





\section{相対的純粋加群}
\label{flat-section-pure}

\noindent
基底に関して純粋な加群という概念は \cite{GruRay} で導入された。

\begin{definition}
\label{flat-definition-pure}
$f : X \to S$ を有限型なスキームの射とする。
$\mathcal{F}$ を有限型準連接 $\mathcal{O}_X$-加群とする。
\begin{enumerate}
\item $s \in S$ とする。$(T, t) \to (S, s)$ が基本エタール近傍となる
$s$ 上の $\mathcal{F}$ の非純性 $(g : T \to S, t' \leadsto t, \xi)$ が
存在しないとき、$\mathcal{F}$ は {\it $X_s$ に沿って純粋}であるという。
\item $s$ 上の $\mathcal{F}$ の非純性が一つも存在しないとき、
$\mathcal{F}$ は {\it $X_s$ に沿って普遍的に純粋}であるという。
\item $\mathcal{O}_X$ が $X_s$ に沿って純粋であるとき、$X$ は
{\it $X_s$ に沿って純粋}であるという。
\item すべての $s \in S$ に対して $\mathcal{F}$ が $X_s$ に沿って
普遍的に純粋であるとき、$\mathcal{F}$ は {\it 普遍的 $S$-純粋}、または
{\it $S$ に関して普遍的に純粋}であるという。
\item すべての $s \in S$ に対して $\mathcal{F}$ が $X_s$ に沿って
純粋であるとき、$\mathcal{F}$ は {\it $S$-純粋}、または
{\it $S$ に関して純粋}であるという。
\item $\mathcal{O}_X$ が $S$ に関して純粋であるとき、$X$ は
{\it $S$-純粋}、または {\it $S$ に関して純粋}であるという。
\end{enumerate}
\end{definition}

\noindent
ここでは意図的に、単に局所有限型な射ではなく有限型な射に限定している。
議論については Remark \ref{flat-remark-discuss-finite-type} を参照せよ。
定義の状況では、Lemma \ref{flat-lemma-pure-along-X-s} により次は同値である。
\begin{enumerate}
\item $\mathcal{F}$ は $X_s$ に沿って純粋である。
\item $g$ が $t$ において準有限となる非純性
$(g : T \to S, t' \leadsto t, \xi)$ は存在しない。
\item $S^h$ を $s$ における $S$ の Hensel 化とするとき、
$(S^h \to S, s' \leadsto s, \xi)$ という形の非純性は存在しない。
\end{enumerate}
$X^h = X \times_S S^h$ と書き、$\mathcal{F}$ の $X^h$ への引き戻しを
$\mathcal{F}^h$ と書けば、最後の条件を次のより肯定的な形で述べられる。
\begin{enumerate}
\item[(4)] $\text{Ass}_{X^h/S^h}(\mathcal{F}^h)$ のすべての点は
$X_s$ の点へ特殊化する。
\end{enumerate}
特に、$\mathcal{F}$ が $X_s$ に沿って純粋であることと、$\mathcal{F}$ の
$X \times_S \Spec(\mathcal{O}_{S, s})$ への引き戻しが
$X_s$ に沿って純粋であることは同値であることが明らかである。

\begin{remark}
\label{flat-remark-discuss-finite-type}
$f : X \to S$ を局所有限型な射とし、$\mathcal{F}$ を有限型準連接
$\mathcal{O}_X$-加群とする。この場合にも上の (1) と (2) は同値である。
実際、Lemma \ref{flat-lemma-quasi-finite-impurity-elementary} の証明では
$f$ が準コンパクトであることを用いない。また (3) と (4) が同値であることも
明らかである。しかし (1) と (3) が同値であるかどうかは分からない。
この場合、\cite{GruRay} で行われているように、同値な条件 (3) と (4) を用いて
純粋性を定義する方が便利なこともある。一方、多くの応用では、真に正しい概念は
普遍的に純粋であることだと思われる。
\end{remark}

\noindent
自然に生じる問いは、基底に関して純粋であるという性質が基底変換で保たれるか、
すなわち純粋であることと普遍的に純粋であることが同じか、というものである。
ネーター基底スキーム上ではこれは成り立つ
（Lemma \ref{flat-lemma-Noetherian-base-change}）。層が平坦な場合にも成り立つ
（Lemmas \ref{flat-lemma-finite-type-flat-pure-along-fibre-is-universal} および
\ref{flat-lemma-finite-type-flat-pure-is-universal}）。射と層が有限表示であっても、
一般には成り立たない。反例については Examples,
Section \ref{examples-section-pure-not-universally} を参照せよ。
% P05-EN-CH38-0177: uniquely determined compound-word repair; see sidecar.
まず、「普遍的」という語のここでの用法を通常の概念と一致させる。

\begin{lemma}
\label{flat-lemma-base-change-universally}
$f : X \to S$ を有限型なスキームの射とする。
$\mathcal{F}$ を有限型準連接 $\mathcal{O}_X$-加群とする。
$s \in S$ とする。次は同値である。
\begin{enumerate}
\item $\mathcal{F}$ は $X_s$ に沿って普遍的に純粋である。
\item 点付きスキームの任意の射 $(S', s') \to (S, s)$ に対して、
引き戻し $\mathcal{F}_{S'}$ は $X_{s'}$ に沿って純粋である。
\end{enumerate}
特に、$\mathcal{F}$ が $S$ に関して普遍的に純粋であることと、
$\mathcal{F}$ のすべての基底変換 $\mathcal{F}_{S'}$ が
$S'$ に関して純粋であることは同値である。
\end{lemma}

\begin{proof}
これは形式的である。
\end{proof}

\begin{lemma}
\label{flat-lemma-quasi-finite-base-change}
$f : X \to S$ を有限型なスキームの射とする。
$\mathcal{F}$ を有限型準連接 $\mathcal{O}_X$-加群とする。
$s \in S$ とし、$(S', s') \to (S, s)$ を点付きスキームの射とする。
$S' \to S$ が $s'$ において準有限であり、$\mathcal{F}$ が
$X_s$ に沿って純粋ならば、$\mathcal{F}_{S'}$ は
$X_{s'}$ に沿って純粋である。
\end{lemma}

\begin{proof}
% P05-EN-CH38-0173: uniquely determined spelling repair; see sidecar.
$(T \to S', t' \leadsto t, \xi)$ が $s'$ 上の $\mathcal{F}_{S'}$ の
非純性であり、$T \to S'$ が $t$ において準有限ならば、
% P05-EN-CH38-0174: source-faithful specialization-arrow mismatch; see sidecar.
$(T \to S, t' \to t, \xi)$ は $s$ 上の $\mathcal{F}$ の非純性であり、
$T \to S$ は $t$ において準有限である。Morphisms,
Lemma \ref{morphisms-lemma-composition-quasi-finite} を参照せよ。
% P05-EN-CH38-0175: uniquely determined missing-preposition repair; see sidecar.
したがって補題は、Definition \ref{flat-definition-pure} に続く議論で与えた
純粋性の特徴づけ (2) から直ちに従う。
\end{proof}

\begin{lemma}
\label{flat-lemma-Noetherian-base-change}
$f : X \to S$ を有限型なスキームの射とする。
$\mathcal{F}$ を有限型準連接 $\mathcal{O}_X$-加群とする。
$s \in S$ とする。$\mathcal{O}_{S, s}$ がネーター環ならば、
$\mathcal{F}$ が $X_s$ に沿って純粋であることと、$\mathcal{F}$ が
$X_s$ に沿って普遍的に純粋であることは同値である。
\end{lemma}

\begin{proof}
まず $S$ を $\Spec(\mathcal{O}_{S, s})$ で置き換えてよい。すなわち
$S$ はネーター的であると仮定してよい。次に
Lemma \ref{flat-lemma-Noetherian-impurity-quasi-finite} と、
% P05-EN-CH38-0176: uniquely determined missing-article repair; see sidecar.
Definition \ref{flat-definition-pure} に続く議論で与えた純粋性の特徴づけ (2) を
用いれば結論を得る。
\end{proof}

\noindent
純粋性は平坦降下を満たす。

\begin{lemma}
\label{flat-lemma-flat-descend-pure}
$f : X \to S$ を有限型なスキームの射とする。
$\mathcal{F}$ を有限型準連接 $\mathcal{O}_X$-加群とする。
$s \in S$ とし、$(S', s') \to (S, s)$ を点付きスキームの射とする。
$S' \to S$ は $s'$ において平坦であると仮定する。
\begin{enumerate}
\item $\mathcal{F}_{S'}$ が $X_{s'}$ に沿って純粋ならば、
$\mathcal{F}$ は $X_s$ に沿って純粋である。
\item $\mathcal{F}_{S'}$ が $X_{s'}$ に沿って普遍的に純粋ならば、
$\mathcal{F}$ は $X_s$ に沿って普遍的に純粋である。
\end{enumerate}
\end{lemma}

\begin{proof}
$(T \to S, t' \leadsto t, \xi)$ を $s$ 上の $\mathcal{F}$ の非純性とする。
$T_1 = T \times_S S'$ とおき、
% P05-EN-CH38-0178: source-faithful false-uniqueness claim; see sidecar.
$t$ と $s'$ に写る $T_1$ の一意な点を $t_1$ とする。
Morphisms, Lemma \ref{morphisms-lemma-base-change-flat} により
$T_1 \to T$ は $t_1$ において平坦なので、Algebra,
Section \ref{algebra-section-going-up} により、$t' \leadsto t$ の上にある
特殊化 $t'_1 \leadsto t_1$ が存在する。
$\Spec(\kappa(t'_1) \otimes_{\kappa(t')} \kappa(\xi))$ の生成点に対応する点
$\xi_1 \in X_{t'_1}$ を選ぶ。Schemes,
Lemma \ref{schemes-lemma-points-fibre-product} を参照せよ。Divisors,
Lemma \ref{divisors-lemma-base-change-relative-assassin} により
$\xi_1 \in \text{Ass}_{X_{T_1}/T_1}(\mathcal{F}_{T_1})$ である。
$X_{T_1}$ における $\{\xi_1\}$ の Zariski 閉包は、$X_T$ における
$\{\xi\}$ の Zariski 閉包内へ写るので、この閉包は $X_{t_1}$ と交わらない。
したがって $(T_1 \to S', t'_1 \leadsto t_1, \xi_1)$ は
$s'$ 上の $\mathcal{F}_{S'}$ の非純性である。言い換えると、
% P05-EN-CH38-0179: uniquely determined wrong-preposition repair; see sidecar.
補題の (2) の対偶を証明した。最後に、$(T, t) \to (S, s)$ が
基本エタール近傍ならば、$(T_1, t_1) \to (S', s')$ も
基本エタール近傍である。このようにして (1) が成り立つことが分かる。
\end{proof}

\begin{lemma}
\label{flat-lemma-supported-on-closed}
$i : Z \to X$ を、スキーム $S$ 上有限型なスキームの閉埋め込みとする。
$s \in S$ とし、$\mathcal{F}$ を $Z$ 上の有限型準連接層とする。
このとき、$\mathcal{F}$ が $Z_s$ に沿って（普遍的に）純粋であることと、
$i_*\mathcal{F}$ が $X_s$ に沿って（普遍的に）純粋であることは同値である。
\end{lemma}

\begin{proof}
これは Divisors,
Lemma \ref{divisors-lemma-relative-weak-assassin-finite} から従う。
\end{proof}








\section{相対的純粋層の例}
\label{flat-section-examples-pure-sheaves}

\noindent
純粋性が何を意味するかを具体的に見て取れる例をいくつか挙げる。

\begin{lemma}
\label{flat-lemma-proper-pure}
$f : X \to S$ を有限型のスキーム射とする。
$\mathcal{F}$ を有限型準連接 $\mathcal{O}_X$-加群とする。
\begin{enumerate}
\item $\mathcal{F}$ の台が $S$ 上固有ならば、$\mathcal{F}$ は
$S$ に関して普遍的に純粋である。
\item $f$ が固有ならば、$\mathcal{F}$ は
$S$ に関して普遍的に純粋である。
\item $f$ が固有ならば、$X$ は $S$ に関して普遍的に純粋である。
\end{enumerate}
\end{lemma}

\begin{proof}
まず (1) を (2) に帰着させる。すなわち、$Z \subset X$ を
$\mathcal{F}$ のスキーム論的台とする。対応する閉埋め込みを
$i : Z \to X$ とし、ある有限型準連接 $\mathcal{O}_Z$-加群
$\mathcal{G}$ を用いて $\mathcal{F} = i_*\mathcal{G}$ と書く。
Morphisms, Section \ref{morphisms-section-support} を参照せよ。
(1) の場合、仮定により $Z \to S$ は固有である。
したがって Lemma \ref{flat-lemma-supported-on-closed} により、(1) は (2) に帰着する。

\medskip\noindent
$f$ が固有であると仮定する。
$(g : T \to S, t' \leadsto t, \xi)$ を、$s \in S$ 上の
$\mathcal{F}$ の非純性とする。$f$ は固有なので普遍閉である。したがって
$f_T : X_T \to T$ は閉写像である。$f_T(\xi) = t'$ であるから、これは
% P05-EN-CH38-0180: source-faithful base-changed-map mismatch; see sidecar.
$t \in f(\overline{\{\xi\}})$ を含意するが、矛盾である。
\end{proof}

\begin{lemma}
\label{flat-lemma-quasi-finite-pure}
$f : X \to S$ を分離有限型のスキーム射とする。
$\mathcal{F}$ を有限型準連接 $\mathcal{O}_X$-加群とする。
すべての $s \in S$ に対して $\text{Supp}(\mathcal{F}_s)$ が有限であると仮定する。
このとき、次は同値である。
\begin{enumerate}
\item $\mathcal{F}$ は $S$ に関して純粋である。
\item $\mathcal{F}$ のスキーム論的台は $S$ 上有限である。
\item $\mathcal{F}$ は $S$ に関して普遍的に純粋である。
\end{enumerate}
特に、準有限分離射 $X \to S$ に対して、$X$ が $S$ に関して
純粋であることと $X \to S$ が有限であることは同値である。
\end{lemma}

\begin{proof}
$Z \subset X$ を $\mathcal{F}$ のスキーム論的台とする。
Morphisms, Definition \ref{morphisms-definition-scheme-theoretic-support}
を参照せよ。このとき $Z \to S$ は有限なファイバーをもつ分離有限型の
スキーム射である。したがってこれは分離かつ準有限である。
Morphisms, Lemma \ref{morphisms-lemma-quasi-finite} を参照せよ。
Lemma \ref{flat-lemma-supported-on-closed}
により、$Z \to S$ と、$Z$ 上の有限型準連接加群とみなした層
$\mathcal{F}$ について補題を証明すれば十分である。したがって
$X \to S$ は分離かつ準有限であり、$\text{Supp}(\mathcal{F}) = X$
であると仮定してよい。

\medskip\noindent
Lemma \ref{flat-lemma-proper-pure}
および
Morphisms, Lemma \ref{morphisms-lemma-finite-proper}
から (2) は (3) を含意する。(3) が (1) を含意することは明らかである。
(1) が成り立つと仮定し、(2) を証明する。$S$ はアフィンであると仮定してよい。
More on Morphisms,
Lemma \ref{more-morphisms-lemma-quasi-finite-separated-pass-through-finite}
により、図式
$$
\xymatrix{
X \ar[rd]_f \ar[rr]_j & & T \ar[ld]^\pi \\
& S &
}
$$
で、$\pi$ が有限、$j$ が準コンパクトな開埋め込みとなるものをとれる。
$j$ が閉であることを示せば、$j$ は閉埋め込みとなり、
$f = \pi \circ j$ が有限であると結論できる。
$j$ が閉であることを示すには、特殊化が $j$ に沿って持ち上がることを
示せば十分である。Schemes, Lemma
\ref{schemes-lemma-quasi-compact-closed} を参照せよ。
$x \in X$ とし、$t' = j(x)$ とおき、$t' \leadsto t$ を特殊化とする。
$t \in j(X)$ を示さなければならない。$s' = f(x)$ および
$s = \pi(t)$ とおくと、$s' \leadsto s$ である。
More on Morphisms, Lemma
\ref{more-morphisms-lemma-etale-splits-off-quasi-finite-part-technical}
により、基本エタール近傍 $(U, u) \to (S, s)$ と、開かつ閉な部分スキームへの分解
$$
T_U = T \times_S U = V \amalg W
$$
で、$V \to U$ が有限であり、$u$ に写る $V$ の点 $v$ が一意に存在し、
しかも $v$ が $T$ の $t$ に写るものをとれる。$V \to T$ はエタールなので、
一般化を持ち上げることができる。Morphisms, Lemmas
\ref{morphisms-lemma-generalizations-lift-flat} および
\ref{morphisms-lemma-etale-flat} を参照せよ。
したがって、特殊化 $v' \leadsto v$ が存在し、$v'$ は $t' \in T$ に写る。
特に $v' \in X_U \subset T_U$ である。
% P05-EN-CH38-0181: source-faithful undefined image of t' in U; see sidecar.
$t'$ の像を $u' \in U$ と記す。
$X_{u'}$ は有限離散集合であり、かつ
$X_{u'} = \text{Supp}(\mathcal{F}_{u'})$ なので、
$v' \in \text{Ass}_{X_U/U}(\mathcal{F})$ であることに注意せよ。
$\mathcal{F}$ は $S$ に関して純粋なので、$v'$ は $X_u$ の点に
特殊化しなければならない。$v$ は $u$ 上にある $V$ の唯一の点であり
（また $W$ の点はどれも $v'$ の特殊化ではありえない）から、
$v \in X_u$ である。したがって $t \in X$ である。
\end{proof}

\begin{lemma}
\label{flat-lemma-flat-geometrically-integral-fibres-pure}
$f : X \to S$ を、幾何学的に整なファイバーをもつ有限型平坦な
スキーム射とする。このとき $X$ は $S$ 上普遍的に純粋である。
\end{lemma}

\begin{proof}
$\xi \in X$ とし、$s' = f(\xi)$ かつ $s' \leadsto s$ を $S$ の特殊化とする。
$\xi$ が $X_{s'}$ の随伴点ならば、$X_{s'}$ は整スキームなので、
$\xi$ はその唯一の生成点である。$\xi_0$ を $X_s$ の唯一の生成点とする。
$X \to S$ は平坦なので、$s' \leadsto s$ を $X$ における特殊化
$\xi' \leadsto \xi_0$ に持ち上げることができる。Morphisms,
Lemma \ref{morphisms-lemma-generalizations-lift-flat} を参照せよ。
% P05-EN-CH38-0182: uniquely determined grammar repair; see sidecar.
$\xi$ は $X_{s'}$ の生成点なので $\xi \leadsto \xi'$ であり、したがって
$\xi \leadsto \xi_0$ である。これは
% P05-EN-CH38-0183: source-faithful specialization-arrow mismatch; see sidecar.
$(\text{id}_S, s' \to s, \xi)$ が $s$ 上の $\mathcal{O}_X$ の
非純性でないことを意味する。$f$ が有限型平坦で幾何学的に整な
ファイバーをもつという仮定は基底変換で保たれるので、任意の基底変換後にも
非純性は存在しない。このようにして、$X$ は普遍的 $S$-純粋である。
\end{proof}

\begin{lemma}
\label{flat-lemma-affine-locally-projective-pure}
$f : X \to S$ を有限型アフィンなスキーム射とする。
$\mathcal{F}$ を有限型準連接 $\mathcal{O}_X$-加群とし、
$f_*\mathcal{F}$ は $S$ 上局所射影的であると仮定する。
Properties, Definition \ref{properties-definition-locally-projective}
を参照せよ。このとき $\mathcal{F}$ は $S$ 上普遍的に純粋である。
\end{lemma}

\begin{proof}
$S$ が Hensel 局所環のスペクトルである場合に帰着した後、これは
Lemma \ref{flat-lemma-explain-why-pure} から従う。
\end{proof}




\section{純粋性の判定条件}
\label{flat-section-criterion-purity}

\noindent
まず、有限型準連接層の平坦族が与えられたとき、相対随伴点が
特殊化するならば、近傍ファイバーの相対随伴点へ特殊化することを証明する。

\begin{lemma}
\label{flat-lemma-associated-point-specializes}
$f : X \to S$ を有限型のスキーム射とする。
$\mathcal{F}$ を有限型準連接 $\mathcal{O}_X$-加群とし、$s \in S$ とする。
$\mathcal{F}$ は $X_s$ のすべての点において $S$ 上平坦であると仮定する。
$f(x') = s'$ を満たす $x' \in \text{Ass}_{X/S}(\mathcal{F})$ をとり、
$s' \leadsto s$ が $S$ における特殊化であるとする。$x'$ が $X_s$ の点に
特殊化するならば、ある $x \in \text{Ass}_{X_s}(\mathcal{F}_s)$ に対して
$x' \leadsto x$ である。
\end{lemma}

\begin{proof}
$t \in X_s$ に対して $x' \leadsto t$ とする。このとき、
$x$ が $\overline{\{x'\}} \cap f^{-1}(\{s\})$ のある既約成分の
生成点に対応するような特殊化 $x' \leadsto x \leadsto t$ をとれる。
仮定により $\mathcal{F}$ は $x$ において $S$ 上平坦である。
More on Morphisms, Lemma
\ref{more-morphisms-lemma-associated-point-specializes}
により、望むとおり $x \in \text{Ass}_{X/S}(\mathcal{F})$ である。
\end{proof}

\begin{lemma}
\label{flat-lemma-criterion}
$f : X \to S$ を有限型のスキーム射とする。
$\mathcal{F}$ を有限型準連接 $\mathcal{O}_X$-加群とし、$s \in S$ とする。
$(S', s') \to (S, s)$ を基本エタール近傍とし、スキーム射の可換図式
$$
\xymatrix{
X \ar[d] & X' \ar[l]^g \ar[d] \\
S & S' \ar[l]
}
$$
が与えられているとする。次を仮定する。
\begin{enumerate}
\item $\mathcal{F}$ は $X_s$ のすべての点において $S$ 上平坦である。
\item $X' \to S'$ は有限型である。
\item $g^*\mathcal{F}$ は $X'_{s'}$ に沿って純粋である。
\item $g : X' \to X$ はエタールである。
\item $g(X')$ は $\text{Ass}_{X_s}(\mathcal{F}_s)$ を含む。
\end{enumerate}
この状況では、$\mathcal{F}$ が $X_s$ に沿って純粋であることと、
$X' \to X \times_S S'$ の像が、$s'$ に特殊化する $S'$ の点の上にある
$\text{Ass}_{X \times_S S'/S'}(\mathcal{F} \times_S S')$ の点を
すべて含むことは同値である。
\end{lemma}

\begin{proof}
$S' \to S$ はエタールなので、$\mathcal{F}$ が $X_s$ に沿って純粋ならば
% P05-EN-CH38-0184: source-faithful base-changed-fibre point mismatch; see sidecar.
$\mathcal{F} \times_S S'$ は $X_s$ に沿って純粋である。
Lemma \ref{flat-lemma-quasi-finite-base-change}.
を参照せよ。純粋性は平坦降下を満たすので、
Lemma \ref{flat-lemma-flat-descend-pure},
を参照せよ、$\mathcal{F} \times_S S'$ が $X_{s'}$ に沿って純粋ならば
$\mathcal{F}$ は $X_s$ に沿って純粋である。したがって $S$ を $S'$ で置き換え、
$S = S'$ かつ $g : X' \to X$ が $S$ 上有限型なスキーム間のエタール射であると
仮定してよい。さらに $S$ を $\Spec(\mathcal{O}_{S, s})$ で置き換え、
$S$ は局所であると仮定してよい。

\medskip\noindent
まず、$\mathcal{F}$ が $X_s$ に沿って純粋であると仮定する。
この場合、純粋性により $\text{Ass}_{X/S}(\mathcal{F})$ のすべての点は
$X_s$ の点に特殊化する。したがって
Lemma \ref{flat-lemma-associated-point-specializes}
により、$\text{Ass}_{X/S}(\mathcal{F})$ のすべての点は
$\text{Ass}_{X_s}(\mathcal{F}_s)$ の点に特殊化する。
ゆえに $\text{Ass}_{X/S}(\mathcal{F})$ のすべての点は $g$ の像に属する
（この像は開であり、$\text{Ass}_{X_s}(\mathcal{F}_s)$ を含むからである）。

\medskip\noindent
逆に、$g(X')$ が $\text{Ass}_{X/S}(\mathcal{F})$ を含むと仮定する。
$S^h = \Spec(\mathcal{O}_{S, s}^h)$ を $s$ における $S$ の Hensel 化とする。
% P05-EN-CH38-0185: uniquely determined denote-construction repair; see sidecar.
$S^h \to S$ による $g$ の基底変換を $g^h : (X')^h \to X^h$ と記し、
$\mathcal{F}$ の $X^h$ への引き戻しを $\mathcal{F}^h$ と記す。
Divisors, Lemma \ref{divisors-lemma-base-change-relative-assassin} および
Remark \ref{divisors-remark-base-change-relative-assassin} により、
相対随伴点集合 $\text{Ass}_{X^h/S^h}(\mathcal{F}^h)$ は、射影
$X^h \to X$ による $\text{Ass}_{X/S}(\mathcal{F})$ の逆像である。
$g(X')$ が $\text{Ass}_{X/S}(\mathcal{F})$ を含むと仮定したので、基底変換
$g^h((X')^h) = g(X') \times_S S^h$ は
$\text{Ass}_{X^h/S^h}(\mathcal{F}^h)$ を含む。こうして、$S$ が Hensel 局所環の
スペクトルである場合に帰着する。$x \in \text{Ass}_{X/S}(\mathcal{F})$ とする。
補題の証明を終えるには、$x$ が $X_s$ の点に特殊化することを示さなければならない。
% P05-EN-CH38-0186: uniquely determined missing-article repair; see sidecar.
Definition \ref{flat-definition-pure} に続く議論の純粋性判定条件 (4) を参照せよ。
% P05-EN-CH38-0187: uniquely determined article-agreement repair; see sidecar.
仮定により、$g(x') = x$ を満たす $x' \in X'$ が存在する。
$g : X' \to X$ はエタールなので、
$x' \in \text{Ass}_{X'/S}(g^*\mathcal{F})$ である。
Lemma \ref{flat-lemma-etale-weak-assassin-up-down} を、$x'$ の像を $w \in S$ としたときの
ファイバーの射 $X'_w \to X_w$ に適用せよ。$g^*\mathcal{F}$ は
$X'_s$ に沿って純粋なので、ある $y \in X'_s$ に対して $x' \leadsto y$ である。
したがって $x = g(x') \leadsto g(y)$ かつ $g(y) \in X_s$ となり、主張が従う。
\end{proof}

\begin{lemma}
\label{flat-lemma-finite-type-flat-pure-along-fibre-is-universal}
$f : X \to S$ をスキーム射とし、$\mathcal{F}$ を準連接
$\mathcal{O}_X$-加群とする。$s \in S$ とし、次を仮定する。
\begin{enumerate}
\item $f$ は有限型である。
\item $\mathcal{F}$ は有限型である。
\item $\mathcal{F}$ は $X_s$ のすべての点において $S$ 上平坦である。
\item $\mathcal{F}$ は $X_s$ に沿って純粋である。
\end{enumerate}
このとき $\mathcal{F}$ は $X_s$ に沿って普遍的に純粋である。
\end{lemma}

\begin{proof}
まず予備的な注意を述べる。$(S', s') \to (S, s)$ を基本エタール近傍とする。
% P05-EN-CH38-0188: uniquely determined denote-construction repair; see sidecar.
$X' = X \times_S S'$ への $\mathcal{F}$ の引き戻しを $\mathcal{F}'$ と記す。
Definition \ref{flat-definition-pure} に続く議論により、$\mathcal{F}'$ は
$X'_{s'}$ に沿って純粋である。さらに $\mathcal{F}'$ は $X'_{s'}$ に沿って
$S'$ 上平坦である。そこで、$\mathcal{F}'$ が $X'_{s'}$ に沿って普遍的に
純粋であることを証明すれば十分である。実際、点付きスキームの任意の射
$(T, t) \to (S, s)$ に対して、ファイバー積
$(T', t') = (T \times_S S', (t, s'))$ は $(T, t)$ 上平坦である。したがって
$\mathcal{F}_{T'}$ が $X_{t'}$ に沿って純粋ならば、
Lemma \ref{flat-lemma-flat-descend-pure} により $\mathcal{F}_T$ は
$X_t$ に沿って純粋である。ゆえに証明中、$(s, S)$ を基本エタール近傍で
いつでも置き換えてよい。問題は局所的なので、$S$ を
$\Spec(\mathcal{O}_{S, s})$ で置き換えてもよい。

\medskip\noindent
基本エタール近傍 $(S', s') \to (S, s)$ と可換図式
$$
\xymatrix{
X \ar[d] & X' \ar[l]^g \ar[d] \\
S & \Spec(\mathcal{O}_{S', s'}) \ar[l]
}
$$
で、$X' \to X \times_S \Spec(\mathcal{O}_{S', s'})$ がエタール、
$X_s = g((X')_{s'})$、スキーム $X'$ がアフィン、かつ
$\Gamma(X', g^*\mathcal{F})$ が自由 $\mathcal{O}_{S', s'}$-加群となるものを選ぶ。
Lemma \ref{flat-lemma-finite-type-flat-along-fibre-free-variant} を参照せよ。
$X' \to \Spec(\mathcal{O}_{S', s'})$ は有限型であることに注意せよ
（これはエタール射と有限型射の基底変換との合成である準コンパクト射だからである）。
証明の第一段落の予備的注意により、$S$ を
$\Spec(\mathcal{O}_{S', s'})$ で置き換えてよい。したがって、$S$ 上有限型な
スキームの可換図式
$$
\xymatrix{
X \ar[dr] & & X' \ar[ll]^g \ar[ld] \\
& S &
}
$$
が存在し、$g$ はエタール、$X_s \subset g(X')$、$S$ は閉点 $s$ をもつ局所スキーム、
$X'$ はアフィン、そして $\Gamma(X', g^*\mathcal{F})$ は自由
$\Gamma(S, \mathcal{O}_S)$-加群であると仮定してよい。この場合、
$g^*\mathcal{F}$ は $S$ 上普遍的に純粋であることに注意せよ。
Lemma \ref{flat-lemma-affine-locally-projective-pure} を参照せよ。

\medskip\noindent
この状況で
Lemma \ref{flat-lemma-criterion}
を適用する。$\mathcal{F}$ が $X_s$ に沿って純粋であり、かつ
$X_s$ に沿って $S$ 上平坦であるという仮定から
$\text{Ass}_{X/S}(\mathcal{F}) \subset g(X')$ が従う。
Divisors, Lemma \ref{divisors-lemma-base-change-relative-assassin} および
Remark \ref{divisors-remark-base-change-relative-assassin} により、点付きスキームの
任意の射 $(T, t) \to (S, s)$ に対して
$$
\text{Ass}_{X_T/T}(\mathcal{F}_T) \subset
(X_T \to X)^{-1}(\text{Ass}_{X/S}(\mathcal{F})) \subset
g(X') \times_S T = g_T(X'_T).
$$
したがって
Lemma \ref{flat-lemma-criterion}
を上の図式の $(T, t)$ への基底変換に適用すると、望むとおり
$\mathcal{F}_T$ は $X_t$ に沿って純粋であると結論できる。
\end{proof}

\begin{lemma}
\label{flat-lemma-finite-type-flat-pure-is-universal}
$f : X \to S$ を有限型のスキーム射とし、$\mathcal{F}$ を有限型準連接
$\mathcal{O}_X$-加群とする。$\mathcal{F}$ は $S$ 上平坦であると仮定する。
このとき、$\mathcal{F}$ が $S$ に関して純粋であることと
$\mathcal{F}$ が $S$ に関して普遍的に純粋であることは同値である。
\end{lemma}

\begin{proof}
これは
Lemma \ref{flat-lemma-finite-type-flat-pure-along-fibre-is-universal}
と諸定義から直ちに従う。
\end{proof}

\begin{lemma}
\label{flat-lemma-limit-purity}
$I$ を有向集合とし、$(S_i, g_{ii'})$ を $I$ 上のアフィンスキームの逆系とする。
$S = \lim_i S_i$ とおき、$s \in S$ とする。射影を $g_i : S \to S_i$ と記し、
$s_i = g_i(s)$ とおく。$f : X \to S$ は有限表示の射であり、
$\mathcal{F}$ は有限表示の準連接 $\mathcal{O}_X$-加群で、$X_s$ に沿って純粋、
かつ $X_s$ のすべての点において $S$ 上平坦であると仮定する。
このとき、ある $i \in I$、有限表示の射 $X_i \to S_i$、および有限表示の
準連接 $\mathcal{O}_{X_i}$-加群 $\mathcal{F}_i$ が存在して、
$(X_i)_{s_i}$ に沿って純粋であり、$(X_i)_{s_i}$ の
すべての点において $S_i$ 上平坦であり、$X \cong X_i \times_{S_i} S$、
かつ $\mathcal{F}_i$ の $X$ への引き戻しは $\mathcal{F}$ と同型である。
\end{lemma}

\begin{proof}
$U \subset X$ を $\mathcal{F}$ が $S$ 上平坦となる点の集合とする。
More on Morphisms, Theorem \ref{more-morphisms-theorem-openness-flatness}
により、これは $X$ の開部分スキームである。仮定により $X_s \subset U$ である。
$X_s$ は準コンパクトなので、$X_s \subset U'$ を満たす準コンパクトな開部分
$U' \subset U$ をとれる。
Limits, Lemma \ref{limits-lemma-descend-finite-presentation}
% P05-EN-CH38-0189: source-faithful descended-morphism identity mismatch; see sidecar.
により、ある $i \in I$ と、$S$ への基底変換が $f_i$ と同型となる
有限表示の射 $f_i : X_i \to S_i$ をとれる。そのような選択を固定し、
$X_{i'} = X_i \times_{S_i} S_{i'}$ とおく。このとき
$X = \lim_{i'} X_{i'}$ であり、遷移射はアフィンである。
Limits, Lemma \ref{limits-lemma-descend-modules-finite-presentation}
により、必要なら $i$ を大きくした後、$S$ への基底変換が $\mathcal{F}$ と
同型となる有限表示の準連接 $\mathcal{O}_{X_i}$-加群 $\mathcal{F}_i$ が
存在すると仮定してよい。
Limits, Lemma \ref{limits-lemma-descend-opens}
により、必要なら $i$ を大きくした後、$X$ における逆像が $U'$ となる開部分
$U'_i \subset X_i$ が存在すると仮定してよい。特に
$(X_i)_{s_i} \subset U'_i$ であることに注意せよ。
Limits, Lemma \ref{limits-lemma-descend-module-flat-finite-presentation}
により（もう一度 $i$ を大きくした後）、$\mathcal{F}_i$ は $U'_i$ 上平坦であると
仮定してよい。特に $\mathcal{F}_i$ は $(X_i)_{s_i}$ に沿って平坦である。

\medskip\noindent
次に
Lemma \ref{flat-lemma-finite-presentation-flat-along-fibre}
を用いて、基本エタール近傍 $(S_i', s_i') \to (S_i, s_i)$ とスキームの可換図式
$$
\xymatrix{
X_i \ar[d] & X_i' \ar[l]^{g_i} \ar[d] \\
S_i & S_i' \ar[l]
}
$$
で、$g_i$ がエタール、$(X_i)_{s_i} \subset g_i(X_i')$、スキーム
$X_i'$ と $S_i'$ がアフィン、かつ $\Gamma(X_i', g_i^*\mathcal{F}_i)$ が
射影的 $\Gamma(S_i', \mathcal{O}_{S_i'})$-加群となるものを選ぶ。
$g_i^*\mathcal{F}_i$ は $S'_i$ 上普遍的に純粋であることに注意せよ。
Lemma \ref{flat-lemma-affine-locally-projective-pure} を参照せよ。
任意の $i' \geq i$ に対し、上の図式を基底変換して、$S_{i'}$ 上の射
$(S'_{i'}, s'_{i'}) \to (S_{i'}, s_{i'})$ と
$g_{i'} : X'_{i'} \to X_{i'}$ をもつ図式を得ることができる。また、この図式を
基底変換して、$S$ 上の射 $(S', s') \to (S, s)$ と $g : X' \to X$ をもつ
図式を得ることができる。

\medskip\noindent
ここで純粋性の判定条件を用いることができる。エタール射
$X'_i \to X_i \times_{S_i} S'_i$ の像を
$W'_i \subset X_i \times_{S_i} S'_i$ とおく。各 $i' \geq i$ について同様に像
$W'_{i'} \subset X_{i'} \times_{S_{i'}} S'_{i'}$ があり、また像
$W' \subset X \times_S S'$ がある。像をとる操作は基底変換と可換なので、
$W'_{i'} = W'_i \times_{S'_i} S'_{i'}$ および
% P05-EN-CH38-0190: source-faithful missing-prime notation; see sidecar.
$W' = W_i \times_{S'_i} S'$ である。$\mathcal{F}$ は $X_s$ に沿って純粋なので、
Lemma \ref{flat-lemma-criterion}
から
\begin{equation}
\label{flat-equation-inclusion}
% P05-EN-CH38-0191: source-faithful base-changed structure-map ambiguity; see sidecar.
f^{-1}(\Spec(\mathcal{O}_{S', s'})) \cap
\text{Ass}_{X \times_S S'/S'}(\mathcal{F} \times_S S') \subset W'
\end{equation}
が従う。
More on Morphisms,
Lemma \ref{more-morphisms-lemma-relative-assassin-constructible}
により、
$$
E = \{t \in S' \mid \text{Ass}_{X_t}(\mathcal{F}_t) \subset W' \}
\quad\text{and}\quad
E_{i'} = \{t \in S'_{i'} \mid
\text{Ass}_{X_t}(\mathcal{F}_{i', t}) \subset W'_{i'} \}
$$
これらはそれぞれ $S'$ と $S'_{i'}$ の局所構成可能部分集合である。
More on Morphisms,
Lemma \ref{more-morphisms-lemma-base-change-assassin-in-U}
により、$E_{i'}$ は射 $S'_{i'} \to S'_i$ による $E_i$ の逆像であり、
$E$ は射 $S' \to S'_i$ による $E_i$ の逆像である。したがって
Equation (\ref{flat-equation-inclusion}) は、$\Spec(\mathcal{O}_{S', s'})$ が
$E_i$ に写るという主張と同値である。
$\mathcal{O}_{S', s'} =
\colim_{i' \geq i} \mathcal{O}_{S'_{i'}, s'_{i'}}$
なので、ある $i' \geq i$ に対して $\Spec(\mathcal{O}_{S'_{i'}, s'_{i'}})$ は
$E_i$ に写る。Limits, Lemma
\ref{limits-lemma-limit-contained-in-constructible} を参照せよ。
そこで $S_{i'}$ 上の状況に Lemma \ref{flat-lemma-criterion} を適用すると、
$\mathcal{F}_{i'}$ は $(X_{i'})_{s_{i'}}$ に沿って純粋であると結論できる。
\end{proof}

\begin{lemma}
\label{flat-lemma-flat-finite-presentation-purity-open}
$f : X \to S$ を有限表示の射とする。$\mathcal{F}$ を $S$ 上平坦な
有限表示の準連接 $\mathcal{O}_X$-加群とする。このとき集合
$$
U = \{s \in S \mid \mathcal{F}\text{ は }X_s\text{ に沿って純粋}\}
$$
は $S$ で開である。
\end{lemma}

\begin{proof}
$s \in U$ とする。
Lemma \ref{flat-lemma-finite-presentation-flat-along-fibre}
を用いて、基本エタール近傍 $(S', s') \to (S, s)$ と可換図式
$$
\xymatrix{
X \ar[d] & X' \ar[l]^g \ar[d] \\
S & S' \ar[l]
}
$$
で、$g$ がエタール、$X_s \subset g(X')$、スキーム $X'$ と $S'$ がアフィン、
かつ $\Gamma(X', g^*\mathcal{F})$ が射影的
$\Gamma(S', \mathcal{O}_{S'})$-加群となるものをとれる。
$g^*\mathcal{F}$ は $S'$ 上普遍的に純粋であることに注意せよ。
Lemma \ref{flat-lemma-affine-locally-projective-pure} を参照せよ。
エタール射 $X' \to X \times_S S'$ の像を
$W' \subset X \times_S S'$ とおく。$W$ は開であり、$S'$ 上準コンパクトである
% P05-EN-CH38-0192: source-faithful missing-prime notation; see sidecar.
ことに注意せよ。さらに
$$
E = \{t \in S' \mid \text{Ass}_{X_t}(\mathcal{F}_t) \subset W' \}.
$$
More on Morphisms,
Lemma \ref{more-morphisms-lemma-relative-assassin-constructible}
により、$E$ は $S'$ の構成可能部分集合である。
Lemma \ref{flat-lemma-criterion}
により、$\Spec(\mathcal{O}_{S', s'}) \subset E$ である。
Morphisms, Lemma \ref{morphisms-lemma-constructible-containing-open}
により、$E$ は $s'$ のある開近傍 $V'$ を含む。もう一度
Lemma \ref{flat-lemma-criterion}
を適用すると、$V'$ の $S$ における像の任意の点 $s_1$ に対して
$\mathcal{F}$ は $X_{s_1}$ に沿って純粋である。$S' \to S$ はエタールなので
$V'$ の $S$ における像は開であり、証明が完了する。
\end{proof}


\section{純粋性の利用法}
\label{flat-section-applications-purity}

\noindent
純粋性の利用例をいくつか挙げる。最初の補題で実際に用いるのは、
純粋性より少し弱い形の条件である。

\begin{lemma}
\label{flat-lemma-injectivity-map-source-flat-pure}
$f : X \to S$ を有限型の射とし、$\mathcal{F}$ を $X$ 上の有限型準連接層とする。
$S$ は閉点 $s$ をもつ局所スキームであると仮定する。
$\mathcal{F}$ は $X_s$ に沿って純粋であり、$\mathcal{F}$ は
$S$ 上平坦であると仮定する。
% P05-EN-CH38-0193: uniquely determined missing-morphism repair; see sidecar.
$\varphi : \mathcal{F} \to \mathcal{G}$ を準連接 $\mathcal{O}_X$-加群の射とする。
このとき、次は同値である。
\begin{enumerate}
\item すべての $x \in \text{Ass}_{X_s}(\mathcal{F}_s)$ に対して
茎上の写像 $\varphi_x$ は単射である。
\item $\varphi$ は単射である。
\end{enumerate}
\end{lemma}

\begin{proof}
$\mathcal{K} = \Ker(\varphi)$ とおく。$\mathcal{K} = 0$ を証明することが目標である。
そのためには $\text{WeakAss}_X(\mathcal{K}) = \emptyset$ を証明すれば十分である。
Divisors, Lemma \ref{divisors-lemma-weakly-ass-zero} を参照せよ。
また $\text{WeakAss}_X(\mathcal{K}) \subset
\text{WeakAss}_X(\mathcal{F})$ である。Divisors,
Lemma \ref{divisors-lemma-ses-weakly-ass} を参照せよ。
$\mathcal{F}$ は平坦なので、
Lemma \ref{flat-lemma-bourbaki-finite-type-general-base}
から $\text{WeakAss}_X(\mathcal{F}) \subset
\text{Ass}_{X/S}(\mathcal{F})$ が従う。純粋性により、
$\text{Ass}_{X/S}(\mathcal{F})$ の任意の点 $x'$ は $X_s$ の点の一般化であり、
% P05-EN-CH38-0194: source-faithful reversed specialization wording; see sidecar.
さらに Lemma \ref{flat-lemma-associated-point-specializes} により、ある
$x \in \text{Ass}_{X_s}(\mathcal{F}_s)$ の特殊化である。
したがって $\varphi_x$ の単射性から $\varphi_{x'}$ の単射性が従い、
$\mathcal{K}_{x'} = 0$ となる。
\end{proof}

\begin{proposition}
\label{flat-proposition-finite-presentation-flat-pure-is-projective}
$f : X \to S$ を有限表示アフィンなスキーム射とする。
$\mathcal{F}$ を、$S$ 上平坦な有限表示の準連接
$\mathcal{O}_X$-加群とする。このとき、次は同値である。
\begin{enumerate}
\item $f_*\mathcal{F}$ は $S$ 上局所射影的である。
\item $\mathcal{F}$ は $S$ に関して純粋である。
\end{enumerate}
特に、有限表示の環準同型 $A \to B$ と、$A$ 上平坦な有限表示
$B$-加群 $N$ が与えられたとき、$N$ が $A$-加群として射影的であることと、
$\Spec(B)$ 上の $\widetilde{N}$ が $\Spec(A)$ に関して純粋であることは同値である。
\end{proposition}

\begin{proof}
(1) $\Rightarrow$ (2) は Lemma \ref{flat-lemma-affine-locally-projective-pure}
である。$\mathcal{F}$ は $S$ に関して純粋であると仮定する。
Lemma \ref{flat-lemma-finite-type-flat-pure-along-fibre-is-universal}
により、$\mathcal{F}$ は任意の基底変換後も純粋であることに注意せよ。
Descent, Lemma \ref{descent-lemma-locally-projective-descends}
により、$f_*\mathcal{F}$ が $S$ 上 fpqc 局所的に射影的であることを証明すれば十分である。
$s \in S$ をとる。$f_*\mathcal{F}$ の $s$ のエタール近傍への制限が
局所射影的であることを証明する。実際、
Lemma \ref{flat-lemma-finite-presentation-flat-along-fibre},
により、$S$ を $s$ のアフィンな基本エタール近傍で置き換えた後、図式
$$
\xymatrix{
X \ar[dr] & & X' \ar[ll]^g \ar[ld] \\
& S &
}
$$
が存在し、そこに現れるスキームは $S$ 上アフィンかつ有限表示、$g$ はエタール、
$X_s \subset g(X')$、かつ $\Gamma(X', g^*\mathcal{F})$ は射影的
$\Gamma(S, \mathcal{O}_S)$-加群であると仮定してよい。この場合
$g^*\mathcal{F}$ は $S$ 上普遍的に純粋であることに注意せよ。
Lemma \ref{flat-lemma-affine-locally-projective-pure} を参照せよ。したがって
Lemma \ref{flat-lemma-criterion}
により、開部分 $g(X')$ は $\Spec(\mathcal{O}_{S, s})$ 上にある
$\text{Ass}_{X/S}(\mathcal{F})$ の点を含む。そこで
$$
E = \{t \in S \mid \text{Ass}_{X_t}(\mathcal{F}_t) \subset g(X') \}.
$$
More on Morphisms,
Lemma \ref{more-morphisms-lemma-relative-assassin-constructible}
により、$E$ は $S$ の構成可能部分集合である。すでに
$\Spec(\mathcal{O}_{S, s}) \subset E$ を示した。
Morphisms, Lemma \ref{morphisms-lemma-constructible-containing-open}
により、$E$ は $s$ のある開近傍を含む。したがって $S$ を $s$ のアフィン近傍で
置き換えた後、$\text{Ass}_{X/S}(\mathcal{F}) \subset g(X')$ と仮定してよい。
Lemma \ref{flat-lemma-base-change-universally-flat}
により、これは
$$
\Gamma(X, \mathcal{F}) \longrightarrow \Gamma(X', g^*\mathcal{F})
$$
が $\Gamma(S, \mathcal{O}_S)$-普遍単射であることを意味する。
Algebra, Lemma \ref{algebra-lemma-pure-submodule-ML} により、
% P05-EN-CH38-0195: uniquely determined article repair; see sidecar.
$\Gamma(X, \mathcal{F})$ は $\Gamma(S, \mathcal{O}_S)$-加群として
Mittag--Leffler である。さらに $\Gamma(X, \mathcal{F})$ は
$\Gamma(S, \mathcal{O}_S)$-加群として可算生成かつ平坦なので、
Algebra, Lemma \ref{algebra-lemma-countgen-projective} により射影的である。
\end{proof}

\noindent
この命題を用いると、先の結果のいくつかを改善できる。次の補題は
Proposition \ref{flat-proposition-finite-presentation-flat-at-point} の改善である。

\begin{lemma}
\label{flat-lemma-flat-finite-presentation-affine-neighbourhood-projective}
$f : X \to S$ を局所有限表示の射とする。$\mathcal{F}$ を有限表示の準連接
$\mathcal{O}_X$-加群とする。$x \in X$ とし、$s = f(x) \in S$ とおく。
% P05-EN-CH38-0196: uniquely determined missing-then repair; see sidecar.
$\mathcal{F}$ が $x$ において $S$ 上平坦ならば、アフィンな基本エタール近傍
$(S', s') \to (S, s)$ と、$x' = (x, s')$ を含むアフィン開部分
$U' \subset X \times_S S'$ が存在して、
$\Gamma(U', \mathcal{F}|_{U'})$ は射影的
$\Gamma(S', \mathcal{O}_{S'})$-加群となる。
\end{lemma}

\begin{proof}
証明中、$X$ を $x$ の開近傍で、$S$ を $s$ の基本エタール近傍で置き換えてよい。
したがって平坦性の開性により、$\mathcal{F}$ は $S$ 上平坦であると仮定してよい
（More on Morphisms, Theorem
\ref{more-morphisms-theorem-openness-flatness} を参照せよ）。
$S$ と $X$ はアフィンであると仮定してよい。$X$ をさらに縮小した後、
$\text{Ass}_{X_s}(\mathcal{F}_s)$ の任意の点は $x$ の一般化であると
仮定してよい。この性質は $(S, s)$ を基本エタール近傍で置き換えても保たれる。
したがって
Lemma \ref{flat-lemma-finite-presentation-flat-along-fibre}
を適用して、図式
$$
\xymatrix{
X \ar[dr] & & X' \ar[ll]^g \ar[ld] \\
& S &
}
$$
が存在し、そこに現れるスキームは $S$ 上アフィンかつ有限表示、$g$ はエタール、
$X_s \subset g(X')$、かつ $\Gamma(X', g^*\mathcal{F})$ は射影的
$\Gamma(S, \mathcal{O}_S)$-加群である状況に到達できる。この場合
$g^*\mathcal{F}$ は $S$ 上普遍的に純粋であることに注意せよ。
Lemma \ref{flat-lemma-affine-locally-projective-pure} を参照せよ。

\medskip\noindent
$U \subset g(X')$ を $x$ のアフィン開近傍とする。
$\mathcal{F}|_U$ は $U_s$ に沿って純粋であると主張する。これを証明すれば、
$S$ を縮小した後 $\mathcal{F}|_U$ は $S$ に関して純粋となり、
Lemma \ref{flat-lemma-flat-finite-presentation-purity-open},
を参照せよ、その後の射影性は
Proposition \ref{flat-proposition-finite-presentation-flat-pure-is-projective}
から従う。主張を証明するには、$(S, s)$ を任意の基本エタール近傍で置き換えた後、
ある $s' \in S$ 上にあり $s' \leadsto s$ を満たす
$\text{Ass}_{U/S}(\mathcal{F}|_U)$ の任意の点 $\xi$ が $U_s$ の点に
特殊化することを示さなければならない。$U \subset g(X')$ なので、
% P05-EN-CH38-0197: uniquely determined article-agreement repair; see sidecar.
$g(\xi') = \xi$ を満たす $\xi' \in X'$ をとれる。$g^*\mathcal{F}$ は
$S$ 上純粋なので、Lemma \ref{flat-lemma-associated-point-specializes} により、
ある特殊化 $\xi' \leadsto x'$ が存在し、
% P05-EN-CH38-0198: source-faithful pullback/fibre namespace mismatch; see sidecar.
$x' \in \text{Ass}_{X'_s}(g^*\mathcal{F}_s)$ である。このとき
$g(x') \in \text{Ass}_{X_s}(\mathcal{F}_s)$ である（例えば、Noether スキームの
エタール射 $X'_s \to X_s$ に Lemma
\ref{flat-lemma-etale-weak-assassin-up-down} を適用せよ）。したがって、上での $X$ の
選び方により $g(x') \leadsto x$ である。$x \in U$ なので
$g(x') \in U$ である。ゆえに望むとおり
$\xi = g(\xi') \leadsto g(x') \in U_s$ である。
\end{proof}

\noindent
次の補題は Lemma \ref{flat-lemma-finite-type-flat-at-point-free-variant}
の改善である。

\begin{lemma}
\label{flat-lemma-flat-finite-type-affine-neighbourhood-projective}
$f : X \to S$ を局所有限型の射とする。$\mathcal{F}$ を有限型準連接
$\mathcal{O}_X$-加群とする。$x \in X$ とし、$s = f(x) \in S$ とおく。
% P05-EN-CH38-0199: uniquely determined missing-then repair; see sidecar.
$\mathcal{F}$ が $x$ において $S$ 上平坦ならば、アフィンな基本エタール近傍
$(S', s') \to (S, s)$ と、$x' = (x, s')$ を含むアフィン開部分
$U' \subset X \times_S \Spec(\mathcal{O}_{S', s'})$ が存在して、
$\Gamma(U', \mathcal{F}|_{U'})$ は自由
$\mathcal{O}_{S', s'}$-加群となる。
\end{lemma}

\begin{proof}
問題は $X$ と $S$ 上 Zariski 局所的である。したがって $X$ と $S$ はアフィンであると
仮定してよい。このとき $S$ 上の閉埋め込み
$i : X \to \mathbf{A}^n_S$ をとれる。$\mathbf{A}^n_S$ 上の層
$i_*\mathcal{F}$ と点 $i(x)$ について補題を証明すれば十分であることは明らかである。
こうして $X \to S$ が有限表示である場合に帰着する。$S$ を
$\Spec(\mathcal{O}_{S', s'})$ で、$X$ を
$X \times_S \Spec(\mathcal{O}_{S', s'})$ の開部分で置き換えた後、
$\mathcal{F}$ は有限表示であると仮定してよい。Proposition
\ref{flat-proposition-finite-type-flat-at-point} を参照せよ。この場合、
Lemma \ref{flat-lemma-flat-finite-presentation-affine-neighbourhood-projective}
および
Algebra, Theorem \ref{algebra-theorem-projective-free-over-local-ring}
から結論が従う。
\end{proof}

\begin{lemma}
\label{flat-lemma-flat-finite-type-local-colimit-free}
$A \to B$ を、本質的有限型な局所環の局所準同型とする。
$N$ を、$A$-加群として平坦な有限 $B$-加群とする。$A$ が Hensel ならば、
$N$ は自由 $A$-加群 $F_i$ のフィルター余極限
$$
N = \colim_i F_i
$$
であり、この系のすべての遷移写像 $u_i : F_i \to F_{i'}$ が誘導する写像
$\overline{u}_i : F_i/\mathfrak m_AF_i \to F_{i'}/\mathfrak m_AF_{i'}$
は単射である。また $N$ は Mittag--Leffler $A$-加群である。
\end{lemma}

\begin{proof}
$S$ の閉点 $s$ 上にある点 $x \in X$ をもつ有限型射
$X \to S = \Spec(A)$ と、$A$-加群として $\mathcal{F}_x \cong N$ を満たす
有限型準連接 $\mathcal{O}_X$-加群 $\mathcal{F}$ をとれる。$X$ を縮小した後、
$\text{Ass}_{X_s}(\mathcal{F}_s)$ の各点は $x$ に特殊化すると仮定してよい。
Lemma \ref{flat-lemma-flat-finite-type-affine-neighbourhood-projective}
により、$x$ のアフィン開近傍からなる基本近傍系 $U_i \subset X$ で、
$\Gamma(U_i, \mathcal{F})$ が自由 $A$-加群 $F_i$ となるものが存在する。
$U_{i'} \subset U_i$ ならば、
$$
F_i/\mathfrak m_AF_i = \Gamma(U_{i, s}, \mathcal{F}_s)
\longrightarrow
\Gamma(U_{i', s}, \mathcal{F}_s) = F_{i'}/\mathfrak m_AF_{i'}
$$
は単射であることに注意せよ。実際、核の切断があれば、その台は $x$ と交わらない
$X_s$ の閉部分集合となり、上での $X$ の選び方に矛盾する。写像
$F_i \to F_{i'}$ は $A$-普遍単射なので
（Lemma \ref{flat-lemma-universally-injective-local}）、Algebra,
Lemma \ref{algebra-lemma-colimit-universally-injective-ML} により
$N$ は Mittag--Leffler である。
\end{proof}

\noindent
初読時には次の補題を飛ばしてよい。

\begin{lemma}
\label{flat-lemma-flat-finite-type-local-valuation-ring-has-content}
$A \to B$ を、本質的有限型な局所環の局所準同型とする。
$N$ を、$A$-加群として平坦な有限 $B$-加群とする。$A$ が付値環ならば、
$N$ の任意の元は内容イデアル $I \subset A$ をもつ
（More on Algebra, Definition
\ref{more-algebra-definition-content-ideal}）。さらに $I$ は単項イデアルである。
\end{lemma}

\begin{proof}
最後の主張は、
More on Algebra, Lemma \ref{more-algebra-lemma-content-finitely-generated}
および Algebra, Lemma \ref{algebra-lemma-characterize-valuation-ring}
により $I$ が有限生成イデアルであることから従う。

\medskip\noindent
次に $I$ の存在を証明する。$A \subset A^h$ を Hensel 化とする。
$B'$ を $B \otimes_A A^h$ の極大イデアル
% P05-EN-CH38-0200: source-faithful missing tensor-base subscripts; see sidecar.
$\mathfrak m_B \otimes A^h + B \otimes \mathfrak m_{A^h}$
における局所化とする。このとき $B \to B'$ は平坦、したがって忠実平坦である。
$N' = N \otimes_B B'$ とおく。$x \in N$ をとり、その像を $x' \in N'$ とする。
イデアル $I \subset A$ に対して
$x \in IN \Leftrightarrow x' \in IN'$ が成り立つと主張する。
実際、$N/IN \to N'/IN'$ は $B \to B'$ と $N/IN$ のテンソル積であり、
Algebra, Lemma \ref{algebra-lemma-faithfully-flat-universally-injective}
により $B \to B'$ は普遍単射である。More on Algebra,
Lemma \ref{more-algebra-lemma-henselization-valuation-ring} および
Algebra, Lemma \ref{algebra-lemma-ideals-valuation-ring} により、写像
$A \to A^h$ はイデアルの集合上に包含関係を保つ全単射
$I \mapsto IA^h$ を定める。したがって $x$ が $A$ において内容イデアルをもつことと、
$x'$ が $A^h$ において内容イデアルをもつことは同値である。
$x' \in N'$ に対する主張は Lemma
\ref{flat-lemma-flat-finite-type-local-colimit-free} および
Algebra, Lemma \ref{algebra-lemma-minimal-contains} から従う。
\end{proof}

\noindent
応用として次を得る。

\begin{lemma}
\label{flat-lemma-proper-flat-over-dvr-reduced-fibre}
$R$ を付値環とし、$X \to \Spec(R)$ を固有平坦射とする。特殊ファイバーが
被約ならば、$X$ および $X \to \Spec(R)$ のすべてのファイバーは被約である。
% P05-EN-CH38-0201: uniquely determined subject-agreement repair; see sidecar.
\end{lemma}

\begin{proof}
特殊ファイバー $X_s$ は被約であると仮定する。任意の点 $x \in X$ をとり、
$\mathcal{O}_{X, x}$ が被約であることを示そう。これにより $X$ が被約であることが
証明される。$x'$ が特殊ファイバーに属するような特殊化
$x \leadsto x'$ をとる。固有射は閉写像なので、そのような特殊化は存在する。
局所環 $A = \mathcal{O}_{X, x'}$ を考える。$\mathcal{O}_{X, x}$ は $A$ の局所化なので、
$A$ が被約であることを示せば十分である。非零元 $a \in A$ をとり、
$I = (\pi) \subset R$ をその内容イデアルとする。Lemma
\ref{flat-lemma-flat-finite-type-local-valuation-ring-has-content} を参照せよ。
このとき $a = \pi a'$ であり、$a'$ は $A/\mathfrak mA$ の非零元に写る。ただし
$\mathfrak m \subset R$ は極大イデアルである。$A$ は $R$ 上平坦なので
$\pi$ は非零因子であり、$a$ が冪零ならば $a'$ も冪零である。
しかし $a'$ は被約環
$A/\mathfrak m A = \mathcal{O}_{X_s, x'}$ の非零元に写る。
$A$ が被約でなければこれは矛盾であり、望む主張が従う。

\medskip\noindent
もちろん、$X$ が被約ならば $R$ 上の $X$ の生成ファイバーも被約である。
$\mathfrak p \subset R$ を素イデアルとすると、Algebra,
Lemma \ref{algebra-lemma-make-valuation-rings} により
$R/\mathfrak p$ は付値環である。したがって $X$ の $R/\mathfrak p$ への
基底変換について同じ議論を行えば、$\mathfrak p$ 上のファイバーは被約である。
\end{proof}






\section{平坦化関手}
\label{flat-section-flattening-functors}

\noindent
スキーム $S$ をとる。関手
$F : (\Sch/S)^{opp} \to \textit{Sets}$ が極限保存であるとは、
極限 $T$ をもつアフィンスキームの任意の有向逆系
$\{T_i\}_{i \in I}$ に対して
$F(T) = \colim_i F(T_i)$ が成り立つことをいう。

\begin{situation}
\label{flat-situation-iso}
スキームの射 $f : X \to S$ をとる。
準連接 $\mathcal{O}_X$ 加群の準同型
$u : \mathcal{F} \to \mathcal{G}$ をとる。$S$ 上の任意のスキーム $T$ に対し、
$u$ の $T$ への基底変換を $u_T : \mathcal{F}_T \to \mathcal{G}_T$ と書く。
すなわち、$u_T$ は射影
$X_T = X \times_S T \to X$ による $u$ の引き戻しである。
この状況では関手
\begin{equation}
\label{flat-equation-iso}
F_{iso} : (\Sch/S)^{opp} \longrightarrow \textit{Sets}, \quad
T \longrightarrow \left\{
\begin{matrix}
\{*\} & \text{if} & u_T \text{ が同型である}, \\
\emptyset & \text{otherwise.} &
\end{matrix}
\right.
\end{equation}
を考えることができる。$u_T$ がそれぞれ単射、全射、または零であることを
要求して得られる変種を $F_{inj}$、$F_{surj}$、$F_{zero}$ と書く。
\end{situation}

\begin{lemma}
\label{flat-lemma-iso-sheaf}
状況 \ref{flat-situation-iso} において、次が成り立つ。
\begin{enumerate}
\item 関手 $F_{iso}$、$F_{inj}$、$F_{surj}$、$F_{zero}$ はいずれも
fpqc 位相に関する層条件を満たす。
\item $f$ が準コンパクトで $\mathcal{G}$ が有限型ならば、
$F_{surj}$ は極限保存である。
% P05-EN-CH38-0202: frozen statement omits “is” before “of finite type”.
\item $f$ が準コンパクトで $\mathcal{F}$ が有限型ならば、
$F_{zero}$ は極限保存である。
\item $f$ が準コンパクト、$\mathcal{F}$ が有限型、かつ
$\mathcal{G}$ が有限表示ならば、$F_{iso}$ は極限保存である。
\end{enumerate}
\end{lemma}

\begin{proof}
$\{T_i \to T\}_{i \in I}$ を $S$ 上のスキームの fpqc 被覆とする。
$X_i = X_{T_i} = X \times_S T_i$ および $u_i = u_{T_i}$ とおく。
$\{X_i \to X_T\}_{i \in I}$ は $X_T$ の fpqc 被覆である
(Topologies, Lemma \ref{topologies-lemma-fpqc} を参照)。
とくに、任意の $x \in X_T$ に対して、$x$ に写る $x_i \in X_i$ と
$i \in I$ が存在する。写像
$\mathcal{O}_{X_T, x} \to \mathcal{O}_{X_i, x_i}$ は平坦であり、したがって
忠実平坦である (Algebra, Lemma \ref{algebra-lemma-local-flat-ff} を参照)。
ゆえに、$(u_i)_{x_i}$ が単射、全射、全単射、または零であることと、
$(u_T)_x$ がそれぞれ単射、全射、全単射、または零であることとは同値である。
これで補題の (1) が従う。

\medskip\noindent
証明の (2)。$f$ が準コンパクトで $\mathcal{G}$ が有限型であると仮定する。
$T = \lim_{i \in I} T_i$ をアフィン $S$-スキームの有向極限とし、
$u_T$ が全射であると仮定する。
$X_i = X_{T_i} = X \times_S T_i$ および
$u_i = u_{T_i} : \mathcal{F}_i = \mathcal{F}_{T_i}
\to \mathcal{G}_i = \mathcal{G}_{T_i}$ とおく。
(2) を証明するには、ある $i$ に対して $u_i$ が全射であることを示せばよい。
$i_0 \in I$ を一つ選び、$I$ を $\{i \mid i \geq i_0\}$ で置き換える。
$f$ は準コンパクトなので、スキーム $X_{i_0}$ は準コンパクトである。
したがって、アフィン開集合 $W_1, \ldots, W_m \subset X$ と、
射影 $X_{i_0} \to X$ によって $U_{j, i_0}$ が $W_j$ の中へ写るような
アフィン開被覆
$X_{i_0} = U_{1, i_0} \cup \ldots \cup U_{m, i_0}$ を選べる。
任意の $i \in I$ に対して $U_{j, i}$ を $U_{j, i_0}$ の逆像とする。
$U_j = \lim_i U_{j, i}$ とおけば、
$X_T = U_1 \cup \ldots \cup U_m$ は $X_T$ のアフィン開被覆である。
したがって、各 $j \in \{1, \ldots, m\}$ に対して、ある
$i = i(j) \in I$ で $u_i|_{U_{j, i}}$ が全射となることを示せば十分である。
Properties, Lemma \ref{properties-lemma-finite-type-module} を用いると、
これは次の代数の問題に翻訳される。
$A$ を環とし、$u : M \to N$ を $A$-加群写像とする。
$R = \colim_{i \in I} R_i$ が $A$-代数の有向余極限であるとする。
$N$ が有限 $A$-加群で、
$u \otimes 1 : M \otimes_A R \to N \otimes_A R$ が全射ならば、
ある $i$ に対して写像
$u \otimes 1 : M \otimes_A R_i \to N \otimes_A R_i$ は全射である。
これは Algebra, Lemma
\ref{algebra-lemma-module-map-property-in-colimit} の (2) である。

\medskip\noindent
証明の (3)。(2) の証明とまったく同じ議論により、
これは次の代数の問題に帰着する。
$A$ を環とし、$u : M \to N$ を $A$-加群写像とする。
$R = \colim_{i \in I} R_i$ が $A$-代数の有向余極限であるとする。
$M$ が有限 $A$-加群で、
$u \otimes 1 : M \otimes_A R \to N \otimes_A R$ が零ならば、
ある $i$ に対して写像
$u \otimes 1 : M \otimes_A R_i \to N \otimes_A R_i$ は零である。
これは Algebra, Lemma
\ref{algebra-lemma-module-map-property-in-colimit} の (1) である。

\medskip\noindent
証明の (4)。
% P05-EN-CH38-0203: the proof strengthens the statement's finite-type hypothesis on F to finite presentation.
$f$ が準コンパクトで $\mathcal{F}, \mathcal{G}$ が有限表示であると仮定する。
直前の段落とまったく同様に議論すると
(さらに Properties, Lemma
\ref{properties-lemma-finite-presentation-module} も用いる)、
% P05-EN-CH38-0204: frozen source says part (3), although this paragraph proves part (4).
(3) は次の代数の主張に翻訳される。
$A$ を環とし、$u : M \to N$ を $A$-加群写像とする。
$R = \colim_{i \in I} R_i$ が $A$-代数の有向余極限であるとする。
$M$ が有限 $A$-加群、$N$ が有限表示の $A$-加群で、
$u \otimes 1 : M \otimes_A R \to N \otimes_A R$ が同型であると仮定する。
このとき、ある $i$ に対して写像
$u \otimes 1 : M \otimes_A R_i \to N \otimes_A R_i$ は同型である。
これは Algebra, Lemma
\ref{algebra-lemma-module-map-property-in-colimit} の (3) である。
\end{proof}

\begin{situation}
\label{flat-situation-flat-at-point}
$(A, \mathfrak m_A)$ を局所環とする。
次の圏を $\mathcal{C}$ と書く。その対象は、局所環である $A$-代数 $A'$ で、
代数構造 $A \to A'$ が局所環の局所準同型となるものである。
$\mathcal{C}$ の対象 $A', A''$ の間の射は、$A$-代数の局所準同型
$A' \to A''$ とする。
$A \to B$ を局所環の局所環写像とし、$M$ を $B$-加群とする。
$A'$ が $\mathcal{C}$ の対象ならば $B' = B \otimes_A A'$ とおき、
$M' = M \otimes_A A'$ を $B'$-加群とみなす。
$A' \in \Ob(\mathcal{C})$ に対し、条件
\begin{equation}
\label{flat-equation-flat-at-primes}
\forall \mathfrak q \in V(\mathfrak m_{A'}B' + \mathfrak m_B B')
\subset \Spec(B') :
M'_{\mathfrak q}\text{ は }A'\text{ 上平坦である}.
\end{equation}
を考える。これは More on Algebra, Equation
(\ref{more-algebra-equation-flat-at-primes}) と類似していることに注意する。
とくに、$A' \to A''$ が $\mathcal{C}$ の射で、
(\ref{flat-equation-flat-at-primes}) が $A'$ に対して成り立つならば、
$A''$ に対しても成り立つ (More on Algebra, Lemma
\ref{more-algebra-lemma-base-change-flat-at-primes} を参照)。
したがって関手
\begin{equation}
\label{flat-equation-flat-at-point}
F_{lf} : \mathcal{C} \longrightarrow \textit{Sets}, \quad
A' \longrightarrow \left\{
\begin{matrix}
\{*\} & \text{if }(\ref{flat-equation-flat-at-primes})\text{ が成り立つ}, \\
\emptyset & \text{otherwise.} &
\end{matrix}
\right.
\end{equation}
を得る。
\end{situation}

\begin{lemma}
\label{flat-lemma-flat-at-point}
状況 \ref{flat-situation-flat-at-point} において、次が成り立つ。
\begin{enumerate}
\item $A' \to A''$ が $\mathcal{C}$ の平坦射ならば、
$F_{lf}(A') = F_{lf}(A'')$ である。
\item $A \to B$ が本質的に有限表示で、$M$ が有限表示の $B$-加群ならば、
$F_{lf}$ は極限保存である。すなわち、$\{A_i\}_{i \in I}$ が
$\mathcal{C}$ の対象の有向系ならば、
$F_{lf}(\colim_i A_i) = \colim_i F_{lf}(A_i)$.
\end{enumerate}
\end{lemma}

\begin{proof}
(1) は More on Algebra, Lemma
\ref{more-algebra-lemma-flat-descent-flat-at-primes} の特別な場合である。
(2) は More on Algebra, Lemma
\ref{more-algebra-lemma-limit-preserving-flat-at-primes} の特別な場合である。
\end{proof}

\begin{lemma}
\label{flat-lemma-flat-at-point-finite}
状況 \ref{flat-situation-flat-at-point} において、
% P05-EN-CH38-0205: frozen sentence has “Let B -> C is” and omits “let” before N.
$B \to C$ を局所 $A$-代数の局所写像とし、$N$ を $C$-加群とする。
組 $(C, N)$ に対応する関手を
$F'_{lf} : \mathcal{C} \to \textit{Sets}$ と書く。
$B$-加群として $M \cong N$ であり $B \to C$ が有限ならば、
$F_{lf} = F'_{lf}$.
\end{lemma}

\begin{proof}
$A'$ を $\mathcal{C}$ の対象とする。状況 \ref{flat-situation-flat-at-point} における
$B'$、$M'$ の定義と同様に、$C' = C \otimes_A A'$ および
$N' = N \otimes_A A'$ とおく。
$B'$-加群として $M' \cong N'$ であることに注意する。
$B \to C$ が有限であるという仮定から、(a)
$\mathfrak m_C = \sqrt{\mathfrak m_B C}$ および (b)
$B' \to C'$ が有限であることが従う。(a) より
$$
V(\mathfrak m_{A'}C' + \mathfrak m_C C')
=
\left(
\Spec(C') \to \Spec(B')
\right)^{-1}V(\mathfrak m_{A'}B' + \mathfrak m_B B').
$$
% P05-EN-CH38-0206: frozen source uses subset where q is a prime point and should use membership.
$\mathfrak q \subset V(\mathfrak m_{A'}B' + \mathfrak m_B B')$ と仮定する。
$M'_{\mathfrak q}$ が $A'$ 上平坦であることと、$C'_{\mathfrak q}$-加群
$N'_{\mathfrak q}$ が $A'$ 上平坦であることとは同値であり
(両者は $A'$-加群として同型だからである)、さらにこれは、
$C'_{\mathfrak q}$ の任意の極大イデアル $\mathfrak r$ に対して
$N'_{\mathfrak r}$ が $A'$ 上平坦であることと同値である
(Algebra, Lemma \ref{algebra-lemma-flat-localization} を参照)。
(b) により $B'_{\mathfrak q} \to C'_{\mathfrak q}$ は有限なので、
$C'_{\mathfrak q}$ の極大イデアルは $\mathfrak q$ の上にある $C'$ の素イデアルと
ちょうど対応する (Algebra, Lemma
\ref{algebra-lemma-integral-going-up} を参照)。また、上の表示式により、
これらの素イデアルはすべて
$V(\mathfrak m_{A'}C' + \mathfrak m_C C')$ に含まれる。
したがって補題の結論が成り立つ。
\end{proof}

\begin{lemma}
\label{flat-lemma-flat-at-point-go-up}
状況 \ref{flat-situation-flat-at-point} において、
$B \to C$ が局所環の平坦な局所準同型であると仮定する。
$N = M \otimes_B C$ とおく。組 $(C, N)$ に対応する関手を
$F'_{lf} : \mathcal{C} \to \textit{Sets}$ と書く。このとき
$F_{lf} = F'_{lf}$ である。
\end{lemma}

\begin{proof}
$A'$ を $\mathcal{C}$ の対象とする。状況 \ref{flat-situation-flat-at-point} における
$B'$、$M'$ の定義と同様に、$C' = C \otimes_A A'$ および
$N' = N \otimes_A A' = M' \otimes_{B'} C'$ とおく。ここで
$$
V(\mathfrak m_{A'}B' + \mathfrak m_B B')
=
\Spec( \kappa(\mathfrak m_B) \otimes_A \kappa(\mathfrak m_{A'}) )
$$
であり、$V(\mathfrak m_{A'}C' + \mathfrak m_C C')$ についても同様である。
環写像
$$
\kappa(\mathfrak m_B) \otimes_A \kappa(\mathfrak m_{A'})
\longrightarrow
\kappa(\mathfrak m_C) \otimes_A \kappa(\mathfrak m_{A'})
$$
は忠実平坦である。したがって
$V(\mathfrak m_{A'}C' + \mathfrak m_C C') \to
V(\mathfrak m_{A'}B' + \mathfrak m_B B')$ は全射である。
最後に、$\mathfrak r \in V(\mathfrak m_{A'}C' + \mathfrak m_C C')$ が
$\mathfrak q \in V(\mathfrak m_{A'}B' + \mathfrak m_B B')$ に写るならば、
$B' \to C'$ は平坦なので、$M'_{\mathfrak q}$ が $A'$ 上平坦であることと
$N'_{\mathfrak r}$ が $A'$ 上平坦であることとは同値である
(Algebra, Lemma \ref{algebra-lemma-flatness-descends-more-general} を参照)。
これらの考察から補題は形式的に従う。
\end{proof}

\begin{situation}
\label{flat-situation-free-at-generic-points}
$f : X \to S$ を幾何学的既約ファイバーをもつ滑らかな射とする。
$\mathcal{F}$ を有限型の準連接 $\mathcal{O}_X$-加群とする。
$S$ 上の任意のスキーム $T$ に対して、$\mathcal{F}$ の $T$ への基底変換を
$\mathcal{F}_T$ と書く。すなわち、$\mathcal{F}_T$ は射影
$X_T = X \times_S T \to X$ による $\mathcal{F}$ の引き戻しである。
$X_T \to T$ は幾何学的既約ファイバーをもつ滑らかな射であることに注意する
(Morphisms, Lemma \ref{morphisms-lemma-base-change-smooth} および
More on Morphisms, Lemma
\ref{more-morphisms-lemma-base-change-fibres-geometrically-irreducible} を参照)。
$p \geq 0$ を整数とする。点 $t \in T$ に対して条件
\begin{equation}
\label{flat-equation-free-at-generic-point-fibre}
\mathcal{F}_T \text{ は }\xi_t\text{ の近傍で階数 }p\text{ の自由加群である}
\end{equation}
を考える。ここで $\xi_t$ はファイバー $X_t$ の生成点である。
すべての $t \in T$ に対するこの条件は基底変換で安定なので、関手
\begin{equation}
\label{flat-equation-free-at-generic-points}
H_p : (\Sch/S)^{opp} \longrightarrow \textit{Sets}, \quad
T \longrightarrow \left\{
\begin{matrix}
\{*\} & \text{if }\mathcal{F}_T\text{ が
(\ref{flat-equation-free-at-generic-point-fibre}) を満たす }\forall t\in T, \\
\emptyset & \text{otherwise.}
\end{matrix}
\right.
\end{equation}
を得る。
\end{situation}

\begin{lemma}
\label{flat-lemma-free-at-generic-points}
状況 \ref{flat-situation-free-at-generic-points} において、次が成り立つ。
\begin{enumerate}
\item 関手 $H_p$ は fpqc 位相に関する層条件を満たす。
% P05-EN-CH38-0207: frozen statement omits the article before “functor H_p”.
\item $\mathcal{F}$ が有限表示ならば、関手 $H_p$ は極限保存である。
\end{enumerate}
\end{lemma}

\begin{proof}
$\{T_i \to T\}_{i \in I}$ を $S$ 上のスキームの fpqc\footnote{有限表示の平坦射が
開であることを用いれば、$H_p$ が fppf 位相に関する層であることは容易に
示せる。後で実際に必要なのはこれだけである。しかし、fpqc 位相に関する
層条件も満たすことを直接証明するのは少し面白い。} 被覆とする。
$X_i = X_{T_i} = X \times_S T_i$ とおき、$\mathcal{F}$ の $X_i$ への
引き戻しを $\mathcal{F}_i$ と書く。すべての $i$ に対して $\mathcal{F}_i$ が
(\ref{flat-equation-free-at-generic-point-fibre}) を満たすと仮定する。
$t \in T$ を選び、$X_t$ の生成点を $\xi_t \in X_T$ と書く。
% P05-EN-CH38-0208: frozen proof says F although the sheaf on X_T is F_T.
$\mathcal{F}$ が $\xi_t$ の近傍で自由であることを示さなければならない。
% P05-EN-CH38-0209: frozen prose uses the malformed article “a t_i”.
ある $i \in I$ に対して、$t$ に写る $t_i \in T_i$ を選べる。
$X_{t_i}$ の生成点を $\xi_i \in X_i$ と書くと、$\xi_i$ は $\xi_t$ に写る。
$\mathcal{F}_i$ が $\xi_i$ の近傍で階数 $p$ の自由加群であることから、
% P05-EN-CH38-0210: frozen formula uses undefined x_i instead of xi_i.
$(\mathcal{F}_i)_{x_i} \cong \mathcal{O}_{X_i, x_i}^{\oplus p}$
が従い、さらに
$\mathcal{F}_{T, \xi_t} \cong \mathcal{O}_{X_T, \xi_t}^{\oplus p}$
が従う。実際、$\mathcal{O}_{X_T, \xi_t} \to \mathcal{O}_{X_i, x_i}$ は
平坦である (例えば Algebra, Lemma
\ref{algebra-lemma-finite-projective-descends} を参照)。
したがって、$X_T$ における $\xi_t$ のアフィン近傍 $U$ と全射
$\mathcal{O}_U^{\oplus p} \to \mathcal{F}_U = \mathcal{F}_T|_U$ が存在する
(Modules, Lemma \ref{modules-lemma-finite-type-surjective-on-stalk} を参照)。
$T$ を縮小することにより、$U \to T$ が全射であると仮定してよい。
したがって $U \to T$ は、幾何学的既約ファイバーをもつアフィンスキーム間の
滑らかな射である。さらに、任意の $t' \in T$ に対して誘導される写像
$$
\alpha :
\mathcal{O}_{U, \xi_{t'}}^{\oplus p}
\longrightarrow
\mathcal{F}_{U, \xi_{t'}}
$$
は同型である (先と同じ議論により、右辺の加群は階数 $p$ の自由加群だからである)。
Lemma \ref{flat-lemma-induction-step} から、任意の $t' \in T$ に対して
$$
\Gamma(U, \mathcal{O}_U^{\oplus p})
\otimes_{\Gamma(T, \mathcal{O}_T)} \mathcal{O}_{T, t'}
\longrightarrow
\Gamma(U, \mathcal{F}_U)
\otimes_{\Gamma(T, \mathcal{O}_T)} \mathcal{O}_{T, t'}
$$
は単射である。したがって全射 $\alpha$ は同型である。これで (1) の証明が完了する。

\medskip\noindent
$\mathcal{F}$ が有限表示であると仮定する。
$T = \lim_{i \in I} T_i$ をアフィン $S$-スキームの有向極限とし、
$\mathcal{F}_T$ が (\ref{flat-equation-free-at-generic-point-fibre}) を
満たすと仮定する。
$X_i = X_{T_i} = X \times_S T_i$ とおき、$\mathcal{F}$ の $X_i$ への
引き戻しを $\mathcal{F}_i$ と書く。
$\mathcal{F}_T$ が $T$ 上平坦となる点からなる開部分スキームを
$U \subset X_T$ と書く (More on Morphisms, Theorem
\ref{more-morphisms-theorem-openness-flatness} を参照)。
仮定により、各ファイバーのすべての生成点は $U$ に属する。すなわち、
$U \to T$ は幾何学的既約ファイバーをもつ滑らかな全射である。
$U$ を少し縮小して、$U$ が準コンパクトであると仮定してよい。
Limits, Lemma \ref{limits-lemma-descend-opens} を用いると、
$X_T$ における逆像が $U$ となる準コンパクト開集合 $U_i \subset X_i$ と
$i \in I$ を見つけられる。$i$ を大きくすることにより、
$\mathcal{F}_i|_{U_i}$ が $T_i$ 上平坦であると仮定してよい
(Limits, Lemma
\ref{limits-lemma-descend-module-flat-finite-presentation} を参照)。
とくに $\mathcal{F}_i|_{U_i}$ は有限局所自由であり、したがって局所定数な
階数関数 $\rho : U_i \to \{0, 1, 2, \ldots \}$ を定める。
$\rho$ が値 $p$ をとる開閉部分集合を $(U_i)_p \subset U_i$ と書く。
$(U_i)_p$ の像を $V_i \subset T_i$ とする。$V_i$ は開かつ準コンパクトである。
仮定により、$T \to T_i$ の像は $V_i$ に含まれる。
したがって Limits, Lemma \ref{limits-lemma-descend-opens} により、
$T_{i'} \to T_i$ が $V_i$ を経由するような $i' \geq i$ が存在する。
すると望んだとおり $\mathcal{F}_{i'}$ は
(\ref{flat-equation-free-at-generic-point-fibre}) を満たす。
いくつかの詳細は省略する。
\end{proof}

\begin{lemma}
\label{flat-lemma-pre-flat-dimension-n}
$f : X \to S$ を局所有限型なスキームの射とする。
$\mathcal{F}$ を有限型の準連接 $\mathcal{O}_X$-加群とする。
$n \geq 0$ とする。次の条件は同値である。
\begin{enumerate}
\item $s \in S$ に対し、$\mathcal{F}$ が $S$ 上平坦でない点からなる
閉部分集合 $Z \subset X_s$ (Lemma \ref{flat-lemma-open-in-fibre-where-flat} を参照)
は $\dim(Z) < n$ を満たす。
\item $\mathcal{F}$ が $x$ で $S$ 上平坦でないような $x \in X$ に対して
$\text{trdeg}_{\kappa(f(x))}(\kappa(x)) < n$ が成り立つ。
\end{enumerate}
これらが成り立つならば、任意の基底変換の後にも成り立つ。
\end{lemma}

\begin{proof}
$x \in X$ を $s \in S$ 上の点とする。
Varieties, Lemma \ref{varieties-lemma-dimension-locally-algebraic} により、
$X_s$ における $\{x\}$ の閉包の次元は
$\text{trdeg}_{\kappa(s)}(\kappa(x))$ である。
逆に、$Z \subset X_s$ が次元 $d$ の閉部分集合ならば、
$\text{trdeg}_{\kappa(s)}(\kappa(x)) = d$ を満たす点 $x \in Z$ が存在する
(同じ参照先)。したがって (1) と (2) は同値である
(実際、ファイバーごとに同値である)。基底変換に関する主張は
Morphisms, Lemmas \ref{morphisms-lemma-base-change-module-flat} および
\ref{morphisms-lemma-dimension-fibre-after-base-change} から従う。
\end{proof}

\begin{definition}
\label{flat-definition-flat-dimension-n}
$f : X \to S$ を局所有限型なスキームの射とする。
$\mathcal{F}$ を有限型の準連接 $\mathcal{O}_X$-加群とする。
$n \geq 0$ とする。Lemma \ref{flat-lemma-pre-flat-dimension-n} の同値な条件が
満たされるとき、{\it $\mathcal{F}$ は次元 $\geq n$ で $S$ 上平坦である}
という。
\end{definition}

\begin{situation}
\label{flat-situation-flat-dimension-n}
$f : X \to S$ を局所有限型なスキームの射とする。
$\mathcal{F}$ を有限型の準連接 $\mathcal{O}_X$-加群とする。
$S$ 上の任意のスキーム $T$ に対して、$\mathcal{F}$ の $T$ への基底変換を
$\mathcal{F}_T$ と書く。すなわち、$\mathcal{F}_T$ は射影
$X_T = X \times_S T \to X$ による $\mathcal{F}$ の引き戻しである。
% P05-EN-CH38-0211: locally finite type is weakened to finite type after base change.
$X_T \to T$ は有限型であり、$\mathcal{F}_T$ は有限型の
$\mathcal{O}_{X_T}$-加群であることに注意する (Morphisms, Lemma
\ref{morphisms-lemma-base-change-finite-type} および
Modules, Lemma \ref{modules-lemma-pullback-finite-type})。
$n \geq 0$ とする。Definition \ref{flat-definition-flat-dimension-n} および
Lemma \ref{flat-lemma-pre-flat-dimension-n} により、関手
\begin{equation}
\label{flat-equation-flat-dimension-n}
F_n : (\Sch/S)^{opp} \longrightarrow \textit{Sets}, \quad
T \longrightarrow \left\{
\begin{matrix}
\{*\} & \text{if }\mathcal{F}_T\text{ は }T\text{ 上 }\dim \geq n\text{ で平坦}, \\
\emptyset & \text{otherwise.}
\end{matrix}
\right.
\end{equation}
を得る。
\end{situation}

\begin{lemma}
\label{flat-lemma-flat-dimension-n}
状況 \ref{flat-situation-flat-dimension-n} において、次が成り立つ。
\begin{enumerate}
\item 関手 $F_n$ は fpqc 位相に関する層条件を満たす。
\item $f$ が準コンパクトかつ局所有限表示で、$\mathcal{F}$ が有限表示ならば、
関手 $F_n$ は極限保存である。
\end{enumerate}
\end{lemma}

\begin{proof}
$\{T_i \to T\}_{i \in I}$ を $S$ 上のスキームの fpqc 被覆とする。
$X_i = X_{T_i} = X \times_S T_i$ とおき、$\mathcal{F}$ の $X_i$ への
引き戻しを $\mathcal{F}_i$ と書く。すべての $i$ に対して
$\mathcal{F}_i$ が次元 $\geq n$ で $T_i$ 上平坦であると仮定する。
$t \in T$ とする。添字 $i$ と、$t$ に写る点 $t_i \in T_i$ を選ぶ。
Cartesian 図式
$$
\xymatrix{
X_{\Spec(\mathcal{O}_{T, t})} \ar[d] &
X_{\Spec(\mathcal{O}_{T_i, t_i})} \ar[d] \ar[l] \\
\Spec(\mathcal{O}_{T, t}) &
\Spec(\mathcal{O}_{T_i, t_i}) \ar[l]
}
$$
を考える。下の水平射は平坦なので、More on Morphisms, Lemma
\ref{more-morphisms-lemma-flat-locus-base-change} により、
$\mathcal{F}_i$ が $T_i$ 上平坦でない点からなる集合
$Z_i \subset X_{t_i}$ と、$\mathcal{F}_T$ が $T$ 上平坦でない点からなる集合
$Z \subset X_t$ は $Z_i = Z_{\kappa(t_i)}$ という関係にある。
したがって Morphisms, Lemma
\ref{morphisms-lemma-dimension-fibre-after-base-change} により、
$\mathcal{F}_T$ は次元 $\geq n$ で $T$ 上平坦である。

\medskip\noindent
$f$ が準コンパクトかつ局所有限表示で、$\mathcal{F}$ が有限表示であると仮定する。
この段落では、まず (2) の証明を $f$ が有限表示である場合に帰着させる。
$T = \lim_{i \in I} T_i$ をアフィン $S$-スキームの有向極限とし、
$\mathcal{F}_T$ が次元 $\geq n$ で平坦であると仮定する。
$X_i = X_{T_i} = X \times_S T_i$ とおき、$\mathcal{F}$ の $X_i$ への
引き戻しを $\mathcal{F}_i$ と書く。ある $i$ に対して $\mathcal{F}_i$ が
次元 $\geq n$ で平坦であることを示さなければならない。
$i_0 \in I$ を選び、$I$ を $\{i \mid i \geq i_0\}$ で置き換える。
$T_{i_0}$ はアフィン (したがって準コンパクト) なので、有限個のアフィン開集合
$W_j \subset S$、$j = 1, \ldots, m$ と、$T_{i_0} \to S$ が各
$V_{j, i_0}$ を $W_j$ の中へ写すようなアフィン開被覆
% P05-EN-CH38-0212: frozen source misspells “covering” as “overing”.
$T_{i_0} = \bigcup_{j = 1, \ldots, m} V_{j, i_0}$ が存在する。
$i \geq i_0$ に対して、$T_i$ における $V_{j, i_0}$ の逆像を
$V_{j, i}$ と書く。各 $j$ に対して、ある $i$ で
% P05-EN-CH38-0213: frozen sheaf restriction uses V_{j,i_0}, not the varying V_{j,i}.
$\mathcal{F}_{V_{j, i_0}}$ が次元 $\geq n$ で平坦となることを示せればよい。
こうして $S$ がアフィンである場合に帰着する。この場合 $X$ は準コンパクトなので、
有限アフィン開被覆 $X = W_1 \cup \ldots \cup W_m$ を選べる。
このとき $(X \to S, \mathcal{F})$ に対する結果は
$(\coprod W_j, \coprod \mathcal{F}|_{W_j})$ に対する結果と同値である。
したがって $f$ が有限表示であると仮定してよい。

\medskip\noindent
$f$ と $\mathcal{F}$ が有限表示であると仮定する。
$\mathcal{F}_T$ が $T$ 上平坦となる点からなる開部分スキームを
$U \subset X_T$ と書く (More on Morphisms, Theorem
\ref{more-morphisms-theorem-openness-flatness} を参照)。
仮定により、$Z = X_T \setminus U$ の $T$ 上の各ファイバーの次元は $< n$ である。
Limits, Lemma \ref{limits-lemma-approximate-given-relative-dimension} により、
すべての $t \in T$ に対して $\dim(Z'_t) < n$ であり、かつ
$Z' \to X_T$ が有限表示となる閉部分スキーム
$Z \subset Z' \subset X_T$ を見つけられる。
Limits, Lemmas \ref{limits-lemma-descend-finite-presentation} および
\ref{limits-lemma-descend-closed-immersion-finite-presentation} により、
$T$ への基底変換が $Z'$ となる有限表示の閉部分スキーム
$Z'_i \subset X_i$ と $i \in I$ が存在する。
Limits, Lemma \ref{limits-lemma-limit-dimension} により、
$Z'_i \to T_i$ のすべてのファイバーの次元が $< n$ であると仮定してよい。
Limits, Lemma \ref{limits-lemma-descend-module-flat-finite-presentation} により、
% P05-EN-CH38-0214: frozen source restricts away from undefined T'_i rather than Z'_i.
$\mathcal{F}_i|_{X_i \setminus T'_i}$ が $T_i$ 上平坦であると仮定してよい。
これは $\mathcal{F}_i$ が次元 $\geq n$ で平坦であることを意味する。
ここで $Z' \to X_T$ が有限表示であり、したがって補集合
$X_T \setminus Z'$ が準コンパクトであることを用いる。
以上で (2) が証明され、補題の証明が完了する。
\end{proof}

\begin{situation}
\label{flat-situation-flat}
$f : X \to S$ をスキームの射とする。
$\mathcal{F}$ を準連接 $\mathcal{O}_X$-加群とする。
$S$ 上の任意のスキーム $T$ に対して、$\mathcal{F}$ の $T$ への基底変換を
$\mathcal{F}_T$ と書く。すなわち、$\mathcal{F}_T$ は射影
$X_T = X \times_S T \to X$ による $\mathcal{F}$ の引き戻しである。
平坦加群の基底変換は平坦なので、関手
\begin{equation}
\label{flat-equation-flat}
F_{flat} : (\Sch/S)^{opp} \longrightarrow \textit{Sets}, \quad
T \longrightarrow \left\{
\begin{matrix}
\{*\} & \text{if } \mathcal{F}_T \text{ は }T\text{ 上平坦}, \\
\emptyset & \text{otherwise.}
\end{matrix}
\right.
\end{equation}
を得る。
\end{situation}

\begin{lemma}
\label{flat-lemma-flat}
状況 \ref{flat-situation-flat} において、次が成り立つ。
\begin{enumerate}
\item 関手 $F_{flat}$ は fpqc 位相に関する層条件を満たす。
\item $f$ が準コンパクトかつ局所有限表示で、$\mathcal{F}$ が有限表示ならば、
関手 $F_{flat}$ は極限保存である。
\end{enumerate}
\end{lemma}

\begin{proof}
(1) は次の主張から従う。$T' \to T$ が $S$ 上のスキームの平坦な全射ならば、
$\mathcal{F}_{T'}$ が $T'$ 上平坦であることと $\mathcal{F}_T$ が $T$ 上平坦で
あることとは同値である (More on Morphisms, Lemma
\ref{more-morphisms-lemma-flat-locus-base-change} を参照)。
(2) は、$X$ と $S$ がアフィンである場合に帰着した後、Limits, Lemma
\ref{limits-lemma-descend-module-flat-finite-presentation} から従う
(Lemma \ref{flat-lemma-flat-dimension-n} の証明と比較せよ)。
\end{proof}




\section{平坦化層別}
\label{flat-section-flattening}

\noindent
ここでは定義だけを述べる。
% P05-EN-CH38-0215: frozen prose has a spurious comma after the quoted object of “looking for”.
「一般平坦性による層別化」を探している読者は More on Morphisms, Section
\ref{more-morphisms-section-generic-flatness-stratification} を参照されたい。

\begin{definition}
\label{flat-definition-flattening}
$X \to S$ をスキームの射とする。
$\mathcal{F}$ を準連接 $\mathcal{O}_X$-加群とする。
状況 \ref{flat-situation-flat} で定義した関手 $F_{flat}$ が $S$ 上のスキーム
$S'$ によって表現可能であるとき、{\it $\mathcal{F}$ の普遍平坦化が存在する}
という。$\mathcal{O}_X$ の普遍平坦化が存在するとき、
{\it $X$ の普遍平坦化が存在する}という。
\end{definition}

\noindent
$\mathcal{F}$ の普遍平坦化 $S'$\footnote{スキーム $S'$ はときに
{\it 普遍平坦化子 (universal flatificator)} と呼ばれる。\cite{GruRay} では
{\it platificateur universel} と呼ばれている。普遍平坦化の存在を、
More on Algebra, Section \ref{more-algebra-section-blowup-flat} で論じる種類の
結果と混同してはならない。} が存在するならば、射 $S' \to S$ はスキームの
全射なモノ射であり、$\mathcal{F}_{S'}$ は $S'$ 上平坦であることに注意する。
さらに、射 $T \to S$ が $S'$ を経由することと、$\mathcal{F}_T$ が
$T$ 上平坦であることとは同値である。

\begin{example}
\label{flat-example-no-universal-flattening}
$k$ を体とし、$X = S = \Spec(k[x, y])$ とする。
$M = k[x, x^{-1}, y]/(y)$ として $\mathcal{F} = \widetilde{M}$ とおく。
$k[x, y]$-代数 $A$ に対して $F_{flat}(A) = F_{flat}(\Spec(A))$ とおく。
このとき、すべての $n$ に対して
$F_{flat}(k[x, y]/(x, y)^n) = \{*\}$ である一方、
$F_{flat}(k[[x, y]]) = \emptyset$ である。これは $F_{flat}$ が表現可能でない
ことを意味する (代数空間によってさえ表現可能でない。Formal Spaces, Lemma
\ref{formal-spaces-lemma-map-into-algebraic-space} を参照)。
したがって、この場合には普遍平坦化は存在しない。
\end{example}

\noindent
(Topology, Remark \ref{topology-remark-locally-finite-stratification} と比較せよ。)
スキーム $S$ の (局所有限でスキーム論的な) {\it 層別化}を、
半順序集合 $I$ で添字づけられた閉部分スキーム $Z_i \subset S$ であって、
$S = \bigcup Z_i$ が集合論的に成り立ち、$S$ の各点が有限個の $Z_i$ としか
交わらない近傍をもち、かつ
$$
Z_i \cap Z_j = \bigcup\nolimits_{k \leq i, j} Z_k.
$$
$S_i = Z_i \setminus \bigcup_{j < i} Z_j$ とおくと、実際の層別化は
局所閉部分スキームへの分解 $S = \coprod S_i$ である。
しばしば層 $S_i$ だけを示し、閉部分スキーム $Z_i$ の構成は読者に委ねる。
層別化が与えられると、モノ射
$$
S' = \coprod\nolimits_{i \in I} S_i \longrightarrow S.
$$
を得る。これを{\it 層別化に付随するモノ射}と呼ぶ。
この用語を用いて、平坦化層別をもつことの意味を定義できる。

\begin{definition}
\label{flat-definition-flattening-stratification}
$X \to S$ をスキームの射とする。
$\mathcal{F}$ を準連接 $\mathcal{O}_X$-加群とする。
状況 \ref{flat-situation-flat} で定義した関手 $F_{flat}$ が、$S$ の局所閉部分スキーム
による層別化に付随するモノ射 $S' \to S$ によって表現可能であるとき、
$\mathcal{F}$ は{\it 平坦化層別}をもつという。
$\mathcal{O}_X$ が平坦化層別をもつとき、$X$ は{\it 平坦化層別}をもつという。
\end{definition}

\noindent
平坦化層別が存在するとき、層を添字づける集合とその半順序を理解することが
しばしば重要である。これは加群の階数と関係することが多い。例えば、
$X = S$ で $\mathcal{F}$ が有限表示の $\mathcal{O}_S$-加群ならば、
平坦化層別は存在し、$\mathcal{F}$ の Fitting イデアルによって与えられる
(Divisors, Lemma \ref{divisors-lemma-finite-presentation-module} を参照)。




\section{Artin 環上の平坦化層別}
\label{flat-section-flattening-artinian}

\noindent
% P05-EN-CH38-0216: frozen source misspells “flattening” as “flatting”.
基底スキームが Artin 環のスペクトルであるとき、平坦化層別は存在する。

\begin{lemma}
\label{flat-lemma-flattening-stratification-artinian}
$S$ を Artin 環のスペクトルとする。
% P05-EN-CH38-0217: frozen statement omits the symbol F after “any quasi-coherent O_X-module”.
$S$ 上の任意のスキーム $X$ と任意の準連接 $\mathcal{O}_X$-加群に対して、
普遍平坦化が存在する。実際、普遍平坦化は閉埋め込み $S' \to S$ によって
与えられ、したがって $\mathcal{F}$ の平坦化層別でもある。
\end{lemma}

\begin{proof}
アフィン開被覆 $X = \bigcup U_i$ を選ぶ。
このとき $F_{flat}$ は各組 $(U_i, \mathcal{F}|_{U_i})$ に対応する関手の
積である。したがって、各 $(U_i, \mathcal{F}|_{U_i})$ について結果を
証明すれば十分である。アフィンの場合、補題は More on Algebra, Lemma
\ref{more-algebra-lemma-flattening-artinian-universal-property} から直ちに従う。
\end{proof}






\section{写像の平坦化}
\label{flat-section-flattening-map}

\noindent
Theorem \ref{flat-theorem-flattening-map} は、この先の平坦化に関する主張の鍵である。

\begin{lemma}
\label{flat-lemma-universally-separating}
$S$ をスキームとする。
$g : X' \to X$ を $S$ 上のスキームの平坦射とし、$X$ は $S$ 上局所有限型で
あるとする。$\mathcal{F}$ を $S$ 上平坦な有限型準連接
$\mathcal{O}_X$-加群とする。$\text{Ass}_{X/S}(\mathcal{F}) \subset g(X')$
ならば、標準写像
$$
\mathcal{F} \longrightarrow g_*g^*\mathcal{F}
$$
は単射であり、任意の基底変換の後にも単射のままである。
\end{lemma}

\begin{proof}
最後の主張は、任意の射 $T \to S$ に対して
$\mathcal{F}_T \to (g_T)_*g_T^*\mathcal{F}_T$ が単射であることを意味する。
仮定 $\text{Ass}_{X/S}(\mathcal{F}) \subset g(X')$ は基底変換で保たれる
(Divisors, Lemma \ref{divisors-lemma-base-change-relative-assassin} および
Remark \ref{divisors-remark-base-change-relative-assassin} を参照)。
平坦性と有限型の仮定についても同様である。
したがって、表示した矢印が単射であることを証明すれば十分である。
$\mathcal{K} = \Ker(\mathcal{F} \to g_*g^*\mathcal{F})$ とおく。
$\mathcal{K} = 0$ を示すことが目標である。そのためには
$\text{WeakAss}_X(\mathcal{K}) = \emptyset$ を示せば十分である
(Divisors, Lemma \ref{divisors-lemma-weakly-ass-zero} を参照)。
Divisors, Lemma \ref{divisors-lemma-ses-weakly-ass} により
$\text{WeakAss}_X(\mathcal{K}) \subset \text{WeakAss}_X(\mathcal{F})$ である。
$\mathcal{F}$ は平坦なので、Lemma
\ref{flat-lemma-bourbaki-finite-type-general-base} により
$\text{WeakAss}_X(\mathcal{F}) \subset \text{Ass}_{X/S}(\mathcal{F})$ である。
仮定により、$\text{Ass}_{X/S}(\mathcal{F})$ の任意の点 $x$ は、ある
$x' \in X'$ の像である。$g$ は平坦なので、局所環写像
$\mathcal{O}_{X, x} \to \mathcal{O}_{X', x'}$ は忠実平坦である。したがって写像
$$
\mathcal{F}_x
\longrightarrow
g^*\mathcal{F}_{x'} =
\mathcal{F}_x \otimes_{\mathcal{O}_{X, x}} \mathcal{O}_{X', x'}
$$
は単射である (Algebra, Lemma
\ref{algebra-lemma-faithfully-flat-universally-injective} を参照)。
これから望んだ $\mathcal{K}_x = 0$ が従う。
\end{proof}

\begin{lemma}
\label{flat-lemma-flattening-module-map}
$A$ を環とする。$u : M \to N$ を $A$-加群の全射とする。
$M$ が $A$-加群として射影的ならば、任意の環写像
$\varphi : A \to B$ に対して次の条件が同値となるようなイデアル
$I \subset A$ が存在する。
\begin{enumerate}
\item $u \otimes 1 : M \otimes_A B \to N \otimes_A B$ は同型である。
\item $\varphi(I) = 0$ である。
\end{enumerate}
\end{lemma}

\begin{proof}
$M$ は射影的なので、$F = M \oplus C$ が自由 $A$-加群となるような射影的
$A$-加群 $C$ を見つけられる。$u$ を
$u \oplus 1 : F = M \oplus C \to N \oplus C$ で置き換えることにより、
$M$ が自由であると仮定してよい。この場合、$M$ のある固定した基底に関する
$\Ker(u)$ のすべての元の係数が生成する $A$ のイデアルを $I$ とする。
これでうまくいく理由は、$u$ が全射で $\otimes_A B$ が右完全なので、
$\Ker(u \otimes 1)$ が $M \otimes_A B$ における
$\Ker(u) \otimes_A B$ の像だからである。
\end{proof}

\begin{theorem}
\label{flat-theorem-flattening-map}
状況 \ref{flat-situation-iso} において、次を仮定する。
\begin{enumerate}
\item $f$ は有限表示である。
\item $\mathcal{F}$ は有限表示で、$S$ 上平坦、かつ $S$ に関して純粋である。
\item $u$ は全射である。
\end{enumerate}
このとき $F_{iso}$ は閉埋め込み $Z \to S$ によって表現可能である。
さらに、$\mathcal{G}$ が有限表示ならば $Z \to S$ は有限表示である。
\end{theorem}

\begin{proof}
Lemma \ref{flat-lemma-finite-type-flat-pure-along-fibre-is-universal} により
$\mathcal{F}$ が $S$ 上普遍的に純粋であることを、以下では断りなく用いる。
Lemma \ref{flat-lemma-iso-sheaf} および Descent, Lemmas
\ref{descent-lemma-closed-immersion}、
\ref{descent-lemma-descent-data-sheaves} により、この問題は $S$ 上の
\'etale 位相に関して局所的である。したがって、$s \in S$ が与えられたとき、
定理が成り立つような $(S, s)$ の \'etale 近傍が存在することを示せば十分である。

\medskip\noindent
Lemma \ref{flat-lemma-finite-presentation-flat-along-fibre} を用い、$S$ を $s$ の
基本 \'etale 近傍で置き換えることにより、可換図式
$$
\xymatrix{
X \ar[dr] & & X' \ar[ll]^g \ar[ld] \\
& S &
}
$$
が存在すると仮定してよい。ここで各スキームは $S$ 上有限表示であり、
$g$ は \'etale、$X_s \subset g(X')$、スキーム $X'$ と $S$ はアフィンである。
% P05-EN-CH38-0218: frozen list omits “is” before the projective-module predicate.
また、$\Gamma(X', g^*\mathcal{F})$ は射影的
$\Gamma(S, \mathcal{O}_S)$-加群である。
Lemma \ref{flat-lemma-affine-locally-projective-pure} により、$g^*\mathcal{F}$ は
$S$ 上普遍的に純粋である。したがって Lemma \ref{flat-lemma-criterion} により、
開集合 $g(X')$ は $\Spec(\mathcal{O}_{S, s})$ 上にある
$\text{Ass}_{X/S}(\mathcal{F})$ の点を含む。
次のようにおく。
$$
E = \{t \in S \mid \text{Ass}_{X_t}(\mathcal{F}_t) \subset g(X') \}.
$$
More on Morphisms, Lemma
\ref{more-morphisms-lemma-relative-assassin-constructible} により、
$E$ は $S$ の構成可能部分集合である。すでに
$\Spec(\mathcal{O}_{S, s}) \subset E$ を示した。Morphisms, Lemma
\ref{morphisms-lemma-constructible-containing-open} により、$E$ は $s$ の
開近傍を含む。したがって $S$ を $s$ のより小さいアフィン近傍で置き換え、
$\text{Ass}_{X/S}(\mathcal{F}) \subset g(X')$ と仮定してよい。

\medskip\noindent
$u$ は全射であると仮定したので $F_{iso} = F_{inj}$ である。
Lemma \ref{flat-lemma-universally-separating} により、
$u : \mathcal{F} \to \mathcal{G}$ が単射であることと
$g^*u : g^*\mathcal{F} \to g^*\mathcal{G}$ が単射であることとは同値であり、
任意の基底変換の後にも同じことが成り立つ。したがって、定理の仮定に加えて、
$X \to S$ がアフィンスキームの射で、$\Gamma(X, \mathcal{F})$ が射影的
$\Gamma(S, \mathcal{O}_S)$-加群である場合に帰着した。
この場合は Lemma \ref{flat-lemma-flattening-module-map} から直ちに従う。

\medskip\noindent
$\mathcal{G}$ が有限表示ならば $Z$ が有限表示であることを示すには、
Lemma \ref{flat-lemma-iso-sheaf} の (4) と Limits, Remark
\ref{limits-remark-limit-preserving} を組み合わせればよい。
\end{proof}

\begin{lemma}
\label{flat-lemma-Weil-restriction-closed-subschemes}
$f:X\to S$ を有限表示、平坦、かつ純粋なスキームの射とする。
$Y$ を $X$ の閉部分スキームとする。$F=f_*Y$ を $f$ に沿う $Y$ の
Weil 制限関手、すなわち
$$
F : (\Sch/S)^{opp} \to \textit{Sets}, \quad
T \mapsto
\left\{
\begin{matrix}
\{*\} & \text{if} & Y_T\to X_T \text{ が同型, }\\
\emptyset & \text{otherwise.} &
\end{matrix}
\right.
$$
で定義される関手とする。このとき $F$ は閉埋め込み $Z\to S$ によって
表現可能である。さらに、$Y\to S$ が有限表示ならば $Z\to S$ も有限表示である。
\end{lemma}

\begin{proof}
$\mathcal{I}$ を $X$ において $Y$ を定義するイデアル層とし、
$u:\mathcal{O}_X\to\mathcal{O}_X/\mathcal{I}$ を全射とする。
$S$-スキーム $T$ に対して、閉埋め込み $Y_T\to X_T$ が同型であることと
$u_T$ が同型であることとは同値である。したがって結果は
Theorem \ref{flat-theorem-flattening-map} から従う。
\end{proof}



\section{局所的な場合の平坦化}
\label{flat-section-flattening-local}

\noindent
この節では、前の内容を適用して平坦化層別の影を得ることから始める。

\begin{theorem}
\label{flat-theorem-flattening-local}
状況 \ref{flat-situation-flat-at-point} において、$A$ は Hensel 的、$B$ は $A$ 上
本質的有限型、$M$ は有限 $B$-加群であると仮定する。このとき、$A/I$ が圏
$\mathcal{C}$ 上の関手 $F_{lf}$ を余表現するようなイデアル $I \subset A$ が
存在する。言い換えると、局所環の局所準同型 $\varphi : A \to A'$ に対し、
$B' = B \otimes_A A'$ および $M' = M \otimes_A A'$ とおけば、次は同値である。
\begin{enumerate}
\item $\forall \mathfrak q \in V(\mathfrak m_{A'}B' + \mathfrak m_B B')
\subset \Spec(B') :
M'_{\mathfrak q}\text{ は }A'\text{ 上平坦である}$。
\item $\varphi(I) = 0$ である。
\end{enumerate}
$B$ が $A$ 上本質的有限表示で、$M$ が $B$ 上有限表示ならば、
$I$ は有限生成イデアルである。
\end{theorem}

\begin{proof}
$B = C_{\mathfrak q}$ および $M = N_{\mathfrak q}$ となるような有限型環写像
$A \to C$、有限 $C$-加群 $N$、および $C$ の素イデアル $\mathfrak q$ を選ぶ。
% P05-EN-CH38-0219: frozen source opens but never closes the quotation around the shorthand phrase.
以下で「定理が $(N/C/A, \mathfrak q)$ に対して成り立つ」とは、
$B = C_{\mathfrak q}$ および $M = N_{\mathfrak q}$ としたときに
$(A \to B, M)$ に対して成り立つことを意味する。
Lemma \ref{flat-lemma-flat-at-point-go-up} により、$B$ を $B$ 上平坦な局所環で
置き換えても関手 $F_{lf}$ は変わらない。したがって $A$ は Hensel 的なので、
Lemma \ref{flat-lemma-existence-algebra} を適用し、$\mathfrak q$ における
$N/C/A$ の完全デヴィサージュが存在すると仮定してよい。

\medskip\noindent
$(A_i, B_i, M_i, \alpha_i, \mathfrak q_i)_{i = 1, \ldots, n}$ をそのような
$\mathfrak q$ における $N/C/A$ の完全デヴィサージュとする。
Definition \ref{flat-definition-complete-devissage-at-x-algebra} のように、
$\mathfrak q_i \subset B_i$ の上にある一意な素イデアルを
$\mathfrak q'_i \subset A_i$ とする。
$C \to A_1$ は全射で、$C$-加群として $N \cong M_1$ なので、
Lemma \ref{flat-lemma-flat-at-point-finite} により、定理を
$(M_1/A_1/A, \mathfrak q'_1)$ に対して証明すれば十分である。
% P05-EN-CH38-0220: frozen prime-over wording and localized finite map reverse B_1 -> A_1.
$B_1 \to A_1$ は有限で、$\mathfrak q_1$ は $\mathfrak q'_1$ の上にある
$B_1$ の唯一の素イデアルなので、
$(A_1)_{\mathfrak q'_1} \to (B_1)_{\mathfrak q_1}$ は有限である
(Algebra, Lemma \ref{algebra-lemma-unique-prime-over-localize-below} または
More on Morphisms, Lemma
\ref{more-morphisms-lemma-finite-morphism-single-point-in-fibre} を参照)。
したがって Lemma \ref{flat-lemma-flat-at-point-finite} により、定理を
$(M_1/B_1/A, \mathfrak q_1)$ に対して証明すれば十分である。

\medskip\noindent
ここで、デヴィサージュの長さ $n$ に関する帰納法により、定理が
$(M_2/B_2/A, \mathfrak q_2)$ に対して成り立つと仮定してよい。
($n = 1$ ならば $M_2 = 0$ であり、これは $A$ 上平坦である。)
前段落の最後の二つの手順を逆向きにたどり、
% P05-EN-CH38-0221: complete devissage gives M_2 congruent Coker(alpha_1), not Coker(alpha_2).
$B_1$-加群として $M_2 \cong \Coker(\alpha_2)$ であることを用いると、
定理は $(\Coker(\alpha_1)/B_1/A, \mathfrak q_1)$ に対して成り立つ。

\medskip\noindent
$A'$ を $\mathcal{C}$ の対象とする。Lemma \ref{flat-lemma-induction-step} により、
$\mathfrak m_{A'}$ の上にある $B_1 \otimes_A A'$ の素イデアル
$\mathfrak q'$ に対して $(M_1 \otimes_A A')_{\mathfrak q'}$ が $A'$ 上平坦ならば、
$(\Coker(\alpha_1) \otimes_A A')_{\mathfrak q'}$ も $A'$ 上平坦である。
したがって $F_{lf}$ は、$(B_1)_{\mathfrak q_1}$ 上の加群
$\Coker(\alpha_1)_{\mathfrak q_1}$ に対応する関手 $F'_{lf}$ の部分関手である。
前段落により、あるイデアル $J \subset A$ に対して $F'_{lf}$ は $A/J$ によって
余表現される。そこで $A$ を $A/J$ で置き換え、
$\Coker(\alpha_1)_{\mathfrak q_1}$ が $A$ 上平坦であると仮定してよい。

\medskip\noindent
% P05-EN-CH38-0222: frozen source invokes a devissage of undefined N_1 instead of Coker(alpha_1), and has subject-verb disagreement.
$\Coker(\alpha_1)$ は、$\mathfrak q_1$ における $N_1/B_1/A$ の完全
デヴィサージュが存在する $B_1$-加群であり、かつ
$\Coker(\alpha_1)_{\mathfrak q_1}$ は $A$ 上平坦なので、Lemma
\ref{flat-lemma-complete-devissage-flat-finite-type-module} により、
$\Coker(\alpha_1)$ は $A$-加群として自由であり、とくに $A$-加群として平坦である。
したがって Lemma \ref{flat-lemma-induction-step} により、$F_{lf}(A')$ が空でない
ことと、
% P05-EN-CH38-0223: frozen proof uses undefined unindexed alpha instead of alpha_1 twice.
$\alpha \otimes 1_{A'}$ が単射であることとは同値である。
$N_1 = \Im(\alpha_1) \subset M_1$ とおくと、完全列
$$
0 \to N_1 \to M_1 \to \Coker(\alpha_1) \to 0
\quad\text{and}\quad
B_1^{\oplus r_1} \to N_1 \to 0
$$
を得る。$\Coker(\alpha_1)$ の平坦性により、第一の列は普遍完全である
(Algebra, Lemma \ref{algebra-lemma-flat-universally-injective} を参照)。
したがって $\alpha \otimes 1_{A'}$ が単射であることと
$B_1^{\oplus r_1} \otimes_A A' \to N_1 \otimes_A A'$ が同型であることとは
同値である。最後に Theorem \ref{flat-theorem-flattening-map} を適用すると、
この関手はあるイデアル $I$ に対して $A/I$ によって余表現可能である。
ゆえに $F_{lf}$ も $A/I$ によって余表現可能である。

\medskip\noindent
最後の主張を証明するため、$A \to B$ が本質的有限表示で、$M$ が $B$ 上
有限表示であると仮定する。$F_{lf}$ が $A/I$ によって余表現されるような
イデアルを $I \subset A$ とする。$I$ に含まれる有限生成イデアル
$I_\lambda$ 全体を動かして $I = \bigcup I_\lambda$ と書く。
$F_{lf}(A/I) = \{*\}$ なので、ある $\lambda$ に対して
$F_{lf}(A/I_\lambda) = \{*\}$ である (Lemma \ref{flat-lemma-flat-at-point} の
(2) を参照)。これは明らかに $I = I_\lambda$ を意味する。
\end{proof}

\begin{remark}
\label{flat-remark-flattening-local-scheme-theoretic}
Theorem \ref{flat-theorem-flattening-local} をスキーム論的に言い換えると次のようになる。
$(X, x) \to (S, s)$ を局所有限型な点付きスキームの射とする。
$\mathcal{F}$ を有限型準連接 $\mathcal{O}_X$-加群とする。
$S$ は閉点 $s$ をもつ Hensel 局所スキームであると仮定する。
次の性質をもつ閉部分スキーム $Z \subset S$ が存在する。点付きスキームの
任意の射 $(T, t) \to (S, s)$ に対して、次は同値である。
\begin{enumerate}
\item $x \in X_s$ に写るファイバー $X_t$ のすべての点で、
$\mathcal{F}_T$ は $T$ 上平坦である。
\item $\Spec(\mathcal{O}_{T, t}) \to S$ は $Z$ を経由する。
\end{enumerate}
さらに、$X \to S$ が $x$ で有限表示で、$\mathcal{F}_x$ が
$\mathcal{O}_{X, x}$ 上有限表示ならば、$Z \to S$ は有限表示である。
\end{remark}

\noindent
ここで、上の結果から極めて一般的な結果をまったく追加の労力なしに得られる。
おそらく最も興味深い場合は $E = X_s$ の場合であることに注意する。

\begin{lemma}
\label{flat-lemma-freebie}
$S$ を閉点 $s$ をもつ Hensel 局所環のスペクトルとする。
$X \to S$ を局所有限型なスキームの射とする。
$\mathcal{F}$ を有限型準連接 $\mathcal{O}_X$-加群とする。
$E \subset X_s$ を部分集合とする。次の性質をもつ閉部分スキーム
$Z \subset S$ が存在する。点付きスキームの任意の射
$(T, t) \to (S, s)$ に対して、次は同値である。
\begin{enumerate}
\item $E \subset X_s$ の点に写るファイバー $X_t$ のすべての点で、
$\mathcal{F}_T$ は $T$ 上平坦である。
\item $\Spec(\mathcal{O}_{T, t}) \to S$ は $Z$ を経由する。
\end{enumerate}
さらに、$X \to S$ が局所有限表示、$\mathcal{F}$ が有限表示、かつ
$E \subset X_s$ が閉かつ準コンパクトならば、$Z \to S$ は有限表示である。
\end{lemma}

\begin{proof}
$x \in X_s$ に対して、Remark
\ref{flat-remark-flattening-local-scheme-theoretic} で得た閉部分スキームを
$Z_x \subset S$ と書く。このとき $Z = \bigcap_{x \in E} Z_x$ が条件を
満たすことは明らかである。

\medskip\noindent
最後の主張を証明するため、$X$ は局所有限表示、$\mathcal{F}$ は有限表示で、
% P05-EN-CH38-0224: frozen proof assumes Z closed/quasi-compact, but the lemma assumes this of E.
$Z$ は閉かつ準コンパクトであると仮定する。まず、
$E \subset \bigcup W_j$ となる有限個のアフィン開集合 $W_j \subset X$ を選ぶ。
% P05-EN-CH38-0225: frozen restriction uses undefined X_j although the opens are W_j.
各射 $W_j \to S$、層 $\mathcal{F}|_{X_j}$、閉部分集合 $E \cap W_j$ に対して
結果を証明すれば十分である。したがって $X$ はアフィンであると仮定してよい。
この場合 More on Algebra, Lemma
\ref{more-algebra-lemma-limit-preserving-flat-at-primes} により、(1) で定義される
関手は「極限保存」である。したがって Theorem
\ref{flat-theorem-flattening-local} の証明の最後の部分とまったく同様に、
$Z \to S$ が有限表示であることを示せる。
\end{proof}

\begin{remark}
\label{flat-remark-flattening-complete-noetherian}
Lemma \ref{flat-lemma-freebie} の証明をその起源までたどると、長く曲がりくねった
道に行き着く。しかし、次を仮定すれば、
\begin{enumerate}
\item $f$ は有限型である。
\item $\mathcal{F}$ は有限型 $\mathcal{O}_X$-加群である。
\item $E = X_s$ である。
\item $S$ は Noether 完備局所環のスペクトルである。
\end{enumerate}
次のように、より初等的な代数だけに依拠する証明がある。まず有限アフィン開被覆を
とり、$X$ がアフィンである場合に帰着する。この場合、More on Algebra, Lemma
\ref{more-algebra-lemma-flattening-complete-local-universal-property} により
$Z$ は存在する。この証明の鍵は、切断
$\Spec(\mathcal{O}_{S, s}/\mathfrak m_s^n)$ の中で閉部分スキーム $Z$ を
段階的に構成することである。これは、基底が Artin 的であるとき平坦化層別が
常に存在するという事実と、
$\mathcal{O}_{S, s} = \lim \mathcal{O}_{S, s}/\mathfrak m_s^n$ という事実に
依拠する。
\end{remark}




\section{ある補題の変種}
\label{flat-section-variants-mod-injective}

\noindent
この節では Algebra, Lemmas \ref{algebra-lemma-mod-injective-general} および
\ref{algebra-lemma-mod-injective} の変種を論じる。
最も一般的な版は Proposition
\ref{flat-proposition-finite-type-injective-into-flat-mod-m} である。
これは \cite[Lemma 4.2.2]{GruRay} として述べられているが、同文献の証明から
得られるのは Lemma
\ref{flat-lemma-upstairs-finite-type-injective-into-flat-mod-m} に述べる弱い結果だけである。
Proposition \ref{flat-proposition-finite-type-injective-into-flat-mod-m} の精緻な証明は
Ofer Gabber による。現時点ではこの命題を応用しないので、Lemma
\ref{flat-lemma-upstairs-finite-type-injective-into-flat-mod-m} の証明を読んだ後は、
次節へ進むことを勧める。この補題は次節で Theorem
\ref{flat-theorem-check-flatness-at-associated-points} を証明するために用いる。

\begin{situation}
\label{flat-situation-mod-injective}
$\varphi : A \to B$ を本質的有限型な局所環の局所準同型とする。
$M$ を平坦 $A$-加群、$N$ を有限 $B$-加群とし、
$u : N \to M$ を $A$-加群写像で、
$\overline{u} : N/\mathfrak m_AN \to M/\mathfrak m_AM$ が単射となるものとする。
\end{situation}

\noindent
この状況で、$u$ が $A$-普遍単射であり、$N$ が $B$ 上有限表示であり、かつ
$N$ が $A$-加群として平坦であることを示すのが目標である。
これが成り立つとき、与えられた状況で{\it 補題が成り立つ}という。

\begin{lemma}
\label{flat-lemma-Noetherian-finite-type-injective-into-flat-mod-m}
状況 \ref{flat-situation-mod-injective} において環 $A$ が Noetherian ならば、補題は成り立つ。
\end{lemma}

\begin{proof}
Algebra, Lemma \ref{algebra-lemma-mod-injective} を適用すると、$u$ は単射で、
% P05-EN-CH38-0226: frozen quotient N/u(M) is ill-typed for u:N->M; the cokernel is M/u(N).
$N/u(M)$ は $A$ 上平坦である。したがって $u$ は $A$-普遍単射
(Algebra, Lemma \ref{algebra-lemma-flat-tor-zero}) であり、$N$ は $A$-平坦
(Algebra, Lemma \ref{algebra-lemma-flat-ses}) である。この場合 $B$ は
Noetherian なので、$N$ は有限表示である。
\end{proof}

\begin{lemma}
\label{flat-lemma-reduce-finite-type-injective-into-flat-mod-m}
$A_0$ を局所環とする。$A = A_0$、$B$ が $A$ 上の多項式代数の局所化、
$N$ が $B$ 上有限表示であるすべての状況 \ref{flat-situation-mod-injective} で
補題が成り立つならば、$A = A_0$ であるすべての状況
\ref{flat-situation-mod-injective} で補題が成り立つ。
\end{lemma}

\begin{proof}
$A \to B$ と $u : N \to M$ を状況 \ref{flat-situation-mod-injective} のとおりとする。
$C$ が $A$ 上の多項式代数を素イデアルで局所化したものであるように
$B = C/I$ と書く。$N$ が $C$-加群として有限表示であることを示せれば、
まして $N$ は $B$-加群として有限表示である (Algebra, Lemma
\ref{algebra-lemma-finitely-presented-over-subring} を参照)。
したがって $B$ は多項式代数の局所化であると仮定してよい。
次に、ある部分加群 $K \subset B^{\oplus n}$ に対して
$N = B^{\oplus n}/K$ と書く。$B/\mathfrak m_AB$ は体上本質的有限型なので
Noetherian である。したがって、$K' = \sum Bk_i$、
$N' = B^{\oplus n}/K'$ とおいたとき標準全射 $N' \to N$ が同型
$N'/\mathfrak m_AN' \cong N/\mathfrak m_AN$ を誘導するような有限個の元
$k_1, \ldots, k_s \in K$ が存在する。合成 $u' : N' \to M$ に対して
補題が成り立つならば $u'$ は単射であり、したがって $N' = N$、$u' = u$ である。
ゆえに元の状況に対して補題が成り立つ。
\end{proof}

\begin{lemma}
\label{flat-lemma-henselian-finite-type-injective-into-flat-mod-m}
状況 \ref{flat-situation-mod-injective} において環 $A$ が Hensel 的ならば、補題は成り立つ。
\end{lemma}

\begin{proof}
Lemma \ref{flat-lemma-reduce-finite-type-injective-into-flat-mod-m} により、
$B$ が $A$ 上本質的有限表示で、$N$ が $B$ 上有限表示である場合を証明すれば十分である。
一時的に $N$ が $A$ 上平坦であるという仮定も加える。このとき $N$ は、遷移写像
$u_{ii'} : F_i \to F_{i'}$ が $\mathfrak m_A$ を法として単射となる自由
$A$-加群 $F_i$ のフィルター余極限 $N = \colim_i F_i$ である。Lemma
\ref{flat-lemma-flat-finite-type-local-colimit-free} を参照せよ。
各合成 $u_i : F_i \to M$ は Lemma
\ref{flat-lemma-universally-injective-local} により $A$-普遍単射であるから、
$u = \colim u_i$ も所望のとおり $A$-普遍単射である。

\medskip\noindent
$A$ を Hensel 局所環、$B$ を $A$ 上本質的有限表示、$N$ を $B$ 上有限表示と仮定する。
Theorem \ref{flat-theorem-flattening-local} により、$N/IN$ が $A/I$ 上平坦であり、
$I = 0$ でない限り $N/I^2N$ が $A/I^2$ 上平坦でないような有限生成イデアル
$I \subset A$ が存在する。前段落の結果により、$A/I$ 上の
$u \bmod I : N/IN \to M/IM$ に対して補題が成り立つ。可換図式
$$
\xymatrix{
0 \ar[r] &
M \otimes_A I/I^2 \ar[r] &
M/I^2M \ar[r] &
M/IM \ar[r] & 0 \\
&
N \otimes_A I/I^2 \ar[r] \ar[u]^u &
N/I^2N \ar[r] \ar[u]^u &
N/IN \ar[r] \ar[u]^u & 0
}
$$
を考える。$\otimes$ の右完全性と $M$ が $A$ 上平坦であることにより、両行は完全である。
左の垂直射は写像
$N/IN \otimes_{A/I} I/I^2 \to M/IM \otimes_{A/I} I/I^2$ であり、したがって単射である。
図式追跡により左下の射が単射であること、すなわち
% P05-EN-CH38-0227: frozen Tor expression uses M/I^2 and does not match the lower-left injectivity or the claimed flatness of N/I^2N.
$\text{Tor}^1_{A/I^2}(I/I^2, M/I^2) = 0$ が分かる。Algebra, Remark
\ref{algebra-remark-Tor-ring-mod-ideal} を参照せよ。ゆえに Algebra, Lemma
\ref{algebra-lemma-what-does-it-mean} により $N/I^2N$ は $A/I^2$ 上平坦であり、
$I = 0$ でなければ矛盾する。
\end{proof}

\noindent
次の補題では、$M$ が $B$-加群構造をもち、$u$ が $B$-線型であるという状況
\ref{flat-situation-mod-injective} の特殊な場合を論じる。これは実際に最も頻繁に用いる場合で、
一般の場合より証明がはるかに容易である。

\begin{lemma}
\label{flat-lemma-upstairs-finite-type-injective-into-flat-mod-m}
$A \to B$ を本質的有限型な局所環の局所準同型とし、
$u : N \to M$ を $B$-加群写像とする。$N$ が有限 $B$-加群、$M$ が $A$ 上平坦で、
$\overline{u} : N/\mathfrak m_A N \to M/\mathfrak m_A M$ が単射ならば、
$u$ は $A$-普遍単射であり、$N$ は $B$ 上有限表示であり、かつ $N$ は $A$ 上平坦である。
\end{lemma}

\begin{proof}
$A \to A^h$ を $A$ の Hensel 化とする。$B'$ を $B \otimes_A A^h$ の極大イデアル
$\mathfrak m_B \otimes A^h + B \otimes \mathfrak m_{A^h}$.
における局所化とする。$B \to B'$ は平坦であり、したがって忠実平坦なので
(Algebra, Lemma \ref{algebra-lemma-local-flat-ff} を参照)、$A \to B$ を
$A^h \to B'$ で、加群 $M$ を $M \otimes_B B'$ で、加群 $N$ を
$N \otimes_B B'$ で、$u$ を $u \otimes \text{id}_{B'}$ で置き換えてよい。
Algebra, Lemmas \ref{algebra-lemma-descend-properties-modules} および
\ref{algebra-lemma-flatness-descends-more-general} を参照せよ。
したがって $A$ が Hensel 局所環であると仮定してよい。この場合、より一般的な
Lemma \ref{flat-lemma-henselian-finite-type-injective-into-flat-mod-m} から本補題が従う。
\end{proof}

\begin{lemma}
\label{flat-lemma-valuation-ring-finite-type-injective-into-flat-mod-m}
状況 \ref{flat-situation-mod-injective} において環 $A$ が付値環ならば、補題は成り立つ。
\end{lemma}

\begin{proof}
$A$-加群が平坦であることと捩れなしであることは同値であることを思い出そう。
More on Algebra, Lemma
\ref{more-algebra-lemma-valuation-ring-torsion-free-flat} を参照せよ。
$T \subset N$ を $A$-捩れ部分とする。このとき $u(T) = 0$ であり、$N/T$ は $A$-平坦である。
したがって $N/T$ は $B$ 上有限表示である。More on Algebra, Lemma
\ref{more-algebra-lemma-flat-finite-type-valuation-ring-finite-presentation} を参照せよ。
ゆえに $T$ は有限 $B$-加群である。Algebra, Lemma \ref{algebra-lemma-extension} を参照せよ。
$N/T$ は $A$-平坦なので
$T/\mathfrak m_A T \subset N/\mathfrak m_A N$ である。Algebra, Lemma
\ref{algebra-lemma-flat-tor-zero} を参照せよ。$\overline{u}$ は単射だが
$u(T) = 0$ なので、$T/\mathfrak m_A T = 0$ と結論できる。したがって Nakayama の補題
(Algebra, Lemma \ref{algebra-lemma-NAK}) により $T = 0$ である。
この時点で三つの主張のうち二つ、すなわち $N$ が $A$-平坦で $B$ 上有限表示であることを
証明した。残るのは $u$ が普遍単射であることの証明である。

\medskip\noindent
Algebra, Theorem \ref{algebra-theorem-universally-exact-criteria} により、
任意の有限表示 $A$-加群 $Q$ に対して $N \otimes_A Q \to M \otimes_A Q$ が
単射であることを示せば十分である。More on Algebra, Lemma
\ref{more-algebra-lemma-generalized-valuation-ring-modules} により、
非零元 $a \in \mathfrak m_A$ に対して $Q = A/(a)$ と仮定してよい。
したがって $N/aN \to M/aM$ が単射であることを示せば十分である。
$u(x) \in aM$ となる $x \in N$ をとる。Lemma
\ref{flat-lemma-flat-finite-type-local-valuation-ring-has-content} により、$x$ は
内容イデアル $I \subset A$ をもつ。$I$ は有限生成であり
(More on Algebra, Lemma \ref{more-algebra-lemma-content-finitely-generated})、
$A$ は付値環なので、ある $b$ に対して $I = (b)$ である
(Algebra, Lemma \ref{algebra-lemma-characterize-valuation-ring})。
More on Algebra, Lemma \ref{more-algebra-lemma-equal-content} により、元 $u(x)$ も
内容イデアル $I$ をもつ。$u(x) \in aM$ なので More on Algebra, Definition
\ref{more-algebra-definition-content-ideal} により $(b) \subset (a)$ である。
$x \in bN$ だから、所望のとおり $x \in aN$ と結論する。
\end{proof}

\noindent
次の状況を考える。
\begin{equation}
\label{flat-equation-star}
\begin{matrix}
A \to B\text{ は有限表示、}S \subset B
\text{ は乗法的部分集合、かつ }\\
N\text{ は有限表示 }S^{-1}B\text{-加群}
\end{matrix}
\end{equation}
この状況で {\it 純粋スプレッドアウト}とは、
$\Spec(S^{-1}B) \subset U$ を満たすアフィン開集合 $U \subset \Spec(B)$ と、
$N$ を延長する有限表示 $\mathcal{O}(U)$-加群 $N'$ で、$N'$ が $A$-射影的であり、
$N' \to N = S^{-1}N'$ が $A$-普遍単射となるものをいう。

\medskip\noindent
(\ref{flat-equation-star}) において $A \to A_1$ が環写像ならば、基底変換を行える。
$B_1 = B \otimes_A A_1$ とし、$S_1 \subset B_1$ を $S$ の像、
$N_1 = N \otimes_A A_1$ とすればよい。これは
$S_1^{-1}B_1 = S^{-1}B \otimes_A A_1$ だからである。以下では断りなくこれを用いる。

\begin{lemma}
\label{flat-lemma-properties-pure-spreadout}
(\ref{flat-equation-star}) において純粋スプレッドアウトが存在するならば、
\begin{enumerate}
\item $N$ の元は $A$ において内容イデアルをもち、
\item $u : N \to M$ が平坦 $A$-加群 $M$ への射で、$A$ のすべての極大イデアル
$\mathfrak m$ に対して $N/\mathfrak m N \to M/\mathfrak m M$ が単射ならば、
$u$ は $A$-普遍単射である。
\end{enumerate}
\end{lemma}

\begin{proof}
純粋スプレッドアウトの定義のとおりに $U$, $N'$ を選ぶ。
$N'$ は $A$-射影的なので、任意の元 $x' \in N'$ は $A$ において内容イデアルをもつ
(これは直接容易に分かるが、More on Algebra, Lemma
\ref{more-algebra-lemma-content-exists-flat-Mittag-Leffler} と Algebra, Example
\ref{algebra-example-ML} からも従う)。$N' \to N$ は $A$-普遍単射なので、任意の
$x' \in N'$ の像 $x \in N$ は $A$ において内容イデアルをもち、それは $x'$ の
内容イデアルと同じである。一般の $x \in N$ に対しては、$s x$ が $N' \to N$ の像に
属するような $s \in S$ を選び、$x$ と $sx$ が同じ内容イデアルをもつことを用いる。

\medskip\noindent
(2) のとおりの $u : N \to M$ をとる。$u$ が $A$-普遍単射であることを示すため、
$A$ を極大イデアルにおける局所化で置き換えてよい (細部は省略する)。
$A$ は極大イデアル $\mathfrak m$ をもつ局所環であると仮定する。
$s \in S$ を選び、合成
$$
N' \to N \xrightarrow{1/s} N \xrightarrow{u} M
$$
を考える。各写像は $\mathfrak m$ を法として単射なので、Lemma
\ref{flat-lemma-universally-injective-local} により合成は $A$-普遍単射である。
$N = \colim_{s \in S} (1/s)N'$ だから、
% P05-EN-CH38-0229: frozen text says A-inversally injective; the established property and argument are A-universally injective.
普遍単射写像の余極限として $u$ は $A$-普遍単射であると結論する。
\end{proof}

\begin{lemma}
\label{flat-lemma-find-pure-spreadout}
(\ref{flat-equation-star}) において、任意の $\mathfrak p \in \Spec(A)$ に対し、
$A_\mathfrak p/I$ 上で純粋スプレッドアウトが存在するような有限生成イデアル
$I \subset \mathfrak pA_\mathfrak p$ が存在する。
\end{lemma}

\begin{proof}
$A$ を $A_\mathfrak p$ で置き換えてよい。したがって $A$ は局所環で、
$\mathfrak p$ は $A$ の極大イデアル $\mathfrak m$ であると仮定してよい。
$S^{-1}B$ 上の $N$ の表示で分母を払うことにより、ある有限表示 $B$-加群 $N'$ に対して
$N = S^{-1}N'$ と書ける。$B/\mathfrak m B$ は Noetherian なので、
$N'/\mathfrak m N' \to N/\mathfrak m N$ の核 $K$ は有限生成である。
したがって $K$ を零化する $s \in S$ を選べる。すなわちこれは $s$ で零化される。
$B$ を $B_s$ で置き換えることは、
$\Spec(B)$ のアフィン開部分スキームへ移るだけなので許される。この置換後、$S$ の元は
$N'/\mathfrak m N'$ 上で単射に作用する。ここで、$\mathbf{Z}$ 上本質的有限型な局所部分環
$A_0 \subset A$、$B = A \otimes_{A_0} B_0$ を満たす有限型環写像
$A_0 \to B_0$、および
$N' = B \otimes_{B_0} N'_0 = A \otimes_{A_0} N'_0$ を満たす有限 $B_0$-加群
$N'_0$ を選ぶ。$I = \mathfrak m_{A_0} A$ が所要のイデアルであると主張する。
実際、
$$
N'/IN' = N'_0/\mathfrak m_{A_0} N'_0 \otimes_{\kappa_{A_0}} A/I
$$
であり、これは $A/I$ 上自由である。$S$ の元による乗法は極大イデアルで割った後に
単射なので、たとえば Lemma \ref{flat-lemma-invert-universally-injective} により
$N'/IN' \to N/IN$ は普遍単射である。
\end{proof}

\begin{lemma}
\label{flat-lemma-universally-injective-if-flat}
(\ref{flat-equation-star}) において $N$ は $A$-平坦、$M$ は平坦 $A$-加群であるとし、
$u : N \to M$ を、すべての $\mathfrak p \in \Spec(A)$ に対して
$u \otimes \text{id}_{\kappa(\mathfrak p)}$ が単射となる $A$-加群写像とする。
このとき $u$ は $A$-普遍単射である。
\end{lemma}

\begin{proof}
Algebra, Lemma \ref{algebra-lemma-check-universally-injective-into-flat} により、
すべてのイデアル $I \subset A$ に対して $N/IN \to M/IM$ が単射であることを
確かめれば十分である。$A$ を $A/I$ で置き換えると、$u$ が単射であることを
証明すれば十分だと分かる。

\medskip\noindent
まず $u$ が単射であることを証明する。$x \in N$ を $u$ の核の非零元とする。
Algebra, Lemma \ref{algebra-lemma-weakly-ass-zero} により、加群 $Ax$ の弱随伴素イデアル
$\mathfrak p$ が存在する。$A$ を $A_\mathfrak p$ で置き換えることにより $A$ は局所環と
仮定でき、零化イデアルの根基が $\mathfrak m_A$ に等しい非零元 $y \in Ax$ が得られる。
Algebra, Lemma \ref{algebra-lemma-weakly-ass-local} を参照せよ。
したがって $\text{Supp}(y) \subset \Spec(S^{-1}B)$ は空でなく、
$\Spec(S^{-1}B) \to \Spec(A)$ の閉ファイバーに含まれる。
Lemma \ref{flat-lemma-find-pure-spreadout} により、$A/I$ 上で純粋スプレッドアウトをもつような
有限生成イデアル $I \subset \mathfrak m_A$ をとる。このとき、ある $n$ に対して
$I^n y = 0$ である。さて、
% P05-EN-CH38-0230: frozen proof puts y in Ann_M although y lies in N, and uses undefined tensor base R in the displayed map.
$y \in \text{Ann}_M(I^n) = \text{Ann}_A(I^n) \otimes_R N$ が平坦性により成り立つ。
したがって、所望の矛盾を得るには、
$$
\text{Ann}_A(I^n) \otimes_R N
\longrightarrow
\text{Ann}_A(I^n) \otimes_R M
$$
が単射であることを示せば十分である。$N$ と $M$ は平坦で、
$\text{Ann}_A(I^n)$ は $I^n$ で零化されるから、$I$ で零化される任意の $A$-加群
$Q$ に対して $Q \otimes_A N \to Q \otimes_A M$ が単射であることを示せば十分である。
これは $I$ の選び方と Lemma \ref{flat-lemma-properties-pure-spreadout} の (2) から従う。
\end{proof}

\begin{lemma}
\label{flat-lemma-big-intersection-is-zero}
$A$ を体でない局所整域とし、$S$ を $A$ の有限生成イデアルからなる集合とする。
$S$ は積について閉じ、$\bigcup_{I \in S} V(I)$ は $\Spec(A)$ の一般点の補集合であると
仮定する。このとき $\bigcap_{I \in S} I = (0)$ である。
\end{lemma}

\begin{proof}
$\mathfrak m_A \subset A$ は $\Spec(A)$ の一般点でないから、少なくとも一つの
$I \in S$ に対して $I \subset \mathfrak m_A$ である。したがって
$\bigcap_{I \in S} I \subset \mathfrak m_A$ である。
非零元 $f \in \mathfrak m_A$ をとる。このとき
$V(f) \subset \bigcup_{I \in S} V(I)$ である。$V(f)$ の構成可能位相は準コンパクトなので
(Topology, Lemma \ref{topology-lemma-constructible-hausdorff-quasi-compact} および
Algebra, Lemma \ref{algebra-lemma-spec-spectral})、ある $I_j \in S$ に対して
$V(f) \subset V(I_1) \cup \ldots \cup V(I_n)$ となる。
$I_1 \ldots I_n \in S$ だから、ある $I$ に対して $V(f) \subset V(I)$ である。
$I$ は有限生成なので、ある $m$ に対して $I^m \subset (f)$ となる。さらに $S$ は
積について閉じているから、ある $I \in S$ に対して $I \subset (f^2)$ となる。
このとき $f \in I$ は不可能である。
\end{proof}

\begin{lemma}
\label{flat-lemma-closed-points-complement}
$A$ を局所環とし、$I, J \subset A$ をイデアルとする。$J$ が有限生成で、
すべての $n \geq 1$ に対して $I \subset J^n$ ならば、$V(I)$ は
$\Spec(A) \setminus V(J)$ の閉点を含む。
\end{lemma}

\begin{proof}
$\mathfrak p \subset A$ を $\Spec(A) \setminus V(J)$ の閉点とする。
$I \subset \mathfrak p$ を示したい。そうでなければ、ある $f \in I$ の
$A/\mathfrak p$ における像は非零である。ここで
$V(J) \cap \Spec(A/\mathfrak p)$ は一般点でない点の集合である。
したがってイデアル族 $J^nA/\mathfrak p$ に Lemma
\ref{flat-lemma-big-intersection-is-zero} を適用すると、$f$ の $A/\mathfrak p$ における像は
零であると結論される。
\end{proof}

\begin{lemma}
\label{flat-lemma-make-smaller-flatness-ideal}
$A$ を局所環、$I \subset A$ をイデアル、$U \subset \Spec(A)$ を準コンパクト開集合、
$M$ を $A$-加群とする。次を仮定する。
\begin{enumerate}
\item $M/IM$ は $A/I$ 上平坦である。
\item $M$ は $U$ 上平坦である。
\end{enumerate}
このとき $I_2 = \Ker(I \to \Gamma(U, I/I^2))$ とおけば、$M/I_2M$ は
$A/I_2$ 上平坦である。
\end{lemma}

\begin{proof}
$M \otimes_A I/I_2 \to IM/I_2M$ が単射であることを示せば十分である。
Algebra, Lemma \ref{algebra-lemma-what-does-it-mean-again} を参照せよ。
仮定 (2) により、これは $U$ 上で成り立つ。したがって
$M \otimes_A I/I_2$ がその $U$ 上の切断へ単射に写ることを示せば十分である。
$M \otimes_A I/I_2 = M/IM \otimes_A I/I_2$ であり、$M/IM$ は有限自由
$A/I$-加群のフィルター余極限である (Algebra, Theorem
\ref{algebra-theorem-lazard})。ゆえに $I/I_2$ がその $U$ 上の切断へ単射に写ることを
示せば十分であり、これは $I_2$ の構成から従う。
\end{proof}

\begin{proposition}
\label{flat-proposition-finite-type-injective-into-flat-mod-m}
$A \to B$ を本質的有限型な局所環の局所準同型とする。$M$ を平坦 $A$-加群、
$N$ を有限 $B$-加群、$u : N \to M$ を
$\overline{u} : N/\mathfrak m_AN \to M/\mathfrak m_AM$ が単射となる
$A$-加群写像とする。このとき $u$ は $A$-普遍単射であり、$N$ は $B$ 上有限表示であり、
かつ $N$ は $A$ 上平坦である。
\end{proposition}

\begin{proof}
Lemma \ref{flat-lemma-reduce-finite-type-injective-into-flat-mod-m} により、
$B$ は有限表示 $A$-代数 $B_0$ の局所化で、$N$ は有限表示 $B_0$-加群 $M_0$ の
局所化であると仮定してよい。More on Morphisms, Lemma
\ref{more-morphisms-lemma-generic-flatness-stratification} により、
$\Spec(A)$ 上の $\Spec(B_0)$ における $\widetilde{M_0}$ に対して
「一般平坦性成層」が存在する。これを $N$ に戻して読み替えると、閉部分スキームの列
$$
S = \Spec(A) \supset S_0 \supset S_1 \supset \ldots \supset S_t = \emptyset
$$
が得られる。ここで $S_i \subset S$ は $A$ の有限生成イデアルで切り出され、
$\widetilde{N}$ の $\Spec(B) \times_S (S_i \setminus S_{i + 1})$ への引き戻しは
$S_i \setminus S_{i + 1}$ 上平坦である。$t$ に関する帰納法で命題を証明する
($t = 1$ の基底の場合も他の段階と並行して証明する)。
$\Spec(A/J_i)$ を $S_i \setminus S_{i + 1}$ のスキーム論的閉包とする。

\medskip\noindent
{\bf 主張 1.} $N/J_iN$ は $A/J_i$ 上平坦である。$i = t - 1$ の場合は直ちに従い、
$i > 0$ の場合は帰納法の仮定から従う。したがって $t > 1$、
$S_{t - 1} \not = \emptyset$、$J_0 = 0$ と仮定して、$N$ が平坦であることを
証明すればよい。$J \subset A$ を $S_1$ を定義するイデアルとする。
再び $t$ に関する帰納法により、$J$ の各冪を法としても平坦性が成り立つ。
$A^h$ を $A$ の Hensel 化とし、$B'$ を $B \otimes_A A^h$ の極大イデアル
$\mathfrak m_B \otimes A^h + B \otimes \mathfrak m_{A^h}$ における局所化とする。
このとき $B \to B'$ は忠実平坦である。$N' = N \otimes_B B'$ とおく。
$N'$ が $A^h$-平坦であることと $N$ が $A$-平坦であることは同値である。
Theorem \ref{flat-theorem-flattening-local} により、$N'/IN'$ が $A^h/I$ 上平坦となる
最小のイデアル $I \subset A^h$ が存在し、$I$ は有限生成である。上で示したことから、
すべての $n \geq 1$ に対して $I \subset J^nA^h$ である。
$S_i^h \subset \Spec(A^h)$ を $S_i \subset \Spec(A)$ の逆像とする。
Lemma \ref{flat-lemma-closed-points-complement} により、$V(I)$ は
$U = \Spec(A^h) - S_1^h$ の閉点を含む。また構成から $N'$ は $U$ 上 $A^h$-平坦である。
Lemma \ref{flat-lemma-make-smaller-flatness-ideal} により、
% P05-EN-CH38-0231: frozen base A/I_2 is inconsistent because I_2 is an ideal of A^h and the module is N' over A^h.
$N'/I_2N'$ は $A/I_2$ 上平坦である。ここで
$I_2 = \Ker(I \to \Gamma(U, I/I^2))$ である。$I$ の最小性により $I = I_2$ である。
したがって $U$ 上局所的に $I = I^2$、すなわちすべての $u \in U$ に対して
$I\mathcal{O}_{U, u} = (0)$ または $I\mathcal{O}_{U, u} = (1)$ である。
$V(I)$ は $U$ の閉点を含むので、$U$ 上 $I = 0$ である。
$U \subset \Spec(A^h)$ はスキーム論的に稠密なので
(この段落の冒頭で $A$ を $A/J_0$ で置き換えたため)、$I = 0$ である。
したがって $N'$ は $A^h$-平坦であり、主張 1 が成り立つ。

\medskip\noindent
主張 1 より前に設定した状況に戻る。$A^h$ を $A$ の Hensel 化、$B'$ を
$B \otimes_A A^h$ の極大イデアル
$\mathfrak m_B \otimes A^h + B \otimes \mathfrak m_{A^h}$ における局所化、
$N' = N \otimes_B B'$ とする。このとき Theorem
\ref{flat-theorem-flattening-local} の平坦化イデアル $I \subset A^h$ は冪零である。
$nil(A^h)$ を冪零元のイデアルとすれば、
$nil(A^h) = nil(A) A^h$
である (More on Algebra, Lemma \ref{more-algebra-lemma-henselization-nil})。
したがって、$N/I_0N$ が $A/I_0$ 上平坦となる有限生成冪零イデアル
$I_0 \subset A$ が存在する。

\medskip\noindent
{\bf 主張 2.} すべての素イデアル $\mathfrak p \subset A$ に対して写像
$\kappa(\mathfrak p) \otimes_A N \to \kappa(\mathfrak p) \otimes_A M$
は単射である。これが偽であるとき $\mathfrak p$ を悪いという。
% P05-EN-CH38-0232: frozen prose says "bad it this is false"; Japanese renders the uniquely determined "if" reading.
悪い素イデアルの空でない鎖を $C$ とし、
$\mathfrak p^* = \bigcup_{\mathfrak p \in C} \mathfrak p$ とおく。
Lemma \ref{flat-lemma-find-pure-spreadout} により、$V(\mathfrak a)$ 上で
純粋スプレッドアウトが存在するような有限生成イデアル
$\mathfrak a \subset \mathfrak p^*A_{\mathfrak p^*}$ が存在する。
$\mathfrak p^*$ が良いならば、Lemma \ref{flat-lemma-properties-pure-spreadout} により
$V(\mathfrak a)$ の点は良い。一方、$\mathfrak a$ は有限生成で、かつ
% P05-EN-CH38-0233: frozen union omits the factors p on its right-hand side.
$\mathfrak p^*A_{\mathfrak p^*} = \bigcup_{\mathfrak p \in C}A_{\mathfrak p^*}$
なので、$V(\mathfrak a)$ はある $\mathfrak p \in C$ を含み、これは矛盾である。
ゆえに $\mathfrak p^*$ は悪い。Zorn の補題により、悪い素イデアルが存在すれば、
極大なもの $\mathfrak p$ が存在する。言い換えると、
$\mathfrak p' \supset \mathfrak p$、$\mathfrak p' \not = \mathfrak p$ を満たすものは
すべて良いと仮定してよい。この場合、$f \not \in \mathfrak p$ である
すべての $f \in A$ に対して、写像 $u \otimes \text{id}_{A/(\mathfrak p + f)}$ は
普遍単射である。Lemma \ref{flat-lemma-universally-injective-if-flat} を参照せよ。
したがって、部分加群 $f(N/\mathfrak pN)$ が定める位相について
$N/\mathfrak p N$ が分離的であることを示せば十分である。
$B \to B'$ は忠実平坦なので、加群 $N'/\mathfrak p N'$ に対して同じことを
証明すれば十分である。Lemma \ref{flat-lemma-flat-finite-type-local-colimit-free} と
More on Algebra, Lemma
\ref{more-algebra-lemma-content-exists-flat-Mittag-Leffler} により、
$N'/\mathfrak pN'$ の元は $A^h/\mathfrak pA^h$ において内容イデアルをもつ。
したがって
$\bigcap_{f \in A, f \not \in \mathfrak p} f(A^h/\mathfrak p A^h) = 0$.
を示せば十分である。さらに、$\mathfrak p A^h$ 上極小なすべての素イデアル
$\mathfrak q \subset A^h$ に対して $A^h/\mathfrak q A^h$ について同じことを
示せば十分である。$A \to A^h$ は Hensel 化なので、各 $\mathfrak q$ は
$\mathfrak p$ に縮約し、$\mathfrak q' \supset \mathfrak q$、
$\mathfrak q' \not = \mathfrak q$ を満たすものはそれぞれ $\mathfrak p$ を真に含む
素イデアル $\mathfrak p'$ に縮約する。したがって Lemma
\ref{flat-lemma-big-intersection-is-zero} から上の共通部分が零であることを得る。

\medskip\noindent
以上をまとめる。主張 1 と主張 2 を用いると、Lemma
\ref{flat-lemma-universally-injective-if-flat} により $N/I_0 N \to M/I_0M$ は
$A/I_0$-普遍単射である。このとき図式
$$
\xymatrix{
N \otimes_A (I_0^n/I_0^{n + 1}) \ar[r] \ar[d] &
M \otimes_A (I_0^n/I_0^{n + 1}) \ar@{=}[d] \\
I_0^n N /I_0^{n + 1} N \ar[r] &
I_0^n M /I_0^{n + 1} M
}
$$
から左の垂直射は単射だと分かる。したがって Algebra, Lemma
\ref{algebra-lemma-what-does-it-mean-again} により $N$ は平坦である。
同様にして、$u$ の普遍単射性は ($N$ の平坦性を先に証明しなくても)
$I_0$ を法とした普遍単射性へ帰着できる。これで証明が完了する。
\end{proof}


\section{有限型平坦加群 III}
\label{flat-section-finite-type-flat-III}

\noindent
次の結果は本章の主要結果の一つである。

\begin{theorem}
\label{flat-theorem-check-flatness-at-associated-points}
$f : X \to S$ を局所有限型とする。$\mathcal{F}$ を有限型準連接
$\mathcal{O}_X$-加群とし、$x \in X$ の像を $s \in S$ とする。
次は同値である。
\begin{enumerate}
\item $\mathcal{F}$ は $x$ において $S$ 上平坦である。
\item $x$ に特殊化するすべての $x' \in \text{Ass}_{X_s}(\mathcal{F}_s)$ に対して、
$\mathcal{F}$ は $x'$ において $S$ 上平坦である。
\end{enumerate}
\end{theorem}

\begin{proof}
$x$ に特殊化する各点について $\mathcal{F}_{x'}$ は
$\mathcal{F}_x$ の局所化なので、(1) が (2) を含意することは明らかである。
$A = \mathcal{O}_{S, s}$、$B = \mathcal{O}_{X, x}$、$N = \mathcal{F}_x$ とおく。
$\Sigma \subset B$ を、$N/\mathfrak m_AN$ 上で非零因子として作用する $B$ の元からなる
乗法的部分集合とする。仮定 (2) と Lemma \ref{flat-lemma-homothety-spectrum} における
% P05-EN-CH38-0234: frozen text takes Spec of the localized module Sigma^{-1}N; the cited lemma describes Spec(Sigma^{-1}B).
$\Spec(\Sigma^{-1}N)$ の記述により、$\Sigma^{-1}N$ は $A$-平坦である。
一方、構成により写像 $N \to \Sigma^{-1}N$ は $\mathfrak m_A$ を法として単射である。
したがって Lemma \ref{flat-lemma-upstairs-finite-type-injective-into-flat-mod-m} を適用すればよい。
\end{proof}

\noindent
これを直接適用して、次の有用な結果を得る。

\begin{lemma}
\label{flat-lemma-check-along-closed-fibre}
$S$ を閉点 $s$ をもつ局所スキーム、$f : X \to S$ を局所有限型、
$\mathcal{F}$ を有限型準連接 $\mathcal{O}_X$-加群とする。次を仮定する。
\begin{enumerate}
\item $\text{Ass}_{X/S}(\mathcal{F})$ の各点は閉ファイバー $X_s$ の点に特殊化する
\footnote{たとえば、$f$ が有限型で $\mathcal{F}$ が $X_s$ に沿って純粋である場合、
または $f$ が固有である場合に成り立つ。}。
\item $\mathcal{F}$ は $X_s$ の各点において $S$ 上平坦である。
\end{enumerate}
このとき $\mathcal{F}$ は $S$ 上平坦である。
\end{lemma}

\begin{proof}
Theorem \ref{flat-theorem-check-flatness-at-associated-points} により、
$S$ 上の $\mathcal{F}$ の相対随伴点集合の点で平坦性を確かめれば十分であることから
直ちに従う。
\end{proof}











\section{普遍平坦化}
\label{flat-section-flattening-final}

\noindent
スキームの射 $f : X \to S$ が固有かつ有限表示ならば、$f$ の普遍平坦化を得られる。
この節ではこれとそのいくつかの変種を論じる。

\begin{lemma}
\label{flat-lemma-free-at-generic-points-representable}
状況 \ref{flat-situation-free-at-generic-points} において、各 $p \geq 0$ に対し関手 $H_p$
(\ref{flat-equation-free-at-generic-points}) は局所閉埋め込み $S_p \to S$ により表現可能である。
$\mathcal{F}$ が有限表示ならば、$S_p \to S$ は有限表示である。
\end{lemma}

\begin{proof}
各 $S$ について、すべての $p \geq 0$ に対する主張を同時に証明する。
Lemma \ref{flat-lemma-free-at-generic-points} により、関手 $H_p$ は fppf 位相の層である。
したがって Descent, Lemma \ref{descent-lemma-descent-data-sheaves}、
More on Morphisms, Lemma
\ref{more-morphisms-lemma-separated-locally-quasi-finite-morphisms-fppf-descend}、および
Descent, Lemma \ref{descent-lemma-descending-fppf-property-immersion} を合わせると、
問題は $S$ 上の \'etale 位相について局所的である。特に $S$ 上 Zariski 局所的である。

\medskip\noindent
$s \in S$ に対し、$\xi_s$ をファイバー $X_s$ の唯一の一般点とする。
各 $s \in S$ に対して $\mathcal{F}$ の制限 $\mathcal{F}_s$ は、$\xi_s$ のある近傍で
ある階数 $p(s) \geq 0$ の局所自由加群であることに注意する
($X_s$ は既約かつ滑らかなので、これは $X_s$ 上の $\mathcal{F}_s$ に対する一般平坦性から
従う。Algebra, Lemma \ref{algebra-lemma-generic-flatness-Noetherian} を参照せよ。
ただしこれは過剰な道具である)。後で用いるため、
$$
p(s) =
\dim_{\kappa(\xi_s)}(
\mathcal{F}_{\xi_s} \otimes_{\mathcal{O}_{X, \xi_s}} \kappa(\xi_s)
).
$$
であることに注意する。特に $H_{p(s)}(s)$ は空でなく、$q \not = p(s)$ ならば
$H_q(s)$ は空である。

\medskip\noindent
$U \subset X$ を開部分スキームとする。$f : X \to S$ は滑らかなので開写像である。
(\ref{flat-equation-free-at-generic-points}) から、組
$(f|_U : U \to f(U), \mathcal{F}|_U)$ に対する関手 $H_p$ と、組
$(f|_{f^{-1}(f(U))}, \mathcal{F}|_{f^{-1}(f(U))})$ に対する関手 $H_p$ は同じである。
したがって $f(U)$ 上で $S_p$ の存在を証明する際には、常に $X$ を $U$ で置き換えてよい。

\medskip\noindent
$s \in S$ をとる。$\mathcal{F}|_U$ が高々 $p(s)$ 個の元で生成されるような
$\xi_s$ のアフィン開近傍 $U$ が存在する。上の議論により、$s$ の近傍で
% P05-EN-CH38-0235: frozen prose says "in an neighbourhood"; Japanese renders the uniquely determined grammar.
$H_{p(s)}$ に対する主張を証明するには、$\mathcal{F}$ が $p(s)$ 個の元で生成される、
すなわち全射
$$
u : \mathcal{O}_X^{\oplus p(s)} \longrightarrow \mathcal{F}
$$
が存在すると仮定してよい。この場合、$H_{p(s)}$ が写像 $u$ に対する $F_{iso}$
(\ref{flat-equation-iso}) に等しいことは明らかである
(Lemma \ref{flat-lemma-injectivity-map-source-flat-pure} から直ちに従うが、$S$ と $X$ がともに
アフィンになるようさらに少し縮小すれば Lemma \ref{flat-lemma-induction-step-fp} からも従う)。
したがって Theorem \ref{flat-theorem-flattening-map} を適用すると、$s$ の近傍で
$H_{p(s)}$ は閉埋め込みにより表現可能である。

\medskip\noindent
以上から形式的に結果が従う。実際、上の議論により関数 $s \mapsto p(s)$ は
$S$ 上局所的に有界である。そこで $p = \max_{s \in S} p(s)$ に関する帰納法を用いる。
上で示したことから、関手 $H_p$ は閉埋め込み $S_p \to S$ により表現可能である。
$S$ を $S \setminus S_p$ で置き換えると最大値は少なくとも一つ減り、帰納法の仮定が使える。

\medskip\noindent
$\mathcal{F}$ が有限表示であると仮定する。Lemma
\ref{flat-lemma-free-at-generic-points} の (2) と Limits, Remark
\ref{limits-remark-limit-preserving} により、$S_p \to S$ は局所有限表示である。
上の帰納法を繰り返すと、$S$ がアフィンならば各 $S_p$ は準コンパクトだと分かる。
まず $p = \max_{s \in S} p(s)$ ならば $S_p \subset S$ は閉であり (上を参照)、
したがって準コンパクトである。次に $S_p \to S$ は有限表示閉埋め込みなので、
$U = S \setminus S_p$ は $S$ の準コンパクト開部分である
(たとえば Morphisms, Section \ref{morphisms-section-constructible} の議論を参照)。
すると $S_{p - 1} \to U$ は有限表示閉埋め込みなので $S_{p - 1}$ は準コンパクトであり、
$U' = S \setminus (S_p \cup S_{p - 1})$ も準コンパクトである。以下同様である。
\end{proof}

\begin{lemma}
\label{flat-lemma-localize-flat-dimension-n}
状況 \ref{flat-situation-flat-dimension-n} において、$h : X' \to X$ を \'etale 射とする。
$\mathcal{F}' = h^*\mathcal{F}$ および $f' = f \circ h$ とおき、$F_n'$ を
$(f' : X' \to S, \mathcal{F}')$ に付随する (\ref{flat-equation-flat-dimension-n}) とする。
このとき $F_n$ は $F_n'$ の部分関手であり、
$h(X') \supset \text{Ass}_{X/S}(\mathcal{F})$ ならば $F_n = F'_n$ である。
\end{lemma}

\begin{proof}
任意の射 $T \to S$ をとる。このとき $h_T : X'_T \to X_T$ は \'etale 射
% P05-EN-CH38-0236: frozen proof calls h_T a base change of undefined g; the lemma defines the etale morphism h.
$g$ の基底変換なので \'etale である。$t \in T$ に対し、$Z \subset X_t$ を
$\mathcal{F}_T$ が $T$ 上平坦でない点の集合とし、同様に $Z' \subset X'_t$ を
$\mathcal{F}'_T$ が $T$ 上平坦でない点の集合とする。
$\mathcal{F}'_T = h_T^*\mathcal{F}_T$ なので、Morphisms, Lemma
\ref{morphisms-lemma-flat-permanence} により $Z' = h_t^{-1}(Z)$ である。
したがって $Z' \to Z$ は \'etale 射であり、$\dim(Z') \leq \dim(Z)$ である
(たとえば Descent, Lemma \ref{descent-lemma-dimension-at-point-local} による。
あるいは単に \'etale 射が相対次元 $0$ の滑らかな射であることによる)。
これは $F_n \subset F_n'$ を含意する。

\medskip\noindent
最後に $h(X') \supset \text{Ass}_{X/S}(\mathcal{F})$ とし、$F_n'(T)$ が空でない、
すなわち $\mathcal{F}'_T$ が次元 $\geq n$ で $T$ 上平坦となる射 $T \to S$ をとる。
$t \in T$ をとり、$Z \subset X_t$ と $Z' \subset X'_t$ を上のとおりとする。
矛盾を得るため $\dim(Z) \geq n$ と仮定する。次元 $\geq n$ の成分に対応する一般点
$\xi \in Z$ を選び、$x \in \text{Ass}_{X_t}(\mathcal{F}_t)$ を $\xi$ の一般化とする。
Divisors, Lemma \ref{divisors-lemma-base-change-relative-assassin} および Remark
\ref{divisors-remark-base-change-relative-assassin} により、$x$ は
$\text{Ass}_{X/S}(\mathcal{F})$ の点へ写る。したがって $x$ は $h_T$ の像に属し、
ある $x' \in X'_T$ に対して $x = h_T(x')$ と書ける。しかし $x \leadsto \xi$ かつ
$\dim(Z') < n$ なので $x' \not \in Z'$ である。したがって $\mathcal{F}'_T$ は
$x'$ において $T$ 上平坦であり、Morphisms, Lemma
\ref{morphisms-lemma-flat-permanence} により $\mathcal{F}_T$ は $x$ において
$T$ 上平坦である。これはそのようなすべての $x$ について成り立つので、Theorem
\ref{flat-theorem-check-flatness-at-associated-points} により $\mathcal{F}_T$ は $\xi$ において
$T$ 上平坦であり、所望の矛盾を得る。
\end{proof}

\begin{lemma}
\label{flat-lemma-compare-H-F}
$X \to S$ を、次元 $d$ の幾何学的既約ファイバーをもつアフィンスキームの滑らかな射とし、
$\mathcal{F}$ を有限表示準連接 $\mathcal{O}_X$-加群とする。このとき、
(\ref{flat-equation-flat-dimension-n}) の $F_d$ と
(\ref{flat-equation-free-at-generic-points}) の $H_p$ に対して、ある $c \geq 0$ が存在して
$F_d = \coprod_{p = 0, \ldots, c} H_p$ である。
\end{lemma}

\begin{proof}
$X$ はアフィンで $\mathcal{F}$ は有限表示準連接なので、$\mathcal{F}$ はある
$c \geq 0$ 個の元で生成される。すると任意の点 $x \in X$ において
$\dim_{\kappa(x)}(\mathcal{F}_x \otimes \kappa(x))$ は $c$ を超えない。
特に $p > c$ に対して $H_p = \emptyset$ である。また明らかに包含
$\coprod H_p \to F_d$ がある。そこで補題の内容は、基底変換 $\mathcal{F}_T$ が
次元 $\geq d$ で $T$ 上平坦で $t \in T$ ならば、$X_t$ の唯一の一般点 $\xi$ の
ある開近傍 $U \subset X_T$ において $\mathcal{F}_T$ がある階数 $r$ の自由加群に
なるということである。実際、このとき $H_r$ は $U$ の像を含み、それは $t$ の
開近傍である。$U$ の存在は More on Morphisms, Lemma
\ref{more-morphisms-lemma-flat-and-free-at-point-fibre} から従う。
\end{proof}

\begin{lemma}
\label{flat-lemma-flat-dimension-n-representable}
状況 \ref{flat-situation-flat-dimension-n} において、$s \in S$ と $d \geq 0$ をとる。
% P05-EN-CH38-0237: frozen statement omits "and" between "Let s in S" and "let d >= 0"; Japanese renders the determined grammar.
次を仮定する。
\begin{enumerate}
\item ある点 $s \in S$ 上で $\mathcal{F}/X/S$ の完全デヴィサージュが存在する。
\item $X$ は $S$ 上有限表示である。
\item $\mathcal{F}$ は有限表示 $\mathcal{O}_X$-加群である。
\item $\mathcal{F}$ は次元 $\geq d + 1$ で $S$ 上平坦である。
\end{enumerate}
このとき、必要ならば $S$ を $s$ の開近傍で置き換えると、関手 $F_d$
(\ref{flat-equation-flat-dimension-n}) は有限表示モノ射 $Z_d \to S$ により表現可能である。
\end{lemma}

\begin{proof}
まず $X$, $S$ はアフィンスキームであり、$S$ 上のアフィンスキームの圏で $F_d$ が
有限表示モノ射 $Z_d \to S$ により表現可能であることを証明すれば十分である
($Z_d$ がアフィンであるとはもちろん仮定しない)。したがって補題の証明を通じて
$S$ 上のアフィンスキームの圏で議論する。

\medskip\noindent
$(Z_k, Y_k, i_k, \pi_k, \mathcal{G}_k, \alpha_k)_{k = 1, \ldots, n}$ を
$s$ 上の $\mathcal{F}/X/S$ の完全デヴィサージュとする。Definition
\ref{flat-definition-complete-devissage} を参照せよ。デヴィサージュの長さ $n$ に関する
帰納法を用いる。Definition \ref{flat-definition-one-step-devissage} により、$Y_k \to S$ は
幾何学的既約ファイバーをもつ滑らかな射である。$d_k$ を $S$ 上の $Y_k$ の相対次元とする。
$i_{k, *}\mathcal{G}_k = \Coker(\alpha_k)$ で、$i_k$ は閉埋め込みであることを思い出そう。
上で参照した定義により、
$d_1 = \dim(\text{Supp}(\mathcal{F}_s))$ であり、
$$
d_k = \dim(\text{Supp}(\Coker(\alpha_{k - 1})_s))
= \dim(\text{Supp}(\mathcal{G}_{k, s}))
$$
が成り立ち、$k = 2, \ldots, n$ についても上の表示が成り立つ。
$\alpha_k$ は $(Y_k)_s$ の一般点において同型なので、
$d_1 > d_2 > \ldots > d_n \geq 0$ が従う。

\medskip\noindent
$i_1$ は閉埋め込みで $\mathcal{F} = i_{1, *}\mathcal{G}_1$ である。
したがって、$T$ がアフィンである任意のスキーム射 $T \to S$ に対して
$\mathcal{F}_T = i_{1, T, *}\mathcal{G}_{1, T}$ であり、$i_{1, T}$ は依然として
$T$ 上のスキームの閉埋め込みである。ゆえに $\mathcal{F}_T$ が次元 $\geq d$ で
$T$ 上平坦であることと、$\mathcal{G}_{1, T}$ が次元 $\geq d$ で $T$ 上平坦であることは
同値である。$\pi_1 : Z_1 \to Y_1$ は有限なので、同様に
$\mathcal{G}_{1, T}$ が次元 $\geq d$ で $T$ 上平坦であることと、
$\pi_{1, T, *}\mathcal{G}_{1, T}$ が次元 $\geq d$ で $T$ 上平坦であることは同値である。
同じ議論は「次元 $\geq d + 1$ で平坦」にも使えるので、$\mathcal{F}$ に関する仮定から、
特に $\pi_{1, *}\mathcal{G}_1$ は次元 $\geq d + 1$ で $S$ 上平坦である。

\medskip\noindent
$d_1 > d$ と仮定する。上の議論から、特に $\pi_{1, *}\mathcal{G}_1$ は
$(Y_1)_s$ の一般点において $S$ 上平坦である。Lemma
\ref{flat-lemma-induction-step-fp} により、$S$ を $s$ のアフィン近傍で置き換えて、
$\alpha_1$ が $S$-普遍単射であると仮定してよい。$\alpha_1$ は $S$-普遍単射なので、
$T$ がアフィンである任意の射 $T \to S$ に対して短完全列
$$
0 \to \mathcal{O}_{Y_{1, T}}^{\oplus r_1}
\to \pi_{1, T, *}\mathcal{G}_{1, T} \to \Coker(\alpha_1)_T \to 0
$$
があり、第一の射は依然として $T$-普遍単射である。したがって、
$(Y_1)_T$ のうち $\pi_{1, T, *}\mathcal{G}_{1, T}$ が $T$ 上平坦である点の集合は、
% P05-EN-CH38-0238: frozen base S is wrong after base change to T; the comparison concerns Coker(alpha_1)_T over T.
$(Y_1)_T$ のうち $\Coker(\alpha_1)_T$ が $S$ 上平坦である点の集合と同じである。
こうして問題は長さ $n - 1$ の完全デヴィサージュをもつ層
$\Coker(\alpha_1)$ に帰着され、帰納法が使える。

\medskip\noindent
$d_1 < d$ ならば $F_d$ は $S$ により表現され、主張が従う。

\medskip\noindent
最後は $d_1 = d$ の場合である。この場合は Lemma \ref{flat-lemma-compare-H-F} と
Lemma \ref{flat-lemma-free-at-generic-points-representable} を合わせれば従う。
\end{proof}

\begin{theorem}
\label{flat-theorem-flat-dimension-n-representable}
状況 \ref{flat-situation-flat-dimension-n} において、さらに $f$ は有限表示、
$\mathcal{F}$ は有限表示 $\mathcal{O}_X$-加群で、$\mathcal{F}$ は $S$ に関して純粋であると
仮定する。このとき $F_n$ は有限表示モノ射 $Z_n \to S$ により表現可能である。
\end{theorem}

\begin{proof}
Lemma \ref{flat-lemma-flat-dimension-n} により、関手 $F_n$ は fppf 位相の層である。
有限表示モノ射は分離かつ準有限であることに注意する (Morphisms, Lemma
\ref{morphisms-lemma-monomorphism-loc-finite-type-loc-quasi-finite})。
したがって Descent, Lemma \ref{descent-lemma-descent-data-sheaves}、
More on Morphisms, Lemma
\ref{more-morphisms-lemma-separated-locally-quasi-finite-morphisms-fppf-descend}、および
Descent, Lemmas \ref{descent-lemma-descending-property-monomorphism} と
\ref{descent-lemma-descending-property-finite-presentation} を合わせると、
問題は $S$ 上の \'etale 位相について局所的である。

\medskip\noindent
特に状況は $S$ 上の Zariski 位相について局所的なので、$S$ はアフィンと仮定してよい。
この場合 $f$ のファイバーの次元は上に有界だから、十分大きい $n$ に対して $F_n$ は
表現可能である。そこで $n$ に関する降下帰納法を用いる。
$F_{n + 1}$ が有限表示モノ射 $Z_{n + 1} \to S$ により表現可能だと仮定する。
基底変換 $X_{n + 1} = Z_{n + 1} \times_S X$ と、$\mathcal{F}$ の $X_{n + 1}$ への
引き戻し $\mathcal{F}_{n + 1}$ を考える。$Z_{n + 1} \to S$ は有限表示モノ射なので
準有限である。したがって Lemma \ref{flat-lemma-quasi-finite-base-change} により
$\mathcal{F}_{n + 1}$ は $Z_{n + 1}$ に関して純粋である。
$F_n$ は $F_{n + 1}$ の部分関手なので、$F_n$ に対する結果を証明するには、状況
$\mathcal{F}_{n + 1}/X_{n + 1}/Z_{n + 1}$ に対応する関手について結果を証明すれば十分である。
このようにして、
% P05-EN-CH38-0239: frozen S_{n+1}=S conflicts with the representing object Z_{n+1} used throughout this induction step.
$S_{n + 1} = S$ の場合に $F_n$ に対する結果を証明することへ帰着される。すなわち、
$\mathcal{F}$ は次元 $\geq n + 1$ で $S$ 上平坦であると仮定してよい。

\medskip\noindent
$n$ を固定し、$\mathcal{F}$ は次元 $\geq n + 1$ で $S$ 上平坦であると仮定する。
証明を終えるには、$F_n$ が有限表示モノ射 $Z_n \to S$ により表現可能であることを
示さなければならない。問題は $S$ 上の \'etale 位相について局所的なので、
各 $s \in S$ に対し、$S'$ への基底変換後に結果が成り立つような基本 \'etale 近傍
$(S', s') \to (S, s)$ が存在することを示せば十分である。
したがって Lemma \ref{flat-lemma-existence-complete} により、各 $j$ に対して $s$ 上の
$\mathcal{F}_j/Y_j/S$ の完全デヴィサージュが存在し、かつ
$X_s \subset \bigcup h_j(Y_j)$ となる \'etale 射 $h_j : Y_j \to X$、
$j = 1, \ldots, m$ が存在すると仮定してよい。ここで $\mathcal{F}_j$ は
$\mathcal{F}$ の $Y_j$ への引き戻しである。Lemma
\ref{flat-lemma-localize-flat-dimension-n} により、層 $\mathcal{F}_j$ は依然として
% P05-EN-CH38-0240: frozen phrase "flat over in dimensions" has a redundant "over"; Japanese renders the determined grammar.
次元 $\geq n + 1$ で $S$ 上平坦である。$W = \bigcup h_j(Y_j)$ とおくと、これは
$X$ の準コンパクト開部分である。$\mathcal{F}$ は $X_s$ に沿って純粋なので、
$$
E = \{t \in S \mid \text{Ass}_{X_t}(\mathcal{F}_t) \subset W \}.
$$
は $s$ のすべての一般化を含む。More on Morphisms, Lemma
\ref{more-morphisms-lemma-relative-assassin-constructible} により、$E$ は $S$ の
構成可能部分集合である。上で $\Spec(\mathcal{O}_{S, s}) \subset E$ を示した。
Morphisms, Lemma \ref{morphisms-lemma-constructible-containing-open} により、$E$ は
$s$ の開近傍を含む。したがって $S$ を縮小した後 $E = S$ と仮定してよい。
Lemma \ref{flat-lemma-localize-flat-dimension-n} により、$X = \coprod Y_j$ および
$\mathcal{F} = \coprod \mathcal{F}_j$ に付随する関手 $F_n$ に対して補題を証明すれば十分である。
$F_{j, n}$ を $Y_j \to S$ と層
% P05-EN-CH38-0241: frozen sheaf index i conflicts with the functor index j and the immediately preceding coproduct F_j.
$\mathcal{F}_i$ に対する関手とすれば、$F_n = \prod F_{j, n}$ である。
したがって各 $F_{j, n}$ がある有限表示モノ射 $Z_{j, n} \to S$ により表現可能であることを
証明すれば十分である。このとき
$$
Z_n = Z_{1, n} \times_S \ldots \times_S Z_{m, n}
$$
だからである。以上により、定理は Lemma
\ref{flat-lemma-flat-dimension-n-representable} で扱った特殊な場合へ帰着された。
\end{proof}

\noindent
次の補題では、普遍平坦化の言葉で定理の意味を明示する。

\begin{lemma}
\label{flat-lemma-when-universal-flattening}
$f : X \to S$ をスキームの射、$\mathcal{F}$ を準連接 $\mathcal{O}_X$-加群とする。
\begin{enumerate}
\item $f$ が有限表示、$\mathcal{F}$ が有限表示 $\mathcal{O}_X$-加群で、
$\mathcal{F}$ が $S$ に関して純粋ならば、$\mathcal{F}$ の普遍平坦化 $S' \to S$ が存在する。
さらに $S' \to S$ は有限表示モノ射である。
\item $f$ が有限表示で $X$ が $S$ に関して純粋ならば、$X$ の普遍平坦化
$S' \to S$ が存在する。さらに $S' \to S$ は有限表示モノ射である。
\item $f$ が固有かつ有限表示で、$\mathcal{F}$ が有限表示 $\mathcal{O}_X$-加群ならば、
$\mathcal{F}$ の普遍平坦化 $S' \to S$ が存在する。さらに $S' \to S$ は
有限表示モノ射である。
\item $f$ が固有かつ有限表示ならば、$X$ の普遍平坦化 $S' \to S$ が存在する。
\end{enumerate}
\end{lemma}

\begin{proof}
これらの主張は、$F_0 = F_{flat}$ に Theorem
\ref{flat-theorem-flat-dimension-n-representable} を適用し、$f$ が固有ならば
$\mathcal{F}$ は基底上自動的に純粋であることを用いれば直ちに従う。
Lemma \ref{flat-lemma-proper-pure} を参照せよ。
\end{proof}







\section{Grothendieck の存在定理 IV}
\label{flat-section-existence}

\noindent
この節では Cohomology of Schemes, Sections \ref{coherent-section-existence},
\ref{coherent-section-existence-proper}, および
\ref{coherent-section-existence-proper-support} の議論を続ける。
以下の状況で考える。

\begin{situation}
\label{flat-situation-existence}
ここでは、全射な遷移写像をもち、その核が局所冪零である環の逆系 $(A_n)$ が
与えられている。$A = \lim A_n$ とおく。また、$A$ 上分離的かつ有限表示な
スキーム $X$ がある。$X_n = X \times_{\Spec(A)} \Spec(A_n)$ とおき、これを
$X$ の閉部分スキームとみなす。さらに系 $(\mathcal{F}_n, \varphi_n)$ が与えられていると仮定する。
ここで $\mathcal{F}_n$ は有限表示 $\mathcal{O}_{X_n}$-加群で、$A_n$ 上平坦であり、
その台は $A_n$ 上固有で、かつ
$$
\varphi_n :
\mathcal{F}_n \otimes_{\mathcal{O}_{X_n}} \mathcal{O}_{X_{n - 1}}
\longrightarrow
\mathcal{F}_{n - 1}
$$
は同型である（記法には Morphisms, Lemma
\ref{morphisms-lemma-i-star-equivalence} の圏同値を用いている）。
\end{situation}

\noindent
目標は、すべての $n$ に対して
$\mathcal{F}_n = \mathcal{F} \otimes_{\mathcal{O}_X} \mathcal{O}_{X_n}$
を満たす $X$ 上の準連接層 $\mathcal{F}$ を見いだせるかを調べることである。

\begin{lemma}
\label{flat-lemma-compute-what-it-should-be}
Situation \ref{flat-situation-existence} において
$$
K = R\lim_{D_\QCoh(\mathcal{O}_X)}(\mathcal{F}_n) =
DQ_X(R\lim_{D(\mathcal{O}_X)}\mathcal{F}_n)
$$
を考える。このとき $K$ は $D^b_{\QCoh}(\mathcal{O}_X)$ に属し、実際、
$K$ のコホモロジー層は次数 $\geq 0$ でのみ非零である。
\end{lemma}

\begin{proof}
Derived Categories of Schemes, Example
\ref{perfect-example-inverse-limit-quasi-coherent} の特殊な場合である。
\end{proof}

\begin{lemma}
\label{flat-lemma-compute-against-perfect}
Situation \ref{flat-situation-existence} において、$K$ を Lemma
\ref{flat-lemma-compute-what-it-should-be} のものとする。$D(\mathcal{O}_X)$ の
任意の完全対象 $E$ に対して次が成り立つ。
\begin{enumerate}
\item $M = R\Gamma(X, K \otimes^\mathbf{L} E)$ は $D(A)$ の完全対象であり、
$R\Gamma(X_n, \mathcal{F}_n \otimes^\mathbf{L} E|_{X_n}) =
M \otimes_A^\mathbf{L} A_n$
という $D(A_n)$ における標準同型がある。
\item $N = R\Hom_X(E, K)$ は $D(A)$ の完全対象であり、
$R\Hom_{X_n}(E|_{X_n}, \mathcal{F}_n) = N \otimes_A^\mathbf{L} A_n$
という $D(A_n)$ における標準同型がある。
\end{enumerate}
いずれの主張でも、$E|_{X_n}$ は $E$ の $X_n$ への導来引き戻しを表す。
\end{lemma}

\begin{proof}
(2) を証明する。$E_n = E|_{X_n}$ および
$N_n = R\Hom_{X_n}(E_n, \mathcal{F}_n)$ と書く。
$R\Hom_{X_n}(-, -)$ は
$R\Gamma(X_n, R\SheafHom(-, -))$ に等しいことを思い出そう。
Cohomology, Section \ref{cohomology-section-global-RHom} を参照せよ。
したがって Derived Categories of Schemes, Lemma
\ref{perfect-lemma-base-change-RHom-perfect} により、$N_n$ は $D(A_n)$ の
完全対象で、その構成は基底変換と可換である。ゆえに $\varphi_n$ から得られる写像
$N_n \otimes_{A_n}^\mathbf{L} A_{n - 1} \to N_{n - 1}$
は同型である。More on Algebra, Lemma
\ref{more-algebra-lemma-Rlim-perfect-gives-perfect} により、$R\lim N_n$ は
完全であり、これを $A_n$ へ基底変換すると $N_n$ が復元される。
一方、三角圏の完全関手
$R\Hom_X(E, -) : D_\QCoh(\mathcal{O}_X) \to D(A)$
は直積と可換であり、したがって導来極限とも可換である。よって
$$
R\Hom_X(E, K) =
R\lim R\Hom_X(E, \mathcal{F}_n) =
R\lim R\Hom_X(E_n, \mathcal{F}_n) =
R\lim N_n
$$
これで (2) が証明された。(1) は Cohomology, Lemma
\ref{cohomology-lemma-dual-perfect-complex} を用いて (2) に帰着すれば従う。
\end{proof}

\begin{lemma}
\label{flat-lemma-relative-pseudo-coherence}
Situation \ref{flat-situation-existence} において、$K$ を Lemma
\ref{flat-lemma-compute-what-it-should-be} のものとする。このとき $K$ は
$A$ に関して擬コヒーレントである。
\end{lemma}

\begin{proof}
% P05-EN-CH38-0242: frozen source has the typographical form "Combinging" for "Combining".
Lemma \ref{flat-lemma-compute-against-perfect} と Derived Categories of Schemes,
Lemma \ref{perfect-lemma-perfect-enough} を組み合わせると、$D(\mathcal{O}_X)$ の
すべての擬コヒーレントな $E$ に対して
$R\Gamma(X, K \otimes^\mathbf{L} E)$ が $D(A)$ で擬コヒーレントであることが分かる。
したがって補題は More on Morphisms, Lemma
\ref{more-morphisms-lemma-characterize-pseudo-coherent} から従う。
\end{proof}

\begin{lemma}
\label{flat-lemma-compute-over-affine}
Situation \ref{flat-situation-existence} において、$K$ を Lemma
\ref{flat-lemma-compute-what-it-should-be} のものとする。任意の準コンパクト開集合
$U \subset X$ に対して
$$
R\Gamma(U, K) \otimes_A^\mathbf{L} A_n =
R\Gamma(U_n, \mathcal{F}_n)
$$
が $D(A_n)$ で成り立つ。ここで $U_n = U \cap X_n$ である。
\end{lemma}

\begin{proof}
$n$ を固定する。Derived Categories of Schemes, Lemma
\ref{perfect-lemma-computing-sections-as-colim} により、$X$ 上の完全複体の系
$E_m$ で
$R\Gamma(U, K) = \text{hocolim} R\Gamma(X, K \otimes^\mathbf{L} E_m)$
を満たすものが存在する。実際、この公式は $K$ に対してだけでなく、
$D_\QCoh(\mathcal{O}_X)$ の任意の対象に対して成り立つ。
これを $\mathcal{F}_n$ に適用すると
\begin{align*}
R\Gamma(U_n, \mathcal{F}_n)
& =
R\Gamma(U, \mathcal{F}_n) \\
& =
\text{hocolim}_m R\Gamma(X, \mathcal{F}_n \otimes^\mathbf{L} E_m) \\
& =
\text{hocolim}_m R\Gamma(X_n, \mathcal{F}_n \otimes^\mathbf{L} E_m|_{X_n})
\end{align*}
を得る。Lemma \ref{flat-lemma-compute-against-perfect} と、関手
$- \otimes_A^\mathbf{L} A_n$ がホモトピー余極限と可換であることを用いれば、結論を得る。
\end{proof}

\begin{lemma}
\label{flat-lemma-finitely-presented}
Situation \ref{flat-situation-existence} において、$K$ を Lemma
\ref{flat-lemma-compute-what-it-should-be} のものとする。
% P05-EN-CH38-0243: frozen source omits "by" in "Denote X_0 ... the closed subset".
$X_0 \subset X$ を、$\Spec(A)$ の閉部分集合
$\Spec(A_1) = \Spec(A_2) = \ldots$ 上にある点からなる閉部分集合とする。
$X_0$ を含む開集合 $W \subset X$ で、次を満たすものが存在する。
\begin{enumerate}
\item $i = 0$ でない限り $H^i(K)|_W$ は零である。
\item $\mathcal{F} = H^0(K)|_W$ は有限表示である。
\item $\mathcal{F}_n = \mathcal{F} \otimes_{\mathcal{O}_X} \mathcal{O}_{X_n}$ である。
\end{enumerate}
\end{lemma}

\begin{proof}
$n \geq 1$ を固定する。構成により $D_\QCoh(\mathcal{O}_X)$ に標準写像
$K \to \mathcal{F}_n$ があり、したがって準連接層の標準写像
$H^0(K) \to \mathcal{F}_n$ がある。これが (3) の意味を説明する。

\medskip\noindent
$x \in X_0$ を一点とする。(1), (2), (3) が成り立つ $x$ の開近傍 $W$ を
見いだそう。$X_0$ は準コンパクトなので、これで補題が証明される。
$U \subset X$ を $x$ のアフィン開近傍とし、$U = \Spec(B)$ と書く。
$P$ が $A$ 上滑らかであるような全射 $P \to B$ を選ぶ。
Lemma \ref{flat-lemma-relative-pseudo-coherence} と相対擬コヒーレンスの定義により、
有限自由 $P$-加群からなる上に有界な複体 $F^\bullet$ で $Ri_*K$ を表すものが存在する。
ここで $i : U \to \Spec(P)$ はこの表示が誘導する閉埋め込みである。
$M_n$ を $\mathcal{F}_n|_U$ に対応する $B$-加群とする。
Lemma \ref{flat-lemma-compute-over-affine} により
$$
H^i(F^\bullet \otimes_A A_n) =
\left\{
\begin{matrix}
0 & \text{if} & i \not = 0 \\
M_n & \text{if} & i = 0
\end{matrix}
\right.
$$
% P05-EN-CH38-0244: frozen source reuses i for the complex degree immediately after naming the closed immersion i.
$F^i$ が非零となる最大の添字を $i$ とする。$i \leq 0$ ならば (1), (2), (3) が成り立つ。
そうでなければ $i > 0$ であり、写像
$$
F^{i - 1} \to F^i
$$
% P05-EN-CH38-0246: frozen source says "in the point x" rather than "at the point x".
の点 $x$ における階数は最大である。したがって $\Spec(P)$ における $x$ の
ある開近傍上で階数は最大である。ゆえに $P$ を単項局所化で置き換えることにより、
表示した写像が全射であると仮定してよい。$F^i$ は有限自由なので、分裂
$F^{i - 1} = F' \oplus F^i$ を選べる。このとき $F^\bullet$ を複体
$$
\ldots \to F^{i - 2} \to F' \to 0 \to \ldots
$$
で置き換えられ、$i$ に関する帰納法で結論を得る。
\end{proof}

\begin{lemma}
\label{flat-lemma-proper-support}
Situation \ref{flat-situation-existence} において、$K$ を Lemma
\ref{flat-lemma-compute-what-it-should-be} のものとし、$W \subset X$ を Lemma
\ref{flat-lemma-finitely-presented} のものとする。
$\mathcal{F} = H^0(K)|_W$ とおく。このとき、必要なら開集合 $W$ を縮小すれば、
$\mathcal{F}$ の台は $A$ 上固有である。
\end{lemma}

\begin{proof}
$n \geq 1$ を固定し、$I_n = \Ker(A \to A_n)$ とおく。
More on Algebra, Lemma \ref{more-algebra-lemma-limit-henselian} により、
対 $(A, I_n)$ は Hensel 的である。$Z \subset W$ を $\mathcal{F}$ の台とする。
$\mathcal{F}$ は有限表示なので、これは閉部分集合である。
Lemma \ref{flat-lemma-finitely-presented} の (3) により、
$Z \times_{\Spec(A)} \Spec(A_n)$ は $\mathcal{F}_n$ の台に等しく、したがって
% P05-EN-CH38-0245: frozen source uses undefined I here after defining I_n.
$\Spec(A/I)$ 上固有である。More on Morphisms, Lemma
\ref{more-morphisms-lemma-split-off-proper-part-henselian} により、
$Z_1, Z_2$ が $Z$ で開かつ閉、$Z_1$ が $A$ 上固有、かつ
$Z_1 \times_{\Spec(A)} \Spec(A/I_n)$ が $\mathcal{F}_n$ の台に等しくなるように
$Z = Z_1 \amalg Z_2$ と書ける。言い換えると、$Z_2$ は $X_0$ と交わらない。
したがって $W$ を $W \setminus Z_2$ で置き換えれば補題を得る。
\end{proof}

\begin{lemma}
\label{flat-lemma-monomorphism-isomorphism}
$A = \lim A_n$ を、遷移写像が全射でその核が局所冪零である環の系の極限とする。
$S = \Spec(A)$ とおく。$T \to S$ を局所有限型なモノ射とする。
すべての $n$ に対して $\Spec(A_n) \to S$ が $T$ を経由するならば、$T = S$ である。
\end{lemma}

\begin{proof}
$S_n = \Spec(A_n)$ とおく。$T_0 \subset T$ を分解 $S_n \to T$ の共通像とする。
このとき $T_0$ は準コンパクトである。$T' \subset T$ を $T_0$ を含む
準コンパクト開集合とする。このとき $S_n \to T$ は $T'$ を経由する。
$T' = S$ を示せば $T' = T = S$ となる。したがって $T$ は準コンパクトと仮定してよい。

\medskip\noindent
$T$ が準コンパクトであると仮定する。この場合 $T \to S$ は分離的かつ準有限である
（Morphisms, Lemma
\ref{morphisms-lemma-monomorphism-loc-finite-type-loc-quasi-finite}）。
Zariski の主定理を More on Morphisms, Lemma
\ref{more-morphisms-lemma-quasi-finite-separated-pass-through-finite} の形で用いて、
$W \to S$ が有限で $T \to W$ が開埋め込みとなる分解 $T \to W \to S$ を選ぶ。
$W = \Spec(B)$ と書く。（一意な）分解 $S_n \to T$ を $W$ への射とみなすと
$$
A \longrightarrow B \longrightarrow \lim A_n = A
$$
を得る。右側の矢印から得られる射
$h : S = \Spec(A) \to \Spec(B) = W$ を考える。このとき
$$
T \times_{W, h} S
$$
は $S$ の開部分スキームで、すべての $n$ に対して $S_n \to S$ の像を含む。
証明を終えるには、ある $n \geq 1$ に対して $S_n \to S$ の像を含む任意の
開集合 $U \subset S$ が $S$ に等しいことを示せば十分である。これは
$(A, \Ker(A \to A_n))$ が Hensel 対であり（More on Algebra, Lemma
\ref{more-algebra-lemma-limit-henselian}）、したがって $S$ のすべての閉点が
$S_n \to S$ の像に含まれることから従う。
\end{proof}

\begin{theorem}[Grothendieck の存在定理]
\label{flat-theorem-existence}
Situation \ref{flat-situation-existence} において、$A$ 上平坦で台が $A$ 上固有な
有限表示 $\mathcal{O}_X$-加群 $\mathcal{F}$ で、すべての $n$ に対して
$\mathcal{F}_n = \mathcal{F} \otimes_{\mathcal{O}_X} \mathcal{O}_{X_n}$
が写像 $\varphi_n$ と両立するものが存在する。
\end{theorem}

\begin{proof}
Lemmas \ref{flat-lemma-compute-what-it-should-be},
\ref{flat-lemma-compute-against-perfect},
\ref{flat-lemma-relative-pseudo-coherence},
\ref{flat-lemma-compute-over-affine},
\ref{flat-lemma-finitely-presented}, および
\ref{flat-lemma-proper-support} を適用すると、すべての $\Spec(A_n)$ 上の点を含む
開部分スキーム $W \subset X$ と、台が $A$ 上固有な有限表示
$\mathcal{O}_W$-加群 $\mathcal{F}$ で
$\mathcal{F}_n = \mathcal{F} \otimes_{\mathcal{O}_W} \mathcal{O}_{X_n}$
がすべての $n \geq 1$ に対して成り立つものを得る（$X_n \subset W$ なのでこれは意味をもつ）。
Lemma \ref{flat-lemma-proper-pure} により、$\mathcal{F}$ は $\Spec(A)$ に関して
普遍的に純粋である。Theorem \ref{flat-theorem-flat-dimension-n-representable}
（説明については Lemma \ref{flat-lemma-when-universal-flattening} を参照）により、
$\mathcal{F}$ の普遍平坦化 $S' \to \Spec(A)$ が存在し、さらに射
$S' \to \Spec(A)$ は有限表示モノ射である。$\mathcal{F}$ の $\Spec(A_n)$ への
基底変換は $\mathcal{F}_n$ なので、各 $n$ に対して
$\Spec(A_n) \to \Spec(A)$ は（一意に）$S'$ を経由する。
Lemma \ref{flat-lemma-monomorphism-isomorphism} により $S' = \Spec(A)$ である。
これは $\mathcal{F}$ が $A$ 上平坦であることを意味する。
最後に、$\mathcal{F}$ のスキーム論的台 $Z$ は $\Spec(A)$ 上固有なので、
射 $Z \to X$ は閉である。したがって押し出し $(W \to X)_*\mathcal{F}$ は
$W$ 上に台をもち、望む性質をすべて備える。
\end{proof}






\section{Grothendieck の存在定理 V}
\label{flat-section-existence-derived}

\noindent
この節では、Section \ref{flat-section-existence} で準連接加群に対して用いた方法に従い、
導来圏における Grothendieck の存在定理の類似を証明する。
古典的な場合は Cohomology of Schemes, Sections
\ref{coherent-section-existence}, \ref{coherent-section-existence-proper}, および
\ref{coherent-section-existence-proper-support} で論じられている。
以下の状況で考える。

\begin{situation}
\label{flat-situation-existence-derived}
ここでは、全射な遷移写像をもち、その核が局所冪零である環の逆系 $(A_n)$ が
与えられている。$A = \lim A_n$ とおく。また、$A$ 上固有、平坦かつ有限表示な
スキーム $X$ がある。$X_n = X \times_{\Spec(A)} \Spec(A_n)$ とおき、これを
$X$ の閉部分スキームとみなす。さらに系 $(K_n, \varphi_n)$ が与えられていると仮定する。
ここで $K_n$ は $D(\mathcal{O}_{X_n})$ の擬コヒーレント対象であり、
$$
\varphi_n : K_n \longrightarrow K_{n - 1}
$$
は $D(\mathcal{O}_{X_n})$ における写像で、$D(\mathcal{O}_{X_{n - 1}})$ に同型
$K_n \otimes_{\mathcal{O}_{X_n}}^\mathbf{L} \mathcal{O}_{X_{n - 1}}
\to K_{n - 1}$ を誘導する。
\end{situation}

\noindent
より正確には、
$\varphi_n : K_n \to Ri_{n - 1, *}K_{n - 1}$
と書くべきである。ここで $i_{n - 1} : X_{n - 1} \to X_n$ は包含射であり、
この記法では条件は随伴写像 $Li_{n - 1}^*K_n \to K_{n - 1}$ が同型であることをいう。
目標は、すべての $n$ に対して
$K_n = K \otimes_{\mathcal{O}_X}^\mathbf{L} \mathcal{O}_{X_n}$ を満たす
擬コヒーレントな $K \in D(\mathcal{O}_X)$ を見いだすことである
（同じ記法の濫用を用いる）。

\begin{lemma}
\label{flat-lemma-compute-what-it-should-be-derived}
Situation \ref{flat-situation-existence-derived} において
$$
K = R\lim_{D_\QCoh(\mathcal{O}_X)}(K_n) =
DQ_X(R\lim_{D(\mathcal{O}_X)} K_n)
$$
を考える。このとき $K$ は $D^-_{\QCoh}(\mathcal{O}_X)$ に属する。
\end{lemma}

\begin{proof}
$X$ は準コンパクトかつ準分離的なので関手 $DQ_X$ は存在する。
Derived Categories of Schemes, Lemma \ref{perfect-lemma-better-coherator} を参照せよ。
$DQ_X$ は右随伴なので直積と可換であり、したがって導来極限とも可換である。
ゆえに補題の主張にある等式を得る。

\medskip\noindent
Derived Categories of Schemes, Lemma
\ref{perfect-lemma-boundedness-better-coherator} により、関手 $DQ_X$ の
コホモロジー次元は有界である。したがって
$R\lim K_n \in D^-(\mathcal{O}_X)$ を示せば十分である。
このために $U \subset X$ をアフィン開集合とする。このとき標準完全列
$$
0 \to
R^1\lim H^{m - 1}(U, K_n) \to H^m(U, R\lim K_n) \to
\lim H^m(U, K_n) \to 0
$$
がある。Cohomology, Lemma \ref{cohomology-lemma-RGamma-commutes-with-Rlim} を参照せよ。
$U$ はアフィンで $K_n$ は擬コヒーレントであり、したがって Derived Categories of Schemes,
Lemma \ref{perfect-lemma-pseudo-coherent} によりそのコホモロジー層は準連接であるから、
Derived Categories of Schemes, Lemma \ref{perfect-lemma-affine-compare-bounded} により
$H^m(U, K_n) = H^m(K_n)(U)$ である。よって $K_n$ が $n$ に依存せず
上に有界であることを示せば十分である。

\medskip\noindent
$K_n$ は擬コヒーレントなので $K_n \in D^-(\mathcal{O}_{X_n})$ である。
$H^{a_n}(K_n)$ が非零となる最大の整数を $a_n$ とする。
もちろん $a_1 \leq a_2 \leq a_3 \leq \ldots$ である。
$H^{a_n}(K_n)$ は有限表示 $\mathcal{O}_{X_n}$-加群であることに注意する
（Cohomology, Lemma \ref{cohomology-lemma-finite-cohomology}）。また
$H^{a_n}(K_{n - 1}) =
H^{a_n}(K_n) \otimes_{\mathcal{O}_{X_n}} \mathcal{O}_{X_{n - 1}}$ である。
$X_{n - 1} \to X_n$ は厚化なので、Nakayama の補題
（Algebra, Lemma \ref{algebra-lemma-NAK}）により、もし
$H^{a_n}(K_n) \otimes_{\mathcal{O}_{X_n}} \mathcal{O}_{X_{n - 1}}$
が零ならば $H^{a_n}(K_n)$ も零である。したがってすべての $n$ に対して
$a_n = a_{n - 1}$ であり、結論を得る。
\end{proof}

\begin{lemma}
\label{flat-lemma-compute-against-perfect-derived}
Situation \ref{flat-situation-existence-derived} において、$K$ を Lemma
\ref{flat-lemma-compute-what-it-should-be-derived} のものとする。
$D(\mathcal{O}_X)$ の任意の完全対象 $E$ に対して、
% P05-EN-CH38-0257: frozen source calls the derived object M = RGamma(...) "the cohomology" rather than an object/derived global sections.
対象
$$
M = R\Gamma(X, K \otimes^\mathbf{L} E)
$$
は $D(A)$ の擬コヒーレント対象であり、標準同型
$$
R\Gamma(X_n, K_n \otimes^\mathbf{L} E|_{X_n}) = M \otimes_A^\mathbf{L} A_n
$$
が $D(A_n)$ において存在する。ここで $E|_{X_n}$ は $E$ の $X_n$ への導来引き戻しを表す。
\end{lemma}

\begin{proof}
$E_n = E|_{X_n}$ および
$M_n = R\Gamma(X_n, K_n \otimes^\mathbf{L} E|_{X_n})$ と書く。
Derived Categories of Schemes, Lemma
\ref{perfect-lemma-flat-proper-pseudo-coherent-direct-image-general} により、
$M_n$ は $D(A_n)$ の擬コヒーレント対象であり、その構成は基底変換と可換である。
したがって $\varphi_n$ から得られる写像
$M_n \otimes_{A_n}^\mathbf{L} A_{n - 1} \to M_{n - 1}$
は同型である。More on Algebra, Lemma
\ref{more-algebra-lemma-Rlim-pseudo-coherent-gives-pseudo-coherent} により、
$R\lim M_n$ は擬コヒーレントであり、これを $A_n$ へ基底変換すると $M_n$ が復元される。
一方、三角圏の完全関手
$R\Gamma(X, -) : D_\QCoh(\mathcal{O}_X) \to D(A)$
は直積と可換であり、したがって導来極限とも可換である。よって
$$
R\Gamma(X, E \otimes^\mathbf{L} K) =
R\lim R\Gamma(X,  E \otimes^\mathbf{L} K_n) =
R\lim R\Gamma(X_n, E_n \otimes^\mathbf{L} K_n) =
R\lim M_n
$$
となり、望む結論を得る。
\end{proof}

\begin{lemma}
\label{flat-lemma-relative-pseudo-coherence-derived}
Situation \ref{flat-situation-existence-derived} において、$K$ を Lemma
\ref{flat-lemma-compute-what-it-should-be-derived} のものとする。このとき $K$ は
$X$ 上擬コヒーレントである。
\end{lemma}

\begin{proof}
% P05-EN-CH38-0247: frozen source has the typographical form "Combinging" for "Combining".
Lemma \ref{flat-lemma-compute-against-perfect-derived} と
Derived Categories of Schemes, Lemma
\ref{perfect-lemma-perfect-enough}
を組み合わせると、$D(\mathcal{O}_X)$ のすべての擬コヒーレントな $E$ に対して
$R\Gamma(X, K \otimes^\mathbf{L} E)$ が $D(A)$ で擬コヒーレントであることが分かる。
したがって More on Morphisms, Lemma
\ref{more-morphisms-lemma-characterize-pseudo-coherent} により、$K$ は $A$ に関して
擬コヒーレントである。
% P05-EN-CH38-0248: frozen source has the redundant construction "is of flat".
$X$ は $A$ 上平坦かつ有限表示なので、これは $X$ 上擬コヒーレントであることと同値である。
More on Morphisms, Lemma
\ref{more-morphisms-lemma-check-relative-pseudo-coherence-on-charts} を参照せよ。
\end{proof}

\begin{lemma}
\label{flat-lemma-compute-over-affine-derived}
Situation \ref{flat-situation-existence-derived} において、$K$ を Lemma
\ref{flat-lemma-compute-what-it-should-be-derived} のものとする。任意の準コンパクト開集合
$U \subset X$ に対して
$$
R\Gamma(U, K) \otimes_A^\mathbf{L} A_n =
R\Gamma(U_n, K_n)
$$
が $D(A_n)$ で成り立つ。ここで $U_n = U \cap X_n$ である。
\end{lemma}

\begin{proof}
$n$ を固定する。Derived Categories of Schemes, Lemma
\ref{perfect-lemma-computing-sections-as-colim} により、$X$ 上の完全複体の系 $E_m$ で
$R\Gamma(U, K) = \text{hocolim} R\Gamma(X, K \otimes^\mathbf{L} E_m)$
を満たすものが存在する。実際、この公式は $K$ に対してだけでなく、
$D_\QCoh(\mathcal{O}_X)$ の任意の対象に対して成り立つ。
これを $K_n$ に適用すると
\begin{align*}
R\Gamma(U_n, K_n)
& =
R\Gamma(U, K_n) \\
& =
\text{hocolim}_m R\Gamma(X, K_n \otimes^\mathbf{L} E_m) \\
& =
\text{hocolim}_m R\Gamma(X_n, K_n \otimes^\mathbf{L} E_m|_{X_n})
\end{align*}
を得る。Lemma \ref{flat-lemma-compute-against-perfect-derived} と、関手
$- \otimes_A^\mathbf{L} A_n$ がホモトピー余極限と可換であることを用いれば、結論を得る。
\end{proof}

\begin{theorem}[導来 Grothendieck 存在定理]
\label{flat-theorem-existence-derived}
Situation \ref{flat-situation-existence-derived} において、$D(\mathcal{O}_X)$ の
擬コヒーレント対象 $K$ で、すべての $n$ に対して
$K_n = K \otimes_{\mathcal{O}_X}^\mathbf{L} \mathcal{O}_{X_n}$ が
写像 $\varphi_n$ と両立して成り立つものが存在する。
\end{theorem}

\begin{proof}
Lemmas \ref{flat-lemma-compute-what-it-should-be-derived},
\ref{flat-lemma-compute-against-perfect-derived},
\ref{flat-lemma-relative-pseudo-coherence-derived}
を適用すると、$D(\mathcal{O}_X)$ の擬コヒーレント対象 $K$ を得る。
Lemma \ref{flat-lemma-compute-over-affine-derived} でアフィン開集合を選べば、
$K$ の $X_n$ への制限が $K_n$ であることが直ちに従う。
\end{proof}

\begin{remark}
\label{flat-remark-correct-generality}
この節の結果は一般化できる。$X \to \Spec(A)$ が分離的かつ有限表示であり、
$K_n$ が $A_n$ に関して擬コヒーレントで、その台が $A_n$ 上固有な
$X_n$ の閉部分集合に含まれることだけを仮定しても、おそらく正しい。
このとき得られる $K$ は $A$ に関して擬コヒーレントで、その台は $A$ 上固有な
閉部分集合に含まれる。この結果が必要になれば、ここで精密な主張を定式化して証明する。
\end{remark}






\section{吹き上げと平坦性}
\label{flat-section-blowup-flat}

\noindent
この節では「吹き上げの後、厳密変換は平坦になる」という形の結果について議論を続ける。
More on Algebra, Section \ref{more-algebra-section-blowup-flat} および
Divisors, Section \ref{divisors-section-blowup-flat} を参照せよ。
この節では次の（ほぼ標準的な）記法を用いる。$X \to S$ をスキームの射、
$\mathcal{F}$ を $X$ 上の準連接加群、$T \to S$ をスキームの射とするとき、
$\mathcal{F}$ の基底変換 $X_T = X \times_S T$ への引き戻しを $\mathcal{F}_T$ と表す。

\begin{remark}
\label{flat-remark-successive-blowups}
$S$ を準コンパクトかつ準分離的なスキーム、$f : X \to S$ をスキームの射、
$\mathcal{F}$ を $X$ 上の準連接加群とする。$U \subset S$ を準コンパクト開部分スキームとする。
$U$-許容吹き上げ $S' \to S$ が与えられたとき、$X$ の厳密変換を $X'$、
$\mathcal{F}$ の厳密変換を $\mathcal{F}'$ と表し、後者を $X'$ 上の準連接加群とみなす
（Divisors, Lemma \ref{divisors-lemma-strict-transform} による）。
$P$ を、$U$-許容吹き上げによる（上記の）厳密変換の下で安定な
$\mathcal{F}/X/S$ の性質とする。この節で扱う一般問題は次である：
$\mathcal{F}/X/S$ に補助条件を課した上で、厳密変換 $\mathcal{F}'/X'/S'$ が
$P$ をもつような $U$-許容吹き上げ $S' \to S$ の存在を示せ。

\medskip\noindent
一般的な方針は、$U$-許容吹き上げの合成が $U$-許容吹き上げであることを用いることである。
Divisors, Lemma \ref{divisors-lemma-composition-admissible-blowups} を参照せよ。
実際には、より精密な Divisors, Lemma
\ref{divisors-lemma-composition-finite-type-blowups} を用い、これを Divisors, Lemma
\ref{divisors-lemma-strict-transform-composition-blowups} と組み合わせる。
その結果、$U$-許容吹き上げの列
$$
S = S_0 \leftarrow S_1 \leftarrow \ldots \leftarrow S_n
$$
で、$\mathcal{F}_0 = \mathcal{F}$、$X_0 = X$ とおき、
$\mathcal{F}_i/X_i$ を $\mathcal{F}_{i - 1}/X_{i - 1}$ の厳密変換としたとき、
$P$ をもつ $\mathcal{F}_n/X_n/S_n$ に到達するものを見いだせば十分である。

\medskip\noindent
特に、$V(\mathcal{I}) = S \setminus U$ となる有限型準連接イデアル層
$\mathcal{I} \subset \mathcal{O}_S$ を選ぶ。
Properties, Lemma \ref{properties-lemma-quasi-coherent-finite-type-ideals} を参照せよ。
$S' \to S$ を $\mathcal{I}$ における吹き上げ、$E \subset S'$ を例外因子とする
（Divisors, Lemma \ref{divisors-lemma-blowing-up-gives-effective-Cartier-divisor}）。
すると問題は、台が
% P05-EN-CH38-0249: frozen source says X minus U although D and U are subschemes of S.
$X \setminus U$ である有効 Cartier 因子 $D \subset S$ が存在する場合に帰着される。
特に $U$ は $S$ でスキーム論的に稠密であると仮定してよい
（Divisors, Lemma \ref{divisors-lemma-complement-effective-Cartier-divisor}）。

\medskip\noindent
$P$ が $S$ 上局所的であると仮定する。すなわち、$S = \bigcup S_i$ が
準コンパクト開集合による有限開被覆で、各 $\mathcal{F}_{S_i}/X_{S_i}/S_i$ に対して
$P$ が成り立つならば、$\mathcal{F}/X/S$ に対しても $P$ が成り立つとする。
この場合、上の一般問題も $S$ 上局所的である。つまり、各 $s \in S$ に対して
$s$ の準コンパクト開近傍 $W$ で $\mathcal{F}_W/X_W/W$ に対する問題が解けるものを
見いだせるならば、$\mathcal{F}/X/S$ に対する問題も解ける。これは Divisors, Lemmas
\ref{divisors-lemma-extend-admissible-blowups} および
\ref{divisors-lemma-dominate-admissible-blowups} から従う。
\end{remark}

\begin{lemma}
\label{flat-lemma-helper-blowup-affine-space}
$R$ を環、$f \in R$ とし、$r\geq 0$ を整数とする。
$R \to S$ を環準同型、$M$ を $S$-加群とする。次を仮定する。
\begin{enumerate}
\item $R \to S$ は有限表示かつ平坦である。
\item 各ファイバー環 $S \otimes_R \kappa(\mathfrak p)$ は
% P05-EN-CH38-0250: frozen source says the fibre ring is geometrically integral over R rather than over kappa(p).
$R$ 上幾何学的整である。
\item $M$ は有限 $S$-加群である。
\item $M_f$ は有限表示 $S_f$-加群である。
\item すべての $\mathfrak p \in R$ で $f \not \in \mathfrak p$ となるものについて、
$\mathfrak q = \mathfrak pS$ とおくと $M_{\mathfrak q}$ は $S_\mathfrak q$ 上階数 $r$ の自由加群である。
\end{enumerate}
このとき $V(f) = V(I)$ を満たす有限生成イデアル $I \subset R$ で、
すべての $a \in I$ に対して $R' = R[\frac{I}{a}]$ とおくと、商
$$
M' = (M \otimes_R R')/a\text{-power torsion}
$$
が $S' = S \otimes_R R'$ 上で次を満たすものが存在する：各素イデアル
$\mathfrak p' \subset R'$ に対して、$g \not \in \mathfrak p'S'$ かつ
$M'_g$ が $S'_g$ 上階数 $r$ の自由加群となる $g \in S'$ が存在する。
\end{lemma}

\begin{proof}
この補題は More on Algebra, Lemma \ref{more-algebra-lemma-blowup-module} の一般化である。
まずそちらの証明を読むことを勧める。
% P05-EN-CH38-0251: frozen source attributes finite generation of M and the resulting surjection to (1), but it is hypothesis (3).
(1) により可能であるとして、全射 $S^{\oplus n} \to M$ を選ぶ。
$f$ を可逆化すると $S^{\oplus n}/K \to M$ が同型になるような有限部分加群
$K \subset \Ker(S^{\oplus n} \to M)$ を選ぶ。これは (4) により可能である。
$M_1 = S^{\oplus n}/K$ とおき、$M_1$ に対して補題を証明できると仮定する。
$I \subset R$ を対応するイデアルとする。このとき $a \in I$ に対して写像
$$
M_1' = (M_1 \otimes_R R')/a\text{-power torsion}
\longrightarrow
M' = (M \otimes_R R')/a\text{-power torsion}
$$
は全射である。また $R'_a = R_f$ なので、$R'$ で $a$ を可逆化すると同型になる。
Algebra, Lemma \ref{algebra-lemma-blowup-in-principal} を参照せよ。
しかし $a$ は $M'_1$ 上の非零因子だから、表示した写像は同型である。
したがって $M$ が有限表示 $S$-加群の場合に補題を証明すれば十分である。

\medskip\noindent
(3) を満たす有限表示 $S$-加群 $M$ を仮定する。このとき
$J = \text{Fit}_r(M) \subset S$ は有限生成イデアルである。
Lemma \ref{flat-lemma-fibres-irreducible-flat-projective-nonnoetherian} により、
$S$ を自由 $R$-加群の直和因子として
$\bigoplus_{\alpha \in A} R = S \oplus C$ と書ける。
任意の元 $h \in S$ を上の分解で $h = \sum a_\alpha$ と書いたとき、
$a_\alpha$ を $h$ の係数と呼ぶ。$I' \subset R$ を $J$ の元の係数からなる
イデアルとする。$S$ の元による乗法は $R$-線形写像 $S \to S$ を定めるので、
$I'$ は $J$ の生成元の係数で生成される。すなわち $I'$ は有限生成イデアルである。
$I = fI'$ が条件を満たすと主張する。

\medskip\noindent
まず $V(f) = V(I)$ を確かめる。包含 $V(f) \subset V(I)$ は明らかである。
逆に $f \not \in \mathfrak p$ ならば、性質 (5) と More on Algebra, Lemma
\ref{more-algebra-lemma-fitting-ideal-generate-locally} により、
$\mathfrak q =  \mathfrak p S$ は $V(J)$ の元でない。
したがって $S \otimes_R \kappa(\mathfrak p)$ で零に写らない $J$ の元がある。
ゆえに $\mathfrak p$ に含まれない $I'$ の元が存在し、
$\mathfrak p \not \in V(fI') = V(I)$ である。

\medskip\noindent
$a \in I$ とし、$R' = R[\frac{I}{a}]$ とおく。ある $a' \in I'$ により
$a = fa'$ と書ける。Algebra, Lemmas \ref{algebra-lemma-affine-blowup} および
\ref{algebra-lemma-blowup-add-principal} により、$I' R' = a'R'$ であり、
$a'$ は $R'$ の非零因子である。
% P05-EN-CH38-0252: frozen source defines S' as S tensor_S R' although the scalar extension is from R to R'.
$S' = S \otimes_S R'$ とおく。$JS' = \text{Fit}_r(M \otimes_S S')$ の各元 $g$ は、
ある $c_\alpha \in I'R'$ により $g = \sum_\alpha c_\alpha$ と書ける。
$I'R' = a'R'$ だから、ある $c'_\alpha \in R'$ により
$c_\alpha = a'c'_\alpha$ と書け、$S'$ で
$g = (\sum c'_\alpha)a' = g' a'$ である。さらに、ある $\alpha$ に対して
$a' = c_\alpha$ となる $g_0 \in J$ がある。この元に対して $S'$ で
$g_0 = g'_0 a'$ であり、$g'_0$ は $S'$ の単元である。
$\mathfrak p' \subset R'$ を素イデアルとし、$\mathfrak q' = \mathfrak p'S'$ とおく。
以上により $JS'_{\mathfrak q'}$ は非零因子 $a'$ で生成される単項イデアルである。
More on Algebra, Lemma \ref{more-algebra-lemma-principal-fitting-ideal} により、
$M'_{\mathfrak q'}$ は $r$ 個の元で生成できる。$M'$ は有限なので、
$m_1, \ldots, m_r \in M'$ と $g \in S'$、$g \not \in \mathfrak q'$ で、
対応する写像 $(S')^{\oplus r} \to M'$ が $g$ を可逆化すると全射になるものが存在する。

\medskip\noindent
最後にイデアル $J' = \text{Fit}_{r - 1}(M')$ を考える。
$J'S'_g$ は $m_1, \ldots, m_r$ の間の関係式の係数で生成されることに注意する
（Fitting イデアルと基底変換の両立性）。したがって $J' = 0$ を示せば十分である。
More on Algebra, Lemma \ref{more-algebra-lemma-fitting-ideal-finite-locally-free} を参照せよ。
$R'_a = R_f$（Algebra, Lemma \ref{algebra-lemma-blowup-in-principal}）かつ
$M'_a = M_f$ なので、(5) により、$\mathfrak p'' \subset R'_a$ に対して
$\mathfrak q'' = \mathfrak p''S'$ の形の任意の素イデアル $\mathfrak q'' \subset S'$ について、
% P05-EN-CH38-0260: frozen source localizes S at q'' even though q'' is a prime of S'; the surrounding localizations are of S'.
$J'_a$ は $S_{\mathfrak q''}$ で零に写る。さらに
$S'_a \subset \prod_{\mathfrak q''\text{ as above}} S'_{\mathfrak q''}$
である（Lemma \ref{flat-lemma-base-change-universally-flat} により
$(S'_a)_{\mathfrak p''} \subset S'_{\mathfrak q''}$ だからである）。したがって
% P05-EN-CH38-0253: frozen source writes J'R'_a although J' is the Fitting ideal in S'.
$J'R'_a = 0$ である。$a$ は $R'$ の非零因子なので $J' = 0$ と結論でき、証明が終わる。
\end{proof}

\begin{lemma}
\label{flat-lemma-flatten-module-pre}
$S$ を準コンパクトかつ準分離的なスキーム、$X \to S$ をスキームの射、
$\mathcal{F}$ を $X$ 上の準連接加群、$U \subset S$ を準コンパクト開集合とする。
次を仮定する。
\begin{enumerate}
\item $X \to S$ はアフィン、有限表示、平坦で、ファイバーは幾何学的整である。
\item $\mathcal{F}$ は有限型加群である。
\item $\mathcal{F}_U$ は有限表示である。
\item $U$ の点上にあるファイバーのすべての生成点で、$\mathcal{F}$ は $S$ 上平坦である。
\end{enumerate}
このとき $U$-許容吹き上げ $S' \to S$ と開部分スキーム $V \subset X_{S'}$ で、
(a) $\mathcal{F}$ の厳密変換 $\mathcal{F}'$ の制限が有限局所自由
$\mathcal{O}_V$-加群であり、(b) $V \to S'$ が全射となるものが存在する。
\end{lemma}

\begin{proof}
補題の仮定を満たす $\mathcal{F}/X/S$ と $U \subset S$ が与えられたとする。
「$\mathcal{F}$ はファイバーのすべての生成点で $S$ 上平坦である」という性質を $P$ と表す。
$P$ が厳密変換の下で保たれることは明らかである。Divisors, Lemma
\ref{divisors-lemma-strict-transform-flat} および Morphisms, Lemma
\ref{morphisms-lemma-base-change-module-flat} を参照せよ。
$P$ が $S$ 上局所的であることも明らかである。したがって Remark
\ref{flat-remark-successive-blowups} のすべての考察を、この補題が提起する問題に適用できる。

\medskip\noindent
各 $u \in U$ に整数
$$
r(u) = \dim_{\kappa(\xi_u)}(\mathcal{F}_{\xi_u} \otimes \kappa(\xi_u))
$$
を対応させる関数 $r : U \to \mathbf{Z}_{\geq 0}$ を考える。ここで $\xi_u$ は
ファイバー $X_u$ の生成点である。More on Morphisms, Lemma
\ref{more-morphisms-lemma-flat-and-free-at-point-fibre} と、
$X_S$ の開集合の $S$ における像が開であることから、$r(u)$ は局所定数である。
したがって $U = U_0 \amalg U_1 \amalg \ldots \amalg U_c$ は開かつ閉な
部分スキームの有限非交和で、$U_i$ 上では $r$ は定数値 $i$ をとる。
Divisors, Lemma \ref{divisors-lemma-separate-disjoint-opens-by-blowing-up} により、
$S$ を二つのスキームの非交和へ分解する $U$-許容吹き上げで、一方が $U_0$ を含み、
他方が $U_1 \cup \ldots \cup U_c$ を含むものを見いだせる。これをさらに $c - 1$ 回
繰り返すと、$S$ は非交和 $S = S_0 \amalg S_1 \amalg \ldots \amalg S_c$ であり、
$U_i \subset S_i$ であると仮定してよい。
したがって上で定めた関数 $r$ は定数、例えば値 $r$ をとると仮定してよい。

\medskip\noindent
Remark \ref{flat-remark-successive-blowups} により、台が $S \setminus U$ である
有効 Cartier 因子 $D \subset S$ があると仮定してよい。
Remark \ref{flat-remark-successive-blowups} をもう一度適用し、Divisors, Lemma
\ref{divisors-lemma-characterize-effective-Cartier-divisor} と組み合わせると、
ある非零因子 $f \in R$ に対して
$S = \Spec(R)$ および $D = \Spec(R/(f))$ と仮定してよい。
この場合は Lemma \ref{flat-lemma-helper-blowup-affine-space} で処理される。
\end{proof}

\begin{lemma}
\label{flat-lemma-trick-fitting-ideal}
$A \to C$ を階数 $d$ の有限局所自由な環準同型とする。
$h \in C$ を、$C_h$ が $A$ 上 \'etale となる元とする。$J \subset C$ をイデアルとする。
$C/J$ を有限 $A$-加群とみなして $I = \text{Fit}_0(C/J)$ とおく。
このとき、あるイデアル $J' \subset C_h$ に対して $IC_h = JJ'$ である。
$J$ が有限生成ならば $I$ と $J'$ も有限生成である。
\end{lemma}

\begin{proof}
Fitting イデアルの基本性質を用いる。More on Algebra, Lemma
\ref{more-algebra-lemma-fitting-ideal-basics} を参照せよ。このとき $IC$ は
$C/J \otimes_A C$ の Fitting イデアルである。
$C \to C \otimes_A C$、$c \mapsto 1 \otimes c$ は切断（乗法写像）をもつことに注意する。
仮定により $C \to C \otimes_A C$ は、この切断による $\Spec(C_h)$ の像にある
すべての素イデアルで \'etale である。したがって乗法写像
$C \otimes_A C_h \to C_h$ は \'etale、特に平坦である。
Algebra, Lemma \ref{algebra-lemma-map-between-etale} を参照せよ。
% P05-EN-CH38-0258: frozen source omits the name C' from "there exists a C_h-algebra such that" although the displayed decomposition introduces C'.
ゆえにある $C_h$-代数で、$C_h$-代数として
$C \otimes_A C_h \cong C_h \oplus C'$ となるものが存在する。
Algebra, Lemma \ref{algebra-lemma-surjective-flat-finitely-presented} を参照せよ。
したがって、あるイデアル $I' \subset C'$ に対して $C_h$-加群として
$(C/J) \otimes_A C_h \cong (C_h/J_h) \oplus C'/I'$ である。
ゆえに $J' = \text{Fit}_0(C'/I')$ とおけば $IC_h = JJ'$ である。
% P05-EN-CH38-0259: frozen source says C'/J' is viewed as a C_h-module, but the Fitting ideal immediately uses the module C'/I'.
ここで $C'/J'$ を $C_h$-加群とみなしている。
\end{proof}

\begin{lemma}
\label{flat-lemma-push-ideal}
$A \to B$ を \'etale な環準同型とする。$a \in A$ を非零因子とし、
$J \subset B$ を $V(J) \subset V(aB)$ を満たす有限型イデアルとする。
各 $\mathfrak q \subset B$ に対して、$V(I) \subset V(a)$ を満たす有限型イデアル
$I \subset A$ と、$g \not \in \mathfrak q$ を満たす $g \in B$ で、ある有限型イデアル
$J' \subset B_g$ に対して $IB_g = JJ'$ となるものが存在する。
\end{lemma}

\begin{proof}
$g \in B$、$g \not \in \mathfrak q$ における単項局所化で $B$ を置き換えてよい。
したがって $B$ は標準 \'etale と仮定してよい。Algebra, Proposition
\ref{algebra-proposition-etale-locally-standard} を参照せよ。
ゆえに、ある次数 $d$ のモニック多項式 $f \in A[x]$ に対する
$C = A[x]/(f)$ の局所化が $B$ であると仮定してよい。
ある $h \in C$ に対して $B = C_h$ とする。$B$ 上で $J$ を生成する元
$h_1, \ldots, h_n \in C$ を選ぶ。条件 $V(J) \subset V(aB)$ は、十分大きい $m$ に対して
$B$ で $a^m = \sum b_i h_i$ となることを意味する。$h_{n + 1} = a^m$ とおく。
Lemma \ref{flat-lemma-trick-fitting-ideal} と同様に
% P05-EN-CH38-0255: frozen source uses h_{r+1} although the generators are indexed through n and h_{n+1} was just defined.
$I = \text{Fit}_0(C/(h_1, \ldots, h_{r + 1}))$ ととる。
加群 $C/(h_1, \ldots, h_{r + 1})$ は $a^m$ で零化されるので、
$a^{dm} \in I$ であり、これは $V(I) \subset V(a)$ を含意する。
\end{proof}

\begin{lemma}
\label{flat-lemma-flatten-module-etale-localize}
$S$ を準コンパクトかつ準分離的なスキーム、$X \to S$ をスキームの射、
$\mathcal{F}$ を $X$ 上の準連接加群、$U \subset S$ を準コンパクト開集合とする。
有限個の可換図式
$$
\xymatrix{
& X_i \ar[r]_{j_i} \ar[d] & X \ar[d] \\
S_i^* \ar[r] & S_i \ar[r]^{e_i} & S
}
$$
が存在し、次を満たすと仮定する。
\begin{enumerate}
\item $e_i : S_i \to S$ は準コンパクト \'etale 射で、$S = \bigcup e_i(S_i)$ である。
\item $j_i : X_i \to X$ は \'etale 射で、$X = \bigcup j_i(X_i)$ である。
\item $S^*_i \to S_i$ は $e_i^{-1}(U)$-許容吹き上げで、
$j_i^*\mathcal{F}$ の厳密変換 $\mathcal{F}_i^*$ は $S^*_i$ 上平坦である。
\end{enumerate}
このとき、$\mathcal{F}$ の厳密変換が $S'$ 上平坦となる $U$-許容吹き上げ
$S' \to S$ が存在する。
\end{lemma}

\begin{proof}
補題の仮定は $U$-許容吹き上げの下で保たれると主張する。実際、
$b : S' \to S$ を準連接イデアル層 $\mathcal{I}$ における $U$-許容吹き上げとする。
さらに $S^*_i \to S_i$ を準連接イデアル層 $\mathcal{J}_i$ における吹き上げとする。
このとき射の族 $e'_i : S'_i = S_i \times_S S' \to S'$ および
$j'_i : X_i' = X_i \times_S S' \to X \times_S S'$ は、吹き上げ $S' \to S$ に関する
$\mathcal{F}$ の厳密変換 $\mathcal{F}'$ に対して条件 (1), (2), (3) を満たす。
まず $S_i'$ は、$\mathcal{I}$ の引き戻しにおける $S_i$ の吹き上げである。
Divisors, Lemma \ref{divisors-lemma-flat-base-change-blowing-up} を参照せよ。
次に $\mathcal{J}_i$ の引き戻しにおける $S_i'$ の吹き上げ
$S_i^{\prime *} \to S_i'$ を考える。Divisors, Lemma
\ref{divisors-lemma-blowing-up-two-ideals} により可換図式
$$
\xymatrix{
S_i^{\prime *} \ar[r] \ar[rd] \ar[d] & S'_i \ar[d] \\
S_i^* \ar[r] & S_i
}
$$
を得て、この図式のすべての射は吹き上げである。したがって Divisors, Lemmas
\ref{divisors-lemma-strict-transform-flat} および
\ref{divisors-lemma-strict-transform-composition-blowups} により
\begin{align*}
& \text{ }(j'_i)^*\mathcal{F}'\text{ の }S_i^{\prime *} \to S_i'
\text{ による厳密変換} \\
= &
\text{ }j_i^*\mathcal{F}\text{ の }S_i^{\prime *} \to S_i
\text{ による厳密変換} \\
= &
\text{ }\mathcal{F}_i'\text{ の }S_i^{\prime *} \to S_i'
\text{ による厳密変換} \\
= &
\text{ }\mathcal{F}_i^*\text{ の }X_i \times_{S_i} S_i^{\prime *} \to X_i
\text{ による引き戻し}
\end{align*}
となる。したがってこれは $S_i^{\prime *}$ 上平坦である
（Morphisms, Lemma \ref{morphisms-lemma-base-change-module-flat}）。
以上により、この補題のような仮定の下で、$\mathcal{F}$ の厳密変換が基底上平坦となる
$U$-許容吹き上げを見いだす問題には Remark \ref{flat-remark-successive-blowups} の
すべての考察が適用される。特に $S \setminus U$ は有効 Cartier 因子
$D \subset S$ の台であると仮定してよい。Remark \ref{flat-remark-successive-blowups} を
もう一度適用し、Divisors, Lemma
\ref{divisors-lemma-characterize-effective-Cartier-divisor} と組み合わせると、
ある非零因子 $a \in A$ に対して $S = \Spec(A)$ および
$D = \Spec(A/(a))$ と仮定してよい。

\medskip\noindent
$i$ と $s \in S_i$ を選ぶ。Lemma \ref{flat-lemma-push-ideal} により、その開近傍
$s \in W_i \subset S_i$ と有限型準連接イデアル
$\mathcal{I} \subset \mathcal{O}_S$ で、ある有限型準連接イデアル
$\mathcal{J}'_i \subset \mathcal{O}_{W_i}$ に対して
$\mathcal{I} \cdot \mathcal{O}_{W_i} = \mathcal{J}_i \mathcal{J}'_i$ となり、かつ
$V(\mathcal{I}) \subset V(a) = S \setminus U$ となるものを見いだせる。
$S_i$ は準コンパクトなので、$S_i$ をこれらの開集合の有限族
$W_1, \ldots, W_n$ で置き換え、各 $i$ に対して準連接イデアル層
$\mathcal{I}_i \subset \mathcal{O}_S$ で、ある有限型準連接イデアル
$\mathcal{J}'_i \subset \mathcal{O}_{S_i}$ に対して
$\mathcal{I}_i \cdot \mathcal{O}_{S_i} = \mathcal{J}_i \mathcal{J}'_i$ となるものが
存在すると仮定してよい。証明の第一段落の議論と同様に、積
$\mathcal{I}_1 \ldots \mathcal{I}_n$ における $S$ の吹き上げ $S'$ を考える
（この吹き上げは構成により $U$-許容である）。$S' \to S$ の $S_i$ への基底変換は
$$
\mathcal{J}_i \cdot
\mathcal{J}'_i \mathcal{I}_1 \ldots \hat{\mathcal{I}_i} \ldots \mathcal{I}_n
$$
における吹き上げであり、与えられた吹き上げ $S_i^* \to S_i$ を経由する
（Divisors, Lemma \ref{divisors-lemma-blowing-up-two-ideals}）。上の図式の記法では、
これは $S_i^{\prime *} = S_i'$ を意味する。したがって $S$ を $S'$ で置き換えた後、
$j_i^*\mathcal{F}$ が $S_i$ 上平坦な状況に到達する。ゆえに
$j_i^*\mathcal{F}$ は $S$ 上平坦である。Lemma \ref{flat-lemma-etale-flat-up-down} を参照せよ。
Morphisms, Lemma \ref{morphisms-lemma-flat-permanence} により、
$\mathcal{F}$ は $S$ 上平坦である。
\end{proof}

\begin{theorem}
\label{flat-theorem-flatten-module}
$S$ を準コンパクトかつ準分離的なスキーム、$X$ を $S$ 上のスキーム、
$\mathcal{F}$ を $X$ 上の準連接加群、$U \subset S$ を準コンパクト開集合とする。
次を仮定する。
\begin{enumerate}
\item $X$ は準コンパクトである。
\item $X$ は $S$ 上局所有限表示である。
\item $\mathcal{F}$ は有限型加群である。
\item $\mathcal{F}_U$ は有限表示である。
\item $\mathcal{F}_U$ は $U$ 上平坦である。
\end{enumerate}
このとき、$\mathcal{F}$ の厳密変換 $\mathcal{F}'$ が有限表示
$\mathcal{O}_{X \times_S S'}$-加群で、かつ $S'$ 上平坦となる
$U$-許容吹き上げ $S' \to S$ が存在する。
\end{theorem}

\begin{proof}
まず、厳密変換が平坦となる $U$-許容吹き上げを見いだせることを証明する。
この問題は始域と終域について \'etale 局所的である。精密な主張については
Lemma \ref{flat-lemma-flatten-module-etale-localize} を参照せよ。
特に $S = \Spec(R)$ と $X = \Spec(A)$ はアフィンであると仮定してよい。
$s \in S$ に対して $\mathcal{F}_s = \mathcal{F}|_{X_s}$ と書く
（$\mathcal{F}$ のファイバーへの引き戻し）。$X \to S$ は有限型なので
$d = \max_{s \in S} \dim(\text{Supp}(\mathcal{F}_s))$ は整数である。
$d$ に関する帰納法を行う。

\medskip\noindent
$x \in X$ を、$X$ の点で $s \in S$ 上にあり
$\dim_x(\text{Supp}(\mathcal{F}_s)) = d$ を満たす点とする。
Lemma \ref{flat-lemma-elementary-devissage} を適用して
$g : X' \to X$、$e : S' \to S$、$i : Z' \to X'$、および
$\pi : Z' \to Y'$ を得る。$Y' \to S'$ はアフィンスキーム間の滑らかな射で、
ファイバーは次元 $d$ の幾何学的既約スキームである。
問題は \'etale 局所的なので、$g^*\mathcal{F}/X'/S'$ に対して定理を証明すれば十分である。
$i : Z' \to X'$ は有限表示閉埋め込みであり、厳密変換はアフィン押し出しと可換なので、
$\mathcal{G}$ に対して平坦化の結果を証明すれば十分である。
Divisors, Lemma \ref{divisors-lemma-strict-transform-affine} を参照せよ。
$\pi$ は有限、したがってアフィンなので、$\pi_*\mathcal{G}/Y'/S'$ に対して
平坦化の結果を証明すれば十分である。ゆえに $X \to S$ はアフィンスキーム間の
滑らかな射で、ファイバーは次元 $d$ の幾何学的既約スキームであると仮定してよい。

\medskip\noindent
次に Lemma \ref{flat-lemma-flatten-module-pre} のような吹き上げを適用する。
すると、$S$ へ全射する開集合 $V \subset X$ で、$\mathcal{F}|_V$ が有限局所自由となるものが
存在する状況に到達する。$\xi \in X$ を $X_s$ の生成点とし、
$r = \dim_{\kappa(\xi)} \mathcal{F}_\xi \otimes \kappa(\xi)$ とおく。
同型 $\kappa(\xi)^{\oplus r} \to \mathcal{F}_\xi \otimes \kappa(\xi)$ を誘導する写像
$\alpha : \mathcal{O}_X^{\oplus r} \to \mathcal{F}$ を選ぶ。
$\mathcal{F}$ は $V$ 上局所自由なので、$\alpha$ が同型となる $\xi$ の開近傍 $W$ がある。
$W \to S$ が全射となるような $s$ のアフィン開近傍へ $S$ を縮小する。
$\mathcal{F}$ は $A$-加群 $N$ に付随する準連接加群であるとする。
$\mathcal{F}$ はファイバーのすべての生成点（実際には $W$ のすべての点）で
$S$ 上平坦なので、
$$
\alpha_\mathfrak p : A_\mathfrak p^{\oplus r} \to N_\mathfrak p
$$
は $R$ のすべての素イデアル $\mathfrak p$ に対して普遍単射である。
Lemma \ref{flat-lemma-induction-step} を参照せよ。したがって $\alpha$ は普遍単射である。
Algebra, Lemma \ref{algebra-lemma-universally-injective-check-stalks} を参照せよ。
$\mathcal{H} = \Coker(\alpha)$ とおく。Divisors, Lemma
\ref{divisors-lemma-strict-transform-universally-injective} により、
$U$-許容吹き上げ $S' \to S$ が与えられたとき、厳密変換
$\mathcal{F}'$ と $\mathcal{H}'$ は完全列
$$
0 \to \mathcal{O}_{X \times_S S'}^{\oplus r} \to \mathcal{F}'
\to \mathcal{H}' \to 0
$$
をなす。したがって Lemma \ref{flat-lemma-induction-step} は、$\mathcal{F}'$ が点 $x'$ で
平坦であることと $\mathcal{H}'$ がその点で平坦であることが同値であることも示す。
特に $\mathcal{H}_U$ は $U$ 上平坦で、$\mathcal{H}_U$ は有限表示加群である。
$\mathcal{H}$ に帰納法の仮定を適用すると、厳密変換 $\mathcal{H}'$ が望みどおり平坦となる
$U$-許容吹き上げが存在することが分かる。

\medskip\noindent
定理の証明を終えるには、必要ならさらに別の $U$-許容吹き上げを行った後、
$\mathcal{F}'$ が有限表示加群であることを示す必要がある。
$U \subset S$ はスキーム論的に稠密であると仮定してよいので、これは Lemma
\ref{flat-lemma-flat-finite-type-finitely-presented-over-dense-open} から従う
（Remark \ref{flat-remark-successive-blowups} の第三段落を参照）。これで定理の証明が終わる。
\end{proof}





\section{応用}
\label{flat-section-applications}

\noindent
この節では上の結果のいくつかを応用する。

\begin{lemma}
\label{flat-lemma-flat-after-blowing-up}
$S$ を準コンパクトかつ準分離的なスキーム、$X$ を $S$ 上のスキーム、
$U \subset S$ を準コンパクト開集合とする。次を仮定する。
\begin{enumerate}
\item $X \to S$ は有限型かつ準分離的である。
\item $X_U \to U$ は平坦かつ局所有限表示である。
\end{enumerate}
このとき、$X$ の厳密変換が $S'$ 上平坦かつ有限表示となる
$U$-許容吹き上げ $S' \to S$ が存在する。
\end{lemma}

\begin{proof}
仮定により $X \to S$ は準コンパクトかつ準分離的なので、吹き上げ
$S' \to S$ に関する $X$ の厳密変換も準コンパクトかつ準分離的である。
したがって補題を証明するには、厳密変換が平坦かつ局所有限表示となる
$U$-許容吹き上げを見いだせば十分である。
$X = W_1 \cup \ldots \cup W_n$ を有限アフィン開被覆とする。
$W_i$ の厳密変換が平坦かつ局所有限表示となる $U$-許容吹き上げ
$S_i \to S$ を見いだせるならば、すべての $S_i \to S$ を支配して条件を満たす
$U$-許容吹き上げ $S' \to S$ が存在する（Divisors, Lemma
\ref{divisors-lemma-dominate-admissible-blowups} および Remark
\ref{flat-remark-successive-blowups} を参照）。したがって $X$ はアフィンと仮定してよい。

\medskip\noindent
$X$ がアフィンであると仮定する。Morphisms, Lemma
\ref{morphisms-lemma-quasi-affine-finite-type-over-S} により、$S$ 上の埋め込み
$j : X \to \mathbf{A}^n_S$ を選べる。$V \subset \mathbf{A}^n_S$ を準コンパクト開部分スキームで、
$j$ が $S$ 上の閉埋め込み $i : X \to V$ を誘導するものとする。
Theorem \ref{flat-theorem-flatten-module} を $V \to S$ と準連接加群
$i_*\mathcal{O}_X$ に適用すると、$i_*\mathcal{O}_X$ の厳密変換が $S'$ 上平坦かつ
$\mathcal{O}_{V \times_S S'}$-加群として有限表示となる $U$-許容吹き上げ
$S' \to S$ を得る。$X'$ を $S' \to S$ に関する $X$ の厳密変換とし、
$i' : X' \to V \times_S S'$ を誘導される射とする。
厳密変換をとる操作はアフィン射に沿う押し出しと可換なので
（Divisors, Lemma \ref{divisors-lemma-strict-transform-affine}）、
% P05-EN-CH38-0261: frozen source says the post-base-change pushforward is flat over S rather than S'.
$i'_*\mathcal{O}_{X'}$ は $S$ 上平坦で、$\mathcal{O}_{V \times_S S'}$-加群として
有限表示である。これは補題を含意する。
\end{proof}

\begin{lemma}
\label{flat-lemma-finite-after-blowing-up}
$S$ を準コンパクトかつ準分離的なスキーム、$X$ を $S$ 上のスキーム、
$U \subset S$ を準コンパクト開集合とする。次を仮定する。
\begin{enumerate}
\item $X \to S$ は固有である。
\item $X_U \to U$ は有限局所自由である。
\end{enumerate}
このとき、$X$ の厳密変換が $S'$ 上有限局所自由となる
$U$-許容吹き上げ $S' \to S$ が存在する。
\end{lemma}

\begin{proof}
Lemma \ref{flat-lemma-flat-after-blowing-up} により、$X \to S$ は平坦かつ有限表示と仮定してよい。
必要なら $S$ を $U$-許容吹き上げで置き換えて、$U \subset S$ はスキーム論的に
稠密であると仮定してよい。このとき
% P05-EN-CH38-0262: frozen source refers to f although the morphism X -> S is not named f in the lemma.
$f$ は Lemma \ref{flat-lemma-proper-flat-finite-over-dense-open} により有限である。
したがって Morphisms, Lemma \ref{morphisms-lemma-finite-flat} により
$f$ は有限局所自由である。
\end{proof}

\begin{lemma}
\label{flat-lemma-zariski-after-blowup}
$\varphi : X \to S$ を有限型分離射とし、$S$ は準コンパクトかつ準分離的であるとする。
$U \subset S$ を、$\varphi^{-1}U \to U$ が同型となる準コンパクト開集合とする。
このとき、$X$ の厳密変換 $X'$ が $S'$ の開部分スキームと同型になる
$U$-許容吹き上げ $S' \to S$ が存在する。
\end{lemma}

\begin{proof}
Remark \ref{flat-remark-successive-blowups} の議論を適用できる。
したがって最初の $U$-許容吹き上げを行い、補集合 $S \setminus U$ が
有効 Cartier 因子 $D$ の台であると仮定してよい。特に $U$ は $S$ で
スキーム論的に稠密である。次にもう一つの $U$-許容吹き上げを行い、
$X \to S$ が平坦かつ有限表示である状況に到達する。
Lemma \ref{flat-lemma-flat-after-blowing-up} を参照せよ。
この場合、結果は Lemma \ref{flat-lemma-zariski} から従う。
\end{proof}

\noindent
次の補題は、固有修正を吹き上げで支配できることを述べる。

\begin{lemma}
\label{flat-lemma-dominate-modification-by-blowup}
$\varphi : X \to S$ を固有射とし、$S$ は準コンパクトかつ準分離的であるとする。
$U \subset S$ を、$\varphi^{-1}U \to U$ が同型となる準コンパクト開集合とする。
このとき $X$ を支配する $U$-許容吹き上げ $S' \to S$、すなわち吹き上げ射の分解
$S' \to X \to S$ が存在するようなものが存在する。
\end{lemma}

\begin{proof}
Remark \ref{flat-remark-successive-blowups} の議論を適用できる。
したがって最初の $U$-許容吹き上げを行い、補集合 $S \setminus U$ が
有効 Cartier 因子 $D$ の台であると仮定してよい。特に $U$ は $S$ で
スキーム論的に稠密である。$X$ の厳密変換 $X'$ が $S'$ の開部分スキームとなる
別の $U$-許容吹き上げ $S' \to S$ を選ぶ。Lemma
\ref{flat-lemma-zariski-after-blowup} を参照せよ。
$X' \to S'$ は固有で $U \subset S'$ は稠密なので、$X' = S'$ である。
細部は省略する。
\end{proof}

\begin{lemma}
\label{flat-lemma-get-section-after-blowup}
$S$ をスキームとし、$U \subset W \subset S$ を開部分スキームとする。
$f : X \to W$ を射、$s : U \to X$ を $f \circ s = \text{id}_U$ を満たす射とする。
次を仮定する。
\begin{enumerate}
\item $f$ は固有である。
\item $S$ は準コンパクトかつ準分離的である。
\item $U$ と $W$ は準コンパクトである。
\end{enumerate}
このとき、$U$-許容吹き上げ $b : S' \to S$ と、$s$ を延長し
$f \circ s' = b|_{b^{-1}(W)}$ を満たす射 $s' : b^{-1}(W) \to X$ が存在する。
\end{lemma}

\begin{proof}
$X$ を $s$ のスキーム論的像で置き換えてよく、実際そうする。
すると $X \to W$ は $U$ 上同型である。Morphisms, Lemma
\ref{morphisms-lemma-scheme-theoretic-image-of-partial-section} を参照せよ。
Lemma \ref{flat-lemma-dominate-modification-by-blowup} により、$U$-許容吹き上げ
$W' \to W$ と $s$ の延長 $W' \to X$ が存在する。
最後に Divisors, Lemma \ref{divisors-lemma-extend-admissible-blowups} を適用して
$W' \to W$ を $S$ の $U$-許容吹き上げへ延長すれば証明が終わる。
\end{proof}









\section{コンパクト化}
\label{flat-section-compactify}

\noindent
$S$ を準コンパクトかつ準分離的なスキームとする。$S$ 上のスキーム $X$ に対し、
$S$ 上固有なスキーム $\overline{X}$ への準コンパクト開埋め込み
$X \to \overline{X}$ が存在するとき、
{\it $S$ 上のコンパクト化をもつ}、または
{\it $S$ 上コンパクト化可能である}という。$X$ が $S$ 上の
コンパクト化をもてば、$X \to S$ は分離的かつ有限型である。逆も成り立つことは
Nagata の定理である。\cite{Lutkebohmert}, \cite{Conrad-Nagata},
\cite{Nagata-1}, \cite{Nagata-2}, \cite{Nagata-3}, \cite{Nagata-4} を参照せよ。
この定理は次節で証明する。Theorem \ref{flat-theorem-nagata} を参照せよ。

\medskip\noindent
$S$ を準コンパクトかつ準分離的なスキームとし、
$X \to S$ をスキームの分離的有限型射とする。
{\it $S$ 上の $X$ のコンパクト化の圏}を次のように定義する。
\begin{enumerate}
\item 対象は、$\overline{X} \to S$ が固有であるような $S$ 上の開埋め込み
$j : X \to \overline{X}$ である。
\item 射 $(j' : X \to \overline{X}') \to (j : X \to \overline{X})$ は、
$f \circ j' = j$ を満たす $S$ 上のスキームの射
$f : \overline{X}' \to \overline{X}$ である。
\end{enumerate}
$j : X \to \overline{X}$ がコンパクト化ならば、$j$ は準コンパクト開埋め込みである。
Schemes, Remark \ref{schemes-remark-quasi-compact-and-quasi-separated} を参照せよ。

\medskip\noindent
{\bf 注意。} コンパクト化 $j : X \to \overline{X}$ の像が稠密であるとは
{\it 仮定しない}。したがって、$f : \overline{X}' \to \overline{X}$ が
コンパクト化の射であっても、$f^{-1}(j(X)) = j'(X)$ とは限らない。

\begin{lemma}
\label{flat-lemma-compactifications-cofiltered}
$S$ を準コンパクトかつ準分離的なスキームとし、
$X$ を $S$ 上コンパクト化可能なスキームとする。
\begin{enumerate}
\item[(a)] $S$ 上の $X$ のコンパクト化の圏は余フィルター的である。
\item[(b)] $j(X)$ が $\overline{X}$ で稠密かつスキーム論的に稠密であるような
コンパクト化 $j : X \to \overline{X}$ からなる充満部分圏は始的である
(Categories, Definition \ref{categories-definition-initial})。
\item[(c)] $f : \overline{X}' \to \overline{X}$ が $X$ のコンパクト化の射で、
$j'(X)$ が $\overline{X}'$ で稠密ならば、
$f^{-1}(j(X)) = j'(X)$ である。
\end{enumerate}
\end{lemma}

\begin{proof}
(a) を証明するには、Categories, Definition
\ref{categories-definition-codirected} の条件 (1), (2), (3) を確認すればよい。
$X$ がコンパクト化可能だと仮定したことにより、条件 (1) はちょうど成り立つ。
$j_i : X \to \overline{X}_i$, $i = 1, 2$ を二つのコンパクト化とする。
このとき $(j_1, j_2) : X \to \overline{X}_1 \times_S \overline{X}_2$ の
スキーム論的像 $\overline{X}$ を考えられる。これにより、両方の $j_i$ を支配する
第三のコンパクト化 $j : X \to \overline{X}$ が定まる。
$$
\xymatrix{
(X, \overline{X}_1) & (X, \overline{X}) \ar[l] \ar[r] & (X, \overline{X}_2)
}
$$
したがって (2) が成り立つ。$j_i : X \to \overline{X}_i$, $i = 1, 2$ の間の
二つの射を $f_1, f_2 : \overline{X}_1 \to \overline{X}_2$ とする。
$\overline{X} \subset \overline{X}_1$ を $f_1$ と $f_2$ の等化子とする。
$\overline{X}_2 \to S$ は分離的なので、$\overline{X}$ は
$\overline{X}_1$ の閉部分スキームであり、したがって $S$ 上固有である。
さらに $f_1|_X = f_2|_X = \text{id}_X$ だから、開埋め込み
$X \to \overline{X}$ を得る。閉埋め込み
$\overline{X} \to \overline{X}_1$ が与える射
$(X \to \overline{X}) \to (j_1 : X \to \overline{X}_1)$ は
$f_1$ と $f_2$ を等化する。これで条件 (3) が証明された。

\medskip\noindent
(b) の証明。$j : X \to \overline{X}$ をコンパクト化とする。
$\overline{X}'$ を $\overline{X}$ における $X$ のスキーム論的閉包とすれば、
Morphisms, Lemma \ref{morphisms-lemma-quasi-compact-immersion} により、
$X$ は $\overline{X}'$ で稠密かつスキーム論的に稠密である。
これで Categories, Definition \ref{categories-definition-initial} の
第一の条件が証明された。$X$ のコンパクト化の圏が余フィルター的であることは
すでに示したので、Categories, Definition
\ref{categories-definition-initial} の第二の条件は第一の条件から従う
（この圏論的演習の解答は省略する）。

\medskip\noindent
(c) の証明。$\overline{X}'$ を $j'(X)$ のスキーム論的閉包で置き換えても
台となる位相空間は変わらない。この置換の後、主張は Morphisms, Lemma
\ref{morphisms-lemma-scheme-theoretic-image-of-partial-section} から従う。
\end{proof}

\noindent
$X$ を動かしてすべてのコンパクト化からなる圏を考えることもできる。
この圏を、内部上に同型を誘導する射の集合で局所化すると、
$S$ 上コンパクト化可能なスキームの圏と同値になる。

\begin{lemma}
\label{flat-lemma-compactifyable}
$S$ を準コンパクトかつ準分離的なスキームとする。
$f : X \to Y$ を $S$ 上のスキームの射とし、$Y$ は $S$ 上分離的かつ有限型、
$X$ は $S$ 上コンパクト化可能であると仮定する。このとき $X$ は $Y$ 上の
コンパクト化をもつ。
\end{lemma}

\begin{proof}
$j : X \to \overline{X}$ を $S$ 上の $X$ のコンパクト化とし、
$\overline{X}'$ を $(j, f) : X \to \overline{X} \times_S Y$ の
スキーム論的像とする。$\overline{X} \times_S Y \to Y$ は
$\overline{X} \to S$ の基底変換として固有なので、射
$\overline{X}' \to Y$ は固有である。一方、$Y$ は $S$ 上分離的なので、射
$(1, f) : X \to X \times_S Y$ は閉埋め込みである
(Schemes, Lemma \ref{schemes-lemma-semi-diagonal})。
したがって、射影 $\overline{X} \times_S Y \to \overline{X}$ の
「部分切断」$s = (j, f)$ に Morphisms, Lemma
\ref{morphisms-lemma-scheme-theoretic-image-of-partial-section} を適用すると、
$X \to \overline{X}'$ は開埋め込みである。
\end{proof}

\noindent
$S$ を準コンパクトかつ準分離的なスキームとする。
{\it コンパクト化の圏}を次の圏として定義する。その対象は、
$\overline{X}$ が $S$ 上固有なスキームで、
$X \subset \overline{X}$ が準コンパクト開部分スキームであるような対
$(X, \overline{X})$ であり、その射は $S$ 上のスキームの射の可換図式
$$
\xymatrix{
X \ar[d] \ar[r]_f & Y \ar[d] \\
\overline{X} \ar[r]^{\overline{f}} & \overline{Y}
}
$$
である。

\begin{lemma}
\label{flat-lemma-right-multiplicative-system}
$S$ を準コンパクトかつ準分離的なスキームとする。
$u$ が同型であるような射
$(u, \overline{u}) : (X', \overline{X}') \to (X, \overline{X})$
の全体は、コンパクト化の圏における射の右乗法系をなす
(Categories, Definition \ref{categories-definition-multiplicative-system})。
\end{lemma}

\begin{proof}
公理 RMS1 の確認は自明である。RMS2 を確認しよう。図式
$$
\xymatrix{
& (X', \overline{X}') \ar[d]_{(u, \overline{u})} \\
(Y, \overline{Y}) \ar[r]^{(f, \overline{f})} & (X, \overline{X})
}
$$
が与えられ、$u : X' \to X$ が同型であるとする。
$Y' = Y \times_X X'$ とおき、射影を $v : Y' \to Y$ とする（これは同型である）。
また $\overline{Y}' = \overline{Y} \times_{\overline{X}} \overline{X}'$ とおき、
射影を $\overline{v} : \overline{Y}' \to \overline{Y}$ とする。
$Y' \to \overline{Y}'$ が開埋め込みであることは明らかである。図式
$$
\xymatrix{
(Y', \overline{Y}') \ar[r]_{(g, \overline{g})} \ar[d]_{(v, \overline{v})} &
(X', \overline{X}') \ar[d]_{(u, \overline{u})} \\
(Y, \overline{Y}) \ar[r]^{(f, \overline{f})} & (X, \overline{X})
}
$$
は公理 RMS2 が成り立つことを示す。

\medskip\noindent
RMS3 を確認しよう。コンパクト化の二つの射
$(f, \overline{f}), (g, \overline{g}) :
(X, \overline{X}) \to (Y, \overline{Y})$
と、射
$(v, \overline{v}) : (Y, \overline{Y}) \to (Y', \overline{Y}')$
が与えられ、$v$ は同型で、かつ
$(v, \overline{v}) \circ (f, \overline{f}) =
(v, \overline{v}) \circ (g, \overline{g})$ を満たすとする。このとき $f = g$ である。
したがって $\overline{X}' \subset \overline{X}$ を $\overline{f}$ と
$\overline{g}$ の等化子とすれば、
$(u, \overline{u}) : (X, \overline{X}') \to (X, \overline{X})$ は
コンパクト化の圏の射であり、所望の通り
$(f, \overline{f}) \circ (u, \overline{u}) =
(g, \overline{g}) \circ (u, \overline{u})$ を満たす。
\end{proof}

\begin{lemma}
\label{flat-lemma-invert-right-multiplicative-system}
$S$ を準コンパクトかつ準分離的なスキームとする。関手
$(X, \overline{X}) \mapsto X$ は、コンパクト化の圏を Lemma
\ref{flat-lemma-right-multiplicative-system} の右乗法系で局所化して得られる圏
(Categories, Lemma \ref{categories-lemma-right-localization}) から、
$S$ 上コンパクト化可能なスキームの圏への同値を定める。
\end{lemma}

\begin{proof}
$\mathcal{C}$ をコンパクト化の圏、
$Q : \mathcal{C} \to \mathcal{C}'$ を Categories, Lemma
\ref{categories-lemma-properties-right-localization} の局所化関手とする。
$\mathcal{D}$ を $S$ 上コンパクト化可能なスキームの圏とする。
今引用した補題と乗法系の選び方から、関手
$\mathcal{C}' \to \mathcal{D}$ が得られることは明らかである。
この関手は明らかに本質的全射である。
$f : X \to Y$ をコンパクト化可能なスキームの射とする。
$S$ 上固有なスキームへの開埋め込み $Y \to \overline{Y}$ を選び、次に
$\overline{Y}$ 上固有なスキーム $\overline{X}$ への埋め込み
$X \to \overline{X}$ を選ぶ（$X \to \overline{Y}$ に Lemma
\ref{flat-lemma-compactifyable} を適用すれば可能である）。これにより、与えられた射
$X \to Y$ を生じるコンパクト化の射
$(X, \overline{X}) \to (Y, \overline{Y})$ を得る。
最後に、局所化圏に同じ始域と終域をもつ二つの射、例えば
$$
a = ((f, \overline{f}) : (X', \overline{X}') \to (Y, \overline{Y}),
(u, \overline{u}) : (X', \overline{X}') \to (X, \overline{X}))
$$
and
$$
b = ((g, \overline{g}) : (X'', \overline{X}'') \to (Y, \overline{Y}),
(v, \overline{v}) : (X'', \overline{X}'') \to (X, \overline{X}))
$$
が、同じ $S$ 上の射 $X \to Y$ を生じると仮定する。すなわち
$f \circ u^{-1} = g \circ v^{-1}$ とする。Categories, Lemma
\ref{categories-lemma-morphisms-right-localization} により、
$(X', \overline{X}') = (X'', \overline{X}'')$ かつ
$(u, \overline{u}) = (v, \overline{v})$ と仮定してよい。この場合、
$\overline{f}$ と $\overline{g}$ の等化子
$\overline{X}''' \subset \overline{X}'$ を考えられる。射
$(w, \overline{w}) : (X', \overline{X}''') \to (X', \overline{X}')$ は
乗法系に属し、$(w, \overline{w})$ と前合成することにより、
局所化圏で $a = b$ であると分かる。
\end{proof}






\section{Nagata コンパクト化}
\label{flat-section-compactifications}

\noindent
本節では Section \ref{flat-section-compactify} で予告した定理を証明する。

\begin{lemma}
\label{flat-lemma-check-separated}
$X \to S$ をスキームの射とする。$X = U \cup V$ が開被覆で、
$U \to S$ と $V \to S$ が分離的、かつ
$U \cap V \to U \times_S V$ が閉ならば、$X \to S$ は分離的である。
\end{lemma}

\begin{proof}
省略する。ヒント：$X \times_S X$ の開被覆
$U \times_S U$, $U \times_S V$, $V \times_S U$, $V \times_S V$ を用いて、
$\Delta : X \to X \times_S X$ が閉であることを確認せよ。
\end{proof}

\begin{lemma}
\label{flat-lemma-separate-disjoint-locally-closed-by-blowing-up}
$X$ を準コンパクトかつ準分離的なスキームとし、
$U \subset X$ を準コンパクト開部分スキームとする。
\begin{enumerate}
\item $Z_1, Z_2 \subset X$ が有限表示閉部分スキームで
$Z_1 \cap Z_2 \cap U = \emptyset$ を満たすならば、$Z_1$ と $Z_2$ の
厳密変換が互いに素となるような $U$-許容吹き上げ $X' \to X$ が存在する。
\item $T_1, T_2 \subset U$ が互いに素な構成可能閉部分集合ならば、
$T_1$ と $T_2$ の閉包が互いに素となるような
$U$-許容吹き上げ $X' \to X$ が存在する。
\end{enumerate}
\end{lemma}

\begin{proof}
(1) の証明。$Z_i \to X$ が有限表示であるという仮定は、$Z_i$ の
準連接イデアル層 $\mathcal{I}_i$ が有限型であることを意味する。
Morphisms, Lemma \ref{morphisms-lemma-closed-immersion-finite-presentation}
を参照せよ。積 $\mathcal{I}_1 \mathcal{I}_2$ で定義される閉部分スキームを
$Z \subset X$ とする。$Z \cap U$ は $Z_1 \cap U$ と $Z_2 \cap U$ の
非交和であることに注意する。Divisors, Lemma
\ref{divisors-lemma-separate-disjoint-opens-by-blowing-up} により、
$Z_1$ と $Z_2$ の厳密変換が互いに素となるような
$U \cap Z$-許容吹き上げ $Z' \to Z$ が存在する。この吹き上げの中心を
$Y \subset Z$ とする。$Y \to X$ は $Y \to Z$ と $Z \to X$ の合成として
有限表示閉埋め込みである (Divisors, Definition
\ref{divisors-definition-admissible-blowup} および Morphisms, Lemma
\ref{morphisms-lemma-composition-finite-presentation})。
したがって $Y$ の吹き上げ $X' \to X$ は $U$-許容吹き上げである。
厳密変換の一般的性質により、$X' \to X$ に関する $Z_1, Z_2$ の厳密変換は、
$Z' \to Z$ に関する $Z_1, Z_2$ の厳密変換と同じである。
Divisors, Lemma \ref{divisors-lemma-strict-transform} を参照せよ。
これで (1) が証明された。

\medskip\noindent
(2) の証明。Properties, Lemma
\ref{properties-lemma-quasi-coherent-finite-type-ideals} により、
$T_i = V(\mathcal{J}_i)$（集合論的に）を満たす有限型準連接イデアル層
$\mathcal{J}_i \subset \mathcal{O}_U$ が存在する。Properties, Lemma
\ref{properties-lemma-extend} により、$U$ への制限が $\mathcal{J}_i$ となる
有限型準連接イデアル層 $\mathcal{I}_i \subset \mathcal{O}_X$ が存在する。
閉部分スキーム $Z_i = V(\mathcal{I}_i)$ に (1) の結果を適用すればよい。
\end{proof}

\begin{lemma}
\label{flat-lemma-blowup-iso-along}
$f : X \to Y$ を準コンパクトかつ準分離的なスキームの固有射とする。
$V \subset Y$ を準コンパクト開部分スキームとし、$U = f^{-1}(V)$ とおく。
$T \subset V$ を閉部分集合とし、$f|_U : U \to V$ は $V$ における
$T$ のある開近傍上で同型であると仮定する。このとき、$f$ の厳密変換
$f' : X' \to Y'$ が $Y'$ における $T$ の閉包のある開近傍上で
同型となるような $V$-許容吹き上げ $Y' \to Y$ が存在する。
\end{lemma}

\begin{proof}
$T' \subset V$ を、$f|_U$ が同型となる最大の開集合の補集合とする。
すると $T', T$ は $V$ で閉であり、$T \cap T' = \emptyset$ である。
$V$ はスペクトル位相空間なので、$T \subset T_c$, $T' \subset T'_c$ かつ
$T_c \cap T'_c = \emptyset$ を満たす構成可能閉部分集合 $T_c, T'_c$ をとれる
（$T'$ を含み $T$ と交わらない $V$ の準コンパクト開集合 $W$ を選んで
$T_c = V \setminus W$ とおき、次に $T_c$ を含み $T'$ と交わらない
$V$ の準コンパクト開集合 $W'$ を選んで $T'_c = V \setminus W'$ とおく）。
Lemma \ref{flat-lemma-separate-disjoint-locally-closed-by-blowing-up} により、
$Y$ を $V$-許容吹き上げで置き換えた後、$T_c$ と $T'_c$ の $Y$ における閉包が
互いに素であると仮定してよい。
$Y_0 = Y \setminus \overline{T}'_c$, $V_0 = V \setminus T'_c$,
$U_0 = U \times_V V_0$, $X_0 = X \times_Y Y_0$ とおく。
$U_0 \to V_0$ は同型なので、$X_0$ の厳密変換 $X'_0$ が $Y'_0$ に
同型に写るような $V_0$-許容吹き上げ $Y'_0 \to Y_0$ をとれる。
Lemma \ref{flat-lemma-zariski-after-blowup} を参照せよ。
Divisors, Lemma \ref{divisors-lemma-extend-admissible-blowups} により、
$Y_0$ への制限が $Y'_0 \to Y_0$ となる $V$-許容吹き上げ
$Y' \to Y$ が存在する。$f' : X' \to Y'$ を $f$ の厳密変換とすれば、
$f'$ の $Y'_0$ 上への制限が同型なので、所望の主張が成り立つ。
\end{proof}

\begin{lemma}
\label{flat-lemma-find-common-blowups}
$S$ を準コンパクトかつ準分離的なスキームとする。
$U \to X_1$ と $U \to X_2$ を $S$ 上のスキームの開埋め込みとし、
$U$, $X_1$, $X_2$ は $S$ 上有限型かつ分離的であると仮定する。
このとき $S$ 上のスキームの可換図式
$$
\xymatrix{
X_1' \ar[d] \ar[r] & X & X_2' \ar[l] \ar[d] \\
X_1 & U \ar[l] \ar[lu] \ar[u] \ar[ru] \ar[r] & X_2
}
$$
で、$X_i' \to X_i$ が $U$-許容吹き上げ、$X_i' \to X$ が開埋め込み、
$X$ が $S$ 上分離的かつ有限型となるものが存在する。
\end{lemma}

\begin{proof}
証明全体を通じて、すべてのスキームは $S$ 上分離的かつ有限型とする。
特に、これらのスキームは準コンパクトかつ準分離的であり、その間の射は
準コンパクトかつ分離的である。Schemes, Sections
\ref{schemes-section-quasi-compact} および
\ref{schemes-section-separation-axioms} を参照せよ。
次の事実を用いる。このような $S$ 上のスキームの埋め込み $U \to W$ に対し、
$W$ における $U$ のスキーム論的像 $Z$ は $W$ の閉部分スキームで、
$U \to Z$ は開埋め込み、$U \subset Z$ はスキーム論的に稠密であり、
$U \subset Z$ は位相的に稠密である。Morphisms, Lemma
\ref{morphisms-lemma-quasi-compact-immersion} を参照せよ。

\medskip\noindent
$X_{12} \subset X_1 \times_S X_2$ を
$U \to X_1 \times_S X_2$ のスキーム論的像とする。Morphisms, Lemma
\ref{morphisms-lemma-scheme-theoretic-image-of-partial-section} により、
射影 $p_i : X_{12} \to X_i$ は同型 $p_i^{-1}(U) \to U$ を誘導する。
$X_{12}$ の厳密変換 $X_{12}^i$ が $X_i^i$ の開部分スキームに同型となるような
$U$-許容吹き上げ $X_i^i \to X_i$ を選ぶ。Lemma
\ref{flat-lemma-zariski-after-blowup} を参照せよ。対応する有限型準連接イデアル層を
$\mathcal{I}_i \subset \mathcal{O}_{X_i}$ とする。
$X_{12}^i \to X_{12}$ は
$p_i^{-1}\mathcal{I}_i \mathcal{O}_{X_{12}}$ における吹き上げであることを
思い出そう。Divisors, Lemma \ref{divisors-lemma-strict-transform} を参照せよ。
$X_{12}'$ を $X_{12}$ の
$p_1^{-1}\mathcal{I}_1 p_2^{-1}\mathcal{I}_2 \mathcal{O}_{X_{12}}$ における
吹き上げとする。その意味については Divisors, Lemma
\ref{divisors-lemma-blowing-up-two-ideals} を参照せよ。特に可換図式
$$
\xymatrix{
X_{12}' \ar[d] \ar[r] & X_{12}^2 \ar[d] \\
X_{12}^1 \ar[r] & X_{12}
}
$$
を得る。ここですべての射は $U$-許容吹き上げである。
$X_{12}^i \subset X_i^i$ は開部分スキームなので、$X_{12}' \to X_{12}^i$ に
制限される $U$-許容吹き上げ $X_i' \to X_i^i$ を選べる。Divisors, Lemma
\ref{divisors-lemma-extend-admissible-blowups} を参照せよ。
すると $X_{12}' \subset X_i'$ は開部分スキームであり、図式
$$
\xymatrix{
X_{12}' \ar[d] \ar[r] & X_i' \ar[d] \\
X_{12}^i \ar[r] & X_i^i
}
$$
は可換で、縦の射は吹き上げ、横の射は開埋め込みである。
$X'_{12} \to X_1' \times_S X_2'$ は埋め込みかつ固有であることに注意する
（$X'_{12} \to X_{12}$ は固有、$X_{12} \to X_1 \times_S X_2$ は閉、
$X_1' \times_S X_2' \to X_1 \times_S X_2$ は分離的であることを用い、
Morphisms, Lemma \ref{morphisms-lemma-image-proper-scheme-closed} を適用する）。
したがって $X'_{12} \to  X_1' \times_S X_2'$ は閉埋め込みである。
共通の開部分スキーム $X_{12}'$ に沿って $X_1'$ と $X_2'$ を貼り合わせて
$X$ を定義すれば、$X \to S$ は有限型かつ分離的である
(Lemma \ref{flat-lemma-check-separated})。
$U$-許容吹き上げの合成は $U$-許容吹き上げなので
(Divisors, Lemma \ref{divisors-lemma-composition-admissible-blowups})、
補題が証明された。
\end{proof}

\begin{lemma}
\label{flat-lemma-replaced-by-strict-transform}
$X \to S$ と $Y \to S$ をスキームの射とし、
$U \subset X$ を開部分スキームとする。$V \to X \times_S Y$ を、
第一射影との合成が $U$ に写る準コンパクト射とする。
$Z \subset X \times_S Y$ を $V \to X \times_S Y$ のスキーム論的像とし、
$X' \to X$ を $U$-許容吹き上げとする。このとき
$V \to X' \times_S Y$ のスキーム論的像は、この吹き上げに関する
$Z$ の厳密変換である。
\end{lemma}

\begin{proof}
$Z' \to Z$ を厳密変換とする。射 $Z' \to X'$ は射
$Z' \to X' \times_S Y$ を誘導し、これは閉埋め込みである
（定義により $Z'$ は $X' \times_X Z$ の閉部分スキームだからである）。
したがって証明を終えるには、$V \to Z'$ のスキーム論的像 $Z''$ が
$Z'$ であることを示せば十分である。$Z'' \subset Z'$ は閉部分スキームで、
$V \to Z'$ は $Z''$ を経由することに注意する。
$V \to X \times_S Y$ と $V \to X' \times_S Y$ はともに準コンパクトである
（後者については Schemes, Lemma
\ref{schemes-lemma-quasi-compact-permanence} と、固有射の基底変換として
$X' \times_S Y \to X \times_S Y$ が分離的であることから従う）。したがって
Morphisms, Lemma \ref{morphisms-lemma-quasi-compact-scheme-theoretic-image}
により、$Z \cap (U \times_S Y) = Z'' \cap (U \times_S Y)$ である。
ゆえに包含射 $Z'' \to Z'$ は、$Z' \to Z$ の例外因子 $E$ の外で同型である。
しかし $Z'$ の構造層には $E$ 上に支持をもつ非零切断が存在しない
（厳密変換の定義による）。よって全射
$\mathcal{O}_{Z'} \to \mathcal{O}_{Z''}$ は同型でなければならない。
\end{proof}

\begin{lemma}
\label{flat-lemma-compactification-dominates}
$S$ を準コンパクトかつ準分離的なスキームとし、$U$ を $S$ 上有限型かつ
分離的なスキームとする。$V \subset U$ を準コンパクト開部分スキームとする。
$V$ が $S$ 上のコンパクト化 $V \subset Y$ をもつならば、
$V$-許容吹き上げ $Y' \to Y$ と開部分スキーム $V \subset V' \subset Y'$ で、
$V \to U$ が固有射 $V' \to U$ に延長されるものが存在する。
\end{lemma}

\begin{proof}
「対角」射 $V \to Y \times_S U$ のスキーム論的像を
$Z \subset Y \times_S U$ とする。$Y$ を $V$-許容吹き上げで置き換えると、
$Z$ はこの吹き上げに関する厳密変換で置き換わる。Lemma
\ref{flat-lemma-replaced-by-strict-transform} を参照せよ。したがって Lemma
\ref{flat-lemma-zariski-after-blowup} により、$Z \to Y$ は開埋め込みであると
仮定してよい。その像を $V' \subset Y$ とする。射影
$Y \times_S U \to U$ は固有で、$V' \cong Z$ は
$Y \times_S U$ の閉部分スキームだから、誘導される射 $V' \to U$ は固有である。
\end{proof}

\noindent
次の補題は Noetherian の場合についてのみ定式化する。
準コンパクトかつ準分離的なスキームに対する版も成り立つが、
それは本節の主定理から直ちに従う。

\begin{lemma}
\label{flat-lemma-two-compactifications}
$S$ を Noetherian スキームとし、$U$ を $S$ 上有限型かつ分離的な
スキームとする。$U = U_1 \cup U_2$ を開被覆とし、$U_1$ と $U_2$ は
$S$ 上のコンパクト化をもち、$U_1 \cap U_2$ は $U$ で稠密であるとする。
このとき $U$ は $S$ 上のコンパクト化をもつ。
\end{lemma}

\begin{proof}
$i = 1, 2$ に対してコンパクト化 $U_i \subset X_i$ を選ぶ。
$U_i$ は $X_i$ でスキーム論的に稠密であると仮定してよい。Lemma
\ref{flat-lemma-compactification-dominates} により、開部分スキーム
$V_i \subset X_i$ と、$\text{id} : U_i \to U_i$ を延長する固有射
$\psi_i : V_i \to U$ が存在すると仮定してよい。図式は
$$
\xymatrix{
U_i \ar[r] \ar[d] & V_i \ar[r] \ar[dl]^{\psi_i} & X_i \\
U
}
$$
で表される。
$\{i, j\} = \{1, 2\}$ のとき、
$Z_i = U \setminus U_j = U_i \setminus (U_1 \cap U_2)$ および
$Z_j = U \setminus U_i = U_j \setminus (U_1 \cap U_2)$ とおく。
すると
$$
U = U_1 \amalg Z_2 = Z_1 \amalg U_2 = Z_1 \amalg (U_1 \cap U_2) \amalg Z_2
$$
$Z_{i, i} \subset V_i$ を $\psi_i$ による $Z_i$ の逆像とする。
$\psi_i$ は $Z_i$ のある開近傍上で同型であることに注意する。
$Z_{i, j} \subset V_i$ を $\psi_i$ による $Z_j$ の逆像とする。
$\psi_i : Z_{i, j} \to Z_j$ は固有射であることに注意する。
$Z_i$ と $Z_j$ は $U$ の互いに素な閉部分集合なので、
$Z_{i, i}$ と $Z_{i, j}$ は $V_i$ の互いに素な閉部分集合である。

\medskip\noindent
$\overline{Z}_{i, i}$ と $\overline{Z}_{i, j}$ を、それぞれ
$Z_{i, i}$ と $Z_{i, j}$ の $X_i$ における閉包とする。
$X_i$ を $V_i$-許容吹き上げで置き換えた後、
$\overline{Z}_{i, i}$ と $\overline{Z}_{i, j}$ は互いに素であると仮定してよい。
Lemma \ref{flat-lemma-separate-disjoint-locally-closed-by-blowing-up} を参照せよ。
$X_1$ と $X_2$ の両方についてこれが成り立つと仮定する。
$X_i$ をさらに $V_i$-許容吹き上げで置き換えても、この性質は保たれる。

\medskip\noindent
$V_{12} = V_1 \times_U V_2$ とおく。埋め込み
$V_{12} \to X_1 \times_S X_2$ があり、これは閉埋め込み
$V_{12} = V_1 \times_U V_2 \to V_1 \times_S V_2$
(Schemes, Lemma \ref{schemes-lemma-fibre-product-after-map}) と開埋め込み
$V_1 \times_S V_2 \to X_1 \times_S X_2$ の合成である。
$X_{12} \subset X_1 \times_S X_2$ を
$V_{12} \to X_1 \times_S X_2$ のスキーム論的像とする。射影
$$
p_1 : X_{12} \to X_1
\quad\text{and}\quad
p_2 : X_{12} \to X_2
$$
は、$X_1$ と $X_2$ が $S$ 上固有なので固有である。$X_1$ を
$V_1$-許容吹き上げで置き換えると、$X_{12}$ はこの吹き上げに関する
厳密変換で置き換わる。Lemma \ref{flat-lemma-replaced-by-strict-transform} を参照せよ。

\medskip\noindent
合成
$\psi = \psi_1 \circ p_1|_{V_{12}} = \psi_2 \circ p_2|_{V_{12}}$ を
$\psi : V_{12} \to U$ と記す。閉部分スキーム
$$
Z_{12, 2} =
(p_1|_{V_{12}})^{-1}(Z_{1, 2}) =
(p_2|_{V_{12}})^{-1}(Z_{2, 2}) =
\psi^{-1}(Z_2) \subset V_{12}
$$
を考える。$\psi_2 : V_2 \to U$ は $Z_2$ のある開近傍上で同型であり、
$V_{12} = V_1 \times_U V_2$ なので、射
$p_1|_{V_{12}} : V_{12} \to V_1$ は $Z_{1, 2}$ のある開近傍上で同型である。
Lemma \ref{flat-lemma-blowup-iso-along} により、$p_1$ の厳密変換
$p'_1 : X'_{12} \to X'_1$ が $X'_1$ における $Z_{1, 2}$ の閉包の
ある開近傍上で同型となるような $V_1$-許容吹き上げ
$X_1' \to X_1$ が存在する。$X_1$ を $X'_1$ で、$X_{12}$ を $X'_{12}$ で
置き換えた後、$p_1$ は $\overline{Z}_{1, 2}$ のある開近傍上で
同型であると仮定してよい。

\medskip\noindent
前段落の帰着により
$$
X_{12} \cap (\overline{Z}_{1, 2} \times_S \overline{Z}_{2, 1}) = \emptyset
$$
であり、ここで交わりは $X_1 \times_S X_2$ の中でとる。実際、$X_{12}$ における
逆像 $p_1^{-1}(\overline{Z}_{1, 2})$ は $\overline{Z}_{1, 2}$ に
同型に写る。特に、$Z_{12, 2}$ は
$p_1^{-1}(\overline{Z}_{1, 2})$ で稠密である。したがって $p_2$ は
$p_1^{-1}(\overline{Z}_{1, 2})$ を $\overline{Z}_{2, 2}$ に写す。
$\overline{Z}_{2, 2} \cap \overline{Z}_{2, 1} = \emptyset$ なので結論を得る。

\medskip\noindent
貼り合わせによって得られるスキーム
$$
W_i = U \coprod\nolimits_{U_i} (X_i \setminus \overline{Z}_{i, j}),
\quad i = 1, 2
$$
を考える。Lemma \ref{flat-lemma-check-separated} を適用して、
$W_i \to S$ が分離的であることを示そう。まず $U \to S$ と
$X_i \to S$ は分離的である。埋め込み
$U_i \to U \times_S (X_i \setminus \overline{Z}_{i, j})$ は閉である。
実際、$u_i \in U_i$ と $u \in U \setminus U_i$ を満たす任意の特殊化
$u_i \leadsto u$ は、固有射 $\psi_i : V_i \to U$ に沿って $V_i$ 内の特殊化
$u_i \leadsto v_i$ に一意に持ち上がり、このとき $v_i$ は
$Z_{i, j}$ に属さなければならない。したがって埋め込みの像は閉であり、
よってこの埋め込みは閉埋め込みである。

\medskip\noindent
一方、$S$ 上の任意の付値環 $A$、その分数体 $K$、および $S$ 上の任意の射
$\gamma : \Spec(K) \to (U_1 \cap U_2)$ に対し、ある $i$ と、$\gamma$ の延長
$h_i : \Spec(A) \to W_i$ が存在する。実際、$i = 1, 2$ のいずれについても、
$X_i$ の $S$ 上の固有性に対する付値判定法により、$\gamma$ を延長する射
$g_i : \Spec(A) \to X_i$ が存在する。Morphisms, Lemma
\ref{morphisms-lemma-characterize-proper} を参照せよ。したがって問題となるのは、
$i = 1, 2$ の両方について
$g_i(\mathfrak m_A) \in \overline{Z}_{i, j}$ となる場合だけである。
しかしこれは、上で証明した $X_{12}$ と
$\overline{Z}_{1, 2} \times_S \overline{Z}_{2, 1}$ の交わりが空であることに反する。

\medskip\noindent
Lemma \ref{flat-lemma-find-common-blowups} におけるような図式
$$
\xymatrix{
W_1' \ar[d] \ar[r] & W & W_2' \ar[l] \ar[d] \\
W_1 & U \ar[l] \ar[lu] \ar[u] \ar[ru] \ar[r] & W_2
}
$$
を考える。前段落により、任意の実線図式
$$
\xymatrix{
\Spec(K) \ar[r]_\gamma  \ar[d] & W \ar[d] \\
\Spec(A) \ar@{..>}[ru] \ar[r] & S
}
$$
で $\Im(\gamma) \subset U_1 \cap U_2$ であるものごとに、ある $i$ と
$\gamma$ の延長 $h_i : \Spec(A) \to W_i$ が存在する。
$W'_i \to W_i$ の固有性に対する付値判定法を用いると、$h_i$ を
$h'_i : \Spec(A) \to W'_i$ に持ち上げられる。したがって図式の点線の射が
存在する。$W$ は $S$ 上分離的なので、この射は一意でもある。
Morphisms, Lemma
\ref{morphisms-lemma-refined-valuative-criterion-universally-closed} により、
これは $W \to S$ が普遍閉であることを意味する。
$W \to S$ はすでに有限型かつ分離的なので、これで証明が完了する。
\end{proof}

\begin{theorem}
\label{flat-theorem-nagata}
\begin{reference}
\cite{Lutkebohmert}, \cite{Conrad-Nagata}, \cite{Nagata-1},
\cite{Nagata-2}, \cite{Nagata-3}, \cite{Nagata-4} を参照せよ。
\end{reference}
$S$ を準コンパクトかつ準分離的なスキームとし、
$X \to S$ を分離的有限型射とする。このとき $X$ は $S$ 上のコンパクト化をもつ。
\end{theorem}

\begin{proof}
まず Noetherian の場合に帰着する。読者にはこの段落を飛ばすことを強く勧める。
$X' \to S$ が有限表示かつ分離的となる閉埋め込み $X \to X'$ が存在する。
Limits, Proposition
\ref{limits-proposition-separated-closed-in-finite-presentation} を参照せよ。
$X'$ の $S$ 上のコンパクト化が得られれば、その中で $X$ の
スキーム論的像をとることで $X$ の $S$ 上のコンパクト化を得る。
したがって $X \to S$ は分離的かつ有限表示であると仮定してよい。
$S = \lim S_i$ を、アフィン遷移射をもつ Noetherian スキームの系の
有向極限として書ける。Limits, Proposition
\ref{limits-proposition-approximate} を参照せよ。ある $i$ と有限表示射
$X_i \to S_i$ で、その $S$ への基底変換が $X \to S$ となるものを選べる。
Limits, Lemma \ref{limits-lemma-descend-finite-presentation} を参照せよ。
$i$ を大きくした後、$X_i \to S_i$ は分離的であると仮定してよい。
Limits, Lemma \ref{limits-lemma-descend-separated-finite-presentation} を参照せよ。
$X_i$ の $S_i$ 上のコンパクト化が得られれば、その $S$ への基底変換は
$X$ の $S$ 上のコンパクト化となる。これで次段落の場合に帰着した。

\medskip\noindent
$S$ は Noetherian であると仮定する。有限アフィン開被覆
$X = \bigcup_{i = 1, \ldots, n} U_i$ で、
$U_1 \cap \ldots \cap U_n$ が $X$ で稠密となるものを選べる。
これは Properties, Lemma
\ref{properties-lemma-point-and-maximal-points-affine} と、$X$ が準コンパクトで
既約成分を有限個しかもたないことから従う。各 $i$ に対し、Morphisms, Lemma
\ref{morphisms-lemma-quasi-affine-finite-type-over-S} により、ある
$n_i \geq 0$ と埋め込み $U_i \to \mathbf{A}^{n_i}_S$ を選べる。
したがって、$\mathbf{P}^{n_i}_S$ におけるスキーム論的像をとることにより、
$i = 1, \ldots, n$ に対して $U_i$ は $S$ 上のコンパクト化をもつ。
Lemma \ref{flat-lemma-two-compactifications} を $(n - 1)$ 回適用すれば、
定理が成り立つと結論できる。
\end{proof}






\section{h 位相}
\label{flat-section-h}

\noindent
大まかにいえば、本章でいう h 層とは、有限表示固有全射に対しても
層条件を満たす Zariski 位相の層である。Lemma
\ref{flat-lemma-characterize-sheaf-h} を参照せよ。ただし、スキームの圏上の
h 位相の定義は文献によって異なることを指摘しておく価値がある。

\medskip\noindent
Voevodsky は当初、h 被覆を、共同で普遍的に submersive な有限型射の有限族として
定義した (Morphisms, Definition \ref{morphisms-definition-submersive})。
\cite[Definition 3.1.2]{Voevodsky} を参照せよ。この定義は、基礎となる
スキームの圏を、固定した Noetherian 基底スキーム上有限型なすべてのスキームに
制限すると最もうまく機能する。この設定で Voevodsky は h 被覆を ph 被覆と
関係づける。ph 位相は Zariski 被覆と固有全射によって生成される。
詳しくは Topologies, Section \ref{topologies-section-ph} を参照せよ。

\medskip\noindent
Topologies, Section \ref{topologies-section-V} では V 位相を研究した。
準コンパクト射 $X \to Y$ は、$Y$ の点の任意の特殊化が $X$ の点の特殊化の像であり、
しかも任意の基底変換後にも同じことが成り立つとき、V 被覆を定める
(Topologies, Lemma \ref{topologies-lemma-refine-qcqs-V})。
この場合 $X \to Y$ は普遍的に submersive である
(Topologies, Lemma \ref{topologies-lemma-V-covering-universally-submersive})。
非 Noetherian スキームを扱う際、V 被覆の概念は
``終域を共有し、共同で普遍的に submersive な射の族'' のよい代替となる。

\medskip\noindent
本章では、まず局所有限表示な射の（無限でもよい）族について、ph 被覆と
V 被覆の同値性を証明する。その後、Stacks project における h 被覆の概念として
これらの族を用いる。Noetherian スキームと有限族について、これらの被覆は
Voevodsky の定義によるものと一致する。Lemma
\ref{flat-lemma-Noetherian-h-covering} を参照せよ。固定した準コンパクトかつ
準分離的なスキーム $S$ 上有限表示なスキームの圏では、これらの被覆は
\cite[Definition 2.7]{Witt-Grass} の位相と同じ位相を定める。

\begin{lemma}
\label{flat-lemma-equivalence-h-v-locally-finite-presentation}
$\{f_i : X_i \to X\}_{i \in I}$ を、終域を共有するスキームの射の族とし、
すべての $i$ に対して $f_i$ は局所有限表示であるとする。次は同値である。
\begin{enumerate}
\item $\{X_i \to X\}$ は ph 被覆である。
\item $\{X_i \to X\}$ は V 被覆である。
\end{enumerate}
\end{lemma}

\begin{proof}
$U \subset X$ をアフィン開部分スキームとする。Topologies, Definitions
\ref{topologies-definition-ph-covering} および
\ref{topologies-definition-V-covering} を見ると、基底変換
$\{X_i \times_X U \to U\}$ が標準 ph 被覆で細分されることと
標準 V 被覆で細分されることが同値だと示せば十分である。したがって
$X$ はアフィンであると仮定してよく、$\{X_i \to X\}$ が標準 ph 被覆で
細分されることと標準 V 被覆で細分されることが同値だと示せばよい。
標準 ph 被覆は標準 V 被覆なので、他方の含意だけを証明すれば十分である。
Topologies, Lemma \ref{topologies-lemma-standard-ph-standard-V} を参照せよ。

\medskip\noindent
$X$ はアフィンで、$\{f_i : X_i \to X\}_{i \in I}$ は標準 V 被覆
$\{g_j : Y_j \to X\}_{j = 1, \ldots, m}$ で細分されると仮定する。
各 $j$ に対し、$g_j = f_{i_j} \circ h_j$ を満たす $i_j$ と射
$h_j : Y_j \to X_{i_j}$ を選ぶ。$Y_j$ はアフィン、したがって準コンパクトなので、
各 $j$ に対し、$\Im(h_j) \subset \bigcup U_{j, k}$ を満たす有限個のアフィン開集合
$U_{j, k} \subset X_{i_j}$ をとれる。すると
$\{U_{j, k} \to X\}_{j, k}$ は $\{X_i \to X\}$ を細分し、標準 V 被覆である
（これはアフィン間の射の有限族であり、$\{Y_j \to X\}$ の付値環に対する
持ち上げ性を継承するからである）。これで次段落の場合に帰着した。

\medskip\noindent
$\{f_i : X_i \to X\}_{i = 1, \ldots, n}$ は標準 V 被覆で、
$f_i$ は有限表示であると仮定する。$\{X_i \to X\}$ が標準 ph 被覆で
細分されることを示さなければならない。一般平坦性成層
$$
X = S \supset S_0 \supset S_1 \supset \ldots \supset S_t = \emptyset
$$
を、アフィンの有限表示射
$$
\coprod\nolimits_{i = 1, \ldots, n} f_i :
\coprod\nolimits_{i = 1, \ldots, n} X_i
\longrightarrow
X
$$
に対して More on Morphisms, Lemma
\ref{more-morphisms-lemma-generic-flatness-stratification-scheme} のように選ぶ。
以下、この成層のすべての性質を断りなく用いる。構成により、各 $f_i$ の
$U_k = S_k \setminus S_{k + 1}$ への基底変換は平坦である。
$Y_k$ を $S_k$ における $U_k$ のスキーム論的閉包とする。
$U_k \to S_k$ は準コンパクト開埋め込みなので (Properties, Lemma
\ref{properties-lemma-quasi-coherent-finite-type-ideals})、
$U_k \subset Y_k$ は準コンパクトで稠密な（かつスキーム論的に稠密な）
開埋め込みである。Morphisms, Lemma
\ref{morphisms-lemma-quasi-compact-scheme-theoretic-image} を参照せよ。
射 $\coprod_{k = 0, \ldots, t - 1} Y_k \to X$ は有限全射なので、
$\{Y_k \to X\}$ は標準 ph 被覆、したがって標準 V 被覆である（上を参照）。
標準 V 被覆の推移性 (Topologies, Lemma
\ref{topologies-lemma-composition-standard-V}) により、被覆
$\{X_i \to X\}$ の各 $Y_k$ への引き戻しが標準 V 被覆で細分されることを
示せば十分である。これで次段落の場合に帰着した。

\medskip\noindent
$\{f_i : X_i \to X\}_{i = 1, \ldots, n}$ は $f_i$ が有限表示である
標準 V 被覆で、稠密な準コンパクト開部分スキーム $U \subset X$ が存在し、
$X_i \times_X U \to U$ は平坦であると仮定する。Theorem
\ref{flat-theorem-flatten-module} により、$f_i$ の厳密変換
$f'_i : X'_i \to X'$ が平坦となるような $U$-許容吹き上げ
$X' \to X$ が存在する。$U \subset X$ は稠密で、$U$ は $X'$ の開部分と
同一視されるので、射影的（したがって閉）な射 $X' \to X$ は全射である。
必要なら $X'$ をさらに $U$-許容吹き上げで置き換えて、
$U \subset X'$ はスキーム論的にも稠密であると仮定してよい
(Remark \ref{flat-remark-successive-blowups})。したがって任意の点 $x \in X'$ に対し、
付値環 $V$ と射 $g : \Spec(V) \to X'$ で、$\Spec(V)$ の一般点が $U$ に写り、
$\Spec(V)$ の閉点が $x$ に写るものが存在する。Morphisms, Lemma
\ref{morphisms-lemma-reach-points-scheme-theoretic-image} を参照せよ。
$\{X_i \to X\}$ は標準 V 被覆なので、付値環の拡大 $V \subset W$、添字 $i$、
および次の図式を可換にする射 $\Spec(W) \to X_i$ を選べる。
$$
\xymatrix{
\Spec(W) \ar[d] \ar[rr] & & X_i \ar[d] \\
\Spec(V) \ar[r] & X' \ar[r] & X
}
$$
$X'_i \subset X' \times_X X_i$ は開部分 $U \times_X X_i$ を含む閉部分スキームで、
$\Spec(W)$ は整スキームであり、誘導される射
$h : \Spec(W) \to X' \times_X X_i$ は $\Spec(W)$ の一般点を
$U \times_X X_i$ に写す。したがって $h$ は閉部分スキーム
$X'_i \subset X' \times_X X_i$ を経由する。よって
$\{f'_i : X'_i \to X'\}$ は V 被覆である。特に $\coprod f'_i$ は全射であり、
$\{X'_i \to X'\}$ は fppf 被覆である。fppf 被覆は ph 被覆なので
(More on Morphisms, Lemma \ref{more-morphisms-lemma-fppf-ph})、
$\{X'_i \to X\}$ を細分する標準 ph 被覆 $\{Y_j \to X'\}$ をとれる。
% P05-EN-CH38-0269: frozen-source pointer; see ERRATA_CANDIDATES.jsonl.
この被覆は、固有全射 $Y \to X'$ と有限アフィン開被覆
$Y = \bigcup Y_j$ で与えられるとする。このとき合成 $Y \to X$ は固有全射であり、
$\{Y_j \to X\}$ は標準 ph 被覆であると結論できる。これで証明が完了した。
\end{proof}

\noindent
以下が我々の定義である。

\begin{definition}
\label{flat-definition-h-covering}
$T$ をスキームとする。{\it $T$ の h 被覆}とは、各 $f_i$ が局所有限表示であり、
Lemma \ref{flat-lemma-equivalence-h-v-locally-finite-presentation} の同値な条件の一つ
（従ってすべて）を満たす射の族
$\{f_i : T_i \to T\}_{i \in I}$ のことである。
\end{definition}

\noindent
Noetherian スキームに対して、これは ph 被覆と同じものであり
（以下の Lemma \ref{flat-lemma-Noetherian-h-ph} に記録する）、
Voevodsky の概念が復元される。

\begin{lemma}
\label{flat-lemma-Noetherian-h-covering}
$X$ を Noetherian スキームとする。$\{X_i \to X\}_{i \in I}$ を
有限型射の有限族とする。このとき次は同値である。
\begin{enumerate}
\item $\coprod_{i \in I} X_i \to X$ は普遍的に submersive である
（Morphisms, Definition
\ref{morphisms-definition-submersive}）。
\item $\{X_i \to X\}_{i \in I}$ は h 被覆である。
\end{enumerate}
\end{lemma}

\begin{proof}
(2) $\Rightarrow$ (1) は、より一般的な Topologies, Lemma
\ref{topologies-lemma-V-covering-universally-submersive} と h 被覆の定義から従う。
$\coprod X_i \to X$ が普遍的に submersive であると仮定する。
$\{X_i \to X\}$ が ph 被覆によって細分できることを示す。
Topologies, Lemma \ref{topologies-lemma-refine-by-ph} と h 被覆の定義により、
これで十分である。議論は Lemma
\ref{flat-lemma-equivalence-h-v-locally-finite-presentation} の証明で用いたものと同じである。

\medskip\noindent
一般平坦性による層別化
$$
X = S \supset S_0 \supset S_1 \supset \ldots \supset S_t = \emptyset
$$
を、有限表示射
$$
\coprod\nolimits_{i = 1, \ldots, n} f_i :
\coprod\nolimits_{i = 1, \ldots, n} X_i
\longrightarrow
X
$$
に対し、More on Morphisms, Lemma
\ref{more-morphisms-lemma-generic-flatness-stratification-scheme} のように選ぶ。
以下ではこの層別化の性質を断りなくすべて用いる。構成により、各 $f_i$ の
$U_k = S_k \setminus S_{k + 1}$ への基底変換は平坦である。
$Y_k$ を $S_k$ における $U_k$ のスキーム論的閉包とする。
$U_k \to S_k$ は準コンパクト開埋め込みであるから（この段落のスキームはすべて
Noetherian である）、$U_k \subset Y_k$ は準コンパクトな稠密
（かつスキーム論的に稠密な）開埋め込みである。Morphisms, Lemma
\ref{morphisms-lemma-quasi-compact-scheme-theoretic-image} を参照せよ。
射 $\coprod_{k = 0, \ldots, t - 1} Y_k \to X$ は有限全射なので、
$\{Y_k \to X\}$ は ph 被覆である。ph 被覆の推移性
（Topologies, Lemma \ref{topologies-lemma-ph}）により、被覆
$\{X_i \to X\}$ の各 $Y_k$ への引き戻しが ph 被覆によって細分できることを
示せば十分である。従って次の段落で述べる場合に帰着する。

\medskip\noindent
$\coprod X_i \to X$ が普遍的に submersive であり、任意の $i$ に対して
$X_i \times_X U \to U$ が平坦となる稠密開部分 $U \subset X$ が存在すると仮定する。
Theorem \ref{flat-theorem-flatten-module} により、各 $i$ について $f_i$ の厳密変換
$f'_i : X'_i \to X'$ が平坦となる $U$-許容吹き上げ $X' \to X$ が存在する。
$U \subset X$ は稠密で、$U$ は $X'$ の開部分と同一視されるから、射影的
（従って閉）な射 $X' \to X$ は全射である。必要なら $X'$ をさらに
$U$-許容吹き上げで置き換えて、$U \subset X'$ が稠密であるとも仮定してよい
（Remark \ref{flat-remark-successive-blowups} を参照）。従って任意の点 $x \in X'$ に対し、
離散付値環 $A$ と射 $g : \Spec(A) \to X'$ で、$\Spec(A)$ の一般点を $U$ に、
$\Spec(A)$ の閉点を $x$ に写すものが存在する。Limits, Lemma
\ref{limits-lemma-reach-point-closure-Noetherian} を参照せよ。
次のようにおく。
$$
W = \Spec(A) \times_X \coprod X_i = \coprod \Spec(A) \times_X X_i
$$
$\coprod X_i \to X$ は普遍的に submersive なので、$W$ に特殊化
$w' \leadsto w$ が存在し、$w'$ は $\Spec(A)$ の一般点に、$w$ は
$\Spec(A)$ の閉点に写る。（さもなければ $W \to \Spec(A)$ の閉ファイバーは
一般化について安定であり、従って開となるが、これは $W \to \Spec(A)$ が
submersive であることに反する。）$w' \in \Spec(A) \times_X X_i$ とする。
するともちろん $w \in \Spec(A) \times_X X_i$ でもある。
$X' \times_X X_i$ における $w' \leadsto w$ の像を
$x'_i \leadsto x_i$ とする。$x'_i \in X'_i$ であり、
$X'_i \subset X' \times_X X_i$ は閉部分スキームなので、$x_i \in X'_i$ である。
$x_i$ は $x \in X'$ に写るから、$\coprod X'_i \to X'$ は全射である！
特に $\{X'_i \to X'\}$ は fppf 被覆である。しかし fppf 被覆は ph 被覆である
（More on Morphisms, Lemma \ref{more-morphisms-lemma-fppf-ph}）。
$X' \to X$ は固有全射なので、$\{X'_i \to X\}$ は ph 被覆である。
これで証明が完了した。
\end{proof}

\begin{lemma}
\label{flat-lemma-Noetherian-h-ph}
$X$ を局所 Noetherian スキームとする。$X$ を終域とする射の族
$\{f_i : X_i \to X\}_{i \in I}$ が h 被覆であるための必要十分条件は、
それが ph 被覆であることである。
\end{lemma}

\begin{proof}
Definition \ref{flat-definition-h-covering} により、h 被覆は ph 被覆である。
逆に $\{f_i : X_i \to X\}$ が ph 被覆なら、射 $f_i$ は局所有限型である
（Topologies, Definition \ref{topologies-definition-ph-covering}）。
$X$ は局所 Noetherian なので各 $f_i$ は局所有限表示であり、定義により
h 被覆を得る。
\end{proof}

\noindent
次の補題と \cite[Theorem 8.4]{rydh_descent} は、我々の定義が David Rydh の論文
\cite{rydh_descent} の定義と一致する（少なくとも密接に関係する）ことを示す。
簡単のためアフィン基底に制限する。

\begin{lemma}
\label{flat-lemma-approximate-h-cover}
$X$ をアフィンスキームとし、$\{X_i \to X\}_{i \in I}$ を h 被覆とする。
このとき有限表示である全射固有射（！）
$$
Y \longrightarrow X
$$
と有限アフィン開被覆 $Y = \bigcup_{j = 1, \ldots, m} Y_j$ で、
$\{Y_j \to X\}_{j = 1, \ldots, m}$ が $\{X_i \to X\}_{i \in I}$ を
細分するものが存在する。
\end{lemma}

\begin{proof}
仮定により、固有全射 $Y \to X$ と有限アフィン開被覆
$Y = \bigcup_{j = 1, \ldots, m} Y_j$ で、
$\{Y_j \to X\}_{j = 1, \ldots, m}$ が $\{X_i \to X\}_{i \in I}$ を
細分するものが存在する。これは各 $j$ に対して添字 $i_j \in I$ と
$X$ 上の射 $h_j : Y_j \to X_{i_j}$ が存在することを意味する。
Definition \ref{flat-definition-h-covering} と Topologies, Definition
\ref{topologies-definition-ph-covering} を参照せよ。
問題は $Y \to X$ が有限表示であるとは分かっていないことである。
Limits, Lemma \ref{limits-lemma-proper-limit-of-proper-finite-presentation} により、
$$
Y = \lim Y_\lambda
$$
と書ける。ここで $Y_\lambda$ は $X$ 上固有かつ有限表示のスキームからなる有向系で、
射 $Y \to Y_\lambda$ および遷移射は閉埋め込みである。各
$Y_\lambda \to X$ が全射であることに注意する。
Limits, Lemma \ref{limits-lemma-descend-opens} により、ある $\lambda$ と、
$Y_\lambda$ を被覆し $Y$ 上で $Y_j$ に制限される準コンパクト開部分
$Y_{\lambda, j} \subset Y_\lambda$（$j = 1, \ldots, m$）をとれる。
このとき $Y_j = \lim Y_{\lambda, j}$ である。
$\lambda$ を大きくすれば、すべての $j$ に対し $Y_{\lambda, j}$ がアフィンであると
仮定してよい。Limits, Lemma \ref{limits-lemma-limit-affine} を参照せよ。
最後に $X_i \to X$ は局所有限表示なので、局所有限表示射の関手的特徴づけ
（Limits, Proposition
\ref{limits-proposition-characterize-locally-finite-presentation}）を用い、各 $j$ に対して
射 $h_{\lambda, j} : Y_{\lambda, j} \to X_{i_j}$ が存在して、その $Y_j$ への制限が
上で選んだ射 $h_j$ となるような $\lambda$ をとれる。
従って $\{Y_{\lambda, j} \to X\}$ は $\{X_i \to X\}$ を細分し、証明が完了する。
\end{proof}

\noindent
h 被覆の一般論の展開に戻る。

\begin{lemma}
\label{flat-lemma-zariski-h}
fppf 被覆は h 被覆である。従って syntomic 被覆、滑らかな被覆、\'etale 被覆、
および Zariski 被覆も h 被覆である。
\end{lemma}

\begin{proof}
これは、fppf 被覆では射が局所有限表示であることを要求され、かつ fppf 被覆が
ph 被覆だからである。More on Morphisms, Lemma
\ref{more-morphisms-lemma-fppf-ph} を参照せよ。第二の主張は第一の主張と
Topologies, Lemma \ref{topologies-lemma-zariski-etale-smooth-syntomic-fppf} から従う。
\end{proof}

\begin{lemma}
\label{flat-lemma-surjective-proper-finite-presentation-h}
$f : Y \to X$ を有限表示であるスキームの全射固有射とする。
このとき $\{Y \to X\}$ は h 被覆である。
\end{lemma}

\begin{proof}
Topologies, Lemmas
\ref{topologies-lemma-zariski-etale-smooth-syntomic-fppf-fpqc-ph-V} and
\ref{topologies-lemma-surjective-proper-ph} を組み合わせればよい。
\end{proof}

\begin{lemma}
\label{flat-lemma-refine-by-h}
$T$ をスキームとする。$\{f_i : T_i \to T\}_{i \in I}$ を、すべての $i$ に対して
$f_i$ が局所有限表示である射の族とする。このとき次は同値である。
\begin{enumerate}
\item $\{T_i \to T\}_{i \in I}$ は h 被覆である。
\item $\{T_i \to T\}_{i \in I}$ を細分する h 被覆が存在する。
\item $\{\coprod_{i \in I} T_i \to T\}$ は h 被覆である。
\end{enumerate}
\end{lemma}

\begin{proof}
これは ph 被覆に対する類似の主張
（Topologies, Lemma \ref{topologies-lemma-refine-by-ph}）、または V 被覆に対する
類似の主張（Topologies, Lemma \ref{topologies-lemma-refine-by-V}）から従う。
\end{proof}

\noindent
次に、我々の h 被覆の概念が Sites, Definition \ref{sites-definition-site} の条件を
満たすことを示す。

\begin{lemma}
\label{flat-lemma-h}
$T$ をスキームとする。
\begin{enumerate}
\item $T' \to T$ が同型なら、$\{T' \to T\}$ は $T$ の h 被覆である。
\item $\{T_i \to T\}_{i\in I}$ が h 被覆であり、各 $i$ に対して
$\{T_{ij} \to T_i\}_{j\in J_i}$ が h 被覆なら、
$\{T_{ij} \to T\}_{i \in I, j\in J_i}$ は h 被覆である。
\item $\{T_i \to T\}_{i\in I}$ が h 被覆であり、$T' \to T$ がスキームの射なら、
$\{T' \times_T T_i \to T'\}_{i\in I}$ は h 被覆である。
\end{enumerate}
\end{lemma}

\begin{proof}
ph 被覆または V 被覆に対する対応する主張
（Topologies, Lemma \ref{topologies-lemma-ph} または
\ref{topologies-lemma-V}）と、局所有限表示射の類が基底変換および合成で
保たれることから直ちに従う。
\end{proof}

\noindent
次に、Stacks project で用いる大 h サイトを定義する。先へ進む前に、
Topologies, Section \ref{topologies-section-procedure} の一般的な議論を
読んでおくのがよい。

\begin{definition}
\label{flat-definition-big-h-site}
{\it 大 h サイト}とは、Sites, Definition \ref{sites-definition-site} のサイト
$\Sch_h$ で、次のように構成されるものをいう。
\begin{enumerate}
\item スキームの任意の集合 $S_0$ と、これらのスキームの間の h 被覆からなる
任意の集合 $\text{Cov}_0$ を選ぶ。
\item 台となる圏として、集合 $S_0$ から出発して Sets, Lemma
\ref{sets-lemma-construct-category} のように構成された任意の圏
$\Sch_\alpha$ をとる。
\item 圏 $\Sch_\alpha$、h 被覆の類、および上で選んだ集合 $\text{Cov}_0$ から出発して、
Sets, Lemma \ref{sets-lemma-coverings-site} のように被覆の任意の集合を選ぶ。
\end{enumerate}
\end{definition}

\noindent
大サイトの定義の動機と説明については、Topologies, Definition
\ref{topologies-definition-big-zariski-site} に続く注意を参照せよ。

\begin{definition}
\label{flat-definition-standard-h}
$T$ をアフィンスキームとする。$T$ の {\it 標準 h 被覆}とは、各 $T_i$ がアフィンで、
$f_i$ が有限表示であり、次の同値な条件のいずれかを満たす族
$\{f_i : T_i \to T\}_{i = 1, \ldots, n}$ をいう。(1) $\{T_i \to T\}$ は
標準 ph 被覆によって細分できる。(2) $\{T_i \to T\}$ は V 被覆である。
\end{definition}

\noindent
これらの条件の同値性は Lemma
\ref{flat-lemma-equivalence-h-v-locally-finite-presentation}、Topologies, Definition
\ref{topologies-definition-ph-covering}、および Lemma
\ref{topologies-lemma-refine-by-ph} から従う。

\medskip\noindent
スキーム $S$ の大 h サイトの導入を続ける前に、大 h サイト $\Sch_h$ 上の位相は、
ある意味ですべてのスキームの圏上の h 位相から誘導されることを指摘しておく。

\begin{lemma}
\label{flat-lemma-h-induced}
$\Sch_h$ を Definition \ref{flat-definition-big-h-site} の大 h サイトとする。
$T \in \Ob(\Sch_h)$ とし、$\{T_i \to T\}_{i \in I}$ を $T$ の任意の h 被覆とする。
\begin{enumerate}
\item サイト $\Sch_h$ における $T$ の被覆 $\{U_j \to T\}_{j \in J}$ で、
$\{T_i \to T\}_{i \in I}$ を細分するものが存在する。
\item $\{T_i \to T\}_{i \in I}$ が標準 h 被覆なら、それは自明な仕方で
$\Sch_h$ の被覆と同値である。
\item $\{T_i \to T\}_{i \in I}$ が Zariski 被覆なら、それは自明な仕方で
$\Sch_h$ の被覆と同値である。
\end{enumerate}
\end{lemma}

\begin{proof}
省略する。ヒント：これは Topologies, Lemma
\ref{topologies-lemma-ph-induced} の証明とまったく同じである。
\end{proof}

\begin{definition}
\label{flat-definition-big-small-h}
$S$ をスキームとし、$\Sch_h$ を $S$ を含む大 h サイトとする。
\begin{enumerate}
\item $S$ の {\it 大 h サイト}は $(\Sch/S)_h$ と記し、Sites, Section
\ref{sites-section-localize} で導入されたサイト $\Sch_h/S$ である。
\item $S$ の {\it 大アフィン h サイト}は $(\textit{Aff}/S)_h$ と記し、
その対象がアフィンな $U/S$ である $(\Sch/S)_h$ の充満部分圏とする。
$(\textit{Aff}/S)_h$ の被覆とは、標準 h 被覆である $(\Sch/S)_h$ の任意の被覆
$\{U_i \to U\}$ のことである。
\end{enumerate}
\end{definition}

\noindent
大アフィン h サイトがサイトであることを明記しておく。

\begin{lemma}
\label{flat-lemma-verify-site-h}
$S$ をスキームとし、$\Sch_h$ を $S$ を含む大 h サイトとする。
このとき $(\textit{Aff}/S)_h$ はサイトである。
\end{lemma}

\begin{proof}
Topologies, Lemma \ref{topologies-lemma-verify-site-etale} の証明と同様に論じれば、
標準 h 被覆の集まりが Sites, Definition \ref{sites-definition-site} の性質
(1)、(2)、(3) を満たすことを示せば十分である。これは明らかである。
例えば標準 h 被覆 $\{T_i \to T\}_{i\in I}$ と、各 $i$ に対する標準 h 被覆
$\{T_{ij} \to T_i\}_{j \in J_i}$ が与えられたとき、
$\{T_{ij} \to T\}_{i \in I, j\in J_i}$ は h 被覆であり
（Lemma \ref{flat-lemma-h}）、$\bigcup_{i\in I} J_i$ は有限で、各 $T_{ij}$ はアフィンである。
従って $\{T_{ij} \to T\}_{i \in I, j\in J_i}$ は標準 h 被覆である。
\end{proof}

\begin{lemma}
\label{flat-lemma-fibre-products-h}
$S$ をスキームとし、$\Sch_h$ を $S$ を含む大 h サイトとする。
サイト $\Sch_h$、$(\Sch/S)_h$、$(\textit{Aff}/S)_h$ の台となる圏は
ファイバー積をもつ。いずれの場合も、すべてのスキームの圏 $\Sch$ への自然な関手は
ファイバー積をとる操作と可換である。圏 $(\Sch/S)_h$ は終対象 $S/S$ をもつ。
\end{lemma}

\begin{proof}
$\Sch_h$ については構成から成り立つ。Sets, Lemma
\ref{sets-lemma-what-is-in-it} を参照せよ。$U, V, W \in \Ob(\Sch_h)$ である
スキームの射 $U \to S$、$V \to U$、$W \to U$ が与えられたとする。
$\Sch_h$ におけるファイバー積 $V \times_U W$ は $\Sch$ におけるファイバー積であり、
$S$ 上のすべてのスキームの圏では $U/S$ 上の $V/S$ と $W/S$ のファイバー積である。
従って $(\Sch/S)_h$ においてもファイバー積である。これで $(\Sch/S)_h$ に対する
主張が示された。$U, V, W$ がアフィンなら $V \times_U W$ もアフィンなので、
$(\textit{Aff}/S)_h$ に対する主張も従う。
\end{proof}

\noindent
次に、大アフィンサイトが大サイトと同じトポスを定めることを確認する。

\begin{lemma}
\label{flat-lemma-affine-big-site-h}
$S$ をスキームとし、$\Sch_h$ を $S$ を含む大 h サイトとする。関手
$(\textit{Aff}/S)_h \to (\Sch/S)_h$ は余連続であり、
$\Sh((\textit{Aff}/S)_h)$ から $\Sh((\Sch/S)_h)$ へのトポスの同値を誘導する。
\end{lemma}

\begin{proof}
特別な余連続関手の概念は Sites, Definition
\ref{sites-definition-special-cocontinuous-functor} で導入されている。
従って Sites, Lemma \ref{sites-lemma-equivalence} の仮定 (1) -- (5) を
検証しなければならない。包含関手を
$u : (\textit{Aff}/S)_h \to (\Sch/S)_h$ と書く。
$T$ がアフィンなら $T/S$ の任意の h 被覆は、例えば Lemma
\ref{flat-lemma-approximate-h-cover} により標準 h 被覆によって細分できるので、
余連続性が従う。従って (1) が成り立つ。標準 h 被覆は h 被覆なので、
$u$ が連続であることも直ちに分かる。従って (2) が成り立つ。
(3) と (4) は $u$ が充満忠実であることから直ちに従う。最後に、すべてのスキームが
アフィン開被覆（これは h 被覆である）をもつことから条件 (5) が従う。
\end{proof}

\begin{lemma}
\label{flat-lemma-characterize-sheaf-h}
$\mathcal{F}$ を $(\Sch/S)_h$ 上の前層とする。このとき $\mathcal{F}$ が層であるための
必要十分条件は次の二条件が成り立つことである。
\begin{enumerate}
\item $\mathcal{F}$ は Zariski 被覆に対する層条件を満たす。
\item $f : V \to U$ が固有、全射、かつ有限表示なら、$\mathcal{F}(U)$ は二つの写像
$\mathcal{F}(V) \to \mathcal{F}(V \times_U V)$ の等化子へ全単射に写る。
\end{enumerate}
さらに (1) のもとで、性質 (2) は次の性質と同値である。
\begin{enumerate}
\item[(2')] $U$ がアフィンである (2) のような $\{V \to U\}$ に対する層条件。
\end{enumerate}
\end{lemma}

\begin{proof}
(1) と (2) が成り立つなら $\mathcal{F}$ が層であることを示す。
$\{T_i \to T\}$ を $(\Sch/S)_h$ の被覆とし、この被覆に対する層条件を検証する。
$s_i \in \mathcal{F}(T_i)$ を、$T_i \times_T T_{i'}$ 上で同じ切断に制限される
切断とする。各 $T_i$ 上で $s_i$ に制限される一意な切断
$s \in \mathcal{F}(T)$ が存在することを示す。
$T = \bigcup U_j$ をアフィン開被覆とする。性質 (1) により、$s$ を作るには
$U_j \cap U_{j'}$ 上で一致する切断 $s_j \in \mathcal{F}(U_j)$ を作れば十分である。
被覆 $\{T_i \times_T U_j \to U_j\}$ を考える。このとき
$s_{ji} = s_i|_{T_i \times_T U_j}$ は
$(T_i \times_T U_j) \times_{U_j} (T_{i'} \times_T U_j)$ 上で一致する切断である。
有限表示の固有全射 $V_j \to U_j$ と有限アフィン開被覆
$V_j = \bigcup V_{jk}$ で、$\{V_{jk} \to U_j\}$ が
$\{T_i \times_T U_j \to U_j\}$ を細分するものを選ぶ。
Lemma \ref{flat-lemma-approximate-h-cover} を参照せよ。
$s_{jk} \in \mathcal{F}(V_{jk})$ を、暗黙の射による $s_{ji}$ の $V_{jk}$ への
引き戻しとすれば、$s_{jk}$ は切断 $s'_j \in \mathcal{F}(V_j)$ に貼り合わさる。
重なり上での一致をもう一度用いると、$s'_j$ は二つの写像
$\mathcal{F}(V_j) \to \mathcal{F}(V_j \times_{U_j} V_j)$ の等化子に属する。
従って (2) により、$s'_j$ は一意な切断 $s_j \in \mathcal{F}(U_j)$ から来る。
これらの切断 $s_j$ が必要な性質をすべてもつことの検証は省略する。

\medskip\noindent
(1) のもとでの (2) と (2') の同値性を証明する。
$V \to U$ を、固有、全射、かつ有限表示である $(\Sch/S)_h$ の射とする。
アフィン開被覆 $U = \bigcup U_i$ を選び、$V_i = V \times_U U_i$ とおく。
(2') により $\mathcal{F}(U_i) \to \mathcal{F}(V_i)$ は単射であり、(1) により
$\mathcal{F}(U) \to \prod \mathcal{F}(U_i)$ は単射なので、
$\mathcal{F}(U) \to \mathcal{F}(V)$ は単射である。
最後に、二つの写像 $\mathcal{F}(V) \to \mathcal{F}(V \times_U V)$ の等化子に属する
$t \in \mathcal{F}(V)$ が与えられたとする。このとき各 $i$ について
$t|_{V_i}$ は二つの写像
$\mathcal{F}(V_i) \to \mathcal{F}(V_i \times_{U_i} V_i)$ の等化子に属する。
従って (2') により、各 $i$ について $t|_{V_i}$ に写る一意な切断
$s_i \in \mathcal{F}(U_i)$ を得る。すべての $i, j$ について
$s_i|_{U_i \cap U_j} = s_j|_{U_i \cap U_j}$ であることの検証は省略する。
これは今示した一意性を用いる。被覆 $U = \bigcup U_i$ に対する層条件により、
切断 $s \in \mathcal{F}(U)$ を得る。$s$ が $\mathcal{F}(V)$ において $t$ に写ることの
証明は省略する。
\end{proof}

\noindent
次に、これらのサイトに付随するトポスの間のいくつかの関係を確立する。

\begin{lemma}
\label{flat-lemma-morphism-big-h}
$\Sch_h$ を大 h サイトとし、$f : T \to S$ を $\Sch_h$ の射とする。関手
$$
u : (\Sch/T)_h \longrightarrow (\Sch/S)_h,
\quad
V/T \longmapsto V/S
$$
は余連続であり、連続な右随伴
$$
v : (\Sch/S)_h \longrightarrow (\Sch/T)_h,
\quad
(U \to S) \longmapsto (U \times_S T \to T).
$$
をもつ。これらは同じトポスの射
$$
f_{big} :
\Sh((\Sch/T)_h)
\longrightarrow
\Sh((\Sch/S)_h)
$$
を誘導する。さらに
$f_{big}^{-1}(\mathcal{G})(U/T) = \mathcal{G}(U/S)$ および
$f_{big, *}(\mathcal{F})(U/S) = \mathcal{F}(U \times_S T/T)$ である。
また $f_{big}^{-1}$ は、ファイバー積および等化子と可換な左随伴 $f_{big!}$ をもつ。
\end{lemma}

\begin{proof}
関手 $u$ は余連続かつ連続であり、ファイバー積および等化子と可換である。
従って Sites, Lemmas \ref{sites-lemma-when-shriek} and
\ref{sites-lemma-preserve-equalizers} を適用でき、$f_{big}^{-1}$ の公式と
$f_{big!}$ の存在が従う。また $U/T$ と $V/S$ が与えられたとき
$\Mor_S(u(U), V) = \Mor_T(U, V \times_S T)$ となるので、関手 $v$ は右随伴である。
従って Sites, Lemmas \ref{sites-lemma-have-functor-other-way} and
\ref{sites-lemma-have-functor-other-way-morphism} を適用して
$f_{big, *}$ の公式を得る。
\end{proof}

\begin{lemma}
\label{flat-lemma-composition-h}
% P05-EN-CH38-0274: frozen-source pointer; see ERRATA_CANDIDATES.jsonl.
$(\Sch/S)_h$ のスキーム $X$, $Y$, $Y$ と射 $f : X \to Y$, $g : Y \to Z$ が
与えられたとき、
$g_{big} \circ f_{big} = (g \circ f)_{big}$.
\end{lemma}

\begin{proof}
これは Lemma \ref{flat-lemma-morphism-big-h} における大サイト上の関手の前進像と
引き戻しの簡明な記述から従う。
\end{proof}










\section{h 位相についてさらに}
\label{flat-section-h-more}

\noindent
この節では h 位相についてさらにいくつかの結果を証明する。
まず非例をいくつか挙げる。

\begin{example}
\label{flat-example-structure-sheaf-not-h-sheaf}
「構造層」$\mathcal{O}$ は h 位相における層ではない。例えば有限表示の全射閉埋め込み
$X \to Y$ を考える。このとき、例えば Lemma
\ref{flat-lemma-surjective-proper-finite-presentation-h} により $\{X \to Y\}$ は h 被覆である。
さらに $X \times_Y X = X$ であることに注意する。従って $\mathcal{O}$ が h 位相における
層なら、$\mathcal{O}_Y(Y) \to \mathcal{O}_X(X)$ は全単射となるはずである。
$X$, $Y$ がアフィンで射 $X \to Y$ が同型でなければ、これは成り立たない。
\end{example}

\begin{example}
\label{flat-example-representable-sheaf-not-h-sheaf}
任意のサイト $(\Sch/S)_h$ 上の位相は準標準ではない。言い換えれば、表現可能層は
% P05-EN-CH38-0276: frozen-source pointer; see ERRATA_CANDIDATES.jsonl.
層ではない。実際、$(\Sch/S)_h$ において
$\mathcal{O}(X) = \Mor_S(X, \mathbf{A}^1_S)$ なので「構造層」$\mathcal{O}$ は表現可能であり、
一方 Example \ref{flat-example-structure-sheaf-not-h-sheaf} で $\mathcal{O}$ が
層ではないことを見た。
\end{example}

\begin{lemma}
\label{flat-lemma-limit-h-topology}
$T$ を、アフィンスキームの有向逆系の極限
$T = \lim_{i \in I} T_i$ と書かれるアフィンスキームとする。
\begin{enumerate}
\item $\mathcal{V} = \{V_j \to T\}_{j = 1, \ldots, m}$ を $T$ の標準 h 被覆とする。
Definition \ref{flat-definition-standard-h} を参照せよ。このとき添字 $i$ と標準 h 被覆
$\mathcal{V}_i = \{V_{i, j} \to T_i\}_{j = 1, \ldots, m}$ で、
その $T$ への基底変換 $T \times_{T_i} \mathcal{V}_i$ が $\mathcal{V}$ と
同型になるものが存在する。
\item $\mathcal{V}_i$, $\mathcal{V}'_i$ を $T_i$ の二つの標準 h 被覆とする。
$f : T \times_{T_i} \mathcal{V}_i \to T \times_{T_i} \mathcal{V}'_i$ が
$T$ の被覆の射なら、添字 $i' \geq i$ と射
% P05-EN-CH38-0277: frozen-source pointer; see ERRATA_CANDIDATES.jsonl.
$f_{i'} : T_{i'} \times_{T_i} \mathcal{V} \to
T_{i'} \times_{T_i} \mathcal{V}'_i$
で、その $T$ への基底変換が $f$ となるものが存在する。
\item
% P05-EN-CH38-0278: frozen-source pointer; see ERRATA_CANDIDATES.jsonl.
$f, g : \mathcal{V} \to \mathcal{V}'_i$
が $T_i$ の標準 h 被覆の射であり、その $T$ への基底変換 $f_T, g_T$ が等しいなら、
$f_{T_{i'}} = g_{T_{i'}}$ となる添字 $i' \geq i$ が存在する。
\end{enumerate}
言い換えれば、$T$ の標準 h 被覆の圏は、$T_i$ の標準 h 被覆の圏の $I$ 上の余極限である。
\end{lemma}

\begin{proof}
Limits, Lemma \ref{limits-lemma-descend-finite-presentation} により、$T$ 上有限表示の
スキームの圏は、$T_i$ 上有限表示の圏の $I$ 上の余極限である。Limits, Lemma
\ref{limits-lemma-descend-affine-finite-presentation} により、$T$ 上アフィンかつ
有限表示のスキームの圏についても同じことが成り立つ。
補題の証明を終えるには、$\{V_{j, i} \to T_i\}_{j = 1, \ldots, m}$ が
$V_{j, i}$ をアフィンとする有限表示射の有限族で、その基底変換族
$\{T \times_{T_i} V_{j, i} \to T\}$ が h 被覆なら、ある $i' \geq i$ に対して族
$\{T_{i'} \times_{T_i} V_{j, i} \to T_{i'}\}$ が h 被覆であることを示せば十分である。
このため Lemma \ref{flat-lemma-approximate-h-cover} を用い、有限表示の固有全射
$Y \to T$ と有限アフィン開被覆 $Y = \bigcup_{k = 1, \ldots, n} Y_k$ で、
$\{Y_k \to T\}_{k = 1, \ldots, n}$ が
$\{T \times_{T_i} V_{j, i} \to T\}$ を細分するものを選ぶ。
上の議論と Limits, Lemmas \ref{limits-lemma-eventually-proper},
\ref{limits-lemma-descend-surjective}, and
\ref{limits-lemma-descend-opens} を用いれば、$i' \geq i$、有限表示の固有全射
$Y_{i'} \to T_{i'}$、およびアフィン開被覆
$Y_{i'} = \bigcup_{k = 1, \ldots, n} Y_{i', k}$ で、さらに
% P05-EN-CH38-0280: frozen-source pointer; see ERRATA_CANDIDATES.jsonl.
$\{Y_{i', k} \to Y_{i'}\}$ が
$\{T_{i'} \times_{T_i} V_{j, i} \to T_{i'}\}$ を細分するものをとれる。
従って最後の族は h 被覆であり、証明が完了する。
\end{proof}

\begin{lemma}
\label{flat-lemma-extend-sheaf-h}
$S$ を大サイト $\Sch_h$ に含まれるスキームとする。
$F : (\Sch/S)_h^{opp} \to \textit{Sets}$ を、$\mathcal{C} = (\Sch/S)_h$ とした
Topologies, Lemma \ref{topologies-lemma-extend} の性質 (b) を満たす h 層とする。
このとき $F$ の $S$ 上のすべてのスキームの圏への拡張 $F'$ は、すべての h 被覆に
対する層条件を満たし、極限を保つ（Limits, Remark
\ref{limits-remark-limit-preserving}）。
\end{lemma}

\begin{proof}
これは Lemmas \ref{flat-lemma-limit-h-topology} and \ref{flat-lemma-h-induced} を用い、
Topologies, Lemmas \ref{topologies-lemma-extend-sheaf-general} and
\ref{topologies-lemma-extend-sheaf} の証明にある議論によって証明される。
詳細は省略する。
\end{proof}








\section{吹き上げ正方形と ph 位相}
\label{flat-section-blow-up-ph}

\noindent
$X$ をスキームとする。$Z \subset X$ を、包含射が有限表示である閉部分スキーム、
すなわち $Z$ に対応するイデアルの準連接層が有限型であるものとする。
$b : X' \to X$ を $Z$ における $X$ の吹き上げとし、$E = b^{-1}(Z)$ を例外因子とする。
Divisors, Section \ref{divisors-section-blowing-up} を参照せよ。この状況で、この節では
\begin{equation}
\label{flat-equation-blow-up-square}
\vcenter{
\xymatrix{
E \ar[d] \ar[r] & X' \ar[d]^b \\
Z \ar[r] & X
}
}
\end{equation}
を {\it 吹き上げ正方形}と呼ぶことにする。

\begin{lemma}
\label{flat-lemma-blow-up-square-ph}
$\mathcal{F}$ をサイト $(\Sch/S)_{ph}$ 上の層とする。Topologies, Definition
\ref{topologies-definition-big-small-ph} を参照せよ。このとき圏 $(\Sch/S)_{ph}$ の
任意の吹き上げ正方形 (\ref{flat-equation-blow-up-square}) に対し、図式
$$
\xymatrix{
\mathcal{F}(E) & \mathcal{F}(X') \ar[l] \\
\mathcal{F}(Z) \ar[u] & \mathcal{F}(X) \ar[u] \ar[l]
}
$$
は集合の圏において Cartesian である。
\end{lemma}

\begin{proof}
$Z \amalg X' \to X$ は固有全射なので、$\{Z \amalg X' \to X\}$ は ph 被覆である
（Topologies, Lemma \ref{topologies-lemma-surjective-proper-ph}）。また
$$
(Z \amalg X') \times_X (Z \amalg X') =
Z \amalg E \amalg E \amalg X' \times_X X'
$$
$\mathcal{F}$ は Zariski 層なので、$\mathcal{F}$ は直和を積へ写す。
従って被覆 $\{Z \amalg X' \to X\}$ に対する層条件によれば、
$\mathcal{F}(X) \to \mathcal{F}(Z) \times \mathcal{F}(X')$ は単射であり、その像は、
(a) $t|_E = s'|_E$ かつ (b) $s'$ が二つの写像
$\mathcal{F}(X') \to \mathcal{F}(X' \times_X X')$ の等化子に属するような組
$(t, s')$ の集合である。次に、自然な射
$$
E \times_Z E \amalg X' \longrightarrow X' \times_X X'
$$
は固有全射であることに注意する。実際 $b$ は同型
$X' \setminus E \to X \setminus Z$ を誘導する。従って
$\mathcal{F}(X' \times_X X') \to
\mathcal{F}(E \times_Z E) \times \mathcal{F}(X')$ は単射である。
これより (a) $\Rightarrow$ (b) が従い、補題が成り立つ。
\end{proof}

\begin{lemma}
\label{flat-lemma-thickening-ph}
$\mathcal{F}$ を Topologies, Definition \ref{topologies-definition-big-small-ph} の
サイト $(\Sch/S)_{ph}$ 上の層とする。$X \to X'$ を、厚化である
$(\Sch/S)_{ph}$ の射とする。このとき
$\mathcal{F}(X') \to \mathcal{F}(X)$ は全単射である。
\end{lemma}

\begin{proof}
$X \to X'$ は固有全射であり、$X \times_{X'} X = X$ であることに注意する。
ph 被覆 $\{X \to X'\}$ に対する層条件
（Topologies, Lemma \ref{topologies-lemma-surjective-proper-ph}）から結論が従う。
\end{proof}








\section{ほとんど吹き上げ正方形と h 位相}
\label{flat-section-blow-up-h}

\noindent
吹き上げ正方形 (\ref{flat-equation-blow-up-square}) を考える。射 $b : X' \to X$ は射影的だが
（Divisors, Lemma \ref{divisors-lemma-blowing-up-projective}）、一般には $b$ が有限表示で
あることを保証する簡明な方法はない。h 被覆は有限表示射を用いて構成されるので、
変形版が必要である。そこで、スキームの可換図式
\begin{equation}
\label{flat-equation-almost-blow-up-square}
\vcenter{
\xymatrix{
E \ar[d] \ar[r] & X' \ar[d]^b \\
Z \ar[r] & X
}
}
\end{equation}
が次の条件を満たすとき、これを {\it ほとんど吹き上げ正方形}と呼ぶことにする。
\begin{enumerate}
\item $Z \to X$ は有限表示の閉埋め込みである。
\item $E = b^{-1}(Z)$ は $X'$ の局所主閉部分スキームである。
\item $b$ は固有かつ有限表示である。
\item $E$ に台をもつ $\mathcal{O}_{X'}$ の切断からなる準連接イデアル
（Properties, Lemma \ref{properties-lemma-sections-supported-on-closed-subset}）で
切り出される閉部分スキーム $X'' \subset X'$ は、$Z$ における $X$ の吹き上げである。
\end{enumerate}
従って射 $b$ は同型 $X' \setminus E \to X \setminus Z$ を誘導する。
ほとんど吹き上げ正方形のごく簡単な例については、Examples
\ref{flat-example-one-generator} and \ref{flat-example-two-generators} を参照せよ。

\medskip\noindent
吹き上げの基底変換は通常は吹き上げではないが、ほとんど吹き上げは基底変換と両立する。

\begin{lemma}
\label{flat-lemma-base-change-almost-blow-up}
ほとんど吹き上げ正方形 (\ref{flat-equation-almost-blow-up-square}) を考える。
$Y \to X$ を任意の射とする。このとき基底変換
$$
\xymatrix{
Y \times_X E \ar[d] \ar[r] & Y \times_X X' \ar[d] \\
Y \times_X Z \ar[r] & Y
}
$$
もほとんど吹き上げ正方形である。
\end{lemma}

\begin{proof}
Morphisms, Lemmas \ref{morphisms-lemma-base-change-proper} and
\ref{morphisms-lemma-base-change-finite-presentation} により、射
$Y \times_X X' \to Y$ は固有かつ有限表示である。射 $Y \times_X Z \to Y$ は有限表示の
閉埋め込みである（Morphisms, Lemma
\ref{morphisms-lemma-base-change-closed-immersion}）。$Y \times_X X'$ における
$Y \times_X Z$ の逆像は、$Y \times_X X'$ における $E$ の逆像に等しく、従って局所主である
（Divisors, Lemma \ref{divisors-lemma-pullback-locally-principal}）。
$X'' \subset X'$、それぞれ $Y'' \subset Y \times_X X'$ を、$E$、それぞれ
$Y \times_X E$ に台をもつ $\mathcal{O}_{X'}$、それぞれ
% P05-EN-CH38-0282: frozen-source pointer; see ERRATA_CANDIDATES.jsonl.
$\mathcal{O}_{Y \times_Y X'}$ の切断からなる準連接イデアルに対応する閉部分スキームとする。
明らかに $Y'' \subset Y \times_X X''$ は、$Y \times_X (E \cap X'')$ に台をもつ
% P05-EN-CH38-0283: frozen-source pointer; see ERRATA_CANDIDATES.jsonl.
$\mathcal{O}_{Y \times_Y X''}$ の切断からなる準連接イデアルに対応する閉部分スキームである。
従って $Y''$ は吹き上げ $X'' \to X$ に関する $Y$ の厳密変換である。Divisors, Definition
\ref{divisors-definition-strict-transform} を参照せよ。従って Divisors, Lemma
\ref{divisors-lemma-strict-transform} により、$Y''$ は $Y$ 上の
% P05-EN-CH38-0284: frozen-source pointer; see ERRATA_CANDIDATES.jsonl.
$Y \times_X Z$ の吹き上げである。
\end{proof}

\noindent
ほとんど吹き上げ正方形は縮小できる。

\begin{lemma}
\label{flat-lemma-shrink-almost-blow-up}
ほとんど吹き上げ正方形 (\ref{flat-equation-almost-blow-up-square}) を考える。
$W \to X'$ を有限表示の閉埋め込みとする。このとき次は同値である。
\begin{enumerate}
\item $X' \setminus E$ はスキーム論的に $W$ に含まれる。
\item $Z$ における $X$ の吹き上げ $X''$ はスキーム論的に $W$ に含まれる。
\item 図式
$$
\xymatrix{
E \cap W \ar[d] \ar[r] & W \ar[d] \\
Z \ar[r] & X
}
$$
はほとんど吹き上げ正方形である。ここで $E \cap W$ はスキーム論的交叉である。
\end{enumerate}
\end{lemma}

\begin{proof}
(1) を仮定する。このとき全射 $\mathcal{O}_{X'} \to \mathcal{O}_W$ は開部分
% P05-EN-CH38-0285: frozen-source pointer; see ERRATA_CANDIDATES.jsonl.
$X' \subset E$ 上で同型である。$X'' \subset X'$ のイデアル層は $E$ に台をもつ
$\mathcal{O}_{X'}$ の切断からなるので（ほとんど吹き上げ正方形の定義による）、
(2) が成り立つ。(2) が成り立てば (3) が成り立つ。(3) が成り立てば、
$X'' \cap (X' \setminus E)$ は $X \setminus Z$ と同型であり、後者はさらに
$X' \setminus E$ と同型なので、(1) が成り立つ。
\end{proof}

\noindent
実際の吹き上げは、任意に与えたほとんど吹き上げの縮小の極限である。

\begin{lemma}
\label{flat-lemma-blow-up-limit-almost-blow-up}
$X$ が準コンパクトかつ準分離であるほとんど吹き上げ正方形
(\ref{flat-equation-almost-blow-up-square}) を考える。このとき $Z$ における $X$ の
吹き上げ $X''$ は
$$
X'' = \lim X'_i
$$
と書ける。ここで極限は、Lemma \ref{flat-lemma-shrink-almost-blow-up} の同値な条件を
満たす有限表示の閉部分スキーム $X'_i \subset X'$ からなる有向系上でとる。
\end{lemma}

\begin{proof}
$\mathcal{I} \subset \mathcal{O}_{X'}$ を $X''$ に対応する準連接イデアル層とする。
Properties, Lemma \ref{properties-lemma-quasi-coherent-colimit-finite-type} により、
$\mathcal{I}$ をその有限型準連接部分加群のフィルター余極限
$\mathcal{I} = \colim \mathcal{I}_i$ と書ける。これらの加群は閉部分スキーム
$X'_i$ と $1$-to-$1$ に対応するので、証明が完了する。
\end{proof}

\noindent
ほとんど吹き上げ正方形は存在する。

\begin{lemma}
\label{flat-lemma-almost-blow-up-square}
$X$ を準コンパクトかつ準分離なスキームとする。$Z \subset X$ を、有限型の
準連接イデアル層で切り出される閉部分スキームとする。このとき
(\ref{flat-equation-almost-blow-up-square}) のようなほとんど吹き上げ正方形が存在する。
\end{lemma}

\begin{proof}
Limits, Proposition \ref{limits-proposition-approximate} により、$X = \lim X_i$ を
アフィン遷移射をもつ Noetherian スキームの逆系の有向極限と書いてよい。
ある添字 $i$ と閉埋め込み $Z_i \to X_i$ で、その $X$ への基底変換が閉埋め込み
$Z \to X$ となるものをとれる。Limits, Lemmas
\ref{limits-lemma-descend-finite-presentation} and
\ref{limits-lemma-descend-closed-immersion-finite-presentation} を参照せよ。
$b_i : X'_i \to X_i$ を中心 $Z_i$ の吹き上げとする。これにより吹き上げ正方形
$$
\xymatrix{
E_i \ar[r] \ar[d] & X'_i \ar[d]^{b_i} \\
Z_i \ar[r] & X_i
}
$$
を得る。ここですべての射は Noetherian スキームの有限型射であり、従って有限表示である。
よってこれはほとんど吹き上げ正方形である。Lemma
\ref{flat-lemma-base-change-almost-blow-up} により、この図式の $X$ への基底変換が求める
ほとんど吹き上げ正方形を与える。
\end{proof}

\noindent
ほとんど吹き上げ正方形は、Lemma \ref{flat-lemma-shrink-almost-blow-up} のような
縮小を除いて一意である。

\begin{lemma}
\label{flat-lemma-almost-blow-up-unique}
$X$ を準コンパクトかつ準分離なスキームとし、$Z \subset X$ を有限型準連接イデアル層で
切り出される閉部分スキームとする。$k = 1, 2$ に対して、ほとんど吹き上げ正方形
(\ref{flat-equation-almost-blow-up-square})
$$
\xymatrix{
E_k \ar[r] \ar[d] & X_k' \ar[d] \\
Z \ar[r] & X
}
$$
が与えられたとする。このとき、ほとんど吹き上げ正方形
$$
\xymatrix{
E \ar[r] \ar[d] & X' \ar[d] \\
Z \ar[r] & X
}
$$
と、$E = i_k^{-1}(E_k)$ を満たす $X$ 上の閉埋め込み
$i_k : X' \to X'_k$ が存在する。
\end{lemma}

\begin{proof}
% P05-EN-CH38-0286: frozen-source pointer; see ERRATA_CANDIDATES.jsonl.
$X'' \to X$ を $X$ における $Z$ の吹き上げと書く。$X''$ を $X'_1$ と $X'_2$ の
双方の閉部分スキームとみなす。Lemma
\ref{flat-lemma-blow-up-limit-almost-blow-up} のように $X'' = \lim X'_{1, i}$ と書く。
Limits, Proposition \ref{limits-proposition-characterize-locally-finite-presentation} により、
ある $i$ と射 $h : X'_{1, i} \to X'_2$ で、包含
$X'' \subset X'_{1, i}$ および $X'' \subset X'_2$ と整合するものが存在する。
Limits, Lemma \ref{limits-lemma-eventually-closed-immersion} により、ある
$i' \geq i$ に対して $h$ の $X'_{1, i'}$ への制限は閉埋め込みである。
これで証明が完了した。
\end{proof}

\noindent
好都合な状況では、吹き上げに対する我々の平坦化手法はほとんど吹き上げにも継承される。

\begin{lemma}
\label{flat-lemma-flat-after-almost-blowing-up}
$Y$ を準コンパクトかつ準分離なスキームとする。$X$ を $Y$ 上有限表示のスキームとする。
$V \subset Y$ を、$X_V \to V$ が平坦となる準コンパクト開部分とする。
このとき可換図式
$$
\xymatrix{
E \ar[ddd] \ar[rd] & & & D \ar[lll] \ar[ddd] \ar[ld] \\
& Y' \ar[d] & X' \ar[l] \ar[d] \\
& Y & X \ar[l] \\
Z \ar[ru] & & & T \ar[lll] \ar[lu]
}
$$
で、右および左の正方形がほとんど吹き上げ正方形、下および上の正方形が
Cartesian であり、$Z \cap V = \emptyset$、かつ $X' \to Y'$ が平坦
（そして有限表示）となるものが存在する。
\end{lemma}

\begin{proof}
$Y$ が Noetherian スキームなら、この場合には吹き上げ正方形がほとんど吹き上げ正方形
なので（厳密変換が吹き上げであることも用いる）、補題は Lemma
\ref{flat-lemma-flat-after-blowing-up} から直ちに従う。一般の場合は絶対 Noetherian 近似により
Noetherian の場合へ帰着する。

\medskip\noindent
Limits, Proposition \ref{limits-proposition-approximate} により、$Y = \lim Y_i$ を
アフィン遷移射をもつ Noetherian スキームの逆系の有向極限と書いてよい。
ある添字 $i$ と有限表示射 $X_i \to Y_i$ で、その $Y$ への基底変換が $X \to Y$ と
なるものをとれる。Limits, Lemma \ref{limits-lemma-descend-finite-presentation} を参照せよ。
$i$ を大きくして、$V$ が開部分スキーム $V_i \subset Y_i$ の逆像であると仮定してよい。
Limits, Lemma \ref{limits-lemma-descend-opens} を参照せよ。最後に、さらに $i$ を大きくして、
$X_{i, V_i} \to V_i$ が平坦であると仮定してよい。Limits, Lemma
\ref{limits-lemma-descend-flat-finite-presentation} を参照せよ。
Noetherian の場合により、$X_i \to Y_i \supset V_i$ に対して補題のような図式を構成できる。
この図式を $Y \to Y_i$ で基底変換すれば解が得られる。基底変換が射の性質を保つこと、
すなわち Morphisms, Lemmas
\ref{morphisms-lemma-base-change-proper},
\ref{morphisms-lemma-base-change-finite-presentation},
\ref{morphisms-lemma-base-change-closed-immersion}, and
\ref{morphisms-lemma-base-change-flat}
を参照せよ。また、ほとんど吹き上げ正方形の基底変換がほとんど吹き上げ正方形であること、
すなわち Lemma \ref{flat-lemma-base-change-almost-blow-up} を用いる。
\end{proof}

\begin{lemma}
\label{flat-lemma-blow-up-square-h}
$\mathcal{F}$ を Definition \ref{flat-definition-big-small-h} で構成されたサイト
$(\Sch/S)_h$ の一つの上の層とする。このとき圏 $(\Sch/S)_h$ の任意の
ほとんど吹き上げ正方形 (\ref{flat-equation-almost-blow-up-square}) に対して図式
$$
\xymatrix{
\mathcal{F}(E) & \mathcal{F}(X') \ar[l] \\
\mathcal{F}(Z) \ar[u] & \mathcal{F}(X) \ar[u] \ar[l]
}
$$
は集合の圏において Cartesian である。
\end{lemma}

\begin{proof}
$Z \amalg X' \to X$ は有限表示の固有全射なので、$\{Z \amalg X' \to X\}$ は h 被覆である
（Lemma \ref{flat-lemma-surjective-proper-finite-presentation-h}）。また
$$
(Z \amalg X') \times_X (Z \amalg X') =
Z \amalg E \amalg E \amalg X' \times_X X'
$$
$\mathcal{F}$ は Zariski 層なので、$\mathcal{F}$ は直和を積へ写す。
従って被覆 $\{Z \amalg X' \to X\}$ に対する層条件によれば、
$\mathcal{F}(X) \to \mathcal{F}(Z) \times \mathcal{F}(X')$ は単射であり、その像は、
(a) $t|_E = s'|_E$ かつ (b) $s'$ が二つの写像
$\mathcal{F}(X') \to \mathcal{F}(X' \times_X X')$ の等化子に属するような組
$(t, s')$ の集合である。次に、自然な射
$$
E \times_Z E \amalg X' \longrightarrow X' \times_X X'
$$
は有限表示の固有全射である。実際 $b$ は同型
$X' \setminus E \to X \setminus Z$ を誘導する。従って
$\mathcal{F}(X' \times_X X') \to
\mathcal{F}(E \times_Z E) \times \mathcal{F}(X')$ は単射である。
これより (a) $\Rightarrow$ (b) が従い、補題が成り立つ。
\end{proof}

\begin{lemma}
\label{flat-lemma-thickening-h}
$\mathcal{F}$ を Definition \ref{flat-definition-big-small-h} で構成されたサイト
$(\Sch/S)_h$ の一つの上の層とする。$X \to X'$ を、厚化かつ有限表示である
$(\Sch/S)_h$ の射とする。このとき $\mathcal{F}(X') \to \mathcal{F}(X)$ は全単射である。
\end{lemma}

\begin{proof}
第一の証明。$X \to X'$ は有限表示の固有全射であり、
$X \times_{X'} X = X$ であることに注意する。h 被覆 $\{X \to X'\}$ に対する層条件
（Lemma \ref{flat-lemma-surjective-proper-finite-presentation-h}）から結論が従う。

\medskip\noindent
第二の証明（馬鹿げた証明）。$X$ における $X'$ の吹き上げは空スキームである。
これは、$a$ が $A$ の冪零元ならアフィン吹き上げ代数 $A[\frac{I}{a}]$
（Algebra, Section \ref{algebra-section-blow-up}）が零だからである。詳細は省略する。
従って次の形のほとんど吹き上げ正方形を得る。
$$
\xymatrix{
\emptyset \ar[r] \ar[d] & \emptyset \ar[d] \\
X \ar[r] & X'
}
$$
$\mathcal{F}$ は層なので $\mathcal{F}(\emptyset)$ は一点集合である。
Lemma \ref{flat-lemma-blow-up-square-h} を適用すれば結論を得る。
\end{proof}

\begin{proposition}
\label{flat-proposition-check-h}
$\mathcal{F}$ を Definition \ref{flat-definition-big-small-h} で構成されたサイト
$(\Sch/S)_h$ の一つの上の前層とする。このとき $\mathcal{F}$ が層であるための
必要十分条件は、次の条件が満たされることである。
\begin{enumerate}
\item $\mathcal{F}$ は Zariski 位相に対する層である。
\item $Y$ がアフィンで、$f$ が全射、平坦、固有、かつ有限表示である
$(\Sch/S)_h$ の射 $f : X \to Y$ が与えられたとき、$\mathcal{F}(Y)$ は二つの写像
$\mathcal{F}(X) \to \mathcal{F}(X \times_Y X)$ の等化子である。
\item 圏 $(\Sch/S)_h$ において $X$ がアフィンであるほとんど吹き上げ正方形
(\ref{flat-equation-almost-blow-up-square}) が与えられたとき、図式
$$
\xymatrix{
\mathcal{F}(E) & \mathcal{F}(X') \ar[l] \\
\mathcal{F}(Z) \ar[u] & \mathcal{F}(X) \ar[u] \ar[l]
}
$$
は集合の圏において Cartesian である。
\end{enumerate}
\end{proposition}

\begin{proof}
$\mathcal{F}$ が層であると仮定する。Zariski 被覆は h 被覆なので条件 (1) が成り立つ。
Lemma \ref{flat-lemma-zariski-h} を参照せよ。(2) のような $f$ に対して
$\{X \to Y\}$ は fppf 被覆であり（これは明らか）、従って h 被覆なので、条件 (2) が
成り立つ。Lemma \ref{flat-lemma-zariski-h} を参照せよ。条件 (3) は Lemma
\ref{flat-lemma-blow-up-square-h} により成り立つ。

\medskip\noindent
逆に $\mathcal{F}$ が (1)、(2)、(3) を満たすと仮定する。Lemma
\ref{flat-lemma-characterize-sheaf-h} を適用して $\mathcal{F}$ が層であることを証明する。
$Y$ をアフィンとする $(\Sch/S)_h$ の有限表示固有全射 $f : X \to Y$ を考える。
$\mathcal{F}(Y)$ が二つの写像
$\mathcal{F}(X) \to \mathcal{F}(X \times_Y X)$ の等化子であることを示せば十分である。

\medskip\noindent
まず $f : X \to Y$ がさらに閉埋め込みであると仮定する（言い換えれば $f$ は厚化である）。
このとき $X$ における $Y$ の吹き上げは空スキームであり、頂点が
$\emptyset, \emptyset, X, Y$ からなるほとんど吹き上げ正方形を与える
（Lemma \ref{flat-lemma-thickening-h} の第二の証明と比較せよ）。従って条件 (3) により
$$
\xymatrix{
\mathcal{F}(\emptyset) & \mathcal{F}(\emptyset) \ar[l] \\
\mathcal{F}(X) \ar[u] & \mathcal{F}(Y) \ar[u] \ar[l]
}
$$
は集合の圏において Cartesian である。$\mathcal{F}$ は Zariski 位相に対する層なので、
$\mathcal{F}(\emptyset)$ は一点集合である。従って
$\mathcal{F}(X) = \mathcal{F}(Y)$ である。

\medskip\noindent
幕間 A：$T \to T'$ を、厚化かつ有限表示である $(\Sch/S)_h$ の射とする。
このとき $\mathcal{F}(T') \to \mathcal{F}(T)$ は全単射である。実際、アフィン開被覆
$T' = \bigcup T'_i$ を選び、$T_i = T \times_{T'} T'_i$ を対応する $T$ のアフィン開部分とする。
前段落の結果により、すべての $i$ に対して
$\mathcal{F}(T'_i) \to \mathcal{F}(T_i)$ は全単射である。Zariski 層の性質を用いると、
$\mathcal{F}(T') \to \mathcal{F}(T)$ が単射であることが分かる。議論を繰り返せば
全単射であることが分かる。細部は省略する。

\medskip\noindent
幕間 B：圏 $(\Sch/S)_h$ のほとんど吹き上げ正方形
(\ref{flat-equation-almost-blow-up-square}) を考える。このとき図式
$$
\xymatrix{
\mathcal{F}(E) & \mathcal{F}(X') \ar[l] \\
\mathcal{F}(Z) \ar[u] & \mathcal{F}(X) \ar[u] \ar[l]
}
$$
は集合の圏において Cartesian であると主張する。これは、$X$ のアフィン開被覆を選び、
幕間 A のように論じることで条件 (3) から従う。詳細は省略する。

\medskip\noindent
次に、$f : X \to Y$ を、$Y$ がアフィンである $(\Sch/S)_h$ の有限表示固有全射とする。
$f : X \to Y$ に対し、More on Morphisms, Lemma
\ref{more-morphisms-lemma-generic-flatness-stratification-scheme} のように
一般平坦性による層別化
$$
Y \supset Y_0 \supset Y_1 \supset \ldots \supset Y_t = \emptyset
$$
を選ぶ。以下ではこの層別化の性質を断りなくすべて用いる。
$X_0 = X \times_Y Y_0$ とおく。幕間 B により
$\mathcal{F}(Y_0) = \mathcal{F}(Y)$、$\mathcal{F}(X_0) = \mathcal{F}(X)$、および
$\mathcal{F}(X_0 \times_{Y_0} X_0) = \mathcal{F}(X \times_Y X)$ である。

\medskip\noindent
$t$ に関する帰納法で結果を証明する。$t = 1$ なら $X_0 \to Y_0$ は全射、固有、平坦、
かつ有限表示なので、性質 (2) により結果が成り立つ。$t > 1$ なら、$Y$ を $Y_0$ で、
$X$ を $X_0$ で置き換え（上を参照）、$Y = Y_0$ と仮定してよい。

\medskip\noindent
準コンパクト開部分スキーム $V = Y \setminus Y_1 = Y_0 \setminus Y_1$ を考える。
$f : X \to Y \supset V$ に対して Lemma
\ref{flat-lemma-flat-after-almost-blowing-up} のような図式
$$
\xymatrix{
E \ar[ddd] \ar[rd] & & & D \ar[lll] \ar[ddd] \ar[ld] \\
& Y' \ar[d] & X' \ar[l] \ar[d] \\
& Y & X \ar[l] \\
Z \ar[ru] & & & T \ar[lll] \ar[lu]
}
$$
を選ぶ。このとき $f' : X' \to Y'$ は平坦かつ有限表示である。また $f'$ は固有である
（これを見るには Morphisms, Lemmas \ref{morphisms-lemma-composition-proper} and
\ref{morphisms-lemma-image-proper-scheme-closed} を用いる）。従って像
$W = f'(X') \subset Y'$ は $Y'$ の開部分スキーム
（Morphisms, Lemma \ref{morphisms-lemma-fppf-open}）かつ閉部分スキームである。
$Y' \setminus E$ は $W$ に含まれることに注意する。Lemma
\ref{flat-lemma-shrink-almost-blow-up} により、上の図式で $Y'$ を $W$ に置き換えてよい。
言い換えれば、$f'$ が全射であると仮定してよく、そう仮定する。この時点で、幕間 B により
$$
\vcenter{
\xymatrix{
\mathcal{F}(E) & \mathcal{F}(Y') \ar[l] \\
\mathcal{F}(Z) \ar[u] & \mathcal{F}(Y) \ar[u] \ar[l]
}
}
\quad\text{and}\quad
\vcenter{
\xymatrix{
\mathcal{F}(D) & \mathcal{F}(X') \ar[l] \\
\mathcal{F}(T) \ar[u] & \mathcal{F}(X) \ar[u] \ar[l]
}
}
$$
はいずれも Cartesian である。$Z \cap Y_1 \to Z$ は有限表示の厚化であることに注意する
（$Z$ は $V$ と交わらない $Y$ の閉部分スキームとして、集合論的に $Y_1$ に含まれる）。
従って、$f$ の $T$ への制限 $T = Z \times_Y X \to Z$ に対し、上のようなフィルトレーション
$$
% P05-EN-CH38-0289: frozen-source pointer; see ERRATA_CANDIDATES.jsonl.
Z \supset Z \cap Y_1 \supset Z \cap Y_2 \subset \ldots \subset
Z \cap Y_t = \emptyset
$$
を得る。従って帰納法の仮定により、$\mathcal{F}(Z) \to \mathcal{F}(T)$ は集合の単射で、
その像は二つの写像 $\mathcal{F}(T) \to
\mathcal{F}(T \times_Z T)$ の等化子である。

\medskip\noindent
$s \in \mathcal{F}(X)$ を二つの写像
$\mathcal{F}(X) \to \mathcal{F}(X \times_Y X)$ の等化子の元とする。
上により、制限 $s|_T$ は一意な元 $t \in \mathcal{F}(Z)$ から来て、同様に制限
$s|_{X'}$ は一意な元 $t' \in \mathcal{F}(Y')$ から来る。上の巨大な可換図式の矢印に
対応する $\mathcal{F}$ の制限写像を用いて切断を追えば、$t$ と $t'$ は
$\mathcal{F}(E)$ の同じ元に制限されることが分かる。実際これらは
$\mathcal{F}(D)$ の同じ元に制限され、(2) が成り立つ。ここで $D \to E$ は
$X' \to Y'$ の制限として全射、平坦、固有、かつ有限表示であることを用いる。
従って上に表示した二つの Cartesian 正方形のうち第一のものにより、$Z$ と $Y'$ 上で
$t$ と $t'$ に制限される一意な切断 $u \in \mathcal{F}(Y)$ を得る。
$u$ が $X$ 上で $s$ に制限されることは第二の図式を用いれば分かる。
\end{proof}

\begin{example}
\label{flat-example-one-generator}
$A$ を環とし、$f \in A$ を元とする。$J \subset A$ を $f$ のある冪で零化される
有限生成イデアルとする。このとき
$$
\xymatrix{
E = \Spec(A/fA + J) \ar[r] \ar[d] & \Spec(A/J) = X' \ar[d] \\
Z = \Spec(A/fA) \ar[r] & \Spec(A) = X
}
$$
はほとんど吹き上げ正方形である。
\end{example}

\begin{example}
\label{flat-example-two-generators}
$A$ を環とし、$f_1, f_2 \in A$ を元とする。
$$
\xymatrix{
E = \text{Proj}(A/(f_1, f_2)[T_0, T_1]) \ar[r] \ar[d] &
% P05-EN-CH38-0291: frozen-source pointer; see ERRATA_CANDIDATES.jsonl.
\text{Proj}(A[T_0, T_1]/(f_2 T_0 - f_1 T_1) = X' \ar[d] \\
Z = \Spec(A/(f_1, f_2)) \ar[r] & \Spec(A) = X
}
$$
はほとんど吹き上げ正方形である。
\end{example}

\begin{lemma}
\label{flat-lemma-refine-check-h}
$\mathcal{F}$ を Definition \ref{flat-definition-big-small-h} で構成されたサイト
$(\Sch/S)_h$ の一つの上の前層とする。このとき $\mathcal{F}$ が層であるための
必要十分条件は、次の条件が満たされることである。
\begin{enumerate}
\item $\mathcal{F}$ は Zariski 位相に対する層である。
\item $Y$ がアフィンで、$f$ が全射、平坦、固有、かつ有限表示である
$(\Sch/S)_h$ の射 $f : X \to Y$ が与えられたとき、$\mathcal{F}(Y)$ は二つの写像
$\mathcal{F}(X) \to \mathcal{F}(X \times_Y X)$ の等化子である。
\item $\mathcal{F}$ は、圏 $(\Sch/S)_h$ の Example
\ref{flat-example-one-generator} のようなほとんど吹き上げ正方形を、集合の Cartesian
図式へ写す。
\item $\mathcal{F}$ は、圏 $(\Sch/S)_h$ の Example
\ref{flat-example-two-generators} のようなほとんど吹き上げ正方形を、集合の Cartesian
図式へ写す。
\end{enumerate}
\end{lemma}

\begin{proof}
Proposition \ref{flat-proposition-check-h} により、圏 $(\Sch/S)_h$ において $X$ がアフィンで
あるほとんど吹き上げ正方形 (\ref{flat-equation-almost-blow-up-square}) が与えられたとき、図式
$$
\xymatrix{
\mathcal{F}(E) & \mathcal{F}(X') \ar[l] \\
\mathcal{F}(Z) \ar[u] & \mathcal{F}(X) \ar[u] \ar[l]
}
$$
が集合の圏において Cartesian であることを示せば十分である。証明の大まかな考えは、
この射を、仮定 (3) と (4) を局所的に適用できる別のほとんど吹き上げ正方形で
支配することである。

\medskip\noindent
圏 $(\Sch/S)_h$ のほとんど吹き上げ正方形
(\ref{flat-equation-almost-blow-up-square})、開被覆 $X = \bigcup U_i$、および開被覆
$U_i \cap U_j = \bigcup U_{ijk}$ で、図式
$$
\vcenter{
\xymatrix{
\mathcal{F}(E \cap b^{-1}(U_i)) & \mathcal{F}(b^{-1}(U_i)) \ar[l] \\
\mathcal{F}(Z \cap U_i) \ar[u] & \mathcal{F}(U_i) \ar[u] \ar[l]
}
}
\quad\text{and}\quad
\vcenter{
\xymatrix{
\mathcal{F}(E \cap b^{-1}(U_{ijk})) &
\mathcal{F}(b^{-1}(U_{ijk})) \ar[l] \\
\mathcal{F}(Z \cap U_{ijk}) \ar[u] &
\mathcal{F}(U_{ijk}) \ar[u] \ar[l]
}
}
$$
が Cartesian となるものが与えられたとする。このとき
$$
\xymatrix{
\mathcal{F}(E) & \mathcal{F}(X') \ar[l] \\
\mathcal{F}(Z) \ar[u] & \mathcal{F}(X) \ar[u] \ar[l]
}
$$
についても同じことが成り立つ。これは $\mathcal{F}$ が Zariski 位相における層である
ことから従う。

\medskip\noindent
特に、$b : X' \to X$ が閉埋め込みで、$Z$ が局所主閉部分スキームである
% P05-EN-CH38-0292: frozen-source pointer; see ERRATA_CANDIDATES.jsonl.
吹き上げ正方形 (\ref{flat-equation-almost-blow-up-square}) が与えられたなら、
$\mathcal{F}(X) = \mathcal{F}(X') \times_{\mathcal{F}(E)} \mathcal{F}(Z)$.
実際、$X$ 上アフィン局所的には (3) のようなほとんど吹き上げ正方形を得る。

\medskip\noindent
$Z \subset X$、$E_k \subset X'_k \to X$、$E \subset X' \to X$、および
$i_k : X' \to X'_k$ を Lemma \ref{flat-lemma-almost-blow-up-unique} の主張のようにとる。
このとき
$$
\xymatrix{
E \ar[d] \ar[r] & X' \ar[d] \\
E_k \ar[r] & X'_k
}
$$
は前段落で論じた種類のほとんど吹き上げ正方形である。従って前段落の結果により
$$
\mathcal{F}(X'_k) = \mathcal{F}(X') \times_{\mathcal{F}(E)} \mathcal{F}(E_k)
$$
が $k = 1, 2$ に対して成り立つ。従って
$$
\mathcal{F}(X) \longrightarrow
\mathcal{F}(X'_k) \times_{\mathcal{F}(E_k)} \mathcal{F}(Z)
$$
が $k = 1$ に対して全単射であることと、$k = 2$ に対して全単射であることは同値である。
従って $X$ が準コンパクトかつ準分離である有限表示の閉埋め込み $Z \to X$ が与えられた
とき、$\mathcal{F}(X) = \mathcal{F}(X') \times_{\mathcal{F}(E)} \mathcal{F}(Z)$ が
成り立つかどうかは、選ぶほとんど吹き上げ正方形
(\ref{flat-equation-almost-blow-up-square}) に依存しない。
（さらに Lemma \ref{flat-lemma-almost-blow-up-square} により、$Z \subset X$ に対する
ほとんど吹き上げ正方形は実際に存在する。）

\medskip\noindent
最後に、$(\Sch/S)_h$ のアフィン対象 $X$ と有限表示の閉埋め込み $Z \to X$ を考える。
$X$ において $Z$ を定義するイデアルの生成元の個数 $r$ に関する帰納法で、組
$(X, Z)$ に対する所望の性質を証明する。生成元の個数が $\leq 2$ なら、
Example \ref{flat-example-two-generators} のようにほとんど吹き上げ正方形を選べるので、
仮定 (4) により結論を得る。

\medskip\noindent
帰納段階。$r > 2$ として、$X = \Spec(A)$ および
$Z = \Spec(A/(f_1, \ldots, f_r))$ と仮定する。組 $(X, Z)$ に対する吹き上げ正方形
(\ref{flat-equation-almost-blow-up-square}) を選ぶ。
$Z_1 = \Spec(A/(f_1, f_2))$ とおき、
$$
\xymatrix{
E_1 \ar[d] \ar[r] & Y \ar[d] \\
Z_1 \ar[r] & X
}
$$
を Example \ref{flat-example-two-generators} で構成されたほとんど吹き上げ正方形とする。
Lemma \ref{flat-lemma-base-change-almost-blow-up} により、基底変換
$$
(I)
\vcenter{
\xymatrix{
Y \times_X E \ar[r] \ar[d] & Y \times_X X' \ar[d] \\
Y \times_X Z \ar[r] & Y
}
}
\quad\text{and}\quad
(II)
\vcenter{
\xymatrix{
E \ar[r] \ar[d] & Z_1 \times_X X' \ar[d] \\
Z \ar[r] & Z_1
}
}
$$
はいずれもほとんど吹き上げ正方形である。$Z_1$ における $Z$ のイデアルは
$r - 2$ 個の元で生成される。$Y \times_X Z$ のイデアルは $f_1, \ldots, f_r$ の
$Y$ への引き戻しで生成される。$Y$ 上局所的に $f_1, f_2$ で生成されるイデアルは
一つの元で生成できるので、$Y \times_X Z$ は $Y$ 上アフィン局所的に高々
$r - 1$ 個の元で切り出される。帰納法の仮定と上の議論により
$$
\mathcal{F}(Y) =
\mathcal{F}(Y \times_X X') \times_{\mathcal{F}(Y \times_X E)}
\mathcal{F}(Y \times_X Z)
$$
and
$$
\mathcal{F}(Z_1) =
\mathcal{F}(Z_1 \times_X X') \times_{\mathcal{F}(E)}
\mathcal{F}(Z)
$$
である。仮定 (4) により
$$
\mathcal{F}(X) =
\mathcal{F}(Y) \times_{\mathcal{F}(E_1)} \mathcal{F}(Z_1)
$$
$s' \in \mathcal{F}(X')$、$t \in \mathcal{F}(Z)$ で、
$\mathcal{F}(E)$ における制限が同じであるような組 $(s', t)$ が与えられたとする。
このとき
% P05-EN-CH38-0293: frozen-source pointer; see ERRATA_CANDIDATES.jsonl.
$(s'|{Z_1 \times_X X'}, t)$ は一意な元 $t_1 \in \mathcal{F}(Z_1)$ の像である。
同様に $(s'|_{Y \times_X X'}, t|_{Y \times_X Z})$ は一意な元
$s_Y \in \mathcal{F}(Y)$ の像である。$s_Y$ と $t_1$ は $\mathcal{F}(E_1)$ の
同じ元に制限されると主張する。これは、ほとんど吹き上げ正方形
$$
\xymatrix{
E_1 \times_X E \ar[r] \ar[d] & E_1 \times_X X' \ar[d] \\
E_1 \times_X Z \ar[r] & E_1
}
$$
が、$E_1 \to Y$ によるほとんど吹き上げ正方形 (I) の基底変換であり、かつ
$E_1 \to Z_1$ によるほとんど吹き上げ正方形 (II) の基底変換であること、そして
$s_Y$ と $t_1$ の構成に用いた切断の組が一致することから分かる。
従って第三のファイバー積等式により、$\mathcal{F}(Y)$ において $s_Y$ に写り、
% P05-EN-CH38-0294: frozen-source pointer; see ERRATA_CANDIDATES.jsonl.
$\mathcal{F}(Z)$ において $t_1$ に写る一意な $s \in \mathcal{F}(X)$ が存在する。
$s$ が $\mathcal{F}(X')$ において $s'$ に、$\mathcal{F}(Z)$ において $t$ に写ることの
検証は省略する。ヒント：今構成した $s$ の一意性を用い、アフィン局所的に作業せよ。
\end{proof}

\begin{lemma}
\label{flat-lemma-refine-check-h-stack}
$p : \mathcal{S} \to (\Sch/S)_h$ を groupoid にファイバー化された圏とする。
このとき $\mathcal{S}$ が groupoid の stack であるための必要十分条件は、次の条件が
満たされることである。
\begin{enumerate}
\item $\mathcal{S}$ は Zariski 位相に対する groupoid の stack である。
\item $Y$ がアフィンで、$f$ が全射、平坦、固有、かつ有限表示である
$(\Sch/S)_h$ の射 $f : X \to Y$ が与えられたとき、
$$
\mathcal{S}_Y \longrightarrow
\mathcal{S}_X \times_{\mathcal{S}_{X \times_Y X}} \mathcal{S}_X
$$
は圏の同値である。
\item 圏 $(\Sch/S)_h$ の Example \ref{flat-example-one-generator} または
\ref{flat-example-two-generators} のようなほとんど吹き上げ正方形に対して、関手
$$
\mathcal{S}_X \longrightarrow
\mathcal{S}_Z \times_{\mathcal{S}_E} \mathcal{S}_{X'}
$$
は圏の同値である。
\end{enumerate}
\end{lemma}

\begin{proof}
この補題は Lemma \ref{flat-lemma-refine-check-h} と我々の groupoid の stack の定義の
形式的帰結である。例えば (1)、(2)、(3) を仮定する。$\mathcal{S}$ が stack であることを
示すには、射と対象に対する降下を証明しなければならない。Stacks, Definition
\ref{stacks-definition-stack-in-groupoids} を参照せよ。

\medskip\noindent
$x, y$ が $(\Sch/S)_h$ の対象 $U$ 上の $\mathcal{S}$ の対象なら、仮定により
$\mathit{Isom}(x, y)$ は $(\Sch/U)_h$ 上の前層で、Lemma
\ref{flat-lemma-refine-check-h} の (1)、(2)、(3)、(4) を満たす。従ってこれは層である。
細部は省略する。

\medskip\noindent
$\{U_i \to U\}_{i \in I}$ を $(\Sch/S)_h$ の被覆とする。
$(x_i, \varphi_{ij})$ を族 $\{U_i \to U\}_{i \in I}$ に関する $\mathcal{S}$ の降下データとする。
Stacks, Definition \ref{stacks-definition-descent-data} を参照せよ。
$(\Sch/U)_h$ の $V/U$ に対し、次の組 $(y, \psi_i)$ の集合を対応させる規則 $F$ を考える。
ここで $y$ は $\mathcal{S}_V$ の対象であり、
$\psi_i : y|_{U_i \times_U V} \to x_i|_{U_i \times_U V}$ は
$U_i \times_U V$ 上の $\mathcal{S}$ の射で、
$$
\varphi_{ij}|_{U_i \times_U U_j \times_U V}
\circ
\psi_i|_{U_i \times_U U_j \times_U V}
=
\psi_j|_{U_i \times_U U_j \times_U V}
$$
が同型を除いて成り立つものとする。射に対する降下はすでに得ているので、
$F(V/U)$ が空集合または一点集合であることは明らかである。一方、
$F(U_{i_0}/U)$ は $(x_{i_0}, \varphi_{i_0i})$ を含むので空でない。
目標は $F(U/U)$ が空でないことを証明することなので、$F$ が $(\Sch/U)_h$ 上の層である
ことを示せば十分である。このため Lemma \ref{flat-lemma-refine-check-h} の判定条件を用いてよい。
ところが仮定 (1)、(2)、(3) により（省略するいくつかの可換図式を描けば）、
Lemma \ref{flat-lemma-refine-check-h} の性質 (1)、(2)、(3)、(4) が $F$ に対して成り立つ。

\medskip\noindent
$\mathcal{S}$ が groupoid の stack なら (1)、(2)、(3) が満たされることの検証は省略する。
\end{proof}





\section{絶対弱正規化と h 被覆}
\label{flat-section-h-and-weak-normalization}

\noindent
この節では Section \ref{flat-section-blow-up-h} で得た判定条件を用いていくつかの h 層を
構成し、構造層の h 層化と絶対弱正規化との関係を与える。そのために次の初等的な補題を
必要とする。

\begin{lemma}
\label{flat-lemma-funny-blow-up}
Example \ref{flat-example-two-generators} のようなほとんど吹き上げ正方形を
$Z, X, X', E$ とする。このとき $p > 0$ に対して
$H^p(X', \mathcal{O}_{X'}) = 0$ であり、
$\Gamma(X, \mathcal{O}_X) \to \Gamma(X', \mathcal{O}_{X'})$
は全射な環準同型で、その核は平方零イデアルである。
\end{lemma}

\begin{proof}
まず $A = \mathbf{Z}[f_1, f_2]$ が多項式環であると仮定する。この場合、考えている
ほとんど吹き上げ正方形は閉部分スキーム $Z$ に沿う $X = \Spec(A)$ の吹き上げである。
実際、$X' \subset \mathbf{P}^1_X$ は
$\mathcal{O}_{\mathbf{P}^1_X}(1)$ の大域切断 $f_2T_0 - f_1 T_1$ で切り出される
有効 Cartier 因子である。従って分解
$$
0 \to
\mathcal{O}_{\mathbf{P}^1_X}(-1) \to
\mathcal{O}_{\mathbf{P}^1_X} \to
\mathcal{O}_{X'} \to
0
$$
を得る。Cohomology of Schemes, Section
\ref{coherent-section-cohomology-projective-space} で与えられるコホモロジーの記述を
用いると、この場合
$\Gamma(X, \mathcal{O}_X) \to \Gamma(X', \mathcal{O}_{X'})$
は同型であり、$H^1(X', \mathcal{O}_{X'}) = 0$ であることが従う。

\medskip\noindent
次に、Example \ref{flat-example-two-generators} のような任意の図式は、前段落の図式を
環準同型 $\mathbf{Z}[f_1, f_2] \to A$ で基底変換したものであることに注意する。
従って More on Morphisms, Lemmas
\ref{more-morphisms-lemma-check-h1-fibre-zero},
\ref{more-morphisms-lemma-h1-fibre-zero}, and
\ref{more-morphisms-lemma-h1-fibre-zero-check-h0-kappa}
により、一般に $H^1(X', \mathcal{O}_{X'})$ は零であり、写像
$H^0(X, \mathcal{O}_X) \to H^0(X', \mathcal{O}_{X'})$ は全射である。

\medskip\noindent
次に一般の場合に核を調べる。$a \in A$ が零に写るとする。アフィンチャート上で見ると、
$$
a = (f_1x - f_2)(a_0 + a_1x + \ldots + a_rx^r)\text{ in }A[x]
$$
となる $r \geq 0$ および $a_0, \ldots, a_r \in A$ が存在し、同様に
$$
a = (f_1 - f_2y)(b_0 + b_1y + \ldots + b_s y^s)\text{ in }A[y]
$$
となる $s \geq 0$ および $b_0, \ldots, b_s \in A$ が存在する。これは
$$
a = f_2 a_0,\ f_1 a_0 = f_2 a_1,\ \ldots,\ f_1 a_r = 0,
\ a = f_1 b_0,\ f_2 b_0 = f_1 b_1,\ \ldots,\ f_2 b_s = 0
$$
を意味する。$(a', r', a'_i, s', b'_j)$ をもう一つのそのような系とすると、
$$
aa' = f_1f_2a_0b'_0 = f_1f_2a_1b'_1 = f_1f_2a_2b'_2 = \ldots = 0
$$
となり、求める結果を得る。
\end{proof}

\noindent
$\mathbf{F}_p$-代数 $A$ に対し、$\colim_F A$ を系
$$
A \xrightarrow{F} A \xrightarrow{F} A \xrightarrow{F} \ldots
$$
の余極限と定める。ここで $F : A \to A$, $a \mapsto a^p$ は Frobenius 自己準同型である。

\begin{lemma}
\label{flat-lemma-h-sheaf-colim-F}
$p$ を素数、$S$ を $\mathbf{F}_p$ 上のスキームとし、$(\Sch/S)_h$ を
Definition \ref{flat-definition-big-small-h} のようなサイトとする。このとき
$$
\mathcal{F}(X) = \colim_F \Gamma(X, \mathcal{O}_X)
$$
が $(\Sch/S)_h$ の任意の準コンパクトかつ準分離な対象 $X$ に対して成り立つような
$(\Sch/S)_h$ 上の層 $\mathcal{F}$ が一意に存在する。
\end{lemma}

\begin{proof}
$\mathcal{F}$ を関手
$$
X \longrightarrow \colim_F \Gamma(X, \mathcal{O}_X)
$$
の Zariski 層化とする。Sheaves, Lemma
\ref{sheaves-lemma-directed-colimits-sections} および $\mathcal{O}$ が Zariski 位相に関する
層であることにより、準コンパクトかつ準分離なスキーム $X$ に対して
$\mathcal{F}(X) = \colim_F \Gamma(X, \mathcal{O}_X)$ である。従って $\mathcal{F}$ が
h 層であることを示せば十分である。このため Lemma \ref{flat-lemma-refine-check-h} の条件
(1)、(2)、(3)、(4) を確認する。条件 (1) は（ほとんど不要な）Zariski 層化を行ったことから
成り立つ。条件 (2) は、$\mathcal{O}$ が fppf 層であり（Descent, Lemma
\ref{descent-lemma-sheaf-condition-holds}）、また $A$ が $\mathbf{F}_p$-代数の二つの写像
$B \to C$ の等化子なら、$\colim_F A$ が二つの写像
$\colim_F B \to \colim_F C$ の等化子だから成り立つ。

\medskip\noindent
条件 (3) を確認する。$A, f, J$ を Example \ref{flat-example-one-generator} のようにとる。
示すべきことは
$$
\colim_F A = \colim_F A/J \times_{\colim_F A/fA + J} \colim_F A/fA
$$
である。これは次の代数的な問題に帰着する。$a', a'' \in A$ が
$F^n(a' - a'') \in fA + J$ を満たすとする。$a - F^m(a') \in J$ および
$a - F^m(a'') \in fA$ を満たす $a \in A$ と $m \geq 0$ を求め、組 $(a, m)$ が
$(a, m) \mapsto (F(a), m + 1)$ という形の置換を除いて一意に定まることを示せばよい。
実際、$h \in A$ と $g \in J$ を用いて $F^n(a' - a'') = f h + g$ と書き、
$a = F^n(a') - g = F^n(a'') + fh$ および $m = n$ とおく。一意性を見るため、
$(a_1, m_1)$ を別の解とする。上の形の置換により $m = m_1$ と仮定してよい。
すると $a - a_1 \in J$ かつ $a - a_1 \in fA$ である。$J$ は $f$ のある冪で
零化されるので、$a - a_1$ は冪零元である。従って十分大きいある $k$ に対して
$F^k(a - a_1)$ は零となる。さらに置換を行えば $a = a_1$ を得る。

\medskip\noindent
条件 (4) を確認する。$X, X', Z, E$ を Example \ref{flat-example-two-generators} のように
とる。Lemma \ref{flat-lemma-funny-blow-up} により
$$
\mathcal{F}(X) = \colim_F \Gamma(X, \mathcal{O}_X)
\longrightarrow
\colim_F \Gamma(X', \mathcal{O}_{X'}) = \mathcal{F}(X')
$$
は全単射である。この場合 $E = \mathbf{P}^1_Z$ なので、
$\mathcal{F}(Z) \to \mathcal{F}(E)$ も全単射である。従ってこの場合にも結論が成り立つ。
\end{proof}

\noindent
$p$ を素数とする。$\mathbf{F}_p$-代数 $A$ に対し、$\lim_F A$ を逆系
$$
\ldots \xrightarrow{F} A \xrightarrow{F} A \xrightarrow{F} A
$$
の極限と定める。ここで $F : A \to A$, $a \mapsto a^p$ は Frobenius 自己準同型である。

\begin{lemma}
\label{flat-lemma-h-sheaf-lim-F}
$p$ を素数、$S$ を $\mathbf{F}_p$ 上のスキームとし、$(\Sch/S)_h$ を
Definition \ref{flat-definition-big-small-h} のようなサイトとする。このとき規則
$$
\mathcal{F}(X) = \lim_F \Gamma(X, \mathcal{O}_X)
$$
は $(\Sch/S)_h$ 上の層を定める。
\end{lemma}

\begin{proof}
$\mathcal{F}$ が層であることを証明するため、Lemma \ref{flat-lemma-refine-check-h} の条件
(1)、(2)、(3)、(4) を確認する。層の極限は層であり、$\mathcal{O}$ は Zariski 層なので
条件 (1) が成り立つ。条件 (2) は、$\mathcal{O}$ が fppf 層であり（Descent, Lemma
\ref{descent-lemma-sheaf-condition-holds}）、また $A$ が $\mathbf{F}_p$-代数の二つの写像
$B \to C$ の等化子なら、$\lim_F A$ が二つの写像
$\lim_F B \to \lim_F C$ の等化子だから成り立つ。

\medskip\noindent
条件 (3) を確認する。$A, f, J$ を Example \ref{flat-example-one-generator} のようにとる。
示すべきことは、写像
\begin{align*}
\lim_F A
& \to
\lim_F A/J \times_{\lim_F A/fA + J} \lim_F A/fA \\
& =
\lim_F (A/J \times_{A/fA + J} A/fA) \\
& =
\lim_F A/(fA \cap J)
\end{align*}
が全単射であることである。$J$ は $f$ のある冪で零化されるので、
$\mathfrak a = fA \cap J$ は冪零イデアル、すなわち $\mathfrak a^n = 0$ を満たす
ある $n$ が存在する。この場合に $\lim_F A \to \lim_F A/\mathfrak a$ が全単射であることは
直ちに確認できる。

\medskip\noindent
条件 (4) を確認する。$X, X', Z, E$ を Example \ref{flat-example-two-generators} のように
とる。Lemma \ref{flat-lemma-funny-blow-up} と上と同じ議論により
$$
\mathcal{F}(X) = \lim_F \Gamma(X, \mathcal{O}_X)
\longrightarrow
\lim_F \Gamma(X', \mathcal{O}_{X'}) = \mathcal{F}(X')
$$
は全単射である。この場合 $E = \mathbf{P}^1_Z$ なので、
$\mathcal{F}(Z) \to \mathcal{F}(E)$ も全単射である。従ってこの場合にも結論が成り立つ。
\end{proof}

\noindent
次の補題ではスキーム $X$ の絶対弱正規化 $X^{awn}$ を用いる。Morphisms, Section
\ref{morphisms-section-seminormalization} を参照せよ。

\begin{lemma}
\label{flat-lemma-weak-normalization-ph-sheaf}
$(\Sch/S)_{ph}$ を Topologies, Definition
\ref{topologies-definition-big-small-ph} のようなサイトとする。このとき規則
$$
X \longmapsto \Gamma(X^{awn}, \mathcal{O}_{X^{awn}})
$$
は $(\Sch/S)_{ph}$ 上の層である。
\end{lemma}

\begin{proof}
$\mathcal{F}$ が層であることを証明するため、Topologies, Lemma
\ref{topologies-lemma-characterize-sheaf} の条件 (1) と (2) を確認する。
$X^{awn}$ の構成は開被覆と可換なので条件 (1) が成り立つ。Morphisms, Lemma
\ref{morphisms-lemma-seminormalization} とその証明を参照せよ。

\medskip\noindent
$\pi : Y \to X$ を全射固有射とする。二つの写像
$$
\Gamma(Y^{awn}, \mathcal{O}_{Y^{awn}}) \to
\Gamma((Y \times_X Y)^{awn}, \mathcal{O}_{(Y \times_X Y)^{awn}})
$$
の等化子が $\Gamma(X^{awn}, \mathcal{O}_{X^{awn}})$ に等しいことを示す必要がある。
$f$ をこの等化子の元とし、射
$$
f : Y^{awn} \longrightarrow \mathbf{A}^1_X
$$
を考える。$Y^{awn} \to X$ は普遍閉なので、$f$ のスキーム論的像 $Z$ は
$X$ 上固有な $\mathbf{A}^1_X$ の閉部分スキームであり、
$f : Y^{awn} \to Z$ は全射である。Morphisms, Lemma
\ref{morphisms-lemma-scheme-theoretic-image-is-proper} を参照せよ。従って
$Z \to X$ は有限（Morphisms, Lemma \ref{morphisms-lemma-finite-proper}）かつ全射である。

\medskip\noindent
$k$ を体とし、$z_1, z_2 : \Spec(k) \to Z$ を $Z \to X$ により同一視される二つの射とする。
$z_1 = z_2$ であると主張する。像 $\lambda_i = z_i^*f \in k$ が一致することを示せば十分である
（$Z$ の構造層は $X$ の構造層上 $f$ で生成されるからである）。これを見るため、体拡大
$K/k$ と射 $y_1, y_2 : \Spec(K) \to Y^{awn}$ で
$z_i \circ (\Spec(K) \to \Spec(k)) = f \circ y_i$ を満たすものを選ぶ。これは写像
$Y^{awn} \to Z$ が全射なので可能である。非常に大きな濃度をもつ代数閉拡大
$\Omega/k$ を選ぶ。任意の $k$-代数準同型 $\sigma_i : K \to \Omega$ に対し
$$
\Spec(\Omega)
\xrightarrow{\sigma_1, \sigma_2}
\Spec(K \otimes_k K)
\xrightarrow{y_1, y_2}
Y^{awn} \times_X Y^{awn}
$$
を得る。標準射 $(Y \times_X Y)^{awn} \to Y^{awn} \times_X Y^{awn}$ は
普遍同相射であり、$\Omega$ は代数閉なので、上の合成は射
$\Spec(\Omega) \to (Y \times_X Y)^{awn}$ に一意に持ち上がる。$f$ は上の等化子に
属するから、$\sigma_1(\lambda_1) = \sigma_2(\lambda_2)$ が従う。体拡大に関する簡単な
補題により、これは $\lambda_1 = \lambda_2$ を含意する。詳細は省略する。

\medskip\noindent
従って $Z \to X$ は普遍単射である。すなわち $Z \to X$ は点上単射であり、純非分離な
剰余体拡大を誘導する（Morphisms, Lemma
\ref{morphisms-lemma-universally-injective}）。以上から $Z \to X$ は普遍同相射である。
Morphisms, Lemma \ref{morphisms-lemma-universal-homeomorphism} を参照せよ。

\medskip\noindent
$g : X^{awn} \to Z$ を $X^{awn}$ の普遍性から得られる写像とする。このとき
$Y^{awn} \to X^{awn} \to Z$ と $f : Y^{awn} \to Z$ は $X$ 上の二つの射である。
$Y^{awn} \to Y$ の普遍性により、対応する $Y$ 上の二つの射
$Y^{awn} \to Y \times_X Z$ は等しくなければならない。従って
% P05-EN-CH38-0299: frozen-source pointer; see ERRATA_CANDIDATES.jsonl.
$g \circ \pi^{wan} = f$ が $\mathbf{A}^1_X$ への射として成り立ち、
$g \in \Gamma(X^{awn}, \mathcal{O}_{X^{awn}})$ が求める元である。
\end{proof}

\begin{lemma}
\label{flat-lemma-weak-normalization-h-sheaf}
$S$ をスキームとし、$(\Sch/S)_h$ を Definition
\ref{flat-definition-big-small-h} のようなサイトとする。このとき規則
$$
X \longmapsto \Gamma(X^{awn}, \mathcal{O}_{X^{awn}})
$$
は $(\Sch/S)_h$ 上の「構造層」$\mathcal{O}$ の層化である。ph 位相についても同様である。
\end{lemma}

\begin{proof}
Lemma \ref{flat-lemma-weak-normalization-ph-sheaf} で、この補題の規則 $\mathcal{F}$ が
ph 位相に関する層を定め、従ってなおさら h 位相に関する層を定めることを見た。
環の前層の標準写像 $\mathcal{O} \to \mathcal{F}$ が存在することは明らかである。
証明を終えるには、次を示せば十分である。
\begin{enumerate}
\item $f \in \mathcal{O}(X)$ が $\mathcal{F}(X)$ で零に写るなら、
$f|_{X_i} = 0$ となる h 被覆 $\{X_i \to X\}$ が存在する。
\item $f \in \mathcal{F}(X)$ が与えられたなら、$f|_{X_i}$ が
$f_i \in \mathcal{O}(X_i)$ の像となる h 被覆 $\{X_i \to X\}$ が存在する。
\end{enumerate}
(1) のような $f$ をとる。このとき $f|_{X^{awn}} = 0$ であり、これは $f$ が局所的に
冪零であることを意味する。従って $X' \subset X$ を $f$ で切り出される閉部分スキームと
すると、$X' \to X$ は有限表示の全射閉埋め込みである。ゆえに $\{X' \to X\}$ が求める
h 被覆である。(2) のような $f$ をとる。$X$ をアフィン開被覆の各要素で置き換え、
$X = \Spec(A)$ がアフィンであると仮定してよい。このとき $f \in A^{awn}$ である。
Morphisms, Lemma \ref{morphisms-lemma-seminormalization-ring} を参照せよ。
Morphisms, Lemma \ref{morphisms-lemma-universal-homeo-limit} により、
$\Spec(B) \to \Spec(A)$ が普遍同相射となる有限表示の環準同型 $A \to B$ で、$f$ が
標準写像 $B \to A^{awn}$ によるある元 $b \in B$ の像となるものをとれる。すると
$\{\Spec(B) \to \Spec(A)\}$ は h 被覆であり、結論を得る。ph 位相についての主張も
同様に従う（または h 位相についての主張から導ける）。
\end{proof}

\noindent
$p$ を素数とする。$\mathbf{F}_p$-代数 $A$ は、写像
$F : A \to A$, $x \mapsto x^p$ が $A$ の自己同型であるとき {\it 完全} と呼ばれる。

\begin{lemma}
\label{flat-lemma-perfect-weankly-normal}
$p$ を素数とする。$\mathbf{F}_p$-代数 $A$ が絶対弱正規であるための必要十分条件は、
完全であることである。
\end{lemma}

\begin{proof}
Morphisms, Definition \ref{morphisms-definition-seminormal-ring} の条件 (2)(b) から、
$A$ が絶対弱正規なら完全であることは直ちに従う。

\medskip\noindent
$A$ が完全であると仮定する。$x^3 = y^2$ を満たす $x, y \in A$ をとる。
$p > 3$ なら、ある $n, m > 0$ によって $p = 2n + 3m$ と書ける。
$a^p = x$ および $b^p = y$ を満たす $a, b \in A$ を選ぶ。
$c = a^n b^m$ とおくと
$$
c^{2p} = x^{2n} y^{2m} = x^{2n + 3m} = x^p
$$
であり、従って $c^2 = x$ である。同様に $c^3 = y$ である。$p = 2$ なら、
$x = a^2$ と書けば $a^6 = y^2$ を得て、これは $a^3 = y$ を含意する。
$p = 3$ なら、$y = a^3$ と書けば $x^3 = a^6$ を得て、これは $x = a^2$ を含意する。

\medskip\noindent
ある素数 $\ell$ に対して $\ell^\ell x = y^\ell$ を満たす $x, y \in A$ をとる。
$\ell \not = p$ なら $a = y/\ell$ は $a^\ell = x$ および $\ell a = y$ を満たす。
$\ell = p$ なら $y = 0$ であり、ある $a$ に対して $x = a^p$ である。
\end{proof}

\begin{lemma}
\label{flat-lemma-char-p}
$p$ を素数とする。
\begin{enumerate}
\item $A$ が $\mathbf{F}_p$-代数なら $\colim_F A = A^{awn}$ である。
\item $S$ が $\mathbf{F}_p$ 上のスキームなら、$\mathcal{O}$ の h 層化は準コンパクトかつ
準分離な $X$ を $\colim_F \Gamma(X, \mathcal{O}_X)$ に写す。
\end{enumerate}
\end{lemma}

\begin{proof}
(1) の証明。Algebra, Lemma \ref{algebra-lemma-p-ring-map} により、
$A \to \colim_F A$ はスペクトル上の普遍同相射を誘導することに注意する。
従って $B = \colim_F A$ が絶対弱正規であることを示せば十分である。Morphisms, Lemma
\ref{morphisms-lemma-seminormalization-ring} を参照せよ。環準同型 $F : B \to B$ は自己同型、
言い換えれば $B$ は完全環であることに注意する。従って Lemma
\ref{flat-lemma-perfect-weankly-normal} を適用できる。

\medskip\noindent
(2) の証明。これはアフィン局所的に考えれば、(1) と Lemmas
\ref{flat-lemma-h-sheaf-colim-F} and \ref{flat-lemma-weak-normalization-h-sheaf} から従う。
\end{proof}






\section{正標数におけるベクトル束の降下}
\label{flat-section-descent-vectorbundles-h}

\noindent
この節の参考文献は \cite{Witt-Grass} である。

\medskip\noindent
スキーム $S$ に対し、有限局所自由 $\mathcal{O}_S$-加群の圏を
$\textit{Vect}(S)$ と書く。$p$ を素数、$S$ を $\mathbf{F}_p$ 上の準コンパクトかつ
準分離なスキームとする。この節では圏
$$
\colim_F \textit{Vect}(S) =
\colim
\left(
\textit{Vect}(S) \xrightarrow{F^*}
\textit{Vect}(S) \xrightarrow{F^*}
\textit{Vect}(S) \xrightarrow{F^*} \ldots
\right)
$$
を扱う。ここで $F : S \to S$ は絶対 Frobenius 射である。具体的には、この圏の対象は
組 $(\mathcal{E}, n)$ であり、$\mathcal{E}$ は有限局所自由 $\mathcal{O}_S$-加群、
$n \geq 0$ は整数である。射は
$$
\Hom_{\colim_F \textit{Vect}(S)}((\mathcal{E}, n), (\mathcal{G}, m)) =
\colim_N \Hom_S(F^{N - n, *}\mathcal{E}, F^{N - m, *}\mathcal{G})
$$
とする。ここで $F : S \to S$ は $S$ の絶対 Frobenius 射である。従って対象
$(\mathcal{E}, n)$ は対象 $(F^*\mathcal{E}, n + 1)$ と同型である。

\begin{lemma}
\label{flat-lemma-colim-F-Vect}
$p$ を素数、$S$ を $\mathbf{F}_p$ 上の準コンパクトかつ準分離なスキームとする。
圏 $\colim_F \textit{Vect}(S)$ は、$S$ 上の環の層 $\colim_F \mathcal{O}_S$ 上の
有限局所自由加群の圏と同値である。
\end{lemma}

\begin{proof}
省略する。
\end{proof}

\begin{lemma}
\label{flat-lemma-vector-bundle-I}
$p$ を素数とする。Example \ref{flat-example-one-generator} のような標数 $p$ の
ほとんど吹き上げ正方形 $X, X', Z, E$ を考える。このとき関手
$$
\colim_F \textit{Vect}(X)
\longrightarrow
\colim_F \textit{Vect}(Z)
\times_{\colim_F \textit{Vect}(E)}
\colim_F \textit{Vect}(X')
$$
は同値である。
\end{lemma}

\begin{proof}
$A, f, J$ を Example \ref{flat-example-one-generator} のようにとる。考えるスキームはすべて
アフィンであり、ベクトル束の圏には内部 Hom があるから、関手の充満忠実性は、有限射影
$A$-加群 $P$ に対して
$$
\colim P \otimes_{A, F^N} A =
\colim P \otimes_{A, F^N} A/J
\times_{\colim P \otimes_{A, F^N} A/fA + J}
\colim P \otimes_{A, F^N} A/fA
$$
を示せば従う。$P$ を有限自由加群の直和因子として書けば、これは $P$ が有限自由である
場合から従う。この場合は直ちに $P = A$ の場合に帰着する。$P = A$ の場合は Lemma
\ref{flat-lemma-h-sheaf-colim-F} から従う（実際、その補題の証明でこの場合を直接証明した）。

\medskip\noindent
本質的全射性。ここでは次の代数的問題を得る。$P_1$ を有限射影 $A/J$-加群、
$P_2$ を有限射影 $A/fA$-加群とし、
$$
\varphi :
P_1 \otimes_{A/J} A/fA + J
\longrightarrow
P_2 \otimes_{A/fA} A/fA + J
$$
を同型とする。目標は、ある $N$、有限射影 $A$-加群 $P$、同型
$\varphi_1 : P \otimes_A A/J \to P_1 \otimes_{A/J, F^N} A/J$,
および同型
$\varphi_2 : P \otimes_A A/fA \to P_2 \otimes_{A/fA, F^N} A/fA$
で、明らかな意味で $\varphi$ と整合するものが存在することを示すことである。これは次の
ように分かる。まず
$$
A/(J \cap fA) = A/J \times_{A/fA + J} A/fA
$$
に注意する。従って More on Algebra, Lemma
\ref{more-algebra-lemma-finitely-presented-module-over-fibre-product} により、
同型 $\varphi'_1 : P' \otimes_A A/J \to P_1$ および
% P05-EN-CH38-0303: frozen-source pointer; see ERRATA_CANDIDATES.jsonl.
$\varphi_2 : P' \otimes_A A/fA \to P_2$
を備えた $A/(J \cap fA)$ 上の有限射影加群 $P'$ が存在し、これらは $\varphi$ と整合する。
$J$ は有限生成イデアルで $f$-冪ねじれなので、$J \cap fA$ は冪零イデアルである。
従ってある $N$ に対して $F^N$ の分解
$$
A \xrightarrow{\alpha} A/(J \cap fA) \xrightarrow{\beta} A
$$
が存在する。$P = P' \otimes_{A/(J \cap fA), \beta} A$ とおけば結論を得る。
\end{proof}

\begin{lemma}
\label{flat-lemma-vector-bundle-II}
$p$ を素数とする。Example \ref{flat-example-two-generators} のような標数 $p$ の
ほとんど吹き上げ正方形 $X, X', Z, E$ を考える。このとき関手
$$
G :
\colim_F \textit{Vect}(X)
\longrightarrow
\colim_F \textit{Vect}(Z)
\times_{\colim_F \textit{Vect}(E)}
\colim_F \textit{Vect}(X')
$$
は同値である。
\end{lemma}

\begin{proof}
充満忠実性。$(\mathcal{E}, n)$ と $(\mathcal{F}, m)$ を
$\colim_F \textit{Vect}(X)$ の対象とし、
$(a, b) : G(\mathcal{E}, n) \to G(\mathcal{F}, m)$ を右辺の射とする。
$N \gg 0$ を選び、$a$ を写像
$a : F^{N - n, *}\mathcal{E}|_Z \to F^{N - m, *}\mathcal{F}|_Z$
と、$b$ を写像
$b : F^{N - n, *}\mathcal{E}|_{X'} \to F^{N - m, *}\mathcal{F}|_{X'}$
とみなして、これらが $E$ 上で一致するとしてよい。$\mathcal{E}|_{X_i}$ と
$\mathcal{F}|_{X_i}$ が有限自由 $\mathcal{O}_{X_i}$-加群となる有限アフィン開被覆
$X = X_1 \cup \ldots \cup X_n$ を選ぶ。各 $i$ に対する基底変換
$$
\xymatrix{
E_i \ar[r] \ar[d] & X'_i \ar[d] \\
Z_i \ar[r] & X_i
}
$$
は Example \ref{flat-example-two-generators} のような別のほとんど吹き上げ正方形である。
これらの正方形に対して Lemma \ref{flat-lemma-h-sheaf-colim-F}（その補題の証明も参照）により
$$
\colim_F H^0(X_i, \mathcal{O}_{X_i}) =
\colim_F H^0(Z_i, \mathcal{O}_{Z_i})
\times_{\colim_F H^0(E_i, \mathcal{O}_{E_i})}
\colim_F H^0(X'_i, \mathcal{O}_{X'_i})
$$
が成り立つ。従って $N$ を大きくすれば、写像 $a|_{Z_i}$ と $b|_{X'_i}$ は写像
$c_i : F^{N - n, *}\mathcal{E}|_{X_i} \to F^{N - m, *}\mathcal{F}|_{X_i}$ から来ると
仮定してよい。必要ならさらに $N$ を大きくし、$c_i$ と $c_j$ が $X_i \cap X_j$ 上で
一致すると仮定してよい。従ってこれらの写像を貼り合わせると、左辺における求める射
$(\mathcal{E}, n) \to (\mathcal{F}, m)$ を得る。

\medskip\noindent
本質的全射性。$(\mathcal{F}, \mathcal{G}, \varphi)$ を、有限局所自由
$\mathcal{O}_Z$-加群 $\mathcal{F}$、有限局所自由 $\mathcal{O}_{X'}$-加群
$\mathcal{G}$、および同型 $\varphi : \mathcal{F}|_E \to \mathcal{G}|_E$ からなる
三つ組とする。この三つ組を Frobenius のある冪による引き戻しで置き換えれば、有限局所自由
$\mathcal{O}_X$-加群から来ることを示す必要がある。

\medskip\noindent
Noetherian 還元。この段落は読み飛ばすことを勧める。
$X = \Spec(A)$、$Z = \Spec(A/(f_1, f_2))$、
$X' = \text{Proj}(A[T_0, T_1]/(f_2T_0 - f_1T_1))$, and
$E = \mathbf{P}^1_Z$ であった。Limits, Lemma
\ref{limits-lemma-descend-invertible-modules} により、$f_1$ と $f_2$ を含む有限生成
$\mathbf{F}_p$-部分代数 $A_0 \subset A$ で、三つ組
$(\mathcal{F}, \mathcal{G}, \varphi)$ が
$X_0 = \Spec(A_0)$ and $Z_0 = \Spec(A_0/(f_1, f_2))$,
$X_0' = \text{Proj}(A_0[T_0, T_1]/(f_2T_0 - f_1T_1))$, and
$E_0 = \mathbf{P}^1_{Z_0}$ に降下するものをとれる。従って考えるスキームは Noetherian
であると仮定してよい。

\medskip\noindent
$X$ が Noetherian であると仮定する。$\mathcal{F}|_{Z \cap X_i}$ が自由となる有限アフィン開被覆
$X = X_1 \cup \ldots \cup X_n$ を選べる。Zariski 位相に関して
$\colim_F \textit{Vect}(X)$ の対象を貼り合わせることができ
（Lemma \ref{flat-lemma-colim-F-Vect}）、$X_i$ と $X_i \cap X_j$ 上の充満忠実性はすでに
分かっているので（証明の第一段落を参照）、各 $X_i$ 上で存在を示せば十分である。
これにより次の段落で扱う場合に帰着する。

\medskip\noindent
$X$ が Noetherian で $\mathcal{F} = \mathcal{O}_Z^{\oplus r}$ であると仮定する。
$\varphi$ により同型 $\mathcal{O}_E^{\oplus r} \to \mathcal{G}|_E$ を得る。
$I = (f_1, f_2) \subset A$ とおく。$\mathcal{I} \subset \mathcal{O}_{X'}$ を $E$ の
イデアル層とする。これは $f_1$ と $f_2$ で大域的に生成される。任意の $n$ に対して全射
$$
(\mathcal{I}^n/\mathcal{I}^{n + 1})^{\oplus r} =
\mathcal{I}^n/\mathcal{I}^{n + 1} \otimes_{\mathcal{O}_E}
\mathcal{G}|_E \longrightarrow
\mathcal{I}^n\mathcal{G}/\mathcal{I}^{n + 1}\mathcal{G}
$$
がある。従ってこの加群の第一コホモロジー群は零である。ここでは
$E = \mathbf{P}^1_Z$ であり、その構造層、さらに実際には任意の大域生成された準連接加群の
$H^1$ が消滅することを用いる。More on Morphisms, Lemma
\ref{more-morphisms-lemma-globally-generated-vanishing} と比較せよ。短完全列
$$
0 \to  \mathcal{I}^n\mathcal{G}/\mathcal{I}^{n + 1}\mathcal{G} \to
\mathcal{G}/\mathcal{I}^{n + 1}\mathcal{G} \to
\mathcal{G}/\mathcal{I}^n\mathcal{G} \to 0
$$
と帰納法を用いると、写像
$$
\lim H^0(X', \mathcal{G}/\mathcal{I}^n\mathcal{G})
\to
H^0(E, \mathcal{G}|_E) = H^0(E, \mathcal{O}_E^{\oplus r}) =
% P05-EN-CH38-0304: frozen-source pointer; see ERRATA_CANDIDATES.jsonl.
A/I^{\oplus r}
$$
は全射である。形式関数定理（Cohomology of Schemes, Theorem
\ref{coherent-theorem-formal-functions}）により、これは写像
$$
H^0(X', \mathcal{G}) \to
H^0(E, \mathcal{G}|_E) = H^0(E, \mathcal{O}_E^{\oplus r}) =
% P05-EN-CH38-0304: frozen-source pointer; see ERRATA_CANDIDATES.jsonl.
A/I^{\oplus r}
$$
が全射であることを含意する。従って、与えられた $\mathcal{G}|_E$ の自明化と整合する写像
$\alpha : \mathcal{O}_{X'}^{\oplus r} \to \mathcal{G}$ を選べる。ゆえに $\alpha$ は
$X'$ における $E$ のある開近傍上で同型である。従って $Z$ の各点は、問題を解ける
アフィン開近傍をもつ。$X' \setminus E \to X \setminus Z$ は同型なので、$Z$ に属さない
$X$ の点についても同じである。最後にもう一度 Zariski 貼り合わせを行えば証明が完了する。
\end{proof}

\begin{proposition}
\label{flat-proposition-h-descent-vector-bundles-p}
$p$ を素数、$S$ を標数 $p$ のスキームとする。このとき groupoid にファイバー化された圏
$$
p : \mathcal{S} \longrightarrow (\Sch/S)_h
$$
で、その $U$ 上のファイバー圏が $U$ 上の有限局所自由
$\colim_F \mathcal{O}_U$-加群の圏であるものは groupoid の stack である。さらに $U$ が
準コンパクトかつ準分離なら、$\mathcal{S}_U$ は $\colim_F \textit{Vect}(U)$ である。
\end{proposition}

\begin{proof}
最後の主張は Lemma \ref{flat-lemma-colim-F-Vect} の内容である。命題を証明するため、
Lemma \ref{flat-lemma-refine-check-h-stack} の条件 (1)、(2)、(3) を確認する。

\medskip\noindent
定義により Zariski 位相に関する貼り合わせがあるので、条件 (1) は成り立つ。

\medskip\noindent
条件 (2) を見るため、$f : X \to Y$ を $S$ 上の有限表示の全射平坦固有射とし、$Y$ は
アフィンであるとする。$Y, X, X \times_Y X$ は準コンパクトかつ準分離なので、命題の主張で
与えたファイバー圏の記述を用いることができる。従って
$$
\textit{Vect}(Y)
\longrightarrow
\textit{Vect}(X)
\times_{\textit{Vect}(X \times_Y X)}
\textit{Vect}(X)
$$
が同値であることを示せば十分である（これは余極限についても同じことを含意する）。
これは有限局所自由加群の fppf 降下から直ちに従う。Descent, Proposition
\ref{descent-proposition-fpqc-descent-quasi-coherent} and
Lemma \ref{descent-lemma-finite-locally-free-descends} を参照せよ。

\medskip\noindent
条件 (3) は Lemmas \ref{flat-lemma-vector-bundle-I} and \ref{flat-lemma-vector-bundle-II} の内容である。
\end{proof}

\begin{lemma}
\label{flat-lemma-trivial-fibres-dvr}
$S$ を離散付値環のスペクトルとし、$f : X \to S$ を幾何学的連結ファイバーをもつ
固有射とする。$\eta \in S$ を生成点とし、$X_n \subset X$ を $S$ 上の一様化元の
$n$ 乗で切り出される閉部分スキームとする。このとき次の性質をもつ整数 $n$ が存在する。
$\mathcal{E}|_{X_\eta}$ と $\mathcal{E}|_{X_n}$ が自由である任意の有限局所自由
$\mathcal{O}_X$-加群 $\mathcal{E}$ は自由である。
\end{lemma}

\begin{proof}
まず $X \to S$ が切断をもつ場合に帰着する。$S = \Spec(A)$ とする。
$X_\eta$ の閉点 $\xi$ を選ぶ。$B$ の分数体が $\kappa(\xi)$ となる離散付値環の拡大
$A \subset B$ を選ぶ。これは Krull--Akizuki（Algebra, Lemma
\ref{algebra-lemma-integral-closure-Dedekind}）と、$\kappa(\xi)$ が $A$ の分数体の有限拡大で
あることから可能である。固有性の付値判定法（Morphisms, Lemma
\ref{morphisms-lemma-characterize-proper}）により、切断
$\sigma : \Spec(B) \to X_B$ を誘導する $B$-値点 $\tau : \Spec(B) \to X$ を得る。
有限局所自由 $\mathcal{O}_X$-加群 $\mathcal{E}$ に対し、$\mathcal{E}_B$ を基底変換
$X_B$ への引き戻しとする。平坦基底変換（Cohomology of Schemes, Lemma
\ref{coherent-lemma-flat-base-change-cohomology}）により
$H^0(X_B, \mathcal{E}_B) = H^0(X, \mathcal{E}) \otimes_A B$ である。従って
$\mathcal{E}_B$ が階数 $r$ の自由加群なら、$H^0(X, \mathcal{E})$ の切断は自由
$B$-加群 $\tau^*\mathcal{E} = \sigma^*\mathcal{E}_B$ を生成する。特に
$\tau^*\mathcal{E}$ を生成する $\mathcal{E}$ の $r$ 個の大域切断
$s_1, \ldots, s_r$ をとれる。このとき
$$
s_1, \ldots, s_r :
\mathcal{O}_X^{\oplus r}
\longrightarrow
\mathcal{E}
$$
は階数 $r$ の有限局所自由 $\mathcal{O}_X$-加群の写像であり、$X_B$ への引き戻しは
自由 $\mathcal{O}_{X_B}$-加群の写像で、各ファイバーのある一点で同型に制限される。
行列式をとると関数
% P05-EN-CH38-0306: frozen-source pointer; see ERRATA_CANDIDATES.jsonl.
$g \in \Gamma(X_\eta, \mathcal{O}_{X_B})$
で、各ファイバーのある一点で可逆なものを得る。ファイバーは固有かつ連結なので、$g$ は
可逆でなければならない（詳細は省略する。ヒント：Varieties, Lemma
\ref{varieties-lemma-proper-geometrically-reduced-global-sections} を用いよ）。従って基底変換
$X_B \to \Spec(B)$ に対して補題を証明すれば十分である。

\medskip\noindent
切断 $\sigma : S \to X$ があると仮定する。$\mathcal{E}$ を生成ファイバー上および
$X_n$ 上で自由である有限局所自由 $\mathcal{O}_X$-加群とする（$n$ は後で選ぶ）。
同型 $\sigma^*\mathcal{E} = \mathcal{O}_S^{\oplus r}$ を選び、$D(A)$ における写像
$$
K = R\Gamma(X, \mathcal{E}) \longrightarrow
R\Gamma(S, \sigma^*\mathcal{E}) = A^{\oplus r}
$$
を考える。上と同様に論じると、誘導写像
$H^0(K) = H^0(X, \mathcal{E}) \to A^{\oplus r}$ が全射であることが
$\mathcal{E}$ が自由であるための必要十分条件である。

\medskip\noindent
$L = R\Gamma(X, \mathcal{O}_X^{\oplus r})$ とおく。対応する写像
$L \to A^{\oplus r}$ は求める性質をもつ。平坦基底変換と $\mathcal{E}$ が生成ファイバー上
自由であるという仮定により
$K \otimes_A Q(A) \cong L \otimes_A Q(A)$ である。$\pi \in A$ を一様化元とする。このとき
$$
K \otimes_A^\mathbf{L} A/\pi^m A =
R\Gamma(X, \mathcal{E} \xrightarrow{\pi^m} \mathcal{E})
$$
であり、$L$ についても同様である。
特殊ファイバー上に台をもつ切断の連接部分層を
$\mathcal{E}_{tors} \subset \mathcal{E}$ と書き、他の $\mathcal{O}_X$-加群についても
同様に書く。
$(\mathcal{O}_X)_{tors} \to \mathcal{O}_X/\pi^k \mathcal{O}_X$
が単射となる $k > 0$ を選ぶ（Cohomology of Schemes, Lemma
\ref{coherent-lemma-Artin-Rees}）。$\mathcal{E}$ は局所自由なので
$\mathcal{E}_{tors} \subset \mathcal{E}/\pi^k\mathcal{E}$ である。
すると $n \geq m + k$ に対して同型
\begin{align*}
(\mathcal{E} \xrightarrow{\pi^m} \mathcal{E})
& \cong
(\mathcal{E}/\pi^k\mathcal{E} \xrightarrow{\pi^m}
\mathcal{E}/\pi^{k + m}\mathcal{E}) \\
& \cong
(\mathcal{O}_X^{\oplus r}/\pi^k\mathcal{O}_X^{\oplus r} \xrightarrow{\pi^m}
\mathcal{O}_X^{\oplus r}/\pi^{k + m}\mathcal{O}_X^{\oplus r}) \\
& \cong
(\mathcal{O}_X^{\oplus r} \xrightarrow{\pi^m} \mathcal{O}_X^{\oplus r})
\end{align*}
が $D(\mathcal{O}_X)$ において成り立つ。これは $D(A)$ における同型
$$
K \otimes_A^\mathbf{L} A/\pi^m A \cong L \otimes_A^\mathbf{L} A/\pi^m A
$$
を定める（$n \geq m + k$ のとき成り立つ）。これらの同型は $\sigma$ による引き戻しと
整合することに注意する。従って特に
$K \otimes_A^\mathbf{L} A/\pi^m A \to (A/\pi^m A)^{\oplus r}$
は次数 $0$ のコホモロジー加群上の全射を定める（$L$ についてそうだからである）。
$A$ は離散付値環なので、$D(A)$ において
$$
K \cong \bigoplus H^i(K)[-i]
\quad\text{and}\quad
L \cong
\bigoplus H^i(L)[-i]
$$
である。More on Algebra, Example
\ref{more-algebra-example-finite-injective-finite-global-dimension} を参照せよ。
コホモロジー群 $H^i(K) = H^i(X, \mathcal{E})$ および
$H^i(L) = H^i(X, \mathcal{O}_X)^{\oplus r}$
は Cohomology of Schemes, Lemma
\ref{coherent-lemma-proper-over-affine-cohomology-finite} により有限 $A$-加群である。
More on Algebra, Lemma \ref{more-algebra-lemma-generalized-valuation-ring-modules} により、
これらの加群は巡回加群の直和である。$H^i(K)$ と $H^i(L)$ の自由部分の階数
$\beta_i$ は等しいことを上で見た。次に
$$
H^i(L \otimes_A^\mathbf{L} A/\pi^m A) =
H^i(L)/\pi^m H^i(L) \oplus H^{i + 1}(L)[\pi^m]
$$
であり、$K$ についても同様であることに注意する。ある $i$ に対して
$A/\pi^eA$ が $H^i(X, \mathcal{O}_X)$、同値に $H^i(L)$ の直和因子として現れるような
最大の整数を $e$ とする。$m = e + 1$ とおけば、
$H^i(L \otimes_A^\mathbf{L} A/\pi^m A)$ は $\beta_i$ 個の $A/\pi^m A$ と、それぞれ
$\pi^e$ で零化されるいくつかの別の巡回加群との直和である。$K$ について同じ議論を行い、
同型
$K \otimes_A^\mathbf{L} A/\pi^m A \cong L \otimes_A^\mathbf{L} A/\pi^m A$
を用いると、$H^i(K)$ は $l > e$ である形 $A/\pi^l A$ の巡回直和因子をもたないことが
従う（さらに $K$ と $L$ が $D(A)$ の対象として同型であることも従うが、これは用いない）。
従って写像
$$
H^0(K \otimes^\mathbf{L}_A A/\pi^{e + 1} A) =
H^0(K)/\pi^{e + 1}H^0(K) \oplus H^1(K)[\pi^{e + 1}]
\longrightarrow
(A/\pi^{e + 1} A)^{\oplus r}
$$
が全射となる唯一の可能性は、第一直和因子上で全射となることである。これが示したかった
ことである。（正確には、切断 $\sigma$ がある場合、補題の主張における整数 $n$ は
$k + e + 1$ に等しくとるべきである。ここで $k$ と $e$ は上のような整数で、$X$ のみに
依存する。）
\end{proof}

\begin{lemma}
\label{flat-lemma-trivial-over-dvrs}
$f : X \to S$ をスキームの射、$\mathcal{E}$ を有限局所自由
$\mathcal{O}_X$-加群とする。次を仮定する。
\begin{enumerate}
\item $f$ は平坦かつ固有で、$\mathcal{O}_S = f_*\mathcal{O}_X$ である。
\item $S$ は正規 Noetherian スキームである。
\item 任意の余次元 $1$ の点 $s \in S$ に対し、$\mathcal{E}$ の
$X \times_S \Spec(\mathcal{O}_{S, s})$ への引き戻しは自由である。
\end{enumerate}
このとき $\mathcal{E}$ は有限局所自由 $\mathcal{O}_S$-加群の引き戻しと同型である。
\end{lemma}

\begin{proof}
標準写像
$$
\Phi : f^*f_*\mathcal{E} \longrightarrow \mathcal{E}
$$
が同型であることを証明する。平坦基底変換（Cohomology of Schemes, Lemma
\ref{coherent-lemma-flat-base-change-cohomology}）と仮定 (1)、(3) により、任意の余次元
$1$ の点 $s \in S$ に対し、これの $X \times_S \Spec(\mathcal{O}_{S, s})$ への引き戻しは
同型である。Divisors, Lemma
\ref{divisors-lemma-check-isomorphism-via-depth-and-ass} により、余次元 $\geq 2$ の点
$s \in S$ に写る任意の点 $x \in X$ に対して
$\text{depth}((f^*f_*\mathcal{E})_x) \geq 2$ を示せば十分である。$f$ は平坦であり、
$(f^*f_*\mathcal{E})_x = (f_*\mathcal{E})_s \otimes_{\mathcal{O}_{S, s}}
\mathcal{O}_{X, x}$
であるから、$\text{depth}((f_*\mathcal{E})_s) \geq 2$ を示せば十分である。Algebra, Lemma
\ref{algebra-lemma-apply-grothendieck} を参照せよ。$S$ は正規 Noetherian スキームで
$\dim(\mathcal{O}_{S, s}) \geq 2$ なので、
$\text{depth}(\mathcal{O}_{S, s}) \geq 2$ である。Properties, Lemma
\ref{properties-lemma-criterion-normal} を参照せよ。従って Divisors, Lemma
\ref{divisors-lemma-depth-pushforward} から求める結果を得る。
\end{proof}

\noindent
上の結果を用いて、次の驚くべき主張を証明できる。

\begin{theorem}
\label{flat-theorem-pullback-trivial-fibres}
$p$ を素数、$Y$ を $\mathbf{F}_p$ 上の準コンパクトかつ準分離なスキームとする。
$f : X \to Y$ を幾何学的連結ファイバーをもつ有限表示の固有全射とする。このとき関手
$$
\colim_F \textit{Vect}(Y) \longrightarrow \colim_F \textit{Vect}(X)
$$
は充満忠実であり、その本質像は次のように記述される。$\mathcal{E}$ を有限局所自由
$\mathcal{O}_X$-加群とする。すべての $y \in Y$ に対し、
$$
F^{n_y, *}\mathcal{E}|_{X_{y, red}}
\cong
\mathcal{O}_{X_{y, red}}^{\oplus r_y}
$$
となる整数 $n_y, r_y \geq 0$ が存在すると仮定する。このとき、ある $n \geq 0$ に対する
$n$ 次 Frobenius 冪引き戻し $F^{n, *}\mathcal{E}$ は、有限局所自由
$\mathcal{O}_Y$-加群の引き戻しである。
\end{theorem}

\begin{proof}
充満忠実性の証明。$Y$ 上のベクトル束は局所的に自明なので、これは写像
$$
\colim_F \Gamma(Y, \mathcal{O}_Y)
\longrightarrow
\colim_F \Gamma(X, \mathcal{O}_X)
$$
が全単射であるという主張に帰着する。$\{X \to Y\}$ は h 被覆なので、二つの写像
$$
\colim_F \Gamma(X, \mathcal{O}_X)
\longrightarrow
\colim_F \Gamma(X \times_Y X, \mathcal{O}_{X \times_Y X})
$$
が等しいことを示せれば、Lemma \ref{flat-lemma-h-sheaf-colim-F} から従う。
$g \in \Gamma(X, \mathcal{O}_X)$ をとり、$g$ の $X \times_Y X$ への二つの引き戻しを
$g_1$ と $g_2$ と書く。$X_{y, red}$ は幾何学的連結なので、
$H^0(X_{y, red}, \mathcal{O}_{X_{y, red}})$ は $\kappa(y)$ の純非分離拡大である。
Varieties, Lemma
\ref{varieties-lemma-proper-geometrically-reduced-global-sections} を参照せよ。従って
$y$ に依存しうるある $p$-冪 $q$ に対し、$g^q|_{X_{y, red}}$ は $\kappa(y)$ の元から来る。
従って $g_1^q$ と $g_2^q$ は、
$(X \times_Y X)_y = X_y \times_y X_y$ の任意の点の剰余体で同じ元に写る。
ゆえに $g_1 - g_2$ は $(X \times_Y X)_{red}$ 上で零に制限される。従ってある $n$ に対して
$(g_1 - g_2)^n = 0$ であり、望むようにこれを $p$-冪にとれる。

\medskip\noindent
本質像の記述。$\mathcal{E}$ を定理の主張のようにとる。まず Noetherian の場合に帰着する。

\medskip\noindent
$y \in Y$ を点とし、剰余体のスペクトルから $Y$ への射 $y \to Y$ とみなす。
$y \to Y$ を、$Y_i$ がアフィンである有限表示射 $Y_i \to Y$ のフィルター極限として
書ける。（自分で証明するのが最善だが、Limits, Lemma
\ref{limits-lemma-relative-approximation} and \ref{limits-lemma-limit-affine} からも形式的に
従う。）各 $i$ に対し $Z_i = Y_i \times_Y X$ とおく。このとき
$X_y = \lim Z_i$ かつ $X_{y, red} = \lim Z_{i, red}$ である。Limits, Lemma
\ref{limits-lemma-descend-modules-finite-presentation} により、
$F^{n_y, *}\mathcal{E}|_{Z_{i, red}} \cong
\mathcal{O}_{Z_{i, red}}^{\oplus r_y}$ となる $i$ をとれる。この $i$ を固定する。
$Z_{i, red} = \lim Z_{i, j}$ と書ける。ここで $Z_{i, j} \to Z_i$ は有限表示の厚化である
（Limits, Lemma
\ref{limits-lemma-closed-is-limit-closed-and-finite-presentation}）。先ほどと同じ補題を用いると、
$F^{n_y, *}\mathcal{E}|_{Z_{i, j}} \cong \mathcal{O}_{Z_{i, j}}^{\oplus r_y}$.
となる $j$ をとれる。従って各 $y \in Y$ に対し、像が $y$ を含む有限表示射
$Y_y \to Y$ と厚化 $Z_y \to Y_y \times_Y X$ で
$F^{n_y, *}\mathcal{E}|_{Z_y} \cong \mathcal{O}_{Z_y}^{\oplus r_y}$.
を満たすものが存在する。$Y_y \to Y$ の像は構成可能であることに注意する
（Morphisms, Lemma \ref{morphisms-lemma-chevalley}）。$Y$ は構成可能位相に関して
準コンパクトなので（Topology, Lemma
\ref{topology-lemma-constructible-hausdorff-quasi-compact} and Properties, Lemma
\ref{properties-lemma-quasi-compact-quasi-separated-spectral}）、有限個の有限表示射
$$
Y_1 \to Y,\ Y_2 \to Y,\ \ldots,\ Y_N \to Y
$$
で、集合論的に $Y = \bigcup \Im(Y_a \to Y)$ となり、各
$a \in \{1, \ldots, N\}$ に対して整数 $n_a, r_a \geq 0$ と有限表示の厚化
$Z_a \subset Y_a \times_Y X$ が存在して
$F^{n_a, *}\mathcal{E}|_{Z_a} \cong \mathcal{O}_{Z_a}^{\oplus r_a}$.
となるものが存在する。

\medskip\noindent
このように定式化すると、条件は絶対 Noetherian 近似に降下する。この段落は読み飛ばすことを
強く勧める。まず $Y = \lim_{i \in I} Y_i$ を、アフィン遷移射をもつ
$\mathbf{F}_p$ 上有限型のスキームの余フィルター極限として書く（Limits, Lemma
\ref{limits-lemma-relative-approximation}）。次に、$Y$ への基底変換が $f : X \to Y$ を
復元する固有射 $f_i : X_i \to Y_i$ があると仮定してよい。Limits, Lemma
\ref{limits-lemma-descend-finite-presentation} を参照せよ。$i$ を大きくすれば、$X$ への
引き戻しが $\mathcal{E}$ と同型となる有限局所自由 $\mathcal{O}_{X_i}$-加群
$\mathcal{E}_i$ が存在すると仮定してよい。Limits, Lemma
\ref{limits-lemma-descend-invertible-modules} を参照せよ。$0 \in I$ を選び、$f_0$ の幾何学的
ファイバーが連結となる構成可能部分集合を $E \subset Y_0$ と書く。More on Morphisms, Lemma
\ref{more-morphisms-lemma-nr-geom-connected-components-constructible} を参照せよ。
このとき $Y \to Y_0$ は $E$ に写る。More on Morphisms, Lemma
\ref{more-morphisms-lemma-base-change-fibres-geometrically-connected} を参照せよ。従って
$i \gg 0$ に対して $Y_i \to Y_0$ は $E$ に写る。Limits, Lemma
\ref{limits-lemma-limit-contained-in-constructible} を参照せよ。ゆえに $i \gg 0$ に対し
$f_i$ のファイバーは幾何学的連結である。Limits, Lemma
\ref{limits-lemma-descend-finite-presentation} により、十分大きい $i$ に対し、$Y$ への
基底変換が $Y_a \to Y$ を復元する有限型射 $Y_{i, a} \to Y_i$ を
$a  \in \{1, \ldots, N\}$ ごとにとれる。必要なら $i$ を大きくし、
$Y_a \times_Y X$ への基底変換が $Z_a$ を復元する厚化
$Z_{i, a} \subset Y_{i, a} \times_{Y_i} X_i$ をとれる（先ほどと同じ参照と Limits, Lemmas
\ref{limits-lemma-descend-closed-immersion-finite-presentation} and
\ref{limits-lemma-descend-surjective} を併用する）。$Z_a = \lim Z_{i, a}$ なので、$i$ を
大きくすれば
$F^{n_a, *}\mathcal{E}_i|_{Z_{i, a}} \cong
\mathcal{O}_{Z_{i, a}}^{\oplus r_a}$ と仮定してよい。Limits, Lemma
\ref{limits-lemma-descend-modules-finite-presentation} を参照せよ。最後にもう一度 $i$ を
大きくすれば、Limits, Lemma \ref{limits-lemma-descend-surjective} により
$\coprod Y_{i, a} \to Y_i$ が全射であると仮定してよい。この時点で $X_i \to Y_i$ と
$\mathcal{E}_i$ はすべての仮定を満たすので、$X_i \to Y_i$ と $\mathcal{E}_i$ に対して
結果を証明すれば十分である。

\medskip\noindent
$Y$ が $\mathbf{F}_p$ 上有限型であると仮定する。結果を証明するため $\dim(Y)$ に関する
帰納法を用いる。$\mathcal{E}$ が定める $\colim_F \textit{Vect}(X)$ の対象へ引き戻される
$\colim_F \textit{Vect}(Y)$ の対象を見つけようとしている。すでに証明した充満忠実性と
Proposition \ref{flat-proposition-h-descent-vector-bundles-p} により、$Y$ を h 被覆の要素で、
$X$ を対応する基底変換で置き換えた後に $\mathcal{E}$ の降下を構成すれば十分である。
従って $Y$ の次元を増加させない限り、$Y$ を $Y$ 上固有全射なスキームで置き換えてよい。
$T \subset T'$ が $\mathbf{F}_p$ 上有限型のスキームの厚化なら
$\colim_F \textit{Vect}(T) = \colim_F \textit{Vect}(T')$
である。実際、$\{T \to T'\}$ は $T \times_{T'} T = T$ を満たす h 被覆である。
$T' \to T$ が $\mathbf{F}_p$ 上有限型のスキームの普遍同相射なら
$\colim_F \textit{Vect}(T) = \colim_F \textit{Vect}(T')$
である。実際、
% P05-EN-CH38-0311: frozen-source pointer; see ERRATA_CANDIDATES.jsonl.
$\{T \to T'\}$ は対角射 $T \subset T \times_{T'} T$ が厚化となる h 被覆である。

\medskip\noindent
上の一般的注意を用い、$X$ をその被約化で置き換えて $X$ が被約であると仮定する。
Stein 分解 $X \to Y' \to Y$ を考える。More on Morphisms, Theorem
\ref{more-morphisms-theorem-stein-factorization-Noetherian} を参照せよ。このとき
$Y' \to Y$ は $\mathbf{F}_p$ 上有限型のスキームの普遍同相射である。上により $Y$ を
$Y'$ で置き換えてよい。従って $f_*\mathcal{O}_X = \mathcal{O}_Y$ かつ $Y$ は被約であると
仮定してよい。これにより次の段落で扱う場合に帰着する。

\medskip\noindent
$Y$ が被約で、$Y$ の稠密開部分スキーム上
$f_*\mathcal{O}_X = \mathcal{O}_Y$ であると仮定する。このとき $X \to Y$ はある稠密開部分
$V \subset Y$ 上で平坦である。Morphisms, Proposition
\ref{morphisms-proposition-generic-flatness-reduced} を参照せよ。Lemma
\ref{flat-lemma-flat-after-blowing-up} により、$X$ の厳密変換 $X'$ が $Y'$ 上平坦となる
$V$-許容吹き上げ $Y' \to Y$ が存在する。$Y$ と $Y'$ は共通の稠密開部分スキームをもつので
$\dim(Y') = \dim(Y)$ であることに注意する。More on Morphisms, Lemma
\ref{more-morphisms-lemma-proper-flat-nr-geom-conn-comps-lower-semicontinuous} と
$V \subset Y'$ が稠密であることにより、$f' : X' \to Y'$ のすべてのファイバーは幾何学的
連結である。依然として $(f'_*\mathcal{O}_{X'})|_V = \mathcal{O}_V$ である。
$$
Y' \times_Y X = X' \cup E \times_Y X
$$
と書く。ここで $E \subset Y'$ は吹き上げの例外因子である。上の一般的注意により、
$Y' \times_Y X \to Y'$ と $\mathcal{E}$ の $Y' \times_Y X$ への制限について存在を
証明すれば十分である。$X'$ への $\mathcal{E}$ の制限（余極限圏の対象とみなす）へ
引き戻される対象 $\xi'$ を $\colim_F \textit{Vect}(Y')$ において見つけたとする。
$\dim(Y)$ に関する帰納法により、$\mathcal{E}$ の $E \times_Y X$ への制限へ引き戻される
対象 $\xi''$ を $\colim_F \textit{Vect}(E)$ においてとれる。すると充満忠実性により、
$E \times_{Y'} X'$ への $\mathcal{E}$ の制限との所与の同一視と整合する一意な同型
$\xi'|_E \to \xi''$ が定まる。
$$
\{E \times_Y X \to Y' \times_Y X, X' \to Y' \times_Y X\}
$$
は二つの閉埋め込みで与えられる h 被覆であり、
$$
(E \times_Y X) \times_{(Y' \times_Y X)} X' = E \times_{Y'} X'
$$
を満たすので、$\xi'$ は $\mathcal{E}$ の $Y' \times_Y X$ への制限へ引き戻される。
従って $\xi'$ を見つければ十分であり、次の段落で扱う場合に帰着する。

\medskip\noindent
$Y$ が被約で、$f$ が平坦であり、$Y$ の稠密開部分スキーム上
$f_*\mathcal{O}_X = \mathcal{O}_Y$ であると仮定する。この場合は正規化
$Y^\nu \to Y$ を考える（Morphisms, Section
\ref{morphisms-section-normalization}）。これは稠密開上で同型となる有限全射である
（Morphisms, Lemma \ref{morphisms-lemma-nagata-normalization} and
\ref{morphisms-lemma-ubiquity-nagata}）。従って一般的注意により、$Y$ を $Y^\nu$ で、
$X$ を $Y^\nu \times_Y X$ で置き換えてよい。この置換の後
$\mathcal{O}_Y = f_*\mathcal{O}_X$ である（この場合 Stein 分解は同型でなければならない
からである。小さな詳細は省略する）。

\medskip\noindent
$Y$ が正規 Noetherian スキームで、$f$ が平坦かつ
$f_*\mathcal{O}_X = \mathcal{O}_Y$ であると仮定する。$\mathcal{E}$ を適当な Frobenius 冪
引き戻しで置き換え、$Y$ の既約成分の生成点における $f$ のスキーム論的ファイバー上で
$\mathcal{E}$ が自明であると仮定してよい（上で見たように
$\colim_F \textit{Vect}(-)$ は普遍同相射上で同値だからである）。上の議論
（Noetherian の場合への帰着）と同様に、$\mathcal{E}|_{f^{-1}(V)}$ が自由となる稠密開部分
$V \subset Y$ が存在する。$Y = V \amalg Z$ が集合論的に成り立つ閉部分スキーム
$Z \subset Y$ をとる。$z_1, \ldots, z_t \in Z$ を、余次元 $1$ である $Z$ の既約成分の
生成点とする。このとき $A_i = \mathcal{O}_{Y, z_i}$ は離散付値環である。
$n_i$ を、$A_i$ 上のスキーム $X_{A_i}$ に対して Lemma
\ref{flat-lemma-trivial-fibres-dvr} で得た整数とする。$\mathcal{E}$ を適当な Frobenius 冪
引き戻しで置き換え、$X_{A_i/\mathfrak m_i^{n_i}}$ 上で $\mathcal{E}$ が自由であると
仮定してよい（上で見たように余極限圏は普遍同相射の下で不変だからである）。すると
Lemma \ref{flat-lemma-trivial-fibres-dvr} により $\mathcal{E}$ は $X_{A_i}$ 上自由である。
最後に Lemma \ref{flat-lemma-trivial-over-dvrs} を適用して結論を得る。
\end{proof}










\section{複体の吹き上げ}
\label{flat-section-blowup-complexes}

\noindent
この節では、吹き上げ後の導来圏の完全対象の標準形を求める。

\begin{lemma}
\label{flat-lemma-fitting-ideals-complex}
$X$ をスキーム、$E \in D(\mathcal{O}_X)$ を擬連接とする。任意の
$p, k \in \mathbf{Z}$ に対し、有限型準連接イデアル層
$\text{Fit}_{p, k}(E) \subset \mathcal{O}_X$ で次の性質をもつものが存在する。
% P05-EN-CH38-0103: source 12965 has the article error ``an finite type''.
$U \subset X$ が開で、$E|_U$ が
$$
\ldots \to
\mathcal{O}_U^{\oplus n_{b - 2}}
\xrightarrow{d_{b - 2}}
\mathcal{O}_U^{\oplus n_{b - 1}}
\xrightarrow{d_{b - 1}}
\mathcal{O}_U^{\oplus n_b} \to 0 \to \ldots
$$
と同型なら、制限 $\text{Fit}_{p, k}(E)|_U$ は $d_p$ の行列のサイズ
$$
- k + n_{p + 1} - n_{p + 2} + \ldots + (-1)^{b - p + 1} n_b
$$
の小行列式で生成される。規約として、$r \leq 0$ なら $r \times r$-小行列式で生成される
イデアルを $\mathcal{O}_U$ とし、$r > \min(n_p, n_{p + 1})$ なら
$r \times r$-小行列式で生成されるイデアルを零とする。
\end{lemma}

\begin{proof}
$X$ 上局所的に $E$ は補題の主張の形をもつことに注意する。More on Algebra, Section
\ref{more-algebra-section-pseudo-coherent}, Cohomology, Section
\ref{cohomology-section-pseudo-coherent}, and Derived Categories of Schemes, Section
\ref{perfect-section-spell-out} を参照せよ。従って小行列式のイデアルが選んだ表示に依存しない
ことを示せば十分であり、これは局所環で確認すればよい。局所環
$(R, \mathfrak m, \kappa)$ 上の上に有界な複体
$$
F^\bullet :
\ldots \to
R^{\oplus n_{b - 2}}
\xrightarrow{d_{b - 2}}
R^{\oplus n_{b - 1}}
\xrightarrow{d_{b - 1}}
R^{\oplus n_b} \to 0 \to \ldots
$$
を考える。
% P05-EN-CH38-0104: source 13005--13006 swaps the Fit indices and reverses the stated minor-size formula.
$\text{Fit}_{k, p}(F^\bullet) \subset R$ を、$d_p$ の行列におけるサイズ
$k - n_{p + 1} + n_{p + 2} - \ldots + (-1)^{b - p} n_b$ の小行列式で生成される
イデアルと書く。$F^\bullet$ のある微分のある行列係数が可逆であるとする。
$d_i$ が可逆な行列係数をもつような最大の整数 $i$ を選ぶ。Algebra, Lemma
\ref{algebra-lemma-add-trivial-complex} により、複体 $F^\bullet$ は、次数 $i$ と $i + 1$ に
非零項をもつ自明な複体 $\ldots \to 0 \to R \to R \to 0 \to \ldots$ と複体
$(F')^\bullet$ との直和に同型である。
$\text{Fit}_{p, k}(F^\bullet) = \text{Fit}_{p, k}((F')^\bullet)$
であることの確認は読者に委ねる。ここで小行列式のサイズの公式を用いる。
$(F')^\bullet$ の別の微分が可逆な行列係数をもつなら、この操作を再び行い、以下同様にする。
このように続けると、最終的にすべての微分の行列係数が $\mathfrak m$ に属する複体
$(F^\infty)^\bullet$ に到達する。ここでは無限回の操作が必要になりうるが、任意の切断次数に
対し、次数 $\geq $ 当該切断次数の複体に影響する操作は有限個だけである。
% P05-EN-CH38-0105: source 13023 leaves the cutoff outside the isolated comparison-symbol math span.
従って「極限」$(F^\infty)^\bullet$ は有限自由加群からなるよく定義された上に有界な複体で、
開始時の複体への擬同型 $(F^\infty)^\bullet \to F^\bullet$ を備え、
$\text{Fit}_{p, k}(F^\bullet) = \text{Fit}_{p, k}((F^\infty)^\bullet)$
である。More on Algebra, Lemma
\ref{more-algebra-lemma-lift-pseudo-coherent-from-residue-field} により複体
$(F^\infty)^\bullet$ は同型を除いて一意なので、証明が完了する。
\end{proof}

\begin{lemma}
\label{flat-lemma-blowup-complex}
$X$ をスキーム、$E \in D(\mathcal{O}_X)$ を完全とする。$U \subset X$ をスキーム論的に
稠密な開部分スキームで、すべての $i \in \mathbf{Z}$ に対して $H^i(E|_U)$ が定数階数
$r_i$ の有限局所自由加群となるものとする。このとき $U$-許容吹き上げ
$b : X' \to X$ で、すべての $i \in \mathbf{Z}$ に対して $H^i(Lb^*E)$ が Tor 次元
$\leq 1$ の完全 $\mathcal{O}_{X'}$-加群となるものが存在する。
% P05-EN-CH38-0106: source 13038, 13084, and 13129 use ``scheme theoretically'' without the compound-adverb hyphen.
\end{lemma}

\begin{proof}
吹き上げをアフィン局所的に構成し調べる。すなわち $V \subset X$ をアフィン開部分スキームで、
$E|_V$ が複体
$$
\mathcal{O}_V^{\oplus n_a} \xrightarrow{d_a}
\ldots \xrightarrow{d_{b - 1}} \mathcal{O}_V^{\oplus n_b}
$$
で表示されるものとする。
$k_ i = r_{i + 1} - r_{i + 2} + \ldots + (-1)^{b - i + 1}r_b$ とおく。
省略する計算により、$U \cap V$ 上で $d_i$ の階数は
% P05-EN-CH38-0107: source 13054 has a stray space in the subscript token k_ i.
% P05-EN-CH38-0108: source 13055 has singular ``A computation'' with bare ``show''.
$$
\rho_i = - k_i + n_{i + 1} - n_{i + 2} + \ldots + (-1)^{b - i + 1}n_b
$$
である。ここで階数とは、$d_i$ の余核が階数 $n_{i + 1} - \rho_i$ の有限局所自由加群で
あるという意味である。$\mathcal{I}_i \subset \mathcal{O}_V$ を $d_i$ の行列の
$\rho_i \times \rho_i$-小行列式で生成されるイデアルとする。

\medskip\noindent
一方、Lemma \ref{flat-lemma-fitting-ideals-complex} と比較すると、イデアル $\mathcal{I}_i$ は
大域イデアル $\text{Fit}_{i, k_i}(E)$ に対応する。このイデアルは $E|_V$ を表示する複体の
選択に依存しないことを示した。他方、$\mathcal{I}_i$ は $\Coker(d_i)$ の
第 $(n_{i + 1} - \rho_i)$ Fitting イデアルである。この点を覚えておいてほしい。

\medskip\noindent
$b : X' \to X$ をイデアル $\text{Fit}_{i, k_i}(E)$ の積に沿う吹き上げとする。$X$ 上
局所的にこれらのイデアルのほとんどすべては単位イデアルに等しいので、これは意味をもつ
（上を参照）。この吹き上げは各イデアル $\text{Fit}_{i, k_i}(E)$ に沿う吹き上げ
$b_i : X'_i \to X$ を支配する。Divisors, Lemma
\ref{divisors-lemma-blowing-up-two-ideals} を参照せよ。Divisors, Lemma
\ref{divisors-lemma-blowup-fitting-ideal} により各 $b_i$ は $U$-許容吹き上げである。
従って $b$ も $U$-許容吹き上げである（小さな詳細は省略する。Divisors, Lemma
\ref{divisors-lemma-dominate-admissible-blowups} の証明と比較せよ）。最後に $U$ は依然として
$X'$ のスキーム論的に稠密な開部分スキームである。従って $X$ を $X'$ で置き換えると、
次の段落で扱う状況になる。

\medskip\noindent
すべての $i$ に対して $\text{Fit}_{i, k_i}(E)$ が可逆イデアルであると仮定する。上のような
アフィン開 $V$ と $E|_V$ を表示する有限自由加群の複体を選ぶ。Divisors, Lemma
\ref{divisors-lemma-blowup-fitting-ideal} により $\Coker(d_i)$ の Tor 次元は $\leq 1$ である。
従って $\Im(d_i)$ は、有限局所自由加群から Tor 次元 $\leq 1$ の有限表示加群への写像の
核として有限局所自由である。ゆえに $\Ker(d_i)$ も有限局所自由である（同じ議論）。
従って短完全列
$$
0 \to \Im(d_{i - 1}) \to \Ker(d_i) \to H^i(E)|_V \to 0
$$
から求めることが従い、証明が完了する。
\end{proof}

\begin{lemma}
\label{flat-lemma-blowup-complex-integral}
$X$ を整スキーム、$E \in D(\mathcal{O}_X)$ を完全とする。このとき空でない開
$U \subset X$ で、すべての $i \in \mathbf{Z}$ に対して $H^i(E|_U)$ が定数階数 $r_i$ の
有限局所自由加群となるものが存在し、さらに $U$-許容吹き上げ $b : X' \to X$ で、
すべての $i \in \mathbf{Z}$ に対して $H^i(Lb^*E)$ が Tor 次元 $\leq 1$ の完全
$\mathcal{O}_{X'}$-加群となるものが存在する。
\end{lemma}

\begin{proof}
$U$ の存在について自分自身の証明を見つけることを強く勧める。$\eta \in X$ を生成点とする。
$E$ の $\eta$ への制限は $D(\kappa(\eta))$ において、微分が零である有限次元ベクトル空間の
有限複体 $V^\bullet$ と同型である。$r_i = \dim_{\kappa(\eta)} V^i$ とおく。各項が
$\mathcal{O}_X^{\oplus r_i}$ で微分が零である複体によって表示される
$D(\mathcal{O}_X)$ の完全対象 $E'$ は、$\eta$ へ引き戻すと $E$ と同型になる。従って
Derived Categories of Schemes, Lemma
\ref{perfect-lemma-descend-relatively-perfect} により、$E|_U$ と $E'|_U$ が同型となる
$\eta$ の開近傍 $U$ が存在する。これで第一の主張が証明された。整スキーム $X$ の任意の
空でない開はスキーム論的に稠密なので、第二の主張は第一の主張と Lemma
\ref{flat-lemma-blowup-complex} から従う。
\end{proof}

\begin{remark}
\label{flat-remark-when-you-have-a-complex}
$X$ をスキームとし、$E \in D(\mathcal{O}_X)$ を、すべての $i \in \mathbf{Z}$ に対して
$H^i(E)$ が Tor 次元 $\leq 1$ の完全 $\mathcal{O}_X$-加群となる完全対象とする。
この性質により、$E$ に関する問題を $H^i(E)$ に関する問題へ帰着できることがある。
例えば
$$
\mathcal{E}^a \xrightarrow{d^a} \ldots
\xrightarrow{d^{b - 2}} \mathcal{E}^{b - 1}
\xrightarrow{d^{b - 1}} \mathcal{E}^b
$$
を $E$ を表示する有限局所自由 $\mathcal{O}_X$-加群の有界複体とする。このときすべての
$i$ に対して $\Im(d^i)$ と $\Ker(d^i)$ は有限局所自由 $\mathcal{O}_X$-加群である。
実際、$i$ より大きいすべての添字についてこれが分かっていると帰納的に仮定する。
まず短完全列
$$
0 \to \Im(d^i) \to \Ker(d^{i + 1}) \to H^{i + 1}(E) \to 0
$$
と、$H^{i + 1}(E)$ が Tor 次元 $\leq 1$ の完全加群であるという仮定を用いて、
$\Im(d^i)$ が有限局所自由であると結論できる。同じ議論を短完全列
$$
0 \to \Ker(d^i) \to \mathcal{E}^i \to \Im(d^i) \to 0
$$
に再び用いると、$\Ker(d^i)$ が有限局所自由であることを得る。従って区別三角形
$$
\tau_{\leq k - 1}E \to \tau_{\leq k}E \to H^k(E)[-k] \to
(\tau_{\leq k - 1}E)[1]
$$
は、有限局所自由加群の有界複体からなる次の短完全列によって表示される。
$$
\begin{matrix}
& &
& &
& &
0 \\
& &
& &
& &
\downarrow \\
\mathcal{E}^a & \to &
\ldots & \to &
\mathcal{E}^{k - 2} & \to &
\Ker(d^{k - 1}) \\
\downarrow & &
& &
\downarrow & &
\downarrow \\
\mathcal{E}^a & \to &
\ldots & \to &
\mathcal{E}^{k - 2} & \to &
\mathcal{E}^{k - 1} & \to &
\Ker(d^k) \\
& &
& &
& &
\downarrow & &
\downarrow \\
& &
& &
& &
\Im(d^{k - 1}) & \to &
\Ker(d^k) \\
& &
& &
& &
\downarrow \\
& &
& &
& &
0
\end{matrix}
$$
ここで複体は行であり、「明らかな」零は表示から省略している。
\end{remark}





\section{完全加群の吹き上げ}
\label{flat-section-blowup-perfect-modules}

\noindent
この節では、Tor 次元 $\leq 1$ の完全加群について、吹き上げ後の標準形を求める。
得られるのは部分的な結果にとどまる。

\begin{lemma}
\label{flat-lemma-blowup-pd1-derived}
$X$ をスキームとし、$\mathcal{F}$ を Tor 次元 $\leq 1$ の完全
$\mathcal{O}_X$-加群とする。任意の吹き上げ $b : X' \to X$ に対して
$Lb^*\mathcal{F} = b^*\mathcal{F}$ であり、$b^*\mathcal{F}$ は Tor 次元
$\leq 1$ の完全 $\mathcal{O}_X$-加群である。
% P05-EN-CH38-0109: source 13228 and 13256 retain \mathcal{O}_X although the modules live on X'.
\end{lemma}

\begin{proof}
$X = \Spec(A)$ はアフィンであり、$\mathcal{F}$ に対応する $A$-加群 $M$ は表示
$$
0 \to A^{\oplus m} \to A^{\oplus n} \to M \to 0
$$
をもつと仮定してよい。$I \subset A$ をイデアル、$a \in I$ とする。アフィン吹き上げ代数
$A[\frac{I}{a}]$ は $A_a$ の部分環であることを想起する。局所化は完全なので
$A_a^{\oplus m} \to A_a^{\oplus n}$ は単射である。従って
$A[\frac{I}{a}]^{\oplus m} \to A[\frac{I}{a}]^{\oplus n}$ も単射である。
これで補題が証明された。
\end{proof}

\begin{lemma}
\label{flat-lemma-blowup-pd1}
$X$ をスキーム、$\mathcal{F}$ を Tor 次元 $\leq 1$ の完全
$\mathcal{O}_X$-加群とする。$U \subset X$ をスキーム論的に稠密な開集合で、
$\mathcal{F}|_U$ が定数階数 $r$ の有限局所自由加群となるものとする。このとき
$U$-許容吹き上げ $b : X' \to X$ と標準的短完全列
% P05-EN-CH38-0110: source 13248 and 13370 add two further unhyphenated ``scheme theoretically dense'' loci.
$$
0 \to \mathcal{K} \to b^*\mathcal{F} \to \mathcal{Q} \to 0
$$
が存在する。ここで $\mathcal{Q}$ は階数 $r$ の有限局所自由加群であり、
$\mathcal{K}$ は Tor 次元 $\leq 1$ の完全 $\mathcal{O}_X$-加群で、その $U$ への制限は零である。
\end{lemma}

\begin{proof}
Divisors, Lemma \ref{divisors-lemma-blowup-fitting-ideal} と Lemma
\ref{flat-lemma-blowup-pd1-derived} を合わせればよい。
\end{proof}

\begin{lemma}
\label{flat-lemma-canonical-blowup-torsion-pd1}
$X$ をスキーム、$\mathcal{F}$ を Tor 次元 $\leq 1$ の完全
$\mathcal{O}_X$-加群とする。$U \subset X$ を $\mathcal{F}|_U = 0$ となる開集合とする。
このとき $U$-許容吹き上げ
$$
b : X' \to X
$$
が存在して、$\mathcal{F}' = b^*\mathcal{F}$ は二つの標準的な局所有限フィルトレーション
$$
0 = F^0 \subset F^1 \subset F^2 \subset \ldots \subset \mathcal{F}'
\quad\text{および}\quad
\mathcal{F}' = F_1 \supset F_2 \supset F_3 \supset \ldots \supset 0
$$
を備える。さらに各 $n \geq 1$ に対し、有効 Cartier 因子 $D_n \subset X'$ が存在して、
$$
F^i/F^{i - 1}
\quad\text{および}\quad
F_i/F_{i + 1}
$$
は $D_i$ 上で階数 $i$ の有限局所自由加群となる。
\end{lemma}

\begin{proof}
$V \subset X$ をアフィン開集合で、ある $n$ と行列 $A$ に対して表示
$$
0 \to \mathcal{O}_V^{\oplus n} \xrightarrow{A} \mathcal{O}_V^{\oplus n} \to
\mathcal{F} \to 0
$$
が存在するものを選ぶ。吹き上げの中心に用いるイデアルは、$k \geq 0$ に対する Fitting
イデアル $\text{Fit}_k(\mathcal{F})$ の積である。上のアフィンな状況では $k > n$ に対して
$\text{Fit}_k(\mathcal{F})|_V = \mathcal{O}_V$ なので、これは意味をもつ。これが
$U$-許容吹き上げであることは明らかである。Divisors, Lemma
\ref{divisors-lemma-blowing-up-two-ideals} により、$X'$ 上でイデアル
$\text{Fit}_k(\mathcal{F})$ は可逆である。従って次の段落で論じる場合に帰着する。

\medskip\noindent
$k \geq 0$ に対して $\text{Fit}_k(\mathcal{F})$ が可逆イデアルであると仮定する。
$E_k \subset X$ を $\text{Fit}_k(\mathcal{F})$ が定める有効 Cartier 因子（$k \geq 0$）とすると、
補題の主張に現れる有効 Cartier 因子 $D_k$ は
$$
E_k = D_{k + 1} + 2 D_{k + 2} + 3 D_{k + 3} + \ldots
$$
を満たす。族 $D_k$ は局所有限になるので、これは意味をもつ。さらにこの等式は有効
Cartier 因子 $D_k$ を一意に定めるから、$D_k$ を局所的に構成すれば十分である。

\medskip\noindent
$V \subset X$ をアフィン開集合とし、上のような $\mathcal{F}|_V$ の表示を選ぶ。
表示に現れる整数 $n$ に関する帰納法で因子とフィルトレーションを構成する。$k > n$ に対して
$D_k|_V = \emptyset$ とおき、$D_n|V = E_{n - 1}|_V$ とおく。
% P05-EN-CH38-0111: source 13323 omits the restriction subscript in D_n|V.
$V$ を縮小して、$\text{Fit}_{n - 1}(\mathcal{F})|_V$ が単一の非零因子
$f \in \Gamma(V, \mathcal{O}_V)$ で生成されると仮定してよい。
$\text{Fit}_{n - 1}(\mathcal{F})|_V$ は $A$ の成分で生成されるイデアルなので、
$\Gamma(V, \mathcal{O}_V)$ の行列 $A'$ で $A = fA'$ となるものが存在する。短完全列
$$
0 \to \mathcal{O}_V^{\oplus n} \xrightarrow{A'} \mathcal{O}_V^{\oplus n} \to
\mathcal{F}' \to 0
$$
によって $V$ 上の $\mathcal{F}'$ を定める。$A'$ の成分は
$\Gamma(V, \mathcal{O}_V)$ の単位イデアルを生成するので、$\mathcal{F}'$ は $V$ 上局所的に
$n$ を $1$ 減らした表示をもつ。Algebra, Lemma
\ref{algebra-lemma-add-trivial-complex} を参照せよ。さらに
$f^{n - k}\text{Fit}_k(\mathcal{F}') = \text{Fit}_k(\mathcal{F})|_V$
が $k = 0, \ldots, n$ に対して成り立つことに注意する。従ってすべての $k$ に対して
$\text{Fit}_k(\mathcal{F}')$ は可逆イデアルである。帰納法により、有効 Cartier 因子
$D'_k \subset V$ が存在し、$\mathcal{F}'$ は補題の主張のような二つの標準的フィルトレーションを
もつ。そこで $k = 1, \ldots, n - 1$ に対して $D_k|_V = D'_k$ とおく。この選択のもとで、
上に表示した有効 Cartier 因子間の等式が成り立つことに注意する。最後にフィルトレーションを
構成する。すなわち短完全列
$$
0 \to
\mathcal{O}_{D_n \cap V}^{\oplus n} \to
\mathcal{F} \to \mathcal{F}' \to 0
\quad\text{および}\quad
0 \to
\mathcal{F}' \to \mathcal{F} \to
\mathcal{O}_{D_n \cap V}^{\oplus n} \to 0
$$
があり、これは $A$ の二つの分解 $A = A'f = f A'$ から得られる。これらの列は標準的である。
実際、第一の列の部分加群は $\Ker(f : \mathcal{F} \to \mathcal{F})$ であり、第二の列の商加群は
$\Coker(f : \mathcal{F} \to \mathcal{F})$ である。
\end{proof}

\begin{lemma}
\label{flat-lemma-blowup-map-pd1}
$X$ をスキームとする。$\varphi : \mathcal{F} \to \mathcal{G}$ を、Tor 次元 $\leq 1$ の
完全 $\mathcal{O}_X$-加群の間の準同型とする。$U \subset X$ をスキーム論的に稠密な開集合で、
$\mathcal{F}|_U = 0$ かつ $\mathcal{G}|_U = 0$ となるものとする。このとき $U$-許容吹き上げ
$b : X' \to X$ で、$b^*\varphi$ の核、像、余核が Tor 次元 $\leq 1$ の完全
$\mathcal{O}_{X'}$-加群となるものが存在する。
% P05-EN-CH38-0112: source 13369 misspells ``homomorphism'' as ``homorphism''.
\end{lemma}

\begin{proof}
仮定により、$D(\mathcal{O}_X)$ の対象 $(\mathcal{F} \to \mathcal{G})$ は完全である。
従って Lemmas \ref{flat-lemma-blowup-complex} および \ref{flat-lemma-blowup-pd1-derived}
（後者は引き戻し後の複体の形を調べるために用いる）により、余核と核の双方に対して機能する
$U$-許容吹き上げが得られる。像は余核への写像の核だから、やはり Tor 次元 $\leq 1$ の
完全加群である。
\end{proof}











\section{Berthelot と Ogus が導入した作用素}
\label{flat-section-eta}

\noindent
まず Cohomology, Section \ref{cohomology-section-eta} を読まれたい。

\medskip\noindent
$X$ をスキーム、$D \subset X$ を有効 Cartier 因子とする。
$\mathcal{I} = \mathcal{I}_D \subset \mathcal{O}_X$ を $D$ のイデアル層とする。
Divisors, Section \ref{divisors-section-effective-Cartier-invertible} を参照せよ。
Cohomology, Section \ref{cohomology-section-eta} の議論を $X$ と $\mathcal{I}$ に
適用できることは明らかである。

\begin{lemma}
\label{flat-lemma-eta-stalks}
$X$ をスキーム、$D \subset X$ をイデアル層 $\mathcal{I} \subset \mathcal{O}_X$ をもつ
有効 Cartier 因子とする。$\mathcal{F}^\bullet$ を準連接 $\mathcal{O}_X$-加群の複体で、
すべての $i$ に対して $\mathcal{F}^i$ が $\mathcal{I}$-捩れなしとなるものとする。
このとき $\eta_\mathcal{I}\mathcal{F}^\bullet$ は準連接 $\mathcal{O}_X$-加群の複体である。
さらに、$U = \Spec(A) \subset X$ がアフィン開集合で $D \cap U = V(f)$ なら、
$\eta_f(\mathcal{F}^\bullet(U))$ は $(\eta_\mathcal{I}\mathcal{F}^\bullet)(U)$ と
標準的に同型である。
\end{lemma}

\begin{proof}
省略する。
\end{proof}

\begin{lemma}
\label{flat-lemma-Leta}
$X$ をスキーム、$D \subset X$ をイデアル層 $\mathcal{I} \subset \mathcal{O}_X$ をもつ
有効 Cartier 因子とする。Cohomology, Lemma \ref{cohomology-lemma-Leta} の函子
$L\eta_\mathcal{I} : D(\mathcal{O}_X) \to D(\mathcal{O}_X)$ は
$D_\QCoh(\mathcal{O}_X)$ をそれ自身へ写す。さらに $X = \Spec(A)$ がアフィンで
$D = V(f)$ なら、More on Algebra, Lemma \ref{more-algebra-lemma-Leta} で定義された
$D(A)$ 上の函子 $L\eta_f$ と $D_\QCoh(\mathcal{O}_X)$ 上の函子
$L\eta_\mathcal{I}$ は、Derived Categories of Schemes, Lemma
\ref{perfect-lemma-affine-compare-bounded} の同値を通じて対応する。
\end{lemma}

\begin{proof}
省略する。
\end{proof}





\section{複体の吹き上げ、II}
\label{flat-section-blowup-complexes-II}

\noindent
この節の内容は、Section \ref{flat-section-blowup-complexes-III} で Macpherson の
グラフ構成の一つの変種を構成するために用いる。
 
\begin{situation}
\label{flat-situation-complex-and-divisor}
ここで $X$ はスキーム、$D \subset X$ はイデアル層
$\mathcal{I} \subset \mathcal{O}_X$ をもつ有効 Cartier 因子、
$M$ は $D(\mathcal{O}_X)$ の完全対象である。
\end{situation}

\noindent
$(X, D, M)$ を Situation \ref{flat-situation-complex-and-divisor} のような三つ組とする。
次を満たすアフィン開集合 $U = \Spec(A) \subset X$ を考える。
\begin{enumerate}
\item ある非零因子 $f \in A$ に対して $D \cap U = V(f)$ である。
\item $M|_U$ を表示する有限自由 $A$-加群の有界複体 $M^\bullet$ が存在する
（Derived Categories of Schemes, Lemma
\ref{perfect-lemma-affine-compare-bounded} の同値を通じて）。
\end{enumerate}
$(U, A, f, M^\bullet)$ を $(X, D, M)$ の{\it アフィンチャート}と呼ぶ。
More on Algebra, Section \ref{more-algebra-section-perfect-eta} で定義されるイデアル
$I_i(M^\bullet, f) \subset A$ を考える。各 $x \in D$ に対して、$x \in U$ であり、
すべての $i \in \mathbf{Z}$ に対して $I_i(M^\bullet, f)$ が単項イデアルとなる
アフィンチャート $(U, A, f, M^\bullet)$ が存在するとき、$(X, S, M)$ を
{\it 良い三つ組}と呼ぶ。
% P05-EN-CH38-0113: source 13482 changes D to the undefined S in the good-triple notation.

\begin{lemma}
\label{flat-lemma-pullback-triple-ideals-good}
Situation \ref{flat-situation-complex-and-divisor} において、$h : Y \to X$ を、$D$ の引き戻し
$E = h^{-1}D$ が定義されるスキームの射とする（Divisors, Definition
\ref{divisors-definition-pullback-effective-Cartier-divisor}）。
$(U, A, f, M^\bullet)$ を $(X, D, M)$ のアフィンチャートとし、
$V = \Spec(B) \subset Y$ を $h(V) \subset U$ となるアフィン開集合とする。
$g \in B$ を $f \in A$ の像と書く。このとき次が成り立つ。
% P05-EN-CH38-0114: source 13493--13494 uses ``Let ... is'' twice instead of ``Let ... be''.
\begin{enumerate}
\item $(V, B, g, M^\bullet \otimes_A B)$ は $(Y, E, Lh^*M)$ のアフィンチャートである。
\item $B$ において $I_i(M^\bullet, f)B = I_i(M^\bullet \otimes_A B, g)$ である。
\item $(X, D, M)$ が良い三つ組なら、$(Y, E, Lh^*M)$ も良い三つ組である。
\end{enumerate}
\end{lemma}

\begin{proof}
第一の主張は次の観察から従う。$g$ は $E \cap V \subset V$ を定義する $B$ の非零因子であり、
$M^\bullet \otimes_A B$ は $M^\bullet \otimes_A^\mathbf{L} B$ を表示するので、
Derived Categories of Schemes, Lemma
\ref{perfect-lemma-quasi-coherence-pullback} により $M$ の $V$ への引き戻しを表示する。
(2) は (1) と More on Algebra, Lemma
\ref{more-algebra-lemma-eta-base-change-pre} から従う。さらに More on Algebra, Lemma
\ref{more-algebra-lemma-eta-base-change-pre} と合わせると、補題の第二の主張が成り立つと分かる。
% P05-EN-CH38-0115: source 13512--13516 cites the same lemma twice, concludes (2) twice, and never proves item (3).
\end{proof}

\begin{lemma}
\label{flat-lemma-good-triple}
$X, D, \mathcal{I}, M$ を Situation \ref{flat-situation-complex-and-divisor} のようなものとする。
$(X, D, M)$ が良い三つ組なら、$L\eta_\mathcal{I}M$ は
$D(\mathcal{O}_X)$ の完全対象である。
\end{lemma}

\begin{proof}
More on Algebra, Lemma \ref{more-algebra-lemma-eta-locally-free} の翻訳である。
この翻訳には Lemma \ref{flat-lemma-Leta} を用いる。
\end{proof}

\begin{lemma}
\label{flat-lemma-good-triple-bdd-loc-free}
$X, D, \mathcal{I}, M$ を Situation \ref{flat-situation-complex-and-divisor} のようなものとし、
$(X, D, M)$ が良い三つ組であると仮定する。$M$ を表示する有限局所自由
$\mathcal{O}_X$-加群の局所有界複体 $\mathcal{M}^\bullet$ が存在するなら、
$L\eta_\mathcal{I}M$ を表示する有限局所自由 $\mathcal{O}_{X'}$-加群の局所有界複体
$\mathcal{Q}^\bullet$ が存在する。
% P05-EN-CH38-0116: source 13542 uses the undefined X'; the eta object still lives on X.
\end{lemma}

\begin{proof}
Cohomology, Lemma \ref{cohomology-lemma-Leta} により、複体
$\mathcal{Q}^\bullet = \eta_\mathcal{I}\mathcal{M}^\bullet$ は
$L\eta_\mathcal{I}M$ を表示する。この複体が局所有界で有限局所自由加群からなることは、
アフィン局所的に確認すればよい。その場合、有界性は明らかである。
$(X, D, M)$ のアフィンチャート $(U, A, f, M^\bullet)$ で、イデアル
$I_i(M^\bullet, f)$ が単項であり、各 $i$ に対して $\mathcal{M}^i|_U$ が有限自由となるものを
選ぶ。$(X, D, M)$ が良い三つ組であるという仮定により、これは可能である。
$N^i = \Gamma(U, \mathcal{M}^i|_U)$ と書けば、有限自由 $A$-加群の有界複体
$N^\bullet$ が得られ、これは複体 $M^\bullet$ と同じ $D(A)$ の対象を表示する
（Derived Categories of Schemes, Lemma
\ref{perfect-lemma-affine-compare-bounded} による）。More on Algebra, Lemma
\ref{more-algebra-lemma-ideal-well-defined} により、すべての $i$ に対して
$I_i(N^\bullet, f)$ は単項イデアルである。従って複体 $\eta_fN^\bullet$ は
有限局所自由 $A$-加群の有界複体である。Lemma \ref{flat-lemma-eta-stalks} により
$\mathcal{Q}^i|_U$ は $\eta_fN^i$ に対応する準連接 $\mathcal{O}_U$-加群なので、結論を得る。
\end{proof}

\begin{lemma}
\label{flat-lemma-complex-and-divisor-eta-pull}
Situation \ref{flat-situation-complex-and-divisor} において、$h : Y \to X$ を、引き戻し
$E = h^{-1}D$ が定義されるスキームの射とする。$(X, D, M)$ が良い三つ組なら、
$$
Lh^*(L\eta_\mathcal{I}M) = L\eta_\mathcal{J}(Lh^*M)
$$
が $D(\mathcal{O}_Y)$ において成り立つ。ここで $\mathcal{J}$ は $E$ のイデアル層である。
\end{lemma}

\begin{proof}
More on Algebra, Lemma \ref{more-algebra-lemma-eta-base-change} の翻訳である。
この翻訳には Lemmas \ref{flat-lemma-eta-stalks} および \ref{flat-lemma-Leta} を用いる。
\end{proof}

\begin{lemma}
\label{flat-lemma-complex-and-divisor-blowup-pre}
Situation \ref{flat-situation-complex-and-divisor} において、次を満たす一意な射
$b : X' \to X$ が存在する。
\begin{enumerate}
\item 引き戻し $D' = b^{-1}D$ が定義され、$M' = Lb^*M$ とおくと
$(X', D', M')$ は良い三つ組である。
\item 引き戻し $E = h^{-1}D$ が定義され、$(Y, E, Lh^*M)$ が良い三つ組となる任意の
スキームの射 $h : Y \to X$ に対して、$h$ は一意に $b$ を経由して分解する。
\end{enumerate}
さらに、任意のアフィンチャート $(U, A, f, M^\bullet)$ に対して、制限
$b^{-1}(U) \to U$ はイデアル $I_i(M^\bullet, f)$ の積に沿う吹き上げであり、任意の
準コンパクト開集合 $W \subset X$ に対して、制限
$b|_{b^{-1}(W)} : b^{-1}(W) \to W$ は $W \setminus D$-許容吹き上げである。
\end{lemma}

\begin{proof}
証明では、任意のアフィンチャート $(U, A, f, M^\bullet)$ に対し、局所的に積イデアル
$I_i(M^\bullet, f)$ に沿って $X$ を吹き上げるだけである。More on Algebra, Section
\ref{more-algebra-section-perfect-eta} の冒頭のいくつかの補題により、これは良定義である。
すると普遍性 (2) は吹き上げの普遍性から従う。以下に詳細を述べる。

\medskip\noindent
$U, A, f, M^\bullet$ を $(X, D, M)$ のアフィンチャートとする。有限個を除くすべての
イデアル $I_i(M^\bullet, f)$ は $A$ に等しいので、
$$
I = \prod\nolimits_i I_i(M^\bullet, f)
$$
を考えることができ、これは $A$ の有限生成イデアルである。
$$
b_U : U' \to U
$$
を $U$ の $I$ に沿う吹き上げと書く。Divisors, Lemma
\ref{divisors-lemma-blow-up-pullback-effective-Cartier} により
$b_U^{-1}(U \cap D)$ は定義される。$f^{r_i} \in I_i(M^\bullet, f)$ であることを
想起すると、$b_U$ は $(U \setminus D)$-許容吹き上げである。Divisors, Lemma
\ref{divisors-lemma-blowing-up-two-ideals} により、各 $i$ に対して射 $b_U$ は
$U' \to U'_i \to U$ と分解する。ここで $U'_i \to U$ は $I_i(M^\bullet, f)$ に沿う
吹き上げであり、$U' \to U'_i$ はもう一つの吹き上げである。従って
$I_i(M^\bullet, f)$ の $U'$ への引き戻し $I_i(M^\bullet, f)\mathcal{O}_{U'}$ は
可逆イデアル層である。Divisors, Lemmas
\ref{divisors-lemma-blow-up-pullback-effective-Cartier} および
\ref{divisors-lemma-blowing-up-gives-effective-Cartier-divisor} を参照せよ。
従って $(U', b^{-1}D, Lb^*M|_U)$ は良い三つ組である。引き戻しに伴うイデアル
$I_i(-,-)$ の振る舞いについては Lemma \ref{flat-lemma-pullback-triple-ideals-good} を参照せよ。
% P05-EN-CH38-0117: source 13638 uses global b inside the local b_U construction before b exists.
最後に、$b_U : U' \to U$ は補題の主張の (2) で述べた普遍性をもつと主張する。
実際、$h : Y \to U$ を、引き戻し $E = h^{-1}(D \cap U)$ が定義され、
$(Y, E, Lh^*M)$ が良い三つ組となるスキームの射とする。このとき $Y$ は、各 $i$ に対して
$I_i(N^\bullet, g)$ が可逆イデアルとなるアフィンチャート $(V, B, g, N^\bullet)$ で覆われる。
$g$ と $f$ の $B$ における像は、どちらも有効 Cartier 因子 $E \cap V$ を切り出すので、
単元倍だけ異なる。従って More on Algebra, Lemma
\ref{more-algebra-lemma-eta-change-unit} により、$g$ は $f$ の像であると仮定してよい。
すると More on Algebra, Lemma \ref{more-algebra-lemma-ideal-well-defined} により、
$I_i(N^\bullet, g)$ は $B$-加群として $I_i(M^\bullet \otimes_A B, g)$ と同型である。
従って $I_i(M^\bullet \otimes_A B, g) = I_i(M^\bullet, f)B$
（Lemma \ref{flat-lemma-pullback-triple-ideals-good}）は可逆 $B$-加群である。
ゆえにイデアル $IB$ は可逆であり、従って $I\mathcal{O}_Y$ も可逆である。Divisors, Lemma
\ref{divisors-lemma-universal-property-blowing-up} により、$h$ の $b_U$ を経由する一意な分解を得る。

\medskip\noindent
$\mathcal{B}$ を、$(X, D, M)$ のアフィンチャート $(U, A, f, M^\bullet)$ が存在する
アフィン開集合 $U \subset X$ の集合とする。すると $\mathcal{B}$ は $X$ の位相の基底である。
詳細は省略する。$U \in \mathcal{B}$ に対し、上で構成した射 $b_U : U' \to U$ は
$U$ 上の普遍性を満たす。$U_1 \subset U_2 \subset X$ がともに $\mathcal{B}$ に属するなら、
$b_{U_1} : U'_1 \to U_1$ は標準的に
$$
b_{U_2}|_{b_{U_2}^{-1}(U_1)} : b_{U_2}^{-1}(U_1) \longrightarrow U_1
$$
と同型である。ここでも普遍性を用いた。換言すれば、$U_1$ 上の同型
$U'_1 \to b_{U_2}^{-1}(U_1)$ を得る。これらの同型はコサイクル条件を満たす
（再び普遍性による）。従って Constructions, Lemma
\ref{constructions-lemma-relative-glueing} により射 $b : X' \to X$ を得て、その
$\mathcal{B}$ の各 $U$ への制限は $U' \to U$ と同型である。これらの性質は局所的に
確認できるので（詳細は省略する）、射 $b : X' \to X$ は補題の主張の性質 (1), (2) を満たす。

\medskip\noindent
補題の最後の主張を示す必要がある。$W \subset X$ を準コンパクト開集合とする。有限被覆
$W = U_1 \cup \ldots \cup U_T$ で、各 $1 \leq t \leq T$ に対してアフィンチャート
$(U_t, A_t, f_t, M_t^\bullet)$ が存在するものを選ぶ。任意のアフィン開集合
$V = \Spec(B) \subset U_t \cap U_{t'}$ に対して、次を用いる。
(a) $f_t$ と $f_{t'}$ の $B$ における像は単元倍だけ異なる。
(b) 複体 $M_t^\bullet \otimes_A B$ と $M_{t'} \otimes_A B$ は $D(B)$ の同型な対象を定める。
% P05-EN-CH38-0119: source 13691 tensors both chart complexes over undefined A instead of A_t and A_{t'}.
$i \in \mathbf{Z}$ に対して
$$
N_i = \max\nolimits_{t = 1, \ldots, T}
\left(\sum\nolimits_{j \geq i} (-1)^{j - i} rk(M_t^j)\right)
$$
とおく。このとき $N_t - \sum\nolimits_{j \geq i} (-1)^{j - i} rk(M_t^j) \geq 0$ であり、
次のイデアルを考えることができる。
% P05-EN-CH38-0118: source 13697 uses N_t although the maximum was defined as N_i.
$$
I_{t, i} =
f_t^{N_i - \sum\nolimits_{j \geq i} (-1)^{j - i} rk(M_t^j)}
I_i(M_t^\bullet, f_t)
\subset A_t
$$
More on Algebra, Lemmas \ref{more-algebra-lemma-eta-change-unit} および
\ref{more-algebra-lemma-ideal-well-defined} により、イデアル $I_{t, i}$ は準連接有限型
イデアル $\mathcal{I}_i \subset \mathcal{O}_W$ に貼り合わさる。さらに有限個を除く
すべてのこれらのイデアルは $\mathcal{O}_W$ に等しい。上で構成した射 $X' \to X$ の
$W$ への制限がイデアル $\mathcal{I}_i$ の積に沿う吹き上げであることは明らかである。
これで証明が完了する。
\end{proof}

\begin{lemma}
\label{flat-lemma-complex-and-divisor-blowup}
Situation \ref{flat-situation-complex-and-divisor} において、$b : X' \to X$ を Lemma
\ref{flat-lemma-complex-and-divisor-blowup-pre} の射とする。イデアル層
$\mathcal{I}' \subset \mathcal{O}_{X'}$ をもつ有効 Cartier 因子 $D' = b^{-1}D$ を考える。
このとき $Q = L\eta_{\mathcal{I}'}Lb^*M$ は $D(\mathcal{O}_{X'})$ の完全対象である。
\end{lemma}

\begin{proof}
Lemmas \ref{flat-lemma-complex-and-divisor-blowup-pre} および \ref{flat-lemma-good-triple} から従う。
\end{proof}

\begin{lemma}
\label{flat-lemma-complex-and-divisor-blowup-base-change}
Situation \ref{flat-situation-complex-and-divisor} において、$h : Y \to X$ を、引き戻し
$E = h^{-1}D$ が定義されるスキームの射とする。$b : X' \to X$ および $c : Y' \to Y$ を、
それぞれ $D \subset X$ と $M$、$E \subset Y$ と $Lh^*M$ に対して Lemma
\ref{flat-lemma-complex-and-divisor-blowup-pre} で構成したものとする。このとき $Y'$ は
$b : X' \to X$ に関する $Y$ の厳密変換であり（正確な定式化は証明を参照せよ）、
$$
L\eta_{\mathcal{J}'}L(h \circ c)^*M = L(Y' \to X')^*Q
$$
が成り立つ。ここで $Q = L\eta_{\mathcal{I}'}Lb^*M$ は Lemma
\ref{flat-lemma-complex-and-divisor-blowup} のものである。特に $(Y, E, Lh^*M)$ が
良い三つ組で、$k : Y \to X'$ が $h = b \circ k$ を満たす一意な射なら、
$L\eta_\mathcal{J}Lh^*M = Lk^*Q$ である。
\end{lemma}

\begin{proof}
$E' = c^{-1}E$ と書く。すると $(Y', E', L(h \circ c)^*M)$ は良い三つ組である。
従って Lemma \ref{flat-lemma-complex-and-divisor-blowup-pre} の普遍性により、一意な射
$$
h' : Y' \longrightarrow X'
$$
で $b \circ h' = h \circ c$ となるものが存在する。特に射
$(h', c) : Y' \to X' \times_X Y$ がある。$W \subset X$ を準コンパクト開集合で、
$b^{-1}(W) \to W$ が吹き上げとなるものとすると、この射は $Y'|_W$ を
$b^{-1}(W) \to W$ に関する $Y_W$ の厳密変換と同一視すると主張する。これは $W$ 上の
局所的な問題なので、アフィンチャート上で主張を証明すればよい。次の段落でそうする。

\medskip\noindent
$(U, A, f, M^\bullet)$ を $(X, D, M)$ のアフィンチャートとする。Lemma
\ref{flat-lemma-complex-and-divisor-blowup} の証明から、$b : X' \to X$ の $U$ への制限は
$U = \Spec(A)$ のイデアル $I_i(M^\bullet, f)$ の積に沿う吹き上げであることを想起する。
% P05-EN-CH38-0120: source 13767 cites blowup rather than the product-ideal construction in blowup-pre.
$V = \Spec(B) \subset Y$ を $h(V) \subset U$ となる任意のアフィン開集合とすると、
$g \in B$ を $f$ の像としたとき、$(V, B, g, M^\bullet \otimes_A B)$ は
$(Y, E, Lh^*M)$ のアフィンチャートである。Lemma
\ref{flat-lemma-pullback-triple-ideals-good} を参照せよ。従って $c : Y' \to Y$ の $V$ への制限は
イデアル $I_i(M^\bullet, f)B$ の積に沿う吹き上げである。すなわち、$h^{-1}(U)$ 上の射
$c : Y' \to Y$ は、$h^{-1}(U)$ のイデアル
$\prod I_i(M^\bullet, f) \mathcal{O}_{h^{-1}(U)}$ に沿う吹き上げである。
厳密変換についても同じことが成り立つので、厳密変換に関する主張を得る。

\medskip\noindent
以上を踏まえると、等式
$L\eta_{\mathcal{J}'}L(h \circ c)^*M = L(Y' \to X')^*Q$ は Lemma
\ref{flat-lemma-complex-and-divisor-eta-pull} から従う。最後の主張はこの特別な場合、すなわち
$c = \text{id}_Y$ かつ $k = h'$ の場合である。
\end{proof}

\begin{lemma}
\label{flat-lemma-complex-and-divisor-blowup-good}
Situation \ref{flat-situation-complex-and-divisor} において、$W \subset X$ を、$M$ のコホモロジー層が
局所自由となる最大の開部分スキームとする。このとき Lemma
\ref{flat-lemma-complex-and-divisor-blowup-pre} の射 $b : X' \to X$ は $W$ 上で同型である。
\end{lemma}

\begin{proof}
$U \subset W$ となる任意のアフィンチャート $(U, A, f, M^\bullet)$ に対して、More on Algebra,
Lemma \ref{more-algebra-lemma-eta-cohomology-free} により $I_i(M^\bullet, f)$ は局所的に
$f$ の冪で生成されるからである。$f$ は非零因子なので、吹き上げ
$b^{-1}(U) \to U$ は同型である。
\end{proof}

\begin{lemma}
\label{flat-lemma-complex-and-divisor-blowup-good-T}
$X, D, \mathcal{I}, M$ を Situation \ref{flat-situation-complex-and-divisor} のようなものとする。
$(X, D, M)$ が良い三つ組なら、有限表示の閉埋め込み
$$
i : T \longrightarrow D
$$
で次の性質をもつものが存在する。
\begin{enumerate}
\item $W \subset X$ を $M$ のコホモロジー層が局所自由となる最大の開集合とすると、
$T$ は $D \cap W$ をスキーム論的に含む。
% P05-EN-CH38-0121: source 13820 and four later loci use ``scheme theoretically'' without the compound-adverb hyphen.
\item $Li^*L\eta_\mathcal{I}M$ のコホモロジー層は局所自由である。
\item 像が $x = i(t) \in W$ となる任意の点 $t \in T$ に対して、
$H^i(M)_x$ の $\mathcal{O}_{X, x}$ 上の階数と
$H^i(Li^*L\eta_\mathcal{I}M)_t$ の $\mathcal{O}_{T, t}$ 上の階数は一致する。
\end{enumerate}
\end{lemma}

\begin{proof}
$(U, A, f, M^\bullet)$ を $(X, D, M)$ のアフィンチャートで、すべての
$i \in \mathbf{Z}$ に対して $I_i(M^\bullet, f)$ が単項イデアルとなるものとする。
$T \cap U \subset D \cap U$ を、イデアル
$$
J(M^\bullet, f) = \sum J_i(M^\bullet, f) \subset A/fA
$$
で定義される閉部分スキームと定める。このイデアルは More on Algebra, Lemmas
\ref{more-algebra-lemma-eta-vanishing-beta-plus-pre} および
\ref{more-algebra-lemma-eta-vanishing-beta-plus} で調べられている。第二の補題の記法では、
$T \cap U \to D \cap U$ はそこで調べた環準同型 $A/fA \to C$ で与えられる。
$(X, D, M)$ は良い三つ組なので、この形のアフィンチャートで $X$ を覆うことができ、
二つの補題の第一のものにより、この構成は貼り合わさる。従って閉部分スキーム
$T \subset D$ を得て、上のような良いアフィンチャート上ではイデアル
$J(M^\bullet, f)$ で与えられる。すると性質 (1), (2) は第二の補題から従う。詳細は省略する。
読者の助けとなる小さな観察を述べる。More on Algebra, Lemma
\ref{more-algebra-lemma-eta-locally-free} により $\eta_fM^\bullet$ は局所自由加群の複体なので、
$Li^*L\eta_\mathcal{I}M|_{T \cap U}$ は $C$-加群の複体
$\eta_fM^\bullet \otimes_A C$ で表示される。階数に関する主張 (3) は Cohomology, Lemma
\ref{cohomology-lemma-eta-cohomology-locally-free} から従う。
\end{proof}

\begin{lemma}
\label{flat-lemma-complex-and-divisor-blowup-T}
Situation \ref{flat-situation-complex-and-divisor} において、$b : X' \to X$ と $D'$ を Lemma
\ref{flat-lemma-complex-and-divisor-blowup-pre} のようなものとし、
$Q = L\eta_{\mathcal{I}'}Lb^*M$ を Lemma \ref{flat-lemma-complex-and-divisor-blowup} の
ようなものとする。$W \subset X$ を、$M$ が局所自由なコホモロジー加群をもつ最大の
開集合とする。このとき有限表示の閉埋め込み $i : T \to D'$ で次を満たすものが存在する。
\begin{enumerate}
\item $D' \cap b^{-1}(W) \subset T$ がスキーム論的に成り立つ。
\item $Li^*Q$ は局所自由なコホモロジー層をもつ。
\item $w \in W$ に写る $t \in T$ に対して、$H^i(Li^*Q)_t$ の
$\mathcal{O}_{T, t}$ 上の階数は $H^i(M)_x$ の $\mathcal{O}_{X, x}$ 上の階数に等しい。
% P05-EN-CH38-0122: source 13872--13874 and 13951--13953 introduce w but use undefined x in the rank comparison.
\end{enumerate}
\end{lemma}

\begin{proof}
Lemma \ref{flat-lemma-complex-and-divisor-blowup-good} により、$b$ は $W$ 上で同型である。
従って $b^{-1}(W) \subset X'$ は、$Lb^*M$ が局所自由なコホモロジー層をもつ最大の
開集合 $W' \subset X'$ に含まれる。そこで補題の実際の主張は、
$(X', D', Lb^*M)$ に適用した Lemma \ref{flat-lemma-complex-and-divisor-blowup-good-T} と、
主張中で言及した他の補題から直ちに従う。
\end{proof}

\begin{lemma}
\label{flat-lemma-complex-and-divisor-blowup-T-ranks}
Situation \ref{flat-situation-complex-and-divisor} において、$b : X' \to X$、$D' \subset X'$、$Q$ を
Lemma \ref{flat-lemma-complex-and-divisor-blowup} のようなものとする。
$\rho = (\rho_i)_{i \in \mathbf{Z}}$ を整数の族とする。$W(\rho) \subset X$ を、すべての
$i$ に対して $H^i(M)$ が階数 $\rho_i$ の局所自由加群となる最大の開部分スキームとする。
$i : T \to D'$ を Lemma \ref{flat-lemma-complex-and-divisor-blowup-T} のようなものとする。
このとき $D' \cap b^{-1}(W(\rho))$ をスキーム論的に含む開かつ閉な部分スキーム
$T(\rho) \subset T$ が存在し、すべての $i$ に対して
$H^i(Li^*Q|_{T(\rho)})$ は階数 $\rho_i$ の局所自由加群である。
\end{lemma}

\begin{proof}
$T(\rho) \subset T$ を、すべての $i$ に対して $H^i(Li^*Q)$ が階数 $\rho_i$ をもつ
開かつ閉な部分スキームとする。このとき主張は、コホモロジー加群の階数に関する Lemma
\ref{flat-lemma-complex-and-divisor-blowup-T} の主張から直ちに従う。
\end{proof}

\begin{lemma}
\label{flat-lemma-complex-and-divisor-blowup-T-represented-bdd-loc-free}
Situation \ref{flat-situation-complex-and-divisor} において、$b : X' \to X$、$D' \subset X'$、$Q$ を
Lemma \ref{flat-lemma-complex-and-divisor-blowup} のようなものとする。$M$ を表示する有限局所自由
$\mathcal{O}_X$-加群の局所有界複体 $\mathcal{M}^\bullet$ が存在するなら、$Q$ を表示する
有限局所自由 $\mathcal{O}_{X'}$-加群の局所有界複体 $\mathcal{Q}^\bullet$ が存在する。
\end{lemma}

\begin{proof}
$Q = L\eta_{\mathcal{I}'}Lb^*M$ であり、$\mathcal{I}'$ は有効 Cartier 因子 $D'$ の
イデアル層であることを想起する。有限局所自由 $\mathcal{O}_{X'}$-加群の局所有界複体
$(\mathcal{M}')^\bullet = b^*\mathcal{M}^\bullet$ は $Lb^*M$ を表示する。従って補題は
Lemma \ref{flat-lemma-good-triple-bdd-loc-free} から従う。
\end{proof}

\begin{lemma}
\label{flat-lemma-complex-and-divisor-derived}
$X$ をスキーム、$D \subset X$ を有効 Cartier 因子とする。
$M \in D(\mathcal{O}_X)$ を完全対象とする。$W \subset X$ を、コホモロジー層
$H^i(M)$ が局所自由となる最大の開集合とする。このとき固有射
$b : X' \longrightarrow X$ と $D(\mathcal{O}_{X'})$ の対象 $Q$ で、次の性質をもつものが存在する。
\begin{enumerate}
\item $b : X' \to X$ は $X \setminus D$ 上で同型である。
\item $b : X' \to X$ は $W$ 上で同型である。
\item $D' = b^{-1}D$ は有効 Cartier 因子である。
\item $\mathcal{I}'$ を $D'$ のイデアル層とすると、$Q = L\eta_{\mathcal{I}'}Lb^*M$ である。
\item $Q$ は $D(\mathcal{O}_{X'})$ の完全対象である。
\item 有限表示の閉埋め込み $i : T \to D'$ で、次を満たすものが存在する。
\begin{enumerate}
\item $D' \cap b^{-1}(W) \subset T$ がスキーム論的に成り立つ。
\item $Li^*Q$ は有限局所自由なコホモロジー層をもつ。
\item 像が $w \in W$ となる $t \in T$ に対して、$H^i(Li^*Q)_t$ の
$\mathcal{O}_{T, t}$ 上の階数は $H^i(M)_x$ の $\mathcal{O}_{X, x}$ 上の階数に等しい。
\end{enumerate}
% P05-EN-CH38-0122: source 13951--13953 again introduces w but uses undefined x in the rank comparison.
\item $(X, D, M)$ の任意のアフィンチャート $(U, A, f, M^\bullet)$ に対して、$b$ の $U$ への
制限は $U = \Spec(A)$ のイデアル $I = \prod I_i(M^\bullet, f)$ に沿う吹き上げである。
\item $(X', D', Lb^*N)$ の任意のアフィンチャート $(V, B, g, N^\bullet)$ で、
$I_i(N^\bullet, g)$ が単項となるものに対して、次が成り立つ。
% P05-EN-CH38-0123: source 13958 uses Lb^*N although M is the pulled-back global object and N is the chart complex.
\begin{enumerate}
\item $Q|_V$ は $\eta_gN^\bullet$ に対応する。
\item $T \subset V \cap D'$ は More on Algebra, Lemma
\ref{more-algebra-lemma-eta-vanishing-beta-plus} で調べたイデアル
$J(N^\bullet, g) = \sum J_i(N^\bullet, g) \subset B/gB$ に対応する。
\end{enumerate}
\item $M$ が有限局所自由 $\mathcal{O}_X$-加群の局所有界複体で表示できるなら、
$Q$ は有限局所自由 $\mathcal{O}_{X'}$-加群の有界複体で表示できる。
% P05-EN-CH38-0124: source 13967--13970 drops ``locally'' only from the conclusion.
\end{enumerate}
\end{lemma}

\begin{proof}
この主張は、Lemmas \ref{flat-lemma-pullback-triple-ideals-good},
\ref{flat-lemma-good-triple},
\ref{flat-lemma-complex-and-divisor-eta-pull},
\ref{flat-lemma-complex-and-divisor-blowup-pre},
\ref{flat-lemma-complex-and-divisor-blowup},
\ref{flat-lemma-complex-and-divisor-blowup-base-change},
\ref{flat-lemma-complex-and-divisor-blowup-good},
\ref{flat-lemma-complex-and-divisor-blowup-good-T},
\ref{flat-lemma-complex-and-divisor-blowup-T}, および
\ref{flat-lemma-complex-and-divisor-blowup-T-represented-bdd-loc-free} で得た情報をまとめたものである。
\end{proof}










\section{複体の吹き上げ、III}
\label{flat-section-blowup-complexes-III}

\noindent
この節では、\cite[Section 18.1]{F} で与えられた Macpherson のグラフ構成の
「代数的版」を与える。
% P05-EN-CH38-0127: source 14001 repeats ``version'' in ``an algebra version of the version''.

\medskip\noindent
$X$ をスキーム、$E$ を $D(\mathcal{O}_X)$ の完全対象とする。
$U \subset X$ を、$E|_U$ が局所自由なコホモロジー層をもつ最大の開部分スキームとする。

\medskip\noindent
可換図式
$$
\xymatrix{
\mathbf{A}^1_X \ar[r] \ar[rd] &
\mathbf{P}^1_X \ar[d]^p &
(\mathbf{P}^1_X)_\infty \ar[l] \ar[ld] \\
& X \ar@/_1em/[ur]_\infty
}
$$
を考える。$\mathbf{A}^1 = D_+(T_0)$ は $\mathbf{P}^1$ の第一標準アフィン開集合であり、
$\infty = V_+(T_0)$ はその補有効 Cartier 因子であることを想起する。上の図式は、これらの
スキームの $X$ への引き戻しである。$\infty : X \to (\mathbf{P}^1_X)_\infty$ は同型である。
すると
$$
(\mathbf{P}^1_X, (\mathbf{P}^1_X)_\infty, Lp^*E)
$$
は Section \ref{flat-section-blowup-complexes-II} の Situation
\ref{flat-situation-complex-and-divisor} のような三つ組である。
$$
b : W \longrightarrow \mathbf{P}^1_X\quad\text{および}\quad
W_\infty = b^{-1}((\mathbf{P}^1_X)_\infty)
$$
を、この三つ組から出発して Lemma \ref{flat-lemma-complex-and-divisor-blowup-pre} で構成した
吹き上げと有効 Cartier 因子とする。また
$$
Q = L\eta_{\mathcal{I}} Lb^*M = L\eta_\mathcal{I} L(p \circ b)^*E
$$
を Lemma \ref{flat-lemma-complex-and-divisor-blowup} で考えた $D(\mathcal{O}_W)$ の
完全対象と書く。ここで $\mathcal{I} \subset \mathcal{O}_W$ は $W_\infty$ のイデアル層である。
% P05-EN-CH38-0125: source 14040 uses undefined M in Lb^*M; the triple contains Lp^*E.

\begin{lemma}
\label{flat-lemma-graph-construction}
上の構成は次の性質をもつ。
\begin{enumerate}
\item $b$ は $\mathbf{P}^1_U \cup \mathbf{A}^1_X$ 上で同型である。
\item $Q$ の $\mathbf{A}^1_X$ への制限は $E$ の引き戻しに等しい。
\item 有限表示の閉埋め込み $i : T \to W_\infty$ で、
$(W_\infty \to X)^{-1}U \subset T$ がスキーム論的に成り立ち、かつ $Li^*Q$ が
局所自由なコホモロジー層をもつものが存在する。
\item 像が $u \in U$ となる $t \in T$ に対して、$H^i(Li^*Q)_t$ の
$\mathcal{O}_{T, t}$ 上の階数は $H^i(M)_u$ の $\mathcal{O}_{U, u}$ 上の階数に等しい。
\item $E$ が有限局所自由 $\mathcal{O}_X$-加群の局所有界複体で表示できるなら、$Q$ は
有限局所自由 $\mathcal{O}_W$-加群の局所有界複体で表示できる。
\end{enumerate}
% P05-EN-CH38-0126: source 14060 uses undefined M in the rank comparison; the section's object is E.
\end{lemma}

\begin{proof}
これは Section \ref{flat-section-blowup-complexes-II} の結果から直ちに従う。必要な内容を
すべてまとめた主張については Lemma \ref{flat-lemma-complex-and-divisor-derived} を参照せよ。
\end{proof}
